\documentclass[10pt]{article}

\usepackage[a4paper,margin=0.88in]{geometry}
\usepackage{amsmath,amssymb,amsthm,mathtools}
\usepackage{bm}
\usepackage{mathrsfs}
\usepackage{enumitem}
\usepackage{microtype}
\usepackage[hidelinks,hypertexnames=false]{hyperref}
\usepackage[nameinlink,capitalize]{cleveref}

\newtheorem{theorem}{Theorem}[section]
\newtheorem{proposition}[theorem]{Proposition}
\newtheorem{lemma}[theorem]{Lemma}
\newtheorem{corollary}[theorem]{Corollary}
\newtheorem{definition}[theorem]{Definition}
\newtheorem{remark}[theorem]{Remark}
\newtheorem{example}[theorem]{Example}
\newtheorem{principle}[theorem]{Principle}
\newtheorem{problem}[theorem]{Problem}
\newtheorem{conjecture}[theorem]{Conjecture}

\newcommand{\R}{\mathbb R}
\newcommand{\Z}{\mathbb Z}
\newcommand{\N}{\mathbb N}
\newcommand{\Sph}{\mathbb S}
\newcommand{\Lag}{\operatorname{Lag}}
\newcommand{\Leg}{\operatorname{Leg}}
\newcommand{\Sat}{\operatorname{Sat}}
\newcommand{\Front}{\operatorname{Front}}
\newcommand{\Graph}{\operatorname{Graph}}
\newcommand{\proj}{\operatorname{pr}}
\newcommand{\rank}{\operatorname{rank}}
\newcommand{\im}{\operatorname{im}}
\newcommand{\dist}{\operatorname{dist}}
\newcommand{\Gr}{\operatorname{Gr}}
\newcommand{\supp}{\operatorname{supp}}
\newcommand{\dimH}{\dim_{\mathrm H}}
\newcommand{\dimP}{\dim_{\mathrm P}}
\newcommand{\dimM}{\dim_{\mathrm M}}
\newcommand{\vol}{\operatorname{vol}}
\newcommand{\Span}{\operatorname{Span}}
\newcommand{\id}{\operatorname{id}}

\title{Sticky Kakeya in $\R^4$ as a Contact--Symplectic Incidence Problem\\
\large Contact Geometry, Maslov Four-Cycles, and Lossless Edge Flow}
\author{Chenxi Cai\\[0.35em]\normalsize\texttt{Chenxi\_Cai@live.com}}
\date{}

\begin{document}
\allowdisplaybreaks[2]
\setlength{\emergencystretch}{3em}
\maketitle

\begin{abstract}
We identify the space of oriented affine lines in $\R^4$ with the symplectic manifold
$T^*S^3$ and marked line segments with Reeb intervals in its contactification
$J^1(S^3)$. Euclidean points become Legendrian probes, so the Sticky Kakeya problem is
an exact contact incidence problem. The infinitesimal Euclidean front is governed by the
matrix pencil $B+sA$. Its determinant satisfies
\[
 |\det(B+iA)|^2=1-\|\omega|_V\|_{\mathrm{op}}^2,
\]
and front collapse is equivalent to simultaneous incidence of the tangent plane $V$ with
the full Lagrangian pencil $L_s=\{(a,-sa)\}$. This isolates three polarized Maslov
degeneracies. At small scales, collision pairs converge to these Maslov incidences;
nondegenerate four-cycles produce physical Lagrangian bushes, while persistent rank loss
is ruled out by the polarized classification.  A mass-preserving edge-flow argument then
propagates the bush alternative across the full stopping tree without squaring the incoming mass.
The key point is to exhaust every fixed-proportion route on its unused occurrence measure before
passing to the next generation.  The resulting absolute flows telescope with coefficient one, so
the apparent $C^L$ loss and the compensation-zone debt disappear.  A deterministic subpower cap
schedule gives a quadratic terminal tail.  The same argument is linear under arbitrary source
restrictions and carries affine fibre and carrier-tree marks as conditional variables; it therefore
yields the uniform shaded, weighted finite-scale estimate required by marked applications.
\end{abstract}

\tableofcontents

\section{Introduction}

Let $\mathcal L^+(\R^4)$ denote the space of oriented affine lines. A full-direction
family contains a line in every direction. Following the continuum formulation of stickiness
\cite{KatzLabaTao2000,KatzTao2002,WangZahlSticky}, the extremal condition is that this family
has packing dimension three, the smallest possible dimension compatible with surjectivity onto
$S^3$.

The purpose of this paper is to exploit the geometry of line space rather than replace it by
an auxiliary finite incidence structure. The key observation is that
\[
 \mathcal L^+(\R^4)\cong T^*S^3
\]
canonically as a symplectic manifold. Adding the marked position along a line produces the
contact manifold $J^1(S^3)$. Translation along a line is its Reeb flow, and incidence with a Euclidean point is
intersection with a distinguished Legendrian sphere; see
\cite{CannasDaSilva2001,Geiges2008,Arnold1990} for the symplectic, contact, and wave-front
background.

This reformulation separates the smooth and rectifiable regimes from the genuinely fractal
one. It also identifies the exact degenerate object: not a small scalar symplectic defect,
but a three-plane incident with every member of a one-parameter Lagrangian pencil. The same
pencil reappears after blowing up small-angle collisions. Thus the infinitesimal and
multiscale parts of the proof are governed by one symplectic object.

\subsection{Main results}

Let $\pi:T^*S^3\to S^3$ be the bundle projection and let
$q:J^1(S^3)\to T^*S^3$ forget the jet coordinate.

\begin{definition}[Sticky Kakeya datum]
A compact marked family $\Gamma\subset J^1(S^3)$ is a \emph{Sticky Kakeya datum} if
\[
 \pi(q(\Gamma))=S^3,
 \qquad
 \dimP q(\Gamma)=3,
\]
and every line in $q(\Gamma)$ carries at least one marked unit interval represented in
$\Gamma$.
\end{definition}

The Euclidean union $K_\Gamma$ is defined precisely in \cref{sec:contact-model}.

\begin{theorem}[Sticky Kakeya theorem in four dimensions]
\label{thm:main}
For every compact Sticky Kakeya datum $\Gamma$ in $\R^4$,
\[
 \boxed{\dimH K_\Gamma=4.}
\]
\end{theorem}

\subsection{Proof architecture}

The proof has four stages.
\begin{enumerate}[label=\textnormal{(\roman*)}]
\item The contact model converts marked segments and Euclidean point incidence into Reeb
intervals and Legendrian intersections.
\item The Reeb-front polynomial $P_V(s)=\det(B+sA)$ detects infinitesimal spreading. If it
vanishes identically, $V$ is trapped by the full Maslov pencil and belongs to one of three
polarized strata.
\item In a slope chart, near-collisions have residual
$r=\beta+s_*\alpha\in\alpha^\perp$. This two-dimensional normal residual is the exact
quantity needed for a Frostman estimate. Angular collapse rescales back to the Maslov pencil;
a nondegenerate four-cycle produces a physical point-probe bush.
\item Bush returns are organized by an incoming edge-flow measure.  Exhausting each hereditary
Markov route makes the absolute flow telescope across generations.  A deterministic cap schedule
then leaves a quadratic terminal tail with no generation factor.
\end{enumerate}

The symplectic argument is therefore not an interpretation added after the incidence proof. It
classifies angular collapse, supplies the four-time obstruction, and keeps each finite physical
bush iteration from losing a power of mass.  The remaining summation is closed by the absolute Markov-flow identity in the final
quantitative section.

\subsection{Where the endpoint is won}

Three points deserve to be visible before the technical argument begins.  First, the relevant
degeneracy is not the vanishing of one chosen symplectic scalar.  It is the intrinsic condition
that the tangent plane meet every member of the Lagrangian pencil \(L_s\).  This is why the
matrix pencil and the contact blow-up lead to the same alternatives.

Second, the graph-theoretic four-cycle is used only as a sampling device.  It is credited only
after the contact equations identify a common Euclidean point and a common Reeb window.  The
resulting object is a physical Lagrangian bush, and its old collision partner and old Maslov flag
remain attached to the source occurrence.

Third, iteration is performed on the ordered edge measure, not on a repeatedly normalized vertex
measure.  Each local alternative is exhausted on the unused part.  Children are therefore
fractional restrictions of one absolute source, and their masses telescope with coefficient one.
This is the mechanism that closes both the compensation regime and the marked finite-scale
interface.

\subsection{Organization}

The first two sections establish the contact model and the selector reduction.  The six
quantitative sections then follow the proof itself: front-null pencils, physical collision
recurrence, angular Maslov extraction, polarized residual routing, bush stopping, and
cross-generation edge-flow closure.  The appendices contain only measurable disintegration
and the uniform finite-scale parameter hierarchy.

\section{The contact model of marked lines}
\label{sec:contact-model}

Every oriented affine line has a unique representation
\[
 \ell_{\theta,y}=\{y+t\theta:t\in\R\},
 \qquad
 \theta\in S^3,\quad y\in\theta^\perp.
\]
Using the Euclidean metric to identify $\theta^\perp$ with $T^*_\theta S^3$, write
$p=y^\flat$. This gives a diffeomorphism
\[
 \mathcal L^+(\R^4)\longrightarrow T^*S^3,
 \qquad
 \ell_{\theta,y}\longmapsto(\theta,p).
\]
Let $\lambda$ be the tautological one-form and $\omega=d\lambda$.

Adjoin the longitudinal variable $z$ and set
\[
 J^1(S^3)=T^*S^3\times\R_z,
 \qquad
 \alpha=dz-\lambda.
\]
The Reeb vector field is $R_\alpha=\partial_z$. The Euclidean evaluation map is
\begin{equation}
\label{eq:evaluation}
 \mathcal E(\theta,p,z)=p^\sharp+z\theta.
\end{equation}

\begin{proposition}[Euclidean incidence form]
Under \cref{eq:evaluation},
\[
 \theta\cdot dx=dz-p\cdot d\theta=\alpha.
\]
In particular, Reeb translation is Euclidean translation along the oriented line.
\end{proposition}

\begin{proof}
Differentiate $x=p^\sharp+z\theta$. Since $p^\sharp\cdot\theta=0$, one has
$\theta\cdot dp^\sharp=-p^\sharp\cdot d\theta=-p\cdot d\theta$. The identities
$\theta\cdot\theta=1$ and $\theta\cdot d\theta=0$ give the claim.
\end{proof}

For $\Gamma\subset J^1(S^3)$ define its unit Reeb saturation and Euclidean front by
\[
 \Sat^R_{1/2}(\Gamma)
 =\{(\theta,p,z+t):(\theta,p,z)\in\Gamma,\ |t|\le\tfrac12\},
 \qquad
 K_\Gamma=\mathcal E(\Sat^R_{1/2}(\Gamma)).
\]

For $x\in\R^4$, let
\[
 h_x(\theta)=x\cdot\theta,
 \qquad
 \Lambda_x=j^1h_x
 =\{(\theta,dh_x(\theta),h_x(\theta)): \theta\in S^3\}.
\]
Since $(j^1h_x)^*\alpha=0$, the sphere $\Lambda_x$ is Legendrian.

\begin{theorem}[Exact Legendrian incidence]
\label{thm:legendrian-incidence}
For every $x\in\R^4$,
\[
 x\in K_\Gamma
 \quad\Longleftrightarrow\quad
 \Lambda_x\cap\Sat^R_{1/2}(\Gamma)\ne\varnothing.
\]
\end{theorem}

\begin{proof}
If $x=p^\sharp+(z+t)\theta$, then $x\cdot\theta=z+t$ and the projection of $x$ to
$\theta^\perp$ is $p^\sharp$. Hence $p=dh_x(\theta)$ and the corresponding point of the
Reeb interval lies on $\Lambda_x$. The converse is the same calculation in reverse.
\end{proof}

Forgetting $z$ turns $\Lambda_x$ into the exact Lagrangian graph
$L_x=\Graph(dh_x)\subset T^*S^3$. The contact condition remembers the additional filter
$|h_x(\theta)-z|\le1/2$. This primitive window is essential: unmarked Lagrangian incidence
alone does not distinguish the chosen unit segment on a line.

\section{Selector reduction and the regular regime}

Compactness and measurable selection allow us to choose one marked line in every direction;
we use the standard Borel uniformization framework of \cite{Kechris1995}.

\begin{proposition}[Borel selector reduction]
\label{prop:selector}
Every compact full-direction marked family contains a Borel subset meeting each direction fibre
in exactly one point. If the original unmarked family has packing dimension three, so does the
selected line graph. Its Euclidean front is contained in the original front.
\end{proposition}

\begin{proof}
The compact fibre relation over the Polish base $S^3$ has a Borel selector. Its graph projects
onto $S^3$, so its packing dimension is at least three; containment in the original line family
gives the reverse inequality. Removing segments can only shrink the Euclidean union.
\end{proof}

Thus it is enough to study a Borel marked section
$s(\theta)=(\theta,p(\theta),z(\theta))$ whose line-space graph has packing dimension three.
If this section were $C^1$, its Reeb sweep would be
\[
 F(\theta,t)=p^\sharp(\theta)+(z(\theta)+t)\theta.
\]

\begin{proposition}[Smooth selectors spread]
\label{prop:smooth}
For a $C^1$ selector, the four-dimensional Jacobian of $F$ is
\[
 J_F(\theta,t)=\det\bigl(\nabla p^\sharp(\theta)+(z(\theta)+t)I\bigr).
\]
For each fixed $\theta$ this is a monic cubic in $t$. Hence the front contains a nonempty
open set.
\end{proposition}

\begin{proof}
For $v\in T_\theta S^3$, the tangential component of $dF(v)$ is
$\nabla_vp^\sharp+(z+t)v$, while $\partial_tF=\theta$ is normal to $T_\theta S^3$.
The determinant is therefore the displayed $3\times3$ determinant. It cannot vanish for all
$t$, and the inverse function theorem applies at a rank-four point.
\end{proof}

The difficult case is consequently neither smooth nor carried by a positive-mass regular
rectifiable sheet. The next section identifies the exact tangent obstruction.

\clearpage
\section*{The complete quantitative argument}
\addcontentsline{toc}{section}{The complete quantitative argument}
The contact model and selector reduction are now fixed.  What follows is the complete
quantitative proof, ordered by its actual geometric flow.  The front-null pencil is treated
first; physical collisions, angular Maslov extraction, polarized residual routing, bush
stopping, and cross-generation conservation follow without a second abbreviated proof.

\section{Reeb-front pencils and polarized null geometry}

\subsection{Matrix pencils and front rank}
The singular-pencil terminology and canonical forms used below follow
\cite{GohbergLancasterRodman2009,VanDooren1979,BergerTrenn2012}.


\begin{theorem}[Exact Reeb-front rank criterion]
\label{thm:v003-rank-criterion}
At Reeb time \(t\), put
\[
s=z+t.
\]
Then the Reeb-saturated four-plane
\[
W\oplus\R R_\alpha
\]
maps with Euclidean rank four exactly when
\begin{equation}
\label{eq:v003-rank-criterion}
\boxed{
P_V(z+t)\neq0.
}
\end{equation}
Consequently a front-regular tangent plane has rank-four Reeb sweep for every \(t\) except at at
most three values.
\end{theorem}

\begin{proof}
Modulo the Reeb image direction \(\theta\), the Euclidean differential sends
\[
(a,b)\longmapsto b+(z+t)a
\in\theta^\perp.
\]
Relative to the chosen basis of \(V\), the tangential \(3\times3\) matrix is therefore
\[
B+(z+t)A.
\]
The Reeb derivative itself maps to \(\theta\), which is transverse to \(\theta^\perp\).  Hence the
full four-dimensional differential has rank four exactly when this tangential matrix is invertible.
Since \(P_V\) has degree at most three, a nonzero \(P_V\) has at most three real roots.
\end{proof}

\begin{theorem}[Complex-evaluation identity]
\label{thm:v003-complex-identity}
For every three-plane \(V\),
\begin{equation}
\label{eq:v003-complex-identity}
\boxed{
|P_V(i)|^2
=
1-\sigma_\omega(V)^2.
}
\end{equation}
\end{theorem}

\begin{proof}
Because the chosen frame of \(V\) is orthonormal,
\[
A^TA+B^TB=I.
\]
Therefore
\[
\begin{aligned}
(B+iA)^*(B+iA)
&=
B^TB+A^TA+i(B^TA-A^TB)\\
&=
I-iC_V.
\end{aligned}
\]
Taking determinants gives
\[
|P_V(i)|^2
=
\det(I-iC_V).
\]
A real \(3\times3\) skew-symmetric matrix has eigenvalues
\[
0,\quad i\sigma,\quad -i\sigma,
\]
where
\[
\sigma=\|C_V\|_{\rm op}.
\]
Hence
\[
\det(I-iC_V)
=
1-\sigma^2.
\]
\end{proof}

\begin{theorem}[Matrix-pencil/Maslov incidence equivalence]
\label{thm:v003-maslov-equivalence}
For every three-plane \(V\),
\begin{equation}
\label{eq:v003-maslov-equivalence}
\boxed{
P_V(s)=0
\iff
V\cap L_s\neq\{0\}.
}
\end{equation}
Consequently,
\begin{equation}
\label{eq:v003-total-trap}
\boxed{
P_V\equiv0
\iff
V\cap L_s\neq\{0\}
\quad\text{for every }s\in\R.
}
\end{equation}
\end{theorem}

\begin{proof}
A nonzero vector of \(V\) has coordinates
\[
(Ac,Bc)
\]
for some \(c\neq0\).  It belongs to \(L_s\) exactly when
\[
Bc=-sAc,
\]
equivalently
\[
(B+sA)c=0.
\]
Thus \(V\cap L_s\neq0\) exactly when \(\det(B+sA)=0\).  The second statement follows because a
polynomial of degree at most three that vanishes for every real \(s\) is identically zero.
\end{proof}

\begin{theorem}[Rectifiable Reeb-front polynomial theorem]
\label{thm:v003-rectifiable-polynomial}
Let
\[
R\subset J^1(S^3)
\]
be countably \(3\)-rectifiable with
\[
0<\mathcal H^3(R)<\infty.
\]
Assume the line projection
\[
q:R\to T^*S^3
\]
has approximate rank three almost everywhere on a measurable subset \(R_0\subset R\) of positive
\(\mathcal H^3\)-measure.  Let \(V_q\) be the projected approximate tangent three-plane.

If
\[
P_{V_q}\not\equiv0
\]
on a positive-measure subset of \(R_0\), then
\begin{equation}
\label{eq:v003-rectifiable-positive-front}
\boxed{
\mathcal L^4
\left(
\mathcal E(\Sat_{1/2}^{R}(R))
\right)
>0.
}
\end{equation}
\end{theorem}

\begin{proof}
Cover \(R\) up to a null set by countably many Lipschitz parametrizations.  On points where the
approximate line projection has rank three, use projected tangent coordinates.  By
\cref{thm:v003-rank-criterion}, for each front-regular tangent the Reeb-sweep differential has rank
four for almost every Reeb time.  Fubini therefore gives a positive-measure subset of parameter
space on which the four-Jacobian is positive.  The area formula then implies that the Euclidean
image has positive four-dimensional Lebesgue measure.
\end{proof}

\begin{lemma}[No constant right kernel]
\label{lem:v004-no-constant-right-kernel}
One has
\begin{equation}
\label{eq:v004-no-common-right}
\boxed{
\ker A\cap\ker B=\{0\}.
}
\end{equation}
Equivalently, the pencil \(B+sA\) has no nonzero constant polynomial right-kernel vector.
\end{lemma}

\begin{proof}
If \(c\in\ker A\cap\ker B\), then
\[
\mathbb Vc=0.
\]
Since the columns of \(\mathbb V\) are a basis of \(V\), the stacked matrix has full column rank and
therefore \(c=0\).
\end{proof}

\begin{theorem}[Three-type Kronecker classification]
\label{thm:v004-three-types}
Let \(V\subset\R_a^3\oplus\R_b^3\) be a front-null three-plane.  Let
\[
\varepsilon
\]
be the unique right minimal index,
\[
\eta
\]
the unique left minimal index, and
\[
r_0
\]
the size of the regular Kronecker part.  Then
\begin{equation}
\label{eq:v004-index-sum}
\boxed{
\varepsilon+\eta+r_0=2,
\qquad
\varepsilon\ge1,
}
\end{equation}
and consequently exactly one of the following occurs:
\begin{equation}
\label{eq:v004-three-types}
\boxed{
(\varepsilon,\eta,r_0)
\in
\{(1,0,1),(1,1,0),(2,0,0)\}.
}
\end{equation}
\end{theorem}

\begin{proof}
Because the normal rank is two, the right and left nullities over \(\R(s)\) are both one.  Hence
there is exactly one right singular block and exactly one left singular block.  The total row and
column counts are
\[
\varepsilon+(\eta+1)+r_0=3,
\]
and
\[
(\varepsilon+1)+\eta+r_0=3,
\]
which both reduce to
\[
\varepsilon+\eta+r_0=2.
\]
A right minimal index \(\varepsilon=0\) would give a nonzero constant right-kernel vector, excluded
by \cref{lem:v004-no-constant-right-kernel}.  Thus \(\varepsilon\ge1\), and the three displayed
nonnegative integer triples are the only possibilities.
\end{proof}

\begin{corollary}[Exact front-null condition in K\"ahler coordinates]
\label{cor:v005-null-condition}
A maximally symplectic three-plane is front-null exactly when
\begin{equation}
\label{eq:v005-null-condition}
\boxed{
(a\times b)\cdot c
=
(a\times b)\cdot d
=
0.
}
\end{equation}
\end{corollary}

\begin{proposition}[Forced common output plane]
\label{prop:v005-forced-output-plane}
If \(V\) is front-null and \(\rho(V)>0\), then
\begin{equation}
\label{eq:v005-cd-in-W}
\boxed{
c,d\in W_V.
}
\end{equation}
Consequently
\begin{equation}
\label{eq:v005-image-in-W}
\boxed{
\operatorname{im}(B+sA)\subseteq W_V
\qquad
\forall s\in\R.
}
\end{equation}
\end{proposition}

\begin{proof}
Since \(a\times b\neq0\), the vector \(a\times b\) spans \(W_V^\perp\).  The front-null equations
\eqref{eq:v005-null-condition} say that both \(c\) and \(d\) are orthogonal to \(a\times b\),
hence belong to \(W_V\).  Every column of \(B+sA\) is then in \(W_V\).
\end{proof}

\begin{theorem}[Transverse branch is exactly \(\mathsf K_{200}\)]
\label{thm:v005-transverse-K200}
If \(V\) is front-null and \(\rho(V)>0\), then
\begin{equation}
\label{eq:v005-transverse-K200}
\boxed{
V\in\mathsf K_{200}.
}
\end{equation}
\end{theorem}

\begin{proof}
By \cref{prop:v005-forced-output-plane}, the left minimal index is zero.  It remains to exclude
right minimal index one.

Rotate the K\"ahler basis inside \(\mathcal C_V\) so that \(a\perp b\), and apply a diagonal
orthogonal transformation so that
\[
a=\alpha e_1,
\qquad
b=\beta e_2,
\qquad
\alpha,\beta>0,
\qquad
\alpha^2+\beta^2=1.
\]
Because \(c,d\in W_V\) and \(w\) is Hermitian-orthogonal to \(u\), there is a phase \(\phi\) with
\[
c=
-\beta\sin\phi\,e_1
+
\alpha\cos\phi\,e_2,
\]
\[
d=
\beta\cos\phi\,e_1
+
\alpha\sin\phi\,e_2.
\]
Delete the zero third output row and take the cross product of the two remaining row vectors of the
resulting \(2\times3\) pencil.  A polynomial generator \(q(s)\) of the right kernel has third
component
\[
-\alpha\beta(1+s^2).
\]
Since \(\alpha\beta\neq0\), no real linear polynomial can divide all components of \(q(s)\).  A
quadratic common divisor would force a constant right-kernel generator, excluded by
\cref{lem:v004-no-constant-right-kernel}.  Thus the primitive right-kernel generator has degree two:
\[
\varepsilon=2.
\]
With \(\eta=0\), the classification \cref{thm:v004-three-types} gives
\(\mathsf K_{200}\).
\end{proof}

\begin{equation}
\label{eq:v005-K101-gauge}
\boxed{
V
=
\operatorname{span}_{\R}
\{
(e_1,0),\,
(0,e_1),\,
(\mu e_2,\nu e_2)
\}.
}
\end{equation}

\begin{theorem}[Front-null direction-rank law]
\label{thm:v005-direction-rank}
Every front-null three-plane satisfies
\begin{equation}
\label{eq:v005-direction-rank}
\boxed{
\rank(d\pi_{\rm dir}|_V)
=
\rank A
\le2.
}
\end{equation}
More precisely:
\begin{enumerate}[label=\textnormal{(\roman*)}]
\item in \(\mathsf K_{110}\), \(\rank A=\rank B=2\);
\item in \(\mathsf K_{200}\), \(\rank A=\rank B=2\);
\item in \(\mathsf K_{101}\), each of \(\rank A,\rank B\) is one or two, with at least one equal to
two.
\end{enumerate}
\end{theorem}

\begin{proof}
The inequality follows already from the vanishing of the leading coefficient of
\(P_V(s)\).  For the sharper statements use the adapted gauges.

In \(\mathsf K_{110}\), \(A\) contains the independent horizontal vectors \(e_1,c\) and \(B\)
contains \(e_1,d\), while \(c,d\neq0\).  In \(\mathsf K_{200}\), the first two horizontal columns
\(a,-b\) are independent and the first two vertical columns \(b,a\) are independent.  In
\(\mathsf K_{101}\), the gauge \eqref{eq:v005-K101-gauge} gives
\[
\rank A=
\begin{cases}
1,&\mu=0,\\
2,&\mu\neq0,
\end{cases}
\qquad
\rank B=
\begin{cases}
1,&\nu=0,\\
2,&\nu\neq0.
\end{cases}
\]
Since \(\mu^2+\nu^2=1\), they cannot both equal one for the same vanishing reason.
\end{proof}

\begin{theorem}[No countably rectifiable Sticky selector with null front]
\label{thm:v005-no-rectifiable-selector}
Assume \(\Lambda_s\) is countably \(3\)-rectifiable.  Let
\[
\Gamma_s
=
\{(\theta,p(\theta),z(\theta)):\theta\in S^3\}
\]
be any Borel marked lift whose relevant rectifiable pieces admit the standard approximate tangent
lifting.  Then the Reeb-saturated Euclidean front cannot be four-dimensional null on all
direction-carrying rectifiable mass.  In particular, whenever the marked graph itself is countably
\(3\)-rectifiable,
\begin{equation}
\label{eq:v005-positive-front-selector}
\boxed{
\mathcal L^4(K_{\Gamma_s})>0.
}
\end{equation}
\end{theorem}

\begin{proof}
Apply the area formula to the Lipschitz direction projection
\[
\pi_{\rm dir}:\Lambda_s\to S^3.
\]
Since the selector graph meets every direction fibre exactly once,
the multiplicity of the projection is one.  Therefore
\[
\int_{\Lambda_s}J_3(\pi_{\rm dir}|_{T_x\Lambda_s})\,d\mathcal H^3(x)
=
\mathcal H^3(S^3)>0.
\]
Hence the approximate direction differential has rank three on a subset of positive
\(\mathcal H^3\)-measure.

At every such point the horizontal tangent matrix \(A\) is invertible, so the leading coefficient
of
\[
P_V(s)=\det(B+sA)
\]
is nonzero.  Thus \(P_V\not\equiv0\).  For a countably rectifiable marked lift, the
rectifiable Reeb-front polynomial theorem
\cref{thm:v003-rectifiable-polynomial} then gives positive four-dimensional front measure.
\end{proof}

\subsection{Polarized normal forms}

\begin{definition}[Collision parameter and secant residual]
\label{def:v006-collision-residual}
For \(\alpha\ne0\), define
\begin{equation}
\label{eq:v006-sstar}
\boxed{
s_*(\alpha,\beta)
=
-\frac{\alpha\cdot\beta}{|\alpha|^2},
}
\end{equation}
and
\begin{equation}
\label{eq:v006-residual}
\boxed{
r(\alpha,\beta)
=
\beta+s_*\alpha.
}
\end{equation}
Then
\[
r\perp\alpha.
\]
\end{definition}

\begin{lemma}[Exact Pythagorean collision identity]
\label{lem:v006-pythagorean}
For every \(s\in\R\),
\begin{equation}
\label{eq:v006-pythagorean}
\boxed{
|\beta+s\alpha|^2
=
|r(\alpha,\beta)|^2
+
|\alpha|^2|s-s_*|^2.
}
\end{equation}
\end{lemma}

\begin{proof}
Write
\[
\beta=r-s_*\alpha.
\]
Then
\[
\beta+s\alpha=r+(s-s_*)\alpha,
\]
and the two terms are orthogonal.
\end{proof}

\begin{proposition}[Exact collision-window formula]
\label{prop:v006-window-formula}
For \(\alpha\ne0\), let \(s_*\) and \(r\) be
\cref{def:v006-collision-residual}.  If
\[
|r|>\delta,
\]
then
\[
m_\delta^I(\alpha,\beta)=0.
\]
If
\[
|r|\le\delta,
\]
then
\begin{equation}
\label{eq:v006-window-formula}
\boxed{
m_\delta^I(\alpha,\beta)
=
\left|
I\cap
\left[
s_*-\frac{\sqrt{\delta^2-|r|^2}}{|\alpha|},
\,
s_*+\frac{\sqrt{\delta^2-|r|^2}}{|\alpha|}
\right]
\right|.
}
\end{equation}
\end{proposition}

\begin{proof}
This is immediate from the Pythagorean identity
\eqref{eq:v006-pythagorean}.
\end{proof}

\begin{definition}[Normalized secant collision energy]
\label{def:v006-energy}
Define
\begin{equation}
\label{eq:v006-energy}
\boxed{
\mathfrak E_\delta^I(\mu)
=
\delta^{-3}
\iint
m_\delta^I(a-a',b-b')\,
d\mu(a,b)\,d\mu(a',b').
}
\end{equation}
\end{definition}

\begin{lemma}[Universal pairwise collision-window bound]
\label{lem:v007-window-upper}
For every pair with \(\alpha\ne0\),
\begin{equation}
\label{eq:v007-window-upper}
\boxed{
m_\delta^J(\alpha,\beta)
\le
\frac{2\delta}{|\alpha|}.
}
\end{equation}
\end{lemma}

\begin{proof}
By the exact Pythagorean formula
\[
|\beta+s\alpha|^2
=
|r|^2+|\alpha|^2|s-s_*|^2,
\]
the collision set is either empty or contained in an interval of radius
\[
\frac{\delta}{|\alpha|}
\]
around \(s_*\).
\end{proof}

\subsection{Rectifiable and selector reductions}

\begin{proposition}[Bush mass gives neighbourhood volume at every size]
\label{prop:v010-bush-volume}
Assume
\[
s_0\in J_\star
\Subset I
\]
and let
\[
H=H_{(c_0,s_0)}(\rho).
\]
Then
\begin{equation}
\label{eq:v010-bush-volume}
\boxed{
\mathcal L^4
\left(
N_\rho(K_b)
\right)
\ge
|U|\,\sigma(H)
\int_I|s-s_0|^3\,ds.
}
\end{equation}
In particular there is \(c_{I,J_\star,U}>0\) such that
\begin{equation}
\label{eq:v010-bush-volume-uniform}
\boxed{
\mathcal L^4(N_\rho(K_b))
\ge
c_{I,J_\star,U}\,
\sigma(H).
}
\end{equation}
\end{proposition}

\begin{proof}
For \(a\in H\), write
\[
b(a)+s_0a=c_0+e(a),
\qquad
|e(a)|\le\rho.
\]
At an arbitrary level \(s\in I\),
\[
b(a)+sa
=
c_0+(s-s_0)a+e(a).
\]
Hence the horizontal slice of \(N_\rho(K_b)\) at height \(s\) contains
\[
c_0+(s-s_0)H.
\]
Its three-dimensional Lebesgue measure is
\[
|s-s_0|^3\,|H|
=
|s-s_0|^3\,|U|\,\sigma(H).
\]
Integrate in \(s\).  The uniform bound follows because
\[
s_0\mapsto\int_I|s-s_0|^3\,ds
\]
has a positive minimum on the compact set \(J_\star\).
\end{proof}

\clearpage
\section{Physical collision, recurrence, and entropy}
\label{sec:v034-entropy-pruning}

\subsection{One-scale collision geometry}

\begin{proposition}[Ideal escape stays near the actual Kakeya slab]
\label{prop:v014-near-K}
For every \(i\),
\begin{equation}
\label{eq:v014-near-K}
\boxed{
E_i
\subset
N_\rho(K_b).
}
\end{equation}
\end{proposition}

\begin{proof}
For \(a\in H_i\), write
\[
b(a)+s_i a=c_i+e_i(a),
\qquad
|e_i(a)|\le\rho.
\]
Then
\[
b(a)+sa
=
c_i+(s-s_i)a+e_i(a).
\]
Thus the actual selected-line point at the same longitudinal coordinate \(s\) lies within
\(\rho\) of the corresponding ideal escaped point.
\end{proof}

\begin{theorem}[Exact one-bush escape volume]
\label{thm:v014-one-volume}
For every \(i\),
\begin{equation}
\label{eq:v014-one-volume}
\boxed{
|E_i|
=
|U|\,\sigma(H_i)
A_I(s_i),
\qquad
A_I(s_i)
:=
\int_I|s-s_i|^3\,ds.
}
\end{equation}
\end{theorem}

\begin{proof}
At fixed \(s\neq s_i\), the horizontal slice is
\[
c_i+(s-s_i)H_i,
\]
whose three-dimensional Lebesgue measure is
\[
|s-s_i|^3|H_i|
=
|U|\,|s-s_i|^3\sigma(H_i).
\]
The exceptional level \(s=s_i\) has measure zero.  Integrate in \(s\).
\end{proof}

\begin{equation}
\label{eq:v014-Psi}
\boxed{
\Psi_{ij,s}(a)
=
\frac{
c_i-c_j+(s-s_i)a
}{
s-s_j
}.
}
\end{equation}

\begin{theorem}[Exact pairwise escaped-overlap formula]
\label{thm:v014-pair-overlap}
For every \(i\neq j\),
\begin{equation}
\label{eq:v014-pair-overlap}
\boxed{
|E_i\cap E_j|
=
|U|
\int_I
|s-s_i|^3
\mathfrak R_{ij}(s)\,ds.
}
\end{equation}
\end{theorem}

\begin{proof}
Fix \(s\) away from \(s_i,s_j\).  The \(s\)-slice of \(E_i\) is
\[
c_i+(s-s_i)H_i.
\]
A point in that slice also belongs to \(E_j\) exactly when its unique \(i\)-slope \(a\in H_i\)
satisfies
\[
\Psi_{ij,s}(a)\in H_j.
\]
Thus
\[
E_i(s)\cap E_j(s)
=
c_i+(s-s_i)
\left(
H_i\cap\Psi_{ij,s}^{-1}(H_j)
\right).
\]
Its three-dimensional measure is
\[
|s-s_i|^3
|U|
\mathfrak R_{ij}(s).
\]
Integrate over \(s\).  The two exceptional longitudinal levels are null.
\end{proof}

\begin{equation}
\label{eq:v014-secondary-collision}
\boxed{
|
[b(a)+sa]
-
[b(a')+sa']
|
\le
2\rho.
}
\end{equation}

\begin{theorem}[Two-time self-homothety]
\label{thm:v016-two-time-self}
Let
\[
D_{ij;s,t}
=
\Psi_{ij,s}^{-1}\circ\Psi_{ij,t}.
\]
Then
\begin{equation}
\label{eq:v016-two-time-self-map}
\boxed{
D_{ij;s,t}
=
\operatorname{Dil}_{a_{ij}}
\bigl(q_{ij}(s,t)\bigr).
}
\end{equation}
Moreover
\begin{equation}
\label{eq:v016-two-time-defect}
\boxed{
\bar\sigma
\left(
H_i
\triangle
D_{ij;s,t}(H_i)
\right)
\le
\Delta_{ij}(s)
+
|q_{ij}(s,t)|^3
\Delta_{ij}(t).
}
\end{equation}
\end{theorem}

\begin{proof}
The map identity is the two-time ratio formula from the preceding stage.

Put
\[
A_s
=
\Psi_{ij,s}^{-1}(H_j),
\qquad
A_t
=
\Psi_{ij,t}^{-1}(H_j).
\]
By definition of \(D_{ij;s,t}\),
\[
D_{ij;s,t}(A_t)=A_s.
\]
Therefore
\[
\begin{aligned}
\bar\sigma
\left(
H_i\triangle D_{ij;s,t}(H_i)
\right)
&\le
\bar\sigma(H_i\triangle A_s)
+
\bar\sigma
\left(
D_{ij;s,t}(A_t)
\triangle
D_{ij;s,t}(H_i)
\right)\\
&=
\Delta_{ij}(s)
+
|q_{ij}(s,t)|^3
\Delta_{ij}(t).
\end{aligned}
\]
\end{proof}

\begin{lemma}[Approximate translation invariance is quantitatively small]
\label{lem:v016-translation-rigidity}
Let
\[
H\subset U
\]
have
\[
m=\bar\sigma(H)>0,
\]
and let
\[
v\in\R^3.
\]
Then
\begin{equation}
\label{eq:v016-translation-rigidity}
\boxed{
\frac{
\bar\sigma
\left(
H\triangle(H+v)
\right)
}{m}
\ge
\min
\left\{
1,\,
\frac{|v|}{D_U}
\right\},
}
\end{equation}
where
\[
D_U=\operatorname{diam}(U).
\]
Consequently, if the left side is at most \(\varepsilon<1\), then
\begin{equation}
\label{eq:v016-translation-small}
\boxed{
|v|
\le
\varepsilon D_U.
}
\end{equation}
\end{lemma}

\begin{proof}
If \(|v|>D_U\), then \(H\) and \(H+v\) are disjoint, so the normalized symmetric difference equals
\(2m\), which is stronger than the claim.

Assume \(0<|v|\le D_U\).  Choose the least positive integer \(n\) such that
\[
n|v|>D_U.
\]
Then
\[
H\cap(H+nv)=\varnothing,
\]
so
\[
\bar\sigma
\left(
H\triangle(H+nv)
\right)
=
2m.
\]
By telescoping and translation invariance of Lebesgue measure,
\[
\bar\sigma
\left(
H\triangle(H+nv)
\right)
\le
n\,
\bar\sigma
\left(
H\triangle(H+v)
\right).
\]
Also
\[
n
\le
\frac{D_U}{|v|}+1
\le
\frac{2D_U}{|v|}.
\]
Hence
\[
\bar\sigma
\left(
H\triangle(H+v)
\right)
\ge
\frac{2m}{n}
\ge
m\frac{|v|}{D_U}.
\]
\end{proof}

\begin{theorem}[Homothety commutator formula]
\label{thm:v016-commutator}
Let
\[
F=D_{a,q},
\qquad
G=D_{b,r}.
\]
Then
\begin{equation}
\label{eq:v016-commutator}
\boxed{
(GF)^{-1}(FG)
=
\tau_v,
\qquad
v
=
\frac{(1-q)(1-r)}{qr}(a-b),
}
\end{equation}
where
\[
\tau_v(x)=x+v.
\]
\end{theorem}

\begin{proof}
The two compositions have the same scalar linear part \(qr\):
\[
FG(x)
=
qrx
+
q(1-r)b
+
(1-q)a,
\]
\[
GF(x)
=
qrx
+
r(1-q)a
+
(1-r)b.
\]
Their translation-part difference is
\[
(1-q)(1-r)(a-b).
\]
Applying \((GF)^{-1}\), whose linear part is \((qr)^{-1}\), gives
\eqref{eq:v016-commutator}.
\end{proof}

\begin{theorem}[Two-axis approximate homothety rigidity]
\label{thm:v016-two-axis}
Let
\[
H\subset U,
\qquad
m=\bar\sigma(H)>0,
\]
and suppose
\[
F=D_{a,q},
\qquad
G=D_{b,r},
\]
with
\[
q,r\ne0.
\]
Assume
\begin{equation}
\label{eq:v016-two-axis-defects}
\boxed{
\bar\sigma(H\triangle F(H))
\le
\varepsilon_Fm,
\qquad
\bar\sigma(H\triangle G(H))
\le
\varepsilon_Gm.
}
\end{equation}
Then
\begin{equation}
\label{eq:v016-commutator-defect}
\boxed{
\bar\sigma
\left(
H\triangle(H+v)
\right)
\le
C(q,r)
(\varepsilon_F+\varepsilon_G)m,
}
\end{equation}
where \(v\) is given by
\eqref{eq:v016-commutator}, and one may take
\[
C(q,r)
=
|qr|^{-3}
\left(
2+|q|^3+|r|^3
\right).
\]
Consequently
\begin{equation}
\label{eq:v016-center-rigidity-general}
\boxed{
\min
\left\{
1,\,
\frac{
|(1-q)(1-r)|\,|a-b|
}{
|qr|D_U
}
\right\}
\le
C(q,r)
(\varepsilon_F+\varepsilon_G).
}
\end{equation}
\end{theorem}

\begin{proof}
First compare the two compositions:
\[
\begin{aligned}
\bar\sigma(FG(H)\triangle GF(H))
&\le
\bar\sigma(FG(H)\triangle F(H))
+
\bar\sigma(F(H)\triangle H)\\
&\quad+
\bar\sigma(H\triangle G(H))
+
\bar\sigma(G(H)\triangle GF(H)).
\end{aligned}
\]
The first term equals
\[
|q|^3\bar\sigma(G(H)\triangle H)
\le
|q|^3\varepsilon_Gm,
\]
and the last equals
\[
|r|^3\bar\sigma(H\triangle F(H))
\le
|r|^3\varepsilon_Fm.
\]
Hence
\[
\bar\sigma(FG(H)\triangle GF(H))
\le
\left(
1+|r|^3
\right)\varepsilon_Fm
+
\left(
1+|q|^3
\right)\varepsilon_Gm.
\]
Apply \((GF)^{-1}\), whose absolute Jacobian is \(|qr|^{-3}\).  By
\cref{thm:v016-commutator}, the two resulting sets are \(H+v\) and \(H\).  This gives
\eqref{eq:v016-commutator-defect}.  Finally use
\cref{lem:v016-translation-rigidity}.
\end{proof}

\begin{proposition}[Sign-reversal strong recurrence remains rigid]
\label{prop:v017-sign-reversal-rigid}
Suppose a fixed direction block \(H_i\) admits two raw small-defect recurrences whose ratio satisfies
\[
q=-1+O(\varepsilon),
\]
and suppose two partner edges produce such sign-reversal self-homotheties with centers \(a,b\).
Then, provided the dilation ratios stay in a fixed compact neighbourhood of \(-1\), the
translation-commutator theorem of the preceding stage still gives
\begin{equation}
\label{eq:v017-sign-reversal-rigid}
\boxed{
|a-b|
\lesssim
D_U\varepsilon.
}
\end{equation}
\end{proposition}

\begin{proof}
Near \(q=r=-1\), the coefficient
\[
\frac{(1-q)(1-r)}{qr}
\]
in the commutator translation stays bounded away from zero, while all Jacobian constants remain
bounded.  Apply
\cref{thm:v016-two-axis}.
\end{proof}

\subsection{Weighted graph regularization and radial recurrence}

\begin{equation}
\label{eq:v029-fixed-angle-overlap}
\boxed{
\mathscr O_{ij}^{+,\tau_0}
=
|U|
\int_{I_+}
|s-s_i|^3
\sigma\left(
\left\{
a\in H_i:
\Psi_{ij,s}(a)\in H_j,
\quad
|a-\Psi_{ij,s}(a)|\ge\tau_0
\right\}
\right)
\,ds.
}
\end{equation}

\begin{equation}
\label{eq:v029-truncated-multiplicity}
\boxed{
\mathcal N_{i_0}^{\tau_0}(a,s)
=
\sum_{j\in J}
\mathbf 1_{H_j}(\Psi_{i_0j,s}(a))
\mathbf 1_{[\tau_0,\infty)}
\bigl(|a-\Psi_{i_0j,s}(a)|\bigr).
}
\end{equation}

\begin{theorem}[Exact truncated star identity and common parameter set]
\label{thm:v029-truncated-star}
If
\[
\mathscr O_{i_0j}^{+,\tau_0}
\ge
\theta R
\qquad
(j\in J),
\]
then
\begin{equation}
\label{eq:v029-truncated-star-identity}
\boxed{
\int_{I_+}\int_{H_{i_0}}
|s-s_{i_0}|^3
\mathcal N_{i_0}^{\tau_0}(a,s)
\,da\,ds
=
\sum_{j\in J}
\mathscr O_{i_0j}^{+,\tau_0}
\ge
D\theta R.
}
\end{equation}
Consequently there is a measurable set
\[
Z_{i_0}
\subset
H_{i_0}\times I_+
\]
such that
\begin{equation}
\label{eq:v029-Z-mass}
\boxed{
(\sigma\times ds)(Z_{i_0})
\gtrsim
\theta R,
}
\end{equation}
and
\begin{equation}
\label{eq:v029-Z-multiplicity}
\boxed{
\mathcal N_{i_0}^{\tau_0}(a,s)
\gtrsim
\theta D
\qquad
((a,s)\in Z_{i_0}).
}
\end{equation}
\end{theorem}

\begin{proof}
Expand
\eqref{eq:v029-truncated-multiplicity} and apply Tonelli.  Since
\(da=|U|\,d\sigma(a)\), each summand is exactly
\eqref{eq:v029-fixed-angle-overlap}.  This proves
\eqref{eq:v029-truncated-star-identity}.  On the outer wing the weight is bounded above and below,
\(0\le\mathcal N_{i_0}^{\tau_0}\le D\), and \(\sigma(H_{i_0})\asymp R\).  Thresholding at a
sufficiently small multiple of \(\theta D\) gives
\eqref{eq:v029-Z-mass}--\eqref{eq:v029-Z-multiplicity}.
\end{proof}

\begin{lemma}[Pointwise heavy longitudinal packet]
\label{lem:v029-heavy-packet}
For every \((a,s)\in Z_{i_0}\), there exists a level interval \(K(a,s)\in\mathcal K_\Delta\) such
that
\begin{equation}
\label{eq:v029-heavy-packet}
\boxed{
\#\left\{
j\in J^{\tau_0}(a,s):
s_j\in K(a,s)
\right\}
\gtrsim
\theta D\Delta.
}
\end{equation}
\end{lemma}

\begin{proof}
The sets in the braces partition \(J^{\tau_0}(a,s)\) into
\(O(\Delta^{-1})\) level bins, while
\eqref{eq:v029-Z-multiplicity} gives
\(\#J^{\tau_0}(a,s)\gtrsim\theta D\).  Pigeonhole.
\end{proof}

\begin{equation}
\label{eq:v029-star-center-line}
\boxed{
c_j
=
q_{i_0}+s_ja^\sharp+e_j,
\qquad
|e_j|
\lesssim
\eta.
}
\end{equation}

\begin{lemma}[Packet centers and packet teeth are close]
\label{lem:v029-packet-cap}
If \(j,k\in J^{\tau_0}(a,s)\) and
\(s_j,s_k\in K(a,s)\), then
\begin{equation}
\label{eq:v029-packet-center-distance}
\boxed{
|x_j-x_k|
\lesssim
d_0^{-1}\Delta+\eta,
}
\end{equation}
and
\begin{equation}
\label{eq:v029-packet-tooth-distance}
\boxed{
|a_j(y)-a_k(y)|
\lesssim
d_0^{-1}\Delta+\eta.
}
\end{equation}
\end{lemma}

\begin{proof}
Subtract
\eqref{eq:v029-star-center-line} for \(j,k\):
\[
c_j-c_k
=
(s_j-s_k)a^\sharp+e_j-e_k.
\]
Since \(|s_j-s_k|\le\Delta\), this proves
\eqref{eq:v029-packet-center-distance}.  At the common escaped point,
\begin{equation}
\label{eq:v029-tooth-difference-identity}
\boxed{
a_j(y)-a_k(y)
=
\frac{
(s_j-s_k)a_j(y)-h_{jk}
}{
s-s_k
}.
}
\end{equation}
The outer-wing denominator is uniformly bounded below, \(a_j(y)\in U\), and
\eqref{eq:v029-packet-center-distance} controls both \(|s_j-s_k|\) and \(|h_{jk}|\).  This gives
\eqref{eq:v029-packet-tooth-distance}.
\end{proof}

\begin{theorem}[Clustered fixed-angle star: degree or angular cap]
\label{thm:v029-star-cap}
A clustered fixed-angle star of degree \(D\) obeys one of the following alternatives:
\begin{enumerate}[label=\textnormal{(\roman*)}]
\item
\begin{equation}
\label{eq:v029-star-degree-bound}
\boxed{
D
\lesssim
\theta^{-1}\Delta^{-1};
}
\end{equation}
\item there is a measurable set \(Y_{i_0}\subset E_{i_0}^+\) such that
\begin{equation}
\label{eq:v029-star-cap-volume}
\boxed{
|Y_{i_0}|
\gtrsim
\theta R,
}
\end{equation}
and every \(y\in Y_{i_0}\) supports at least
\begin{equation}
\label{eq:v029-star-cap-multiplicity}
\boxed{
\gtrsim
\theta D\Delta
}
\end{equation}
simultaneous neighbor teeth in one cap of radius
\begin{equation}
\label{eq:v029-star-cap-radius}
\boxed{
C(d_0^{-1}\Delta+\eta).
}
\end{equation}
\end{enumerate}
\end{theorem}

\begin{proof}
If the right side of
\eqref{eq:v029-star-cap-multiplicity} is smaller than a sufficiently large fixed constant, then
\eqref{eq:v029-star-degree-bound} holds.  Otherwise map \(Z_{i_0}\) into the central escaped bush by
\[
\Phi_{i_0}(a,s)
=
\bigl(c_{i_0}+(s-s_{i_0})a,s\bigr).
\]
This map is injective on the outer wing and has Jacobian \(|s-s_{i_0}|^3\asymp1\), so its image
\(Y_{i_0}\) satisfies
\eqref{eq:v029-star-cap-volume}.  At each image point, choose the variable packet from
\cref{lem:v029-heavy-packet}; its cardinality is
\eqref{eq:v029-star-cap-multiplicity}, and
\cref{lem:v029-packet-cap} puts all its teeth in the cap
\eqref{eq:v029-star-cap-radius}.
\end{proof}

\begin{theorem}[Far-chart packing-or-cap decomposition]
\label{thm:v029-far-decomposition}
There is an edge-disjoint decomposition
\begin{equation}
\label{eq:v029-far-edge-decomposition}
\boxed{
\mathcal E_{\theta,\tau_0}^{\mathrm{far}}
=
\mathcal E_{\theta,\tau_0}^{\mathrm{far,pack}}
\mathbin{\dot\cup}
\mathop{\dot\bigcup}_{\nu}
\mathcal S_{\nu}^{\mathrm{far,cap}}
}
\end{equation}
such that
\begin{equation}
\label{eq:v029-far-pack-count}
\boxed{
\#\mathcal E_{\theta,\tau_0}^{\mathrm{far,pack}}
\lesssim
d_R^{-3}\theta^{-1}R^{-3}
=
\theta^{-1}R^{-3+o(1)},
}
\end{equation}
and every removed star satisfies
\eqref{eq:v029-star-cap-volume}--\eqref{eq:v029-star-cap-radius} with
\begin{equation}
\label{eq:v029-far-cap-radius}
\boxed{
C(d_R+\eta_R)
\longrightarrow
0.
}
\end{equation}
\end{theorem}

\begin{proof}
For far edges, \(|a_{ij}|\lesssim d_R^{-1}\).  Partition this radial-center region into
\(\eta_R\)-cubes.  Their number is
\[
O(d_R^{-3}\eta_R^{-3})
=
O(d_R^{-1}R^{-1}).
\]
Inside each cell graph, repeatedly remove the full star of a maximum-degree vertex whenever its
degree exceeds
\[
C\theta^{-1}\Delta_R^{-1}
=
C\theta^{-1}d_R^{-2}.
\]
By
\cref{thm:v029-star-cap}, every removed star has the asserted cap output.  The residual graph in one
cell has at most
\[
O(N\theta^{-1}d_R^{-2})
=
O(\theta^{-1}d_R^{-2}R^{-2})
\]
edges.  Multiplication by the cell count gives
\[
O(d_R^{-1}R^{-1})
O(\theta^{-1}d_R^{-2}R^{-2})
=
O(d_R^{-3}\theta^{-1}R^{-3}),
\]
proving
\eqref{eq:v029-far-pack-count}.  The cap radius is
\eqref{eq:v029-star-cap-radius} with
\(d_0=d_R\), \(\Delta_R=d_R^2\), and
\(\eta_R=(R/d_R^2)^{1/3}\).
\end{proof}

\begin{lemma}[An oriented near-diagonal displacement star has close centers]
\label{lem:v029-near-center-cluster}
Suppose all edges of a star use one orientation and their ordered displacements lie in one
\(\zeta_R\)-cube.  Then its neighbor centers satisfy
\begin{equation}
\label{eq:v029-near-center-cluster}
\boxed{
|x_j-x_k|
\lesssim
d_R+\zeta_R.
}
\end{equation}
At any common escaped point, the corresponding neighbor teeth satisfy the same bound.
\end{lemma}

\begin{proof}
For the orientation \(i_0<j\), one has
\(c_j=c_{i_0}-h_{i_0j}\); for the opposite orientation,
\(c_j=c_{i_0}+h_{ji_0}\).  Thus the horizontal neighbor diameter is \(O(\zeta_R)\).  Their levels
are each within \(d_R\) of \(s_{i_0}\), so their longitudinal diameter is at most \(2d_R\).  This
proves the center bound.  Apply the exact identity
\eqref{eq:v029-tooth-difference-identity} and the outer-wing denominator bound to obtain the tooth
bound.
\end{proof}

\begin{theorem}[Near-chart packing-or-cap decomposition]
\label{thm:v029-near-decomposition}
There is an edge-disjoint decomposition
\begin{equation}
\label{eq:v029-near-edge-decomposition}
\boxed{
\mathcal E_{\theta,\tau_0}^{\mathrm{near}}
=
\mathcal E_{\theta,\tau_0}^{\mathrm{near,pack}}
\mathbin{\dot\cup}
\mathop{\dot\bigcup}_{\nu}
\mathcal S_{\nu}^{\mathrm{near,cap}}
}
\end{equation}
such that
\begin{equation}
\label{eq:v029-near-pack-count}
\boxed{
\#\mathcal E_{\theta,\tau_0}^{\mathrm{near,pack}}
\lesssim
\theta^{-1}R^{-3},
}
\end{equation}
and every removed star has a common physical region of volume \(\gtrsim\theta R\) on which at least
\(\gtrsim\theta D\) simultaneous neighbor teeth lie in a cap of radius
\begin{equation}
\label{eq:v029-near-cap-radius}
\boxed{
C(d_R+R^{1/3})
\longrightarrow
0.
}
\end{equation}
\end{theorem}

\begin{proof}
Work in one displacement-cell graph.  While a vertex has degree larger than
\(C\theta^{-1}\), choose a maximum-degree vertex and split its incident edges according to whether
the vertex is the first or second endpoint in the fixed ordering.  One orientation retains at least
half the degree.  Apply
\cref{thm:v029-truncated-star} to that oriented star.  On its set \(Z_{i_0}\), at least
\(\gtrsim\theta D\) neighbors occur simultaneously.  Map this set into \(E_{i_0}^+\) by the
central escaped-bush parametrization; the image has volume \(\gtrsim\theta R\).  By
\cref{lem:v029-near-center-cluster}, all simultaneous neighbor teeth lie in the cap
\eqref{eq:v029-near-cap-radius}.  Record and remove the oriented star.

When the process stops, the residual graph in each cell has maximum degree
\(O(\theta^{-1})\), hence at most \(O(N\theta^{-1})\) edges.  There are \(O(R^{-1})\) cells and
\(N\asymp R^{-2}\), giving
\eqref{eq:v029-near-pack-count}.
\end{proof}

\subsection{Entropy pruning and common-neighbor structure}

\begin{lemma}[The fixed-angle collision time is nondegenerate]
\label{lem:v042-time-bilipschitz}
For a fixed edge \(ij\) and fixed source direction \(a\),
\begin{equation}
\label{eq:v042-time-derivative}
\boxed{
\partial_\xi\Psi_{ij,\xi}(a)
=
\frac{
a-\Psi_{ij,\xi}(a)
}{
\xi-s_j
}.
}
\end{equation}
On the outer wing and the fixed-angle set,
\[
c_{\tau_0}
\le
\left|
\partial_\xi\Psi_{ij,\xi}(a)
\right|
\le
C.
\]
Moreover the trajectory lies on one affine line and is monotone there.  Consequently the set of
times for which \(\Psi_{ij,\xi}(a)\) belongs to one direction \(R\)-cube has length \(O_{\tau_0}(R)\).
\end{lemma}

\begin{proof}
Differentiate
\[
\Psi_{ij,\xi}(a)
=
\frac{c_i-c_j+(\xi-s_i)a}{\xi-s_j}.
\]
This gives
\[
\partial_\xi\Psi_{ij,\xi}(a)
=
\frac{(s_i-s_j)a-(c_i-c_j)}{(\xi-s_j)^2}
=
\frac{a-\Psi_{ij,\xi}(a)}{\xi-s_j}.
\]
The outer-wing denominator is bounded above and below, while the truncation gives
\(|a-\Psi_{ij,\xi}(a)|\ge\tau_0\).  Also
\[
\Psi_{ij,\xi}(a)-a
=
\frac{
c_i-c_j-(s_i-s_j)a
}{
\xi-s_j
},
\]
so the trajectory has a fixed direction and a monotone scalar parameter.  Its intersection with
an \(R\)-cube therefore has time length \(O_{\tau_0}(R)\).
\end{proof}

\begin{corollary}[One edge contributes only \(R^4\) to one double direction code]
\label{cor:v042-one-edge-double-code}
For every fixed edge \(ij\),
\begin{equation}
\label{eq:v042-one-edge-double-code}
\boxed{
(\sigma\otimes d\xi)
\left\{
(a,\xi):
a\in\tau,\ 
\Psi_{ij,\xi}(a)\in\tau',\
|a-\Psi_{ij,\xi}(a)|\ge\tau_0
\right\}
\lesssim_{\tau_0}
R^4.
}
\end{equation}
The same bound holds after imposing the two intercept cells and both center-level intervals.
\end{corollary}

\begin{proof}
For each fixed \(a\in\tau\), \cref{lem:v042-time-bilipschitz} bounds the admissible
\(\xi\)-length by \(CR\).  Integrate over the direction cube, whose \(\sigma\)-measure is
\(O(R^3)\).  All remaining restrictions only decrease the measure.
\end{proof}

\begin{theorem}[A double graph code has bounded endpoints and \(R^4\) mass]
\label{thm:v042-double-code-capacity}
For one nonempty double code \(\kappa^{(2)}\), all source endpoints \(x_i=(c_i,s_i)\) lie in one
\(O(R)\)-ball and all output endpoints \(x_j=(c_j,s_j)\) lie in another \(O(R)\)-ball.  Hence an
\(R\)-separated simple center graph contributes only \(O(1)\) ordered edges to the code, and
\begin{equation}
\label{eq:v042-double-code-capacity}
\boxed{
\sum_{i,j}
(\sigma\otimes d\xi)
\left(
\mathcal W_{ij}^{\tau_0}
\cap
\{\kappa^{(2)}\}
\right)
\lesssim_{\tau_0}
R^4.
}
\end{equation}
\end{theorem}

\begin{proof}
Choose actual graph points
\((a_0,b(a_0))\in\chi\) and
\((u_0,b(u_0))\in\chi'\) occurring in the code, and representatives
\(s_0\in K\), \(t_0\in K'\).  Since \(a\in H_{x_i}(R)\),
\[
|c_i-(b(a)+s_i a)|
\le
R.
\]
All \(a,b(a),s_i\) in the source code differ from \(a_0,b(a_0),s_0\) by \(O(R)\), so all
\(x_i\) lie in an \(O(R)\)-ball about
\((b(a_0)+s_0a_0,s_0)\).  The same argument with \(u\in H_{x_j}(R)\) places all \(x_j\) in an
\(O(R)\)-ball about
\((b(u_0)+t_0u_0,t_0)\).

Each ball contains \(O(1)\) vertices of an \(R\)-separated family.  Simplicity of the graph gives
only \(O(1)\) ordered edges between the two endpoint sets.  Apply
\eqref{eq:v042-one-edge-double-code} to each edge and sum.
\end{proof}

\begin{lemma}[One physical edge has universal one-dimensional code valence]
\label{lem:v043-one-edge-code-valence}
Fix a physical edge \(ij\) and a source endpoint code \(v=(\chi,K)\).  The fixed-angle
outer-wing incidences of this edge meet at most
\begin{equation}
\label{eq:v043-one-edge-code-valence}
\boxed{
B_R
\lesssim
R^{-1}
}
\end{equation}
output endpoint codes, with no goodness assumption on either direction cube.
\end{lemma}

\begin{proof}
The output level interval \(K'\) is determined by the fixed level \(s_j\).  For a representative
source direction \(a_0\) of \(\tau\), the curve
\[
\xi
\longmapsto
\Psi_{ij,\xi}(a_0)
\]
is a line segment of bounded length inside the bounded direction domain \(U\), by
\cref{lem:v042-time-bilipschitz}.  It meets \(O(R^{-1})\) direction \(R\)-cubes.  For
\(a\in\tau\), the affine homothety formula for \(\Psi_{ij,\xi}(a)\) has a uniformly bounded
coefficient of \(a\) on the outer wing.  Hence allowing all \(a\in\tau\) only thickens that segment
by \(O(R)\), and the same \(O(R^{-1})\) direction-cube bound remains valid.

Fix one output direction cube \(\tau'\) met by this tube and choose \(u_{\tau'}\in\tau'\).
Every incidence counted for the physical endpoint \(j\) satisfies
\[
|b(u)+s_ju-c_j|
\lesssim
R.
\]
Since \(u\in\tau'\) and the center levels remain in a fixed bounded interval,
\[
|b(u)-(c_j-s_ju_{\tau'})|
\lesssim
R.
\]
Thus all output graph points of this fixed edge above \(\tau'\) meet only \(O(1)\) intercept
\(R\)-cubes \(\zeta'\), regardless of how many other selector-graph sheets lie above
\(\tau'\).  The fixed level \(s_j\) determines one interval \(K'\).  Multiplying the
\(O(R^{-1})\) direction-cube count by this \(O(1)\) anchored intercept count proves
\eqref{eq:v043-one-edge-code-valence}.
\end{proof}

\begin{theorem}[Universal high endpoint-code degree gives a genuine physical star]
\label{thm:v043-code-degree-to-star}
Let \(\mathfrak D\) be any double-code family, and let
\(\deg_{\mathfrak D}(v)\) denote endpoint-code degree.  If
\[
\deg_{\mathfrak D}(v)>L,
\]
then one of the \(O(1)\) physical centers localized by \(v\) is incident to at least
\begin{equation}
\label{eq:v043-code-degree-to-star}
\boxed{
\gtrsim
RL
}
\end{equation}
distinct original fixed-angle edges.  These edges form a genuine physical star and retain their
outer-wing collision witnesses.
\end{theorem}

\begin{proof}
The proof of \cref{thm:v042-double-code-capacity} shows that one endpoint code \(v=(\chi,K)\)
contains only \(O(1)\) possible physical centers.  For one such center \(i\), every fixed physical
edge \(ij\) accounts for at most \(B_R\) incident double codes by
\eqref{eq:v043-one-edge-code-valence}.  Therefore the \(L\) code edges incident to \(v\) require at
least \(cL/B_R\) distinct physical edges after pigeonholing the \(O(1)\) centers.  If \(v\) was
the output endpoint, reverse the physical edges; the fixed-angle outer-wing relation and
\eqref{eq:v043-one-edge-code-valence} are symmetric under this reversal.  Use
\eqref{eq:v043-one-edge-code-valence}.
\end{proof}

\subsection{Endpoint codes and the transition to a Maslov shell}

\begin{lemma}[One double code determines its the preceding stage extraction cell]
\label{lem:v045-double-code-extraction-cell}
All physical edges contained in one double code satisfy
\begin{equation}
\label{eq:v045-double-code-extraction-localization}
\boxed{
\operatorname{diam}\{a_{ij}\}
\lesssim
Rd_R^{-2}
\ll
\eta_R
\quad\text{in the far chart},
\qquad
\operatorname{diam}\{h_{ij}\}
\lesssim
R
\ll
\zeta_R
\quad\text{in the near chart}.
}
\end{equation}
Consequently one double code meets only \(O(1)\) of the extraction cells used in
\cref{thm:v029-far-decomposition,thm:v029-near-decomposition}.
\end{lemma}

\begin{proof}
By \cref{thm:v042-double-code-capacity}, the two endpoint codes localize \(x_i=(c_i,s_i)\) and
\(x_j=(c_j,s_j)\) to respective \(O(R)\)-balls.  Hence the ordered displacement
\(h_{ij}=c_i-c_j\) varies by \(O(R)\).  In the far chart
\(|s_i-s_j|\ge d_R\), and the quotient map
\[
(h,d)\longmapsto h/d
\]
has Lipschitz constant \(O(d_R^{-2})\) on the bounded physical parameter region.  Thus
\(a_{ij}=h_{ij}/(s_i-s_j)\) varies by \(O(Rd_R^{-2})\).  Finally,
\[
Rd_R^{-2}
=o\!\left((R/d_R^2)^{1/3}\right)
=o(\eta_R),
\qquad
R=o(R^{1/3})=o(\zeta_R).
\]
A set of diameter smaller than the grid scale meets only boundedly many half-open grid cells.
\end{proof}

\begin{corollary}[A nonsummable high code has the prescribed clustered/separated center split]
\label{cor:v045-extraction-cell-split}
Let an endpoint code \(v\) satisfy
\(\deg(v)>AL_*\).  Assign each incident double code to one of the \(O(1)\) compatible
the preceding stage extraction cells furnished by
\cref{lem:v045-double-code-extraction-cell}, and write \(\deg_\alpha(v)\) for the resulting
cell-degree.  After adjusting absolute constants, at least one of the following holds:
\begin{enumerate}[label=\textnormal{(\roman*)}]
\item some extraction cell satisfies \(\deg_\alpha(v)>cL_*\), and the physical edges in that cell
contain a clustered star of degree \(\gtrsim D_0\);
\item one of the \(O(1)\) physical centers localized by \(v\) is the center of a genuine star
containing at least \(cA\) physical edges whose radial centers, respectively displacements, occupy
extraction-scale-separated cells.
\end{enumerate}
Equivalently,
\begin{equation}
\label{eq:v045-extraction-cell-split}
\boxed{
\max_\alpha\deg_\alpha(v)
\gtrsim
L_*
\quad\text{or}\quad
\#\{\alpha:\deg_\alpha(v)>0\}
\gtrsim
A.
}
\end{equation}
Thus, after choosing any slowly divergent \(A_R=R^{-o(1)}\), a fixed-angle contribution outside
the dimension budget can leave the double-code stage only through the clustered alternative or
through a high endpoint code meeting at least \(cA_R\) separated distinguished-center cells,
exactly as prescribed in the preceding stage.
\end{corollary}

\begin{proof}
The cell-degrees sum to \(\deg(v)\), up to the harmless bounded assignment multiplicity.  If their
maximum is \(\gtrsim L_*\), then
\cref{lem:v043-one-edge-code-valence,thm:v043-code-degree-to-star} give at least
\[
R L_*
\asymp
D_0
\]
distinct physical edges inside that single extraction cell.  In the near chart first split the
edges according to the two ordered displacement orientations, losing only an absolute factor.
Otherwise division of \(\deg(v)>AL_*\) by the maximum cell-degree gives at least \(cA\) occupied
cells.  For every incident double code choose one realizing ordered physical edge and assign the
code to the unique half-open extraction cell containing that edge's distinguished center or
displacement.  Color the extraction grid by coordinate residues modulo \(3\); one color retains a
fixed fraction and makes the chosen cells extraction-scale separated.  Choose one realizing edge
from each retained cell.  Distinct retained cells give distinct physical edges.  Their endpoint
belonging to \(v\) is one of only \(O(1)\) physical centers by
\cref{thm:v042-double-code-capacity}; pigeonholing that bounded set leaves \(cA\) edges incident to
one center.  Reverse all edges first if \(v\) is their output endpoint, and in the near chart split
the two ordered displacement orientations.  Both operations lose only an absolute factor.
\end{proof}

\begin{lemma}[The distinguished center is selected by the contact secant]
\label{lem:v045-contact-secant-identity}
Let a double-code incidence come from a physical edge \(ij\), source direction \(a\), output
direction \(u=\Psi_{ij,\xi}(a)\), and outer-wing parameter \(\xi\).  Then
\begin{equation}
\label{eq:v045-contact-secant-identity}
\boxed{
\begin{aligned}
b(u)-b(a)
&=
\xi(a-u)+O(R),\\
h_{ij}
&=
(\xi-s_j)u-(\xi-s_i)a,\\
a_{ij}
&=
a+\frac{\xi-s_j}{s_i-s_j}(u-a)
\qquad (|s_i-s_j|\ge d_R).
\end{aligned}
}
\end{equation}
Consequently, on the fixed-angle piece \(|a-u|\ge\tau_0\), the two endpoint graph cells
determine \(\xi\) up to \(O_{\tau_0}(R)\).  Thus the separated-center alternative in
\cref{cor:v045-extraction-cell-split} is not a free-center family: its center cells are the ones
selected by the contact secants of the same graph incidences.
\end{lemma}

\begin{proof}
Choose errors \(e_i,e_j\) of size \(O(R)\) so that
\[
c_i=b(a)+s_i a+e_i,
\qquad
c_j=b(u)+s_j u+e_j.
\]
The collision relation
\[
c_i+(\xi-s_i)a
=
c_j+(\xi-s_j)u
\]
gives the first two identities in \eqref{eq:v045-contact-secant-identity}.  Since
\(h_{ij}=c_i-c_j\), the second identity can also be written
\[
h_{ij}
=
(s_i-s_j)a+(\xi-s_j)(u-a).
\]
Division by \(s_i-s_j\) gives the far-chart formula.  Finally, taking the scalar product of
the first identity with \(a-u\) yields
\[
\xi
=
\frac{(b(u)-b(a))\mathbin{\cdot}(a-u)}{|a-u|^2}
+O_{\tau_0}(R),
\]
which proves the asserted selection statement.
\end{proof}

\clearpage
\section{Angular blow-up and Maslov four-cycles}

\subsection{Darboux blow-up and four-cycle extraction}

\begin{equation}
\label{eq:v046-newton-near-diagonal}
\boxed{
\iint_{|a-u|\le\tau}
\frac{d\sigma(a)d\sigma(u)}{|a-u|}
\lesssim_U \tau^2.
}
\end{equation}

\begin{proposition}[Critical mass commutes with the Darboux blow-up]
\label{prop:v046-critical-renormalization}
Let \(q_R(P)\) be the collision mass assigned to \(P\).  Up to the fixed enlargement used to hold
neighboring grid cells,
\begin{equation}
\label{eq:v046-local-critical-renormalization}
\boxed{
q_R(P)
=
m_P^2\widehat\tau_R^{-1}
M_{\epsilon_R}^{J,\mathrm{sh}}(\mu_P^\sharp).
}
\end{equation}
On the density-good patches
\(m_P\asymp_U\widehat\tau_R^3\), which carry all but \(o(1)\) of the direction mass after passing
to a dyadic subsequence, this gives
\begin{equation}
\label{eq:v046-kappa-local-average}
\boxed{
\kappa_R
\asymp_U
\sum_P
m_P\,M_{\epsilon_R}^{J,\mathrm{sh}}(\mu_P^\sharp).
}
\end{equation}
Thus \(\kappa_R\) is precisely the direction-mass average of the critical collision masses in the
unit Darboux blow-ups.
\end{proposition}

\begin{proof}
Under the simultaneous dilation
\[
(a-a_P,b-b_P)
=
\widehat\tau_R(\bar a,\bar b),
\]
the collision inequality at thickness \(R\) becomes the same inequality at thickness
\(\epsilon_R\), while the collision-time length is unchanged.  Hence the normalized weight scales
as
\[
w_R^J
=
\widehat\tau_R^{-1}w_{\epsilon_R}^J.
\]
Replacing each copy of \(\sigma|_P\) by \(m_P\sigma_P\) proves
\eqref{eq:v046-local-critical-renormalization}.  Summing in \(P\), dividing by
\(\widehat\tau_R^2\), and using
\(m_P/\widehat\tau_R^3\asymp_U1\) gives
\eqref{eq:v046-kappa-local-average}.

For the density-good assertion, use the dyadic martingale differentiation theorem on
\(\mathbf 1_U\): the \(U\)-mass in dyadic cells whose relative density stays below any fixed
positive threshold tends to zero.  The bounded coloring and enlargement change only fixed
constants.
\end{proof}

\begin{lemma}[Rescaled contact-secant identity]
\label{lem:v046-rescaled-contact-secant}
Every physical re-collision represented in
\eqref{eq:v045-contact-secant-identity} satisfies
\begin{equation}
\label{eq:v046-rescaled-contact-secant}
\boxed{
\Pi_\xi(\alpha_R,\beta_R)
=
\beta_R+\xi\alpha_R
=
O\!\left(\frac{R}{\tau_R}\right),
\qquad
\Pi_\xi(\alpha,\beta)=\beta+\xi\alpha.
}
\end{equation}
Thus the rescaled secant lies within \(O(R/\tau_R)\) of
\[
L_\xi=\ker\Pi_\xi=\{(\alpha,-\xi\alpha):\alpha\in\R^3\},
\]
the Lagrangian member of the Reeb/Maslov pencil.
\end{lemma}

\begin{proof}
The first line of \eqref{eq:v045-contact-secant-identity} is
\[
b(u)-b(a)+\xi(u-a)=O(R).
\]
Divide by \(\tau_R\).  The description of the kernel is immediate, and it is Lagrangian because
it is the graph of the symmetric map \(-\xi I\).
\end{proof}

\begin{equation}
\label{eq:v046-limit-maslov-support}
\boxed{
\supp\mathfrak q_\infty
\subset
\{(a,\alpha,\beta,\xi):
(\alpha,\beta)\in L_\xi\}.
}
\end{equation}

\begin{equation}
\label{eq:v046-direction-determinant-sublevel}
\boxed{
\rho^{\otimes4}
\left\{
(a_0,a_1,a_2,a_3):
\left|
\det(a_1-a_0,a_2-a_0,a_3-a_0)
\right|
\le\Delta
\right\}
\lesssim \Delta.
}
\end{equation}

\begin{equation}
\label{eq:v046-approx-edge-relation}
|b_y-b_x+s_{xy}(a_y-a_x)|\le\epsilon.
\end{equation}

\begin{theorem}[Quantitative Maslov four-cycle saving]
\label{thm:v046-quantitative-four-cycle}
With
\[
e(h)=\iint h(x,y)\,d\nu(x)d\nu(y),
\]
one has, for every \(0<\Delta<1\),
\begin{equation}
\label{eq:v046-four-cycle-saving}
\boxed{
e(h)^4
\lesssim_{U,J}
\Delta
+
\mathfrak B_{C_{U,J}\epsilon/\Delta}(h).
}
\end{equation}
In particular, with \(\Delta=\sqrt\epsilon\),
\[
e(h)
\lesssim_{U,J}
\epsilon^{1/8}
+
\mathfrak B_{C_{U,J}\sqrt\epsilon}(h)^{1/4}.
\]
\end{theorem}

\begin{proof}
Twice applying Cauchy--Schwarz to the common-neighbor kernel gives the four-cycle inequality
\[
e(h)^4
\le
\int h_{01}h_{12}h_{23}h_{30}\,d\nu^4.
\]
Split the integral according to whether the horizontal tetrahedron determinant is at most
\(\Delta\).  The first part is \(O(\Delta)\) by
\eqref{eq:v046-direction-determinant-sublevel}.

For a nondegenerate four-cycle, orient its horizontal edge vectors cyclically as
\(v_1,v_2,v_3,v_4\), so \(v_1+v_2+v_3+v_4=0\).  Summing
\eqref{eq:v046-approx-edge-relation} around the cycle and subtracting the fourth edge label gives
\[
(s_1-s_4)v_1
+(s_2-s_4)v_2
+(s_3-s_4)v_3
=
O(\epsilon).
\]
The three vectors \(v_1,v_2,v_3\) have bounded norm and determinant comparable, up to a fixed
linear change, to the horizontal tetrahedron determinant.  Cramer's rule therefore yields
\[
\max_{1\le i\le4}|s_i-s_4|
\lesssim_{U,J}
\frac{\epsilon}{\Delta}.
\]
Choose \(s=s_4\) and \(c=b_{x_0}+sa_{x_0}\).  Propagating the four approximate edge relations
around the cycle gives
\[
|b_{x_i}+sa_{x_i}-c|
\lesssim_{U,J}
\frac{\epsilon}{\Delta}
\qquad(0\le i\le3).
\]
Thus every nondegenerate cycle is counted by
\(\mathfrak B_{C\epsilon/\Delta}(h)\), proving
\eqref{eq:v046-four-cycle-saving}.
\end{proof}

\begin{theorem}[Quantitative Maslov common-neighbor dichotomy]
\label{thm:v046-common-neighbor-dichotomy}
Under the hypotheses of \cref{thm:v046-quantitative-four-cycle}, put
\[
\eta=e(h).
\]
There are constants depending only on the bounded coordinate region and on \(J\) such that at
least one of the following holds:
\begin{enumerate}[label=\textnormal{(\roman*)}]
\item
\begin{equation}
\label{eq:v046-common-neighbor-small}
\boxed{
\eta\lesssim \epsilon^{1/4}.
}
\end{equation}
\item there are \(s_0\in J\) and \(c_0\in\R^3\) for which
\begin{equation}
\label{eq:v046-common-neighbor-bush}
\boxed{
\nu\left\{(a,b):|b+s_0a-c_0|
\le C\epsilon\eta^{-4}\right\}
\gtrsim \eta^3.
}
\end{equation}
\end{enumerate}
Thus a collision graph of mass larger than the fourth-root error scale produces, at the same
finite scale, one positive-mass approximate affine Lagrangian point-probe bush.
\end{theorem}

\begin{proof}
Write
\[
K(p,q)=\int h(p,y)h(q,y)\,d\nu(y),
\qquad
d\lambda_{p,q}(y)=h(p,y)h(q,y)\,d\nu(y).
\]
The four-cycle integral is
\(
\iint K(p,q)^2\,d\nu(p)d\nu(q)
\)
and is at least \(\eta^4\).  Choose
\[
\Delta=c_*\eta^4
\]
with \(c_*>0\) sufficiently small.  By
\eqref{eq:v046-direction-determinant-sublevel}, the determinant-degenerate cycles have total mass
at most \(\eta^4/2\).  Hence the nondegenerate common-neighbor pairs have total mass
\begin{equation}
\label{eq:v046-good-common-neighbor-mass}
G\gtrsim\eta^4.
\end{equation}

For fixed \(p,q\), let \(G(p,q)\) be their contribution to \(G\).  Since
\[
\iint K(p,q)\,d\nu(p)d\nu(q)
=\int\left(\int h(x,y)\,d\nu(x)\right)^2d\nu(y)
\le\iint h\,d\nu^2
=\eta,
\]
there is a pair \((p,q)\) with
\begin{equation}
\label{eq:v046-good-over-common-neighbor}
\frac{G(p,q)}{K(p,q)}\gtrsim\eta^3.
\end{equation}
For a nondegenerate pair of common neighbors \(y,y'\), the Cramer-rule argument in the proof of
\cref{thm:v046-quantitative-four-cycle} gives
\[
|s_{py}-s_{py'}|
\lesssim
r,
\qquad
r:=C\epsilon\eta^{-4}.
\]
If \(r\gtrsim1\), then (i) holds.  Otherwise partition \(J\) into half-open intervals of length
comparable to \(r\).  A good pair \(y,y'\) can occur only in the same interval or in a bounded
number of neighboring intervals.  Therefore
\[
G(p,q)
\lesssim
K(p,q)
\sup_I\lambda_{p,q}\{y:s_{py}\in I\},
\]
where the supremum is over intervals \(I\) of length \(Cr\).  By
\eqref{eq:v046-good-over-common-neighbor}, some such interval \(I_0\), centered at \(s_0\), has
\[
\lambda_{p,q}\{y:s_{py}\in I_0\}
\gtrsim\eta^3.
\]
Because \(0\le h\le1\), the same set has \(\nu\)-measure at least \(c\eta^3\).  For every \(y\) in
this set the edge relation between \(p\) and \(y\) gives
\[
|b_y+s_0a_y-(b_p+s_0a_p)|
\le
\epsilon+C_U|s_{py}-s_0|
\lesssim r.
\]
Taking \(c_0=b_p+s_0a_p\) proves (ii).
\end{proof}

\begin{lemma}[A dyadic code-edge layer retains the joint-box mass]
\label{lem:v046-weighted-code-edge-layer}
For fixed \(L\), let \(\omega_L(e)\) be the total assigned incidence mass from
\(\mathcal B_L\) carried by the double-code edge \(e\).  Then
\[
0\le\omega_L(e)\lesssim R^4
\]
and, with dyadic \(R^3\le\lambda\lesssim1\),
\begin{equation}
\label{eq:v046-weighted-code-layer}
\boxed{
\sum_{B\in\mathcal B_L}m_B
\lesssim
R^{-1-o(1)}
+
(\log R^{-1})
\max_\lambda
\lambda R^4\,|E(\mathscr G_{L,\lambda})|,
}
\end{equation}
where
\[
E(\mathscr G_{L,\lambda})
=
\{e:\lambda R^4\le\omega_L(e)<2\lambda R^4\}.
\]
Thus any contribution above the target \(R^{-2+o(1)}\) survives in a simple code graph whose edges
all have comparable positive incidence mass.
\end{lemma}

\begin{proof}
The half-open center and time assignments partition the incidence mass, so summing
\(\omega_L(e)\) over \(e\) gives the left side.  The double-code capacity
\eqref{eq:v042-double-code-capacity} gives the pointwise upper bound.  There are at most
\(N_R^2\) possible code edges, where
\[
N_R=|\mathscr V_R|\lesssim R^{-4-o(1)}.
\]
Therefore the edges with \(\omega_L(e)<R^7\) carry total mass at most
\[
R^7N_R^2\lesssim R^{-1-o(1)},
\]
which is already below the required \(R^{-2+o(1)}\) budget.  The remaining normalized weights
\(\omega_L(e)/R^4\) occupy \(O(\log R^{-1})\) dyadic classes.  Pigeonholing proves
\eqref{eq:v046-weighted-code-layer}.
\end{proof}

\begin{lemma}[Power excess creates many code four-cycles]
\label{lem:v046-code-cycle-supersaturation}
Let \(\mathscr G\) be any simple graph on the endpoint-code vertex set and put
\(N_R=|\mathscr V_R|\).  If, for some \(A\ge A_0\),
\[
|E(\mathscr G)|
\ge
A N_R^{3/2},
\]
then \(\mathscr G\) contains at least
\begin{equation}
\label{eq:v046-code-cycle-supersaturation}
\boxed{
cA^4N_R^2
}
\end{equation}
unoriented simple four-cycles.
\end{lemma}

\begin{proof}
Write \(E=|E(\mathscr G)|\).  If \(d(x,y)\) denotes the number of common neighbors of \(x,y\), then
\[
\#C_4(\mathscr G_L)
=\frac12\sum_{\{x,y\}}\binom{d(x,y)}2.
\]
Also
\[
\sum_{\{x,y\}}d(x,y)=\sum_z\binom{\deg z}2
\gtrsim E^2/N_R.
\]
Twice using convexity, and taking \(A_0\) large enough to absorb the linear terms in the binomial
coefficients, gives
\[
\#C_4(\mathscr G_L)
\gtrsim E^4/N_R^4
\ge cA^4N_R^2.
\]
\end{proof}

\begin{lemma}[A nondegenerate code four-cycle is one near-Lagrangian bush]
\label{lem:v046-code-cycle-maslov-bush}
Let
\[
v_0v_1v_2v_3v_0
\]
be a four-cycle of endpoint codes.  Choose once and for all one graph representative
\[
z_k=(a_k,b_k)\in\chi_k\cap G_b
\]
for each \(v_k=(\chi_k,K_k)\).  Then every edge \(v_kv_{k+1}\), with indices modulo four, has a
collision label \(\xi_k\in I+O(R)\) satisfying
\begin{equation}
\label{eq:v046-code-cycle-contact}
\boxed{
b_{k+1}-b_k+\xi_k(a_{k+1}-a_k)=O(R).
}
\end{equation}
If
\[
\left|
\det(a_1-a_0,a_2-a_1,a_3-a_2)
\right|
\ge\Delta,
\]
then there are \(s_0\in I+O(R)\) and \(c_0\in\R^3\) such that
\begin{equation}
\label{eq:v046-code-cycle-bush}
\boxed{
|b_k+s_0a_k-c_0|
\lesssim
R/\Delta
\qquad(0\le k\le3).
}
\end{equation}
Thus every horizontally nondegenerate code four-cycle lies in one
\(O(R/\Delta)\)-approximate affine Lagrangian point-probe bush.
\end{lemma}

\begin{proof}
For an edge \(v_kv_{k+1}\), choose one realizing physical incidence.  Its actual graph endpoints
belong to \(\chi_k,\chi_{k+1}\), so replacing them by the fixed representatives changes the
contact-secant identity \eqref{eq:v045-contact-secant-identity} by \(O(R)\).  This proves
\eqref{eq:v046-code-cycle-contact}.

Put \(w_k=a_{k+1}-a_k\).  Summing \eqref{eq:v046-code-cycle-contact} around the cycle gives
\[
\xi_0w_0+\xi_1w_1+\xi_2w_2+\xi_3w_3=O(R),
\qquad
w_0+w_1+w_2+w_3=0.
\]
After subtracting \(\xi_3\) times the second identity, Cramer's rule yields
\[
\max_{0\le k\le3}|\xi_k-\xi_3|
\lesssim R/\Delta.
\]
Take \(s_0=\xi_3\) and \(c_0=b_0+s_0a_0\).  Propagating
\eqref{eq:v046-code-cycle-contact} around the four edges proves
\eqref{eq:v046-code-cycle-bush}.
\end{proof}

\subsection{Time synchronization and physical bush cells}

\begin{lemma}[Four-time Maslov interpolation on one contact cycle]
\label{lem:v046-four-time-maslov-cycle}
For a code four-cycle as in \cref{lem:v046-code-cycle-maslov-bush}, put
\[
d_k=(a_{k+1}-a_k,b_{k+1}-b_k)\in\R_a^3\oplus\R_b^3
\qquad(k\ {\rm mod}\ 4).
\]
Then
\[
d_0+d_1+d_2+d_3=0.
\]
If \(\Span\{d_0,d_1,d_2\}\) has dimension at most two, the cycle already has full secant rank at
most two.  Otherwise let
\[
V_C=\Span\{d_0,d_1,d_2\}
\]
and let \(P_{V_C}(s)\) be its Reeb-front polynomial.  For the four collision labels
\(\xi_0,\ldots,\xi_3\),
\begin{equation}
\label{eq:v046-cycle-polynomial-values}
\boxed{
|P_{V_C}(\xi_k)|
\lesssim_{\tau_0,I}
R
\qquad(0\le k\le3).
}
\end{equation}
Consequently, if
\[
\mathfrak V(\xi_0,\xi_1,\xi_2,\xi_3)
:=
\prod_{0\le i<j\le3}|\xi_i-\xi_j|
\ge\upsilon,
\]
then
\begin{equation}
\label{eq:v046-four-time-maslov-interpolation}
\boxed{
\sup_{s\in J}|P_{V_C}(s)|
\lesssim_{I,J,\tau_0}
\frac{R}{\upsilon}.
}
\end{equation}
In particular \cref{lem:v046-polynomial-separation} gives
\[
|\det A_C|+|P_{V_C}(i)|
\lesssim
R/\upsilon,
\]
so a four-time cycle with \(R/\upsilon\to0\) is quantitatively horizontal-rank deficient and
maximally symplectic.
\end{lemma}

\begin{proof}
The telescoping identity gives \(\sum_kd_k=0\), hence \(d_3\in V_C\).  Choose an orthonormal frame
\[
T_Cc=(A_Cc,B_Cc)
\]
for \(V_C\).  Write \(d_k=T_Cc_k\).  Fixed-angle separation gives
\[
|c_k|=|d_k|\ge|a_{k+1}-a_k|\ge\tau_0.
\]
By \eqref{eq:v046-code-cycle-contact},
\[
|(B_C+\xi_kA_C)c_k|
=
|b_{k+1}-b_k+\xi_k(a_{k+1}-a_k)|
\lesssim R.
\]
Thus the smallest singular value of \(B_C+\xi_kA_C\) is \(O_{\tau_0}(R)\).  The other two
singular values are uniformly bounded on the fixed interval containing the \(\xi_k\)'s, proving
\eqref{eq:v046-cycle-polynomial-values}.

Lagrange interpolation of the cubic \(P_{V_C}\) at the four nodes \(\xi_k\) gives
\[
P_{V_C}(s)
=
\sum_{k=0}^3
P_{V_C}(\xi_k)
\prod_{\ell\ne k}\frac{s-\xi_\ell}{\xi_k-\xi_\ell}.
\]
All numerator factors are bounded on \(J\).  Since all pairwise distances are also bounded above,
the lower bound on their total product implies
\[
\prod_{\ell\ne k}|\xi_k-\xi_\ell|
\gtrsim_I\upsilon
\]
for each \(k\).  Insert \eqref{eq:v046-cycle-polynomial-values} to obtain
\eqref{eq:v046-four-time-maslov-interpolation}, and then apply
\cref{lem:v046-polynomial-separation}.
\end{proof}

\begin{lemma}[Two-time common neighbors are bushes or angular packets]
\label{lem:v046-two-time-common-neighbors}
Fix two code vertices \(p,q\) and a family \(Y\) of their common neighbors.  For each \(y\in Y\)
choose collision labels \(s_y\) on \(py\) and \(t_y\) on \(qy\), so that
\[
\begin{aligned}
b_y-b_p+s_y(a_y-a_p)&=O(R),\\
b_y-b_q+t_y(a_y-a_q)&=O(R).
\end{aligned}
\]
Suppose all \(s_y,t_y\) lie in the union of two intervals \(I_1,I_2\) of length \(\eta_s\), whose
distance is at least \(\delta_s\ge4\eta_s\).  Then \(Y\) is the union of four pieces with the
following disposition:
\begin{enumerate}[label=\textnormal{(\roman*)}]
\item if \(s_y,t_y\in I_\alpha\), the corresponding graph points \(y\) lie in one
\(O(R+\eta_s)\)-approximate Lagrangian bush at the representative time of \(I_\alpha\);
\item if \(s_y\in I_\alpha\), \(t_y\in I_\beta\), \(\alpha\ne\beta\), their horizontal directions
lie in one ball of radius
\begin{equation}
\label{eq:v046-two-time-angular-radius}
\boxed{
C_{I,U}\frac{R+\eta_s}{\delta_s^2}.
}
\end{equation}
\end{enumerate}
Consequently any positive common-neighbor mass carried by two separated Reeb-time packets has a
fixed positive fraction in a physical near-bush or in a shrinking angular packet.
\end{lemma}

\begin{proof}
If both labels lie in \(I_\alpha\), with representative \(s_\alpha\), the first edge equation gives
\[
b_y+s_\alpha a_y
=
b_p+s_\alpha a_p+O(R+\eta_s).
\]
Thus all such \(y\)'s belong to the displayed bush.  The second equation shows, whenever the piece
is nonempty, that the bush center defined from \(q\) differs from the one defined from \(p\) by the
same error.

For the off-diagonal pieces, subtract the two edge equations to obtain
\begin{equation}
\label{eq:v046-two-anchor-rational}
(t_y-s_y)a_y
=
b_q-b_p+t_ya_q-s_ya_p+O(R).
\end{equation}
Here \(|t_y-s_y|\ge\delta_s/2\).  Replacing \(s_y,t_y\) by the representatives
\(s_\alpha,t_\beta\) changes the rational expression on the right of
\eqref{eq:v046-two-anchor-rational}, after division by \(t_y-s_y\), by at most
\[
C_{I,U}(R+\eta_s)\delta_s^{-2}.
\]
Hence every \(a_y\) lies in the asserted ball about the value obtained from
\((s_\alpha,t_\beta)\).  There are two diagonal and two off-diagonal packet pairs, completing the
four-piece decomposition.
\end{proof}

\begin{lemma}[Anchored time-packet trichotomy]
\label{lem:v046-anchored-time-packet-trichotomy}
Let \(Y=N(p)\cap N(q)\) be a common-neighbor set of cardinality \(T\), equipped with the two edge
labels \(s_y,t_y\) from \cref{lem:v046-two-time-common-neighbors}.  Fix
\[
R\le\delta_s\ll1,
\qquad
\tau_1>0.
\]
At least one of the following holds, with constants depending only on the bounded coordinate
region and on \(\tau_1\):
\begin{enumerate}[label=\textnormal{(\roman*)}]
\item \(|a_p-a_q|<\tau_1\);
\item at least \(cT^2\) pairs \((y,y')\in Y^2\) have three labels among
\[
s_y,t_y,s_{y'},t_{y'}
\]
which are pairwise \(c\delta_s\)-separated;
\item there are two intervals \(I_1,I_2\) of length \(C\delta_s\) such that
\[
\#\{y\in Y:s_y,t_y\in I_1\cup I_2\}
\ge cT.
\]
\end{enumerate}
In (iii), \cref{lem:v046-two-time-common-neighbors} sends a fixed fraction of the retained
neighbors to one near-bush or to at most two angular packets.
\end{lemma}

\begin{proof}
Assume (i) fails.  Call two neighbors \(y,y'\) compatible when their four labels can be covered by
two intervals of length \(10\delta_s\).  An elementary one-dimensional packing argument says that
an incompatible pair contains three \(c\delta_s\)-separated labels.  Hence, if a fixed positive
fraction of \(Y^2\) is incompatible, (ii) holds.

Otherwise discard at most a fixed small fraction of \(Y\) so that every remaining neighbor is
compatible with a fixed positive fraction of the others.  If one such \(y_0\) has
\[
|s_{y_0}-t_{y_0}|>20\delta_s,
\]
then compatibility with \(y_0\) forces both labels of every compatible neighbor into the two
\(O(\delta_s)\)-intervals about \(s_{y_0}\) and \(t_{y_0}\).  This gives (iii).

It remains to consider the neighbors with
\(|s_y-t_y|\le20\delta_s\).  Subtracting their two contact-edge equations as in
\eqref{eq:v046-two-anchor-rational}, and putting \(r_y=(s_y+t_y)/2\), gives
\[
|b_q-b_p+r_y(a_q-a_p)|
\lesssim R+\delta_s.
\]
Since \(|a_p-a_q|\ge\tau_1\), all such \(r_y\)'s lie in one interval of length
\(C_{\tau_1}(R+\delta_s)\lesssim_{\tau_1}\delta_s\).  Thus their two labels lie in one
\(C_{\tau_1}\delta_s\)-interval.  If they form a fixed fraction of \(Y\), (iii) holds with the
second interval empty.  Otherwise the preceding off-diagonal argument applies to a fixed fraction
of the remaining neighbors and again gives (iii).
\end{proof}

\begin{lemma}[Nondegenerate cycle density concentrates in one bush cell]
\label{lem:v046-cycle-bush-cell-aggregation}
Fix \(0<\Delta\le1\), and let \(\mathcal Z_\Delta\) be any family of code four-cycles satisfying the
determinant lower bound in \cref{lem:v046-code-cycle-maslov-bush}.  Then there are
\(s_Q\in I+O(R)\), \(c_Q\in\R^3\), and a set \(\mathscr W_Q\subset\mathscr V_R\) such that every
representative \(z_v=(a_v,b_v)\), \(v\in\mathscr W_Q\), satisfies
\[
|b_v+s_Qa_v-c_Q|
\lesssim R/\Delta,
\]
and
\begin{equation}
\label{eq:v046-cycle-bush-cell-vertices}
\boxed{
|\mathscr W_Q|
\gtrsim
|\mathcal Z_\Delta|^{1/4}\frac{R}{\Delta}.
}
\end{equation}
Consequently, if the nondegenerate cycles form a fixed proportion of the output in
\eqref{eq:v046-code-cycle-supersaturation}, then
\begin{equation}
\label{eq:v046-cycle-bush-cell-power}
\boxed{
|\mathscr W_Q|
\gtrsim
A N_R^{1/2}\frac{R}{\Delta}.
}
\end{equation}
\end{lemma}

\begin{proof}
The relevant \((s,c)\)-parameters stay in a fixed bounded subset of \(\R^4\).  Cover that set by
half-open cells of side comparable to \(r=R/\Delta\).  There are \(O(r^{-4})\) cells.  Assign each
cycle in \(\mathcal Z_\Delta\) to a cell containing one bush parameter supplied by
\cref{lem:v046-code-cycle-maslov-bush}.  One cell \(Q\) receives at least
\[
c|\mathcal Z_\Delta|r^4
\]
cycles.  All vertices of every assigned cycle satisfy the enlarged bush inequality for the center
of \(Q\).  If their union has \(T\) vertices, it supports at most \(T^4\) ordered four-cycles.
Thus
\[
T^4\gtrsim|\mathcal Z_\Delta|r^4,
\]
which proves \eqref{eq:v046-cycle-bush-cell-vertices}.  Insert
\eqref{eq:v046-code-cycle-supersaturation} for the last assertion.
\end{proof}

\begin{corollary}[Bush-cell aggregation retains comparable edge mass]
\label{cor:v046-weighted-cycle-bush-cell}
Apply \cref{lem:v046-cycle-bush-cell-aggregation} inside one graph
\(\mathscr G_{L,\lambda}\).  If a bush cell \(Q\) receives \(Z_Q\) four-cycles, and
\(\mathscr E_Q\) denotes the union of their code edges, then
\begin{equation}
\label{eq:v046-weighted-cycle-bush-mass}
\boxed{
\sum_{e\in\mathscr E_Q}\omega_L(e)
\gtrsim
\lambda R^4 Z_Q^{1/2}.
}
\end{equation}
In particular, if
\[
|E(\mathscr G_{L,\lambda})|
\ge A N_R^{3/2}
\]
and a fixed proportion of the resulting four-cycles are \(\Delta\)-nondegenerate, some
\(O(R/\Delta)\)-Lagrangian bush cell carries incidence mass at least
\begin{equation}
\label{eq:v046-power-bush-cell-mass}
\boxed{
c\lambda A^2R^4N_R(R/\Delta)^2.
}
\end{equation}
Since \(N_R\gtrsim R^{-4}\), the last quantity is
\[
\gtrsim \lambda A^2R^2\Delta^{-2}.
\]
\end{corollary}

\begin{proof}
A simple four-cycle is determined, up to a fixed number of choices, by two of its opposite edges.
Hence \(Z_Q\lesssim|\mathscr E_Q|^2\), so
\[
|\mathscr E_Q|\gtrsim Z_Q^{1/2}.
\]
Every edge in the dyadic layer has mass at least \(\lambda R^4\), proving
\eqref{eq:v046-weighted-cycle-bush-mass}.  The proof of
\cref{lem:v046-cycle-bush-cell-aggregation}, combined with
\cref{lem:v046-code-cycle-supersaturation}, gives
\[
Z_Q\gtrsim A^4N_R^2(R/\Delta)^4,
\]
which yields \eqref{eq:v046-power-bush-cell-mass}.

Finally, the projection of \(G_b\) onto \(U\) is all of \(U\), so at least \(c_UR^{-3}\)
direction cubes contain an occupied graph product cell.  The definition of \(\mathscr V_R\)
pairs every such graph cell with \(cR^{-1}\) level intervals.  Thus \(N_R\gtrsim R^{-4}\).
\end{proof}

\begin{proposition}[Double-cell de-multiplication produces a physical bush]
\label{prop:v046-double-cell-demultiplication}
Let \(Q\) receive \(Z_Q\) nondegenerate four-cycles from one dyadic edge layer
\(\mathscr G_{L,\lambda}\), and let \(\mathscr E_Q\) be the union of their code edges.  Then there is
a physical \(O(R/\Delta)\)-bush \(H_Q\) satisfying
\begin{equation}
\label{eq:v046-double-cell-bush-mass}
\boxed{
\sigma(H_Q)
\gtrsim
\lambda R^4 Z_Q^{1/4}.
}
\end{equation}
If \(Q\) is the cell supplied by power supersaturation, then
\begin{equation}
\label{eq:v046-power-direction-bush}
\boxed{
\sigma(H_Q)
\gtrsim
\lambda A R^4N_R^{1/2}\frac{R}{\Delta}
\gtrsim
\lambda A R^3\Delta^{-1}.
}
\end{equation}
The same quantities are lower bounds, up to fixed constants, for
\(\mathcal L^4(N_{CR/\Delta}(K_b))\).
\end{proposition}

\begin{proof}
Let \(\mathscr X_Q\) be the set of selector graph product cells \(\chi\) occurring in endpoint
codes incident to \(\mathscr E_Q\), and put \(D_Q=|\mathscr X_Q|\).  Each \(\chi\) supports only
\(O(R^{-1})\) endpoint codes \((\chi,K)\).  Since a simple graph on
\(O(D_QR^{-1})\) vertices has at most \(O(D_Q^2R^{-2})\) edges,
\[
D_Q
\gtrsim
R|\mathscr E_Q|^{1/2}
\gtrsim
RZ_Q^{1/4},
\]
where the last inequality is the first step in the proof of
\cref{cor:v046-weighted-cycle-bush-cell}.

Now fix \(e\in\mathscr E_Q\) incident to a graph cell \(\chi\).  The double code restricts its
collision time to an interval of length \(O(R)\), and contains only \(O(1)\) physical edges.
Therefore its lower mass bound
\[
\omega_L(e)\ge\lambda R^4
\]
forces the set of source directions in \(\chi\) that actually support this code to have
\(\sigma\)-mass at least \(c\lambda R^3\).  On the outer wing the source--output homothety
Jacobian is bounded above and below, so the same conclusion holds at either endpoint of \(e\).

Choose, for each \(\chi\in\mathscr X_Q\), one incident edge and the corresponding positive-mass
direction subset.  Distinct half-open graph product cells give disjoint subsets of the selector
graph, hence disjoint direction subsets because \(G_b\) is a graph over \(U\).  All their graph
points are within \(O(R/\Delta)\) of the affine Lagrangian leaf
\[
b+s_Qa=c_Q
\]
assigned to \(Q\).  Their union is therefore contained in one physical
\(O(R/\Delta)\)-bush and has mass at least
\[
c\lambda R^3D_Q
\gtrsim
\lambda R^4Z_Q^{1/4},
\]
proving \eqref{eq:v046-double-cell-bush-mass}.  Insert the lower bound for \(Z_Q\) from the proof of
\cref{lem:v046-cycle-bush-cell-aggregation} and use \(N_R\gtrsim R^{-4}\) to obtain
\eqref{eq:v046-power-direction-bush}.  The volume assertion is
\cref{prop:v010-bush-volume}.
\end{proof}

\subsection{Rank loss, interpolation, and packet reuse}

\begin{lemma}[Uniform separation from total Maslov trapping]
\label{lem:v046-polynomial-separation}
For every nontrivial compact interval \(J\subset\R\), there is \(C_J<\infty\) such that
\begin{equation}
\label{eq:v046-polynomial-separation}
\boxed{
|\det A|+|P_V(i)|
\le
C_J\sup_{s\in J}|P_V(s)|.
}
\end{equation}
Consequently, if
\(\sup_{s\in J}|P_V(s)|\le\varepsilon\), then
\begin{equation}
\label{eq:v046-near-maximal-defect}
\boxed{
|\det A|\lesssim_J\varepsilon,
\qquad
\|\omega|_V\|_{\rm op}
\ge
\sqrt{1-C_J^2\varepsilon^2}.
}
\end{equation}
Moreover the distance of \(V\), in the Grassmannian, from the front-null locus
\[
\mathfrak N
=
\{W:P_W\equiv0\}
=
\mathsf K_{101}\cup\mathsf K_{110}\cup\mathsf K_{200}
\]
tends uniformly to zero with \(\varepsilon\).
\end{lemma}

\begin{proof}
On the four-dimensional vector space of real polynomials of degree at most three, coefficient norm,
evaluation at \(i\), and the supremum norm on \(J\) are equivalent.  The leading coefficient of
\(P_V\) is \(\det A\), proving \eqref{eq:v046-polynomial-separation}.  The complex-evaluation
identity \eqref{eq:v003-complex-identity} gives
\[
1-\|\omega|_V\|_{\rm op}^2=|P_V(i)|^2,
\]
and hence \eqref{eq:v046-near-maximal-defect}.  For the last assertion, otherwise there would be
\(\varepsilon_n\downarrow0\) and planes \(V_n\) staying a fixed distance from \(\mathfrak N\).
Compactness of the Grassmannian gives \(V_n\to V\); continuity yields
\(P_V|_J=0\), hence \(P_V\equiv0\), a contradiction.  The equality with the three strata is the
classification of the preceding flag construction.
\end{proof}

\begin{lemma}[Rank-adapted Maslov interpolation]
\label{lem:v046-rank-adapted-interpolation}
Let \(V\) have orthonormal frame matrices \(A,B\), let \(r\in\{0,1,2,3\}\), and suppose
\[
\operatorname{dist}\!\left(A,\{A_0:\rank A_0\le r\}\right)
\le\delta_A.
\]
Let \(s_0,\ldots,s_r\) lie in a fixed bounded interval and satisfy
\[
\prod_{0\le i<j\le r}|s_i-s_j|\ge\upsilon
\]
(with the empty product equal to one).  If
\[
|P_V(s_j)|\le\epsilon
\qquad(0\le j\le r),
\]
then
\begin{equation}
\label{eq:v046-rank-adapted-interpolation}
\boxed{
\sup_{s\in J}|P_V(s)|
\lesssim_{I,J}
\frac{\epsilon+\delta_A}{\upsilon}.
}
\end{equation}
\end{lemma}

\begin{proof}
Choose \(A_0\) of rank at most \(r\) with \(\|A-A_0\|\le2\delta_A\), and put
\[
Q(s)=\det(B+sA_0).
\]
Multilinearity of the determinant shows that \(Q\) has degree at most \(r\), while its local
Lipschitz dependence on the matrix entries gives
\[
\sup_{s\in I\cup J}|P_V(s)-Q(s)|
\lesssim_{I,J}\delta_A.
\]
Hence \(|Q(s_j)|\lesssim\epsilon+\delta_A\).  Lagrange interpolation of the degree-\(r\)
polynomial \(Q\) at \(s_0,\ldots,s_r\), followed by the boundedness of all numerator factors,
gives
\[
\sup_J|Q|
\lesssim_{I,J}
(\epsilon+\delta_A)\upsilon^{-1}.
\]
Add the preceding approximation error.
\end{proof}

\begin{corollary}[Three-time trapping closes a horizontally planar code cycle]
\label{cor:v046-three-time-planar-cycle}
Use the notation of \cref{lem:v046-four-time-maslov-cycle} and suppose
\[
\mathfrak g_C:=|d_0\wedge d_1\wedge d_2|\ge\gamma,
\qquad
\left|\det(a_1-a_0,a_2-a_1,a_3-a_2)\right|\le\Delta.
\]
If three of the four collision labels, say \(\xi_{j_0},\xi_{j_1},\xi_{j_2}\), satisfy
\[
\prod_{0\le p<q\le2}|\xi_{j_p}-\xi_{j_q}|
\ge\upsilon,
\]
then
\begin{equation}
\label{eq:v046-three-time-planar-cycle}
\boxed{
\sup_{s\in J}|P_{V_C}(s)|
\lesssim
\frac{R+(\Delta/\gamma)^{1/3}}{\upsilon}.
}
\end{equation}
Thus, provided the right side tends to zero, the cycle plane converges to the front-null,
maximally symplectic K\"ahler locus.  If \(\mathfrak g_C<\gamma\), the full secant triple itself is
quantitatively rank deficient.
\end{corollary}

\begin{proof}
Let \(C\) be the coordinate matrix determined by
\[
[d_0\ d_1\ d_2]=T_CC.
\]
Because \(T_C\) is orthonormal, \(|\det C|=\mathfrak g_C\).  Taking horizontal components gives
\[
[a_1-a_0,\ a_2-a_1,\ a_3-a_2]=A_CC,
\]
and hence
\[
|\det A_C|
\le
\Delta/\gamma.
\]
The distance from a \(3\times3\) matrix to the rank-two locus is its smallest singular value, which
is at most \(|\det A_C|^{1/3}\).  Use
\eqref{eq:v046-cycle-polynomial-values} at the three selected labels and apply
\cref{lem:v046-rank-adapted-interpolation} with \(r=2\).
\end{proof}

\begin{lemma}[A shared two-edge Maslov flag fixes the horizontal plane]
\label{lem:v046-shared-two-edge-flag}
Let \(u,v\in\R^3\) be linearly independent, let \(s,t\in\R\), and put
\[
W
=
\Span\{(u,-su),(v,-tv)\}
\subset\R_a^3\oplus\R_b^3.
\]
If a three-plane \(V\supset W\) is front-null, then
\begin{equation}
\label{eq:v046-shared-flag-plane}
\boxed{
V
\subset
E\oplus E,
\qquad
E=\Span\{u,v\}\subset\R^3.
}
\end{equation}
Equivalently, once two horizontally independent contact edges are fixed, every front-null
three-dimensional extension has the same horizontal and vertical polarization plane.
\end{lemma}

\begin{proof}
Choose a third vector \((w,z)\) so that
\[
V=\Span\{(u,-su),(v,-tv),(w,z)\}.
\]
In this, not necessarily orthonormal, frame the Reeb-front determinant is a nonzero constant
multiple of
\[
\det\bigl((r-s)u,(r-t)v,z+rw\bigr)
=
(r-s)(r-t)
\left[
\det(u,v,z)+r\det(u,v,w)
\right].
\]
Front-nullity says that this polynomial vanishes identically.  Hence
\[
\det(u,v,z)=\det(u,v,w)=0,
\]
so \(w,z\in E\).  The first two spanning vectors already belong to \(E\oplus E\), proving
\eqref{eq:v046-shared-flag-plane}.
\end{proof}

\begin{corollary}[Quantitative coherence of anchored cycle planes]
\label{cor:v046-anchored-plane-coherence}
Suppose a three-plane \(V\) contains unit-scale vectors
\[
d_1=(u,\beta_1),
\qquad
d_2=(v,\beta_2)
\]
such that
\[
|u\times v|\ge\kappa,
\qquad
|\beta_1+su|+|\beta_2+tv|\le\varepsilon,
\qquad
\sup_{r\in J}|P_V(r)|\le\varepsilon.
\]
Then
\begin{equation}
\label{eq:v046-anchored-plane-coherence}
\boxed{
\operatorname{dist}_{\mathrm{Gr}}
\left(
V,\,
\{V'\subset E\oplus E:\dim V'=3\}
\right)
\le
C_{I,J}\frac{\varepsilon}{\kappa^2},
\qquad
E=\Span\{u,v\}.
}
\end{equation}
Thus all near-front-null cycle planes sharing the same well-conditioned two-edge anchor are
quantitatively concentrated near one fixed polarized four-space \(E\oplus E\); their horizontal
directions lie near the same affine two-plane.
\end{corollary}

\begin{proof}
Choose a unit vector \(d_3=(w,z)\in V\) orthogonal to \(\Span\{d_1,d_2\}\).  The three vectors
\(d_1,d_2,d_3\) have volume at least \(c\kappa\), since their first two horizontal components have
cross product at least \(\kappa\).  Put
\[
Q(r)
=
\det(\beta_1+ru,\beta_2+rv,z+rw).
\]
Changing from this bounded basis to an orthonormal frame of \(V\) gives
\[
\sup_{r\in J}|Q(r)|\lesssim\varepsilon.
\]
Write \(e_1=\beta_1+su\) and \(e_2=\beta_2+tv\).  Multilinearity and the bounds on all data give
\[
\sup_{r\in J}
\left|
(r-s)(r-t)
\bigl[\det(u,v,z)+r\det(u,v,w)\bigr]
\right|
\lesssim
\varepsilon.
\]
On the finite-dimensional space of polynomials of degree at most three, multiplication of a linear
polynomial by the monic quadratic \((r-s)(r-t)\) has a uniformly bounded inverse on its image when
\(s,t\) range in the fixed bounded interval.   Hence
\[
|\det(u,v,z)|+|\det(u,v,w)|
\lesssim_{I,J}
\varepsilon.
\]
The distance of a vector \(x\in\R^3\) from \(E\) is
\(|\det(u,v,x)|/|u\times v|\), so both \(w\) and \(z\) lie within
\(C_{I,J}\varepsilon/\kappa\) of \(E\).  The first two spanning vectors lie within
\(\varepsilon\) of \(E\oplus E\).  Finally the condition number of
\([d_1\ d_2\ d_3]\) is \(O(\kappa^{-1})\), giving the Grassmannian bound
\(C_{I,J}\varepsilon/\kappa^2\).
\end{proof}

\begin{proposition}[Anchored determinant-small cycles form one geometric packet]
\label{prop:v046-anchored-rank-packet}
Let \(Y=N(p)\cap N(q)\) have cardinality \(T\), assume
\[
|a_p-a_q|\ge\tau_1>0,
\]
and suppose all edges from \(p,q\) to \(Y\) belong to one weight layer with lower weight
\(\lambda R^4\).  Let \(\varnothing\ne\Gamma\subset\binom Y2\), put
\[
Z_\Gamma:=|\Gamma|,
\qquad
m_\Gamma:=\max\{1,Z_\Gamma/T\},
\]
For every \(\{y,y'\}\in\Gamma\), suppose the contact cycle
\[
p,y,q,y',p
\]
has horizontal determinant at most \(\Delta\), and three of its four collision labels have
Vandermonde product at least \(\upsilon\).
For such a pair put
\[
\mathfrak g(y,y')
=
\left|
(z_y-z_p)\wedge(z_q-z_y)\wedge(z_{y'}-z_q)
\right|.
\]

For every \(0<\gamma,\kappa<1\), there are \(y_0\in Y\) and
\(Y_0\subset Y\setminus\{y_0\}\), with
\[
|Y_0|\gtrsim m_\Gamma,
\]
for which at least one of the
following holds:
\begin{enumerate}[label=\textnormal{(\roman*)}]
\item with
\[
F
=
\Span\{z_{y_0}-z_p,z_q-z_{y_0}\}
\subset\R^6,
\]
one has \(\dim F=2\), \(\dim\pi_a(F)\le2\), and every \(y\in Y_0\) belongs to the same
full-secant planar packet
\begin{equation}
\label{eq:v046-anchored-full-rank-collapse}
\operatorname{dist}(z_y,z_p+F)
\lesssim
\gamma/\kappa;
\end{equation}
\item every \(y\in Y_0\) belongs to the horizontal line packet
\begin{equation}
\label{eq:v046-anchored-line-packet}
\operatorname{dist}
\bigl(a_y,\operatorname{aff}\{a_p,a_q\}\bigr)
\lesssim
\kappa/\tau_1;
\end{equation}
\item putting
\[
E
=
\Span\{a_{y_0}-a_p,a_q-a_{y_0}\},
\]
one has
\[
|(a_{y_0}-a_p)\times(a_q-a_p)|\ge\kappa,
\]
and, for every \(y\in Y_0\), the cycle secant plane \(V_{y_0,y}\) satisfies
\begin{equation}
\label{eq:v046-anchored-polarized-packet}
\boxed{
\begin{aligned}
\sup_{s\in J}|P_{V_{y_0,y}}(s)|
&\lesssim
\varepsilon_{\Delta,\gamma,\upsilon},\\
\operatorname{dist}_{\mathrm{Gr}}
\left(
V_{y_0,y},
\{V'\subset E\oplus E:\dim V'=3\}
\right)
&\le
C\varepsilon_{\Delta,\gamma,\upsilon}/\kappa^2,\\
\operatorname{dist}(a_y,a_p+E)
&\lesssim
\varepsilon_{\Delta,\gamma,\upsilon}/\kappa^2,
\end{aligned}
}
\end{equation}
where
\begin{equation}
\label{eq:v046-anchored-packet-error}
\varepsilon_{\Delta,\gamma,\upsilon}
=
R+
\frac{R+(\Delta/\gamma)^{1/3}}{\upsilon}.
\end{equation}
\end{enumerate}
In every case the retained endpoint codes de-multiply to an actual selector-direction subset
\(H_{Y_0}\) of mass
\begin{equation}
\label{eq:v046-anchored-packet-mass}
\boxed{
\sigma(H_{Y_0})
\gtrsim
\lambda R^4m_\Gamma.
}
\end{equation}
In every case this subset lies in the corresponding line or plane packet enlarged by \(O(R)\).
\end{proposition}

\begin{proof}
Suppose first that at least a fixed fraction of the edges of \(\Gamma\) have
\(\mathfrak g(y,y')<\gamma\); denote that subgraph by \(\Gamma_{\rm rk}\).  Put
\[
Y_{\rm full-line}
=
\left\{
y\in Y:
|(z_y-z_p)\wedge(z_q-z_y)|<\kappa
\right\}.
\]
If \(|Y_{\rm full-line}|\gtrsim m_\Gamma\), then, since
\(|z_q-z_p|\ge|a_q-a_p|\ge\tau_1\), all these graph points lie within
\(C\kappa/\tau_1\) of the affine line through \(z_p,z_q\).  Horizontal projection gives (ii).
Otherwise deletion of \(Y_{\rm full-line}\) removes fewer than half the edges of
\(\Gamma_{\rm rk}\).  The remaining graph has a vertex \(y_0\) with \(\gtrsim m_\Gamma\)
neighbors \(Y_0\).  Its two base secants have wedge norm at least \(\kappa\); dividing the
triple-volume inequality \(\mathfrak g(y_0,y)<\gamma\) by that base area gives
\[
\operatorname{dist}(z_y,z_p+F)
\lesssim
\gamma/\kappa,
\qquad
F=\Span\{z_{y_0}-z_p,z_q-z_{y_0}\}.
\]
This is (i).

Otherwise a subgraph \(\Gamma_{\rm fr}\subset\Gamma\), with
\(|\Gamma_{\rm fr}|\gtrsim Z_\Gamma\), has full secant volume at least \(\gamma\).  Put
\[
Y_{\rm line}
=
\left\{
y\in Y:
|(a_y-a_p)\times(a_q-a_p)|<\kappa
\right\}.
\]
If \(|Y_{\rm line}|\gtrsim m_\Gamma\), the elementary point-to-line formula and
\(|a_q-a_p|\ge\tau_1\) give (ii).  Otherwise, choosing the implicit small constant in this last
alternative absolutely small, deletion of \(Y_{\rm line}\) removes fewer than half the edges of
\(\Gamma_{\rm fr}\), because it removes at most \(|Y_{\rm line}|T\) edges.  The remaining graph has
a vertex \(y_0\notin Y_{\rm line}\) with \(\gtrsim m_\Gamma\) neighbors (using one retained edge
when \(Z_\Gamma<T\)); call that neighbor set \(Y_0\).

For every \(y\in Y_0\), the cycle \(p,y_0,q,y,p\) has full secant volume at least \(\gamma\),
horizontal determinant at most \(\Delta\), and three collision labels with Vandermonde at least
\(\upsilon\).  Therefore \cref{cor:v046-three-time-planar-cycle} gives the first line of
\eqref{eq:v046-anchored-polarized-packet}.  All these cycle planes contain the two fixed contact
secants
\[
z_{y_0}-z_p,
\qquad
z_q-z_{y_0},
\]
whose horizontal parts have cross product at least \(\kappa\); their contact errors are \(O(R)\)
by \eqref{eq:v046-code-cycle-contact}.  Apply
\cref{cor:v046-anchored-plane-coherence}, absorbing the bounded normalization constants into its
modulus, to obtain the second line.  Since \(z_y-z_q\in V_{y_0,y}\), horizontal projection gives
the third line.  This proves (iii).

It remains to verify that the packet has physical mass.  Let \(\mathscr X_0\) be the set of selector
graph product cells underlying the endpoint codes in \(Y_0\).  A fixed graph cell supports only
\(O(R^{-1})\) level codes, hence
\[
|\mathscr X_0|\gtrsim R|Y_0|\gtrsim Rm_\Gamma.
\]
For each \(\chi\in\mathscr X_0\), choose one code \(y\) over \(\chi\) and one of the weighted edges
\(py\) or \(qy\).  Exactly as in
\cref{prop:v046-double-cell-demultiplication}, its weight lower bound \(\lambda R^4\) forces at
least \(c\lambda R^3\) selector-direction mass inside \(\chi\).  Distinct half-open graph cells are
disjoint on \(G_b\), so summation proves \eqref{eq:v046-anchored-packet-mass}.  The graph cell
diameter accounts for the final \(O(R)\) enlargement.
\end{proof}

\begin{proposition}[Exact graph-cell reuse ledger for anchored Maslov packets]
\label{prop:v046-anchor-cell-reuse}
Let a family of anchored packets in one \(\lambda\)-weight layer contain a total of \(S\) retained
endpoint-code occurrences.  For a packet \(Q\), let \(\mathscr X_Q\) be the distinct selector graph
cells supporting its retained endpoint codes.  Define
\[
r(\chi)
=
\#\{Q:\chi\in\mathscr X_Q\},
\qquad
\mathcal R_{\rm anc}
=
\sum_\chi r(\chi)^2.
\]
Then
\begin{equation}
\label{eq:v046-anchor-reuse-L1}
\sum_\chi r(\chi)
\gtrsim
RS,
\end{equation}
and the union \(H_{\rm anc}\) of the corresponding actual selector-direction subsets obeys
\begin{equation}
\label{eq:v046-anchor-reuse-L2}
\boxed{
\sigma(H_{\rm anc})
\gtrsim
\lambda R^5
\frac{S^2}{\mathcal R_{\rm anc}}.
}
\end{equation}
More explicitly, for every \(K\ge1\), at least one of the following holds:
\begin{enumerate}[label=\textnormal{(\roman*)}]
\item the cells with \(r(\chi)\le K\) carry a selector-direction union of mass
\begin{equation}
\label{eq:v046-anchor-low-reuse-mass}
\boxed{
\sigma(H_{\rm anc}^{\rm low})
\gtrsim
\lambda R^4S/K;
}
\end{equation}
\item some endpoint code belongs to at least \(cRK\) distinct anchored packets, and is incident to
a code star which produces a genuine physical star of degree
\begin{equation}
\label{eq:v046-anchor-high-reuse-star}
\boxed{
D_{\rm phys}
\gtrsim
R^{3/2}K^{1/2}.
}
\end{equation}
\end{enumerate}
Thus packet reuse is not a loss of information: its high part returns to a higher-order common
contact star at one endpoint code.
\end{proposition}

\begin{proof}
A selector graph cell supports only \(O(R^{-1})\) level codes.  Thinning separately inside each
packet therefore gives
\[
|\mathscr X_Q|\gtrsim R|Y_Q|.
\]
Summation proves \eqref{eq:v046-anchor-reuse-L1}.  Each occupied graph cell contains, by the same
weighted double-cell de-multiplication as in
\cref{prop:v046-double-cell-demultiplication}, an actual direction subset of mass at least
\(c\lambda R^3\).  Distinct cells are disjoint on the selector graph.  Cauchy--Schwarz gives
\[
\#\{\chi:r(\chi)>0\}
\ge
\frac{(\sum_\chi r(\chi))^2}{\sum_\chi r(\chi)^2},
\]
which proves \eqref{eq:v046-anchor-reuse-L2}.

If at least half of \(\sum r(\chi)\) lies on \(r(\chi)\le K\), those cells number at least
\(cRS/K\), and the preceding \(c\lambda R^3\) mass per cell proves
\eqref{eq:v046-anchor-low-reuse-mass}.  Otherwise some cell has \(r(\chi)>K\).  At most
\(O(R^{-1})\) endpoint codes lie over that graph cell, so one endpoint code \(y\) occurs in at least
\(cRK\) of the anchored packets.  Their anchors are distinct unordered pairs \(\{p,q\}\) with
\(p,q\in N(y)\); hence
\[
\binom{\deg(y)}2\gtrsim RK,
\qquad
\deg(y)\gtrsim(RK)^{1/2}.
\]
Apply \eqref{eq:v043-code-degree-to-star} to obtain
\eqref{eq:v046-anchor-high-reuse-star}.
\end{proof}

\subsection{Front seeds and Reeb-time completion}

\begin{lemma}[Four-cluster Reeb-front seed]
\label{lem:v046-four-cluster-seed}
Fix a compact interval \(J\) with nonempty interior and a bounded set
\(E\subset\R_a^3\oplus\R_b^3\).  For four points
\[
x_j=(a_j,b_j)\in E,
\qquad 0\le j\le3,
\]
put
\[
A_\Delta=[a_1-a_0,\ a_2-a_0,\ a_3-a_0],
\qquad
B_\Delta=[b_1-b_0,\ b_2-b_0,\ b_3-b_0].
\]
If \(|\det A_\Delta|\ge\kappa>0\), then
\begin{equation}
\label{eq:v046-four-cluster-polynomial}
\boxed{
\sup_{s\in J}
\left|
\det
\bigl[
\Pi_s(x_1)-\Pi_s(x_0),
\Pi_s(x_2)-\Pi_s(x_0),
\Pi_s(x_3)-\Pi_s(x_0)
\bigr]
\right|
\ge c_J\kappa.
}
\end{equation}
There are constants \(c_{J}>0\) and \(c'_{J}>0\), independent of \(E\) and \(\kappa\), such that
\begin{equation}
\label{eq:v047-uniform-good-time}
\boxed{
\left|
\left\{
s\in J:
\left|
\det
\bigl[
\Pi_s(x_1)-\Pi_s(x_0),
\Pi_s(x_2)-\Pi_s(x_0),
\Pi_s(x_3)-\Pi_s(x_0)
\bigr]
\right|
\ge c_J\kappa
\right\}
\right|
\ge c'_J.
}
\end{equation}

If a probability measure \(\nu\) gives mass at least \(m\) to each of four sufficiently small
balls of radius \(r\le c_{J,E}\kappa\) about the \(x_j\), then for every such good \(s\), the
projected measure \(\Pi_s{}_{\#}\nu\) gives mass at least \(m\) to four disjoint output balls whose
centers form a tetrahedron with smallest altitude at least \(c_{J,E}\kappa\).  Thus every
horizontally nondegenerate four-cluster produces a definite coarse projection-entropy gain on a
longitudinal set of fixed positive measure.
\end{lemma}

\begin{proof}
The determinant in \eqref{eq:v046-four-cluster-polynomial} is
\[
P_\Delta(s)=\det(B_\Delta+sA_\Delta),
\]
a cubic polynomial whose leading coefficient is \(\det A_\Delta\).  Equivalence of polynomial
norms on \(J\) gives
\[
\|P_\Delta\|_{L^\infty(J)}\ge c_J\kappa.
\]
Apply the degree-three Remez inequality.  If
\[
\bigl|\{s\in J:|P_\Delta(s)|<c_J^*\kappa\}\bigr|>|J|-c'_J
\]
for a sufficiently small fixed \(c_J^*>0\), Remez would bound the supremum on all of \(J\) by
\(C_Jc_J^*\kappa<c_J\kappa\), a contradiction.  Choose \(c'_J>0\) first, then choose
\(c_J^*\) accordingly and rename it as the constant in
\eqref{eq:v047-uniform-good-time}.  This is stronger than the derivative estimate: no factor
\(\kappa\) is lost in the length of the good-time set.

On that set the three projected difference vectors have determinant at least \(c_J\kappa\) and,
because \(E\) and \(J\) are bounded, norm at most \(C_{J,E}\).  The smallest altitude of the
tetrahedron is therefore at least \(c_{J,E}\kappa\).  Input balls of radius
\(r\le c_{J,E}\kappa\) project into disjoint output balls about the four vertices and retain their
input mass.
\end{proof}

\begin{lemma}[Sticky graph mass contains a quantitative Reeb-front simplex]
\label{lem:v046-sticky-four-cluster-extraction}
Assume the quantitative packing-piece estimate
\[
N_r(G_b)\le C_\varepsilon r^{-3-\varepsilon}.
\]
Let \(H\subset U\) be measurable with \(\sigma(H)=m>0\), and let
\[
0<r\le c_Um.
\]
Then there are graph points
\[
z_j=(a_j,b(a_j)),
\qquad
0\le j\le3,
\]
and graph product cells \(\chi_j\) of diameter \(O(r)\), containing \(z_j\), such that
\begin{equation}
\label{eq:v046-sticky-four-cluster-input}
\boxed{
\left|
\det(a_1-a_0,a_2-a_0,a_3-a_0)
\right|
\gtrsim_U
m,
\qquad
\mu_H(\chi_j)
\gtrsim_{\varepsilon,U}
m r^{3+\varepsilon},
}
\end{equation}
where \(\mu_H=(\id,b)_\#(\sigma|_H)\).  Consequently there is a set
\(G_H\subset J\) with \(|G_H|\gtrsim_J1\) such that, for every \(s\in G_H\), the projected measure
\(\Pi_s{}_\#\mu_H\) has four \(O(r)\)-clusters, each of mass
\(\gtrsim_{\varepsilon,U}mr^{3+\varepsilon}\), whose centers form a tetrahedron of determinant
\(\gtrsim_Um\).
\end{lemma}

\begin{proof}
Use the half-open product grid of side \(r\) from \cref{sec:v034-entropy-pruning}.  There are at
most \(C_\varepsilon r^{-3-\varepsilon}\) cells meeting \(G_b\).  Call a cell heavy when
\[
\mu_H(\chi)\ge c_\varepsilon m r^{3+\varepsilon},
\]
where \(c_\varepsilon\) is chosen so that the union of the nonheavy cells has
\(\mu_H\)-mass at most \(m/2\).  Let \(H_{\rm hv}\) be the horizontal preimage of the heavy cells;
then
\[
\sigma(H_{\rm hv})\ge m/2.
\]

We use the elementary maximal-simplex bound: every bounded measurable
\(F\subset\R^3\) of Lebesgue measure \(v\) contains four points whose tetrahedron determinant is at
least \(c_Uv\).  Indeed choose successively a diameter, a point of maximal distance from its line,
and a point of maximal distance from the resulting plane.  Up to a fixed enlargement, the set is
contained in a box whose three side lengths are those successive widths; its volume is bounded by
their product, which is comparable to the selected determinant.

Apply this bound to \(H_{\rm hv}\), recalling that \(\sigma\) is normalized Lebesgue measure on
\(U\).  It gives four points with horizontal determinant at least \(c_Um\).  Moving each point by
\(O(r)\) changes that determinant by \(O_U(r)\), so the choice \(r\le c_Um\) keeps the lower bound
after replacing the points by representatives of their heavy graph cells.  The four cells are
necessarily distinct.  Their defining heaviness gives the second bound in
\eqref{eq:v046-sticky-four-cluster-input}.

Apply \cref{lem:v046-four-cluster-seed} with horizontal tetrahedron constant \(c_Um\).  Its
good-parameter set has measure at least \(c_J'>0\); product cells of diameter \(O(r)\) are contained
in \(O(r)\)-balls, and every one retains its \(\mu_H\)-mass under projection.  This proves the final
assertion.
\end{proof}

\begin{corollary}[Uniform longitudinal persistence of shrinking simplices]
\label{cor:v046-longitudinal-limsup}
Let \(R_n\downarrow0\).  At the \(n\)-th scale suppose
\cref{lem:v046-four-cluster-seed} applies with an arbitrary \(\kappa_n>0\), and choose the input
cluster radius at most \(c_{J,E}\kappa_n\).  If \(G_n\subset J\) is the corresponding set of good
longitudinal parameters, then
\[
\left|\limsup_{n\to\infty}G_n\right|
\ge c'_J.
\]
Thus a fixed positive-measure set of Reeb parameters sees quantitatively nondegenerate projection
seeds on infinitely many scales even when \(\kappa_n\downarrow0\).
\end{corollary}

\begin{proof}
The preceding lemma gives \(|G_n|\ge c'_J\).  Apply reverse Fatou to the bounded functions
\(\mathbf 1_{G_n}\):
\[
\int_J\limsup_n\mathbf 1_{G_n}
\ge
\limsup_n\int_J\mathbf 1_{G_n}.
\]
\end{proof}

\begin{lemma}[Reeb-time integration lifts transverse Frostman measures]
\label{lem:v047-slice-frostman-lift}
Let
\[
K_s:=\{\Pi_s(a,b(a)):a\in U\},
\qquad
K:=\{(y,s):s\in J,\ y\in K_s\}.
\]
Suppose \(G\subset J\) has positive Lebesgue measure and, for almost every \(s\in G\), there is a
weakly measurable probability kernel \(s\mapsto\nu_s\) carried by \(K_s\) such that, uniformly in
\(s\),
\begin{equation}
\label{eq:v047-slice-frostman}
\nu_s(B(y,\rho))
\le C_\varepsilon\rho^{3-\varepsilon}
\qquad
(y\in\R^3,\ 0<\rho<1).
\end{equation}
Then the measure
\begin{equation}
\label{eq:v047-integrated-front-measure}
\nu(A)
:=
\frac1{|G|}
\int_G
\nu_s\{y:(y,s)\in A\}\,ds
\end{equation}
is carried by \(K\) and satisfies
\begin{equation}
\label{eq:v047-integrated-front-frostman}
\boxed{
\nu(B((y_0,s_0),\rho))
\le
\frac{2C_\varepsilon}{|G|}\rho^{4-\varepsilon}.
}
\end{equation}
Consequently \(\dim_HK\ge4-\varepsilon\).
\end{lemma}

\begin{proof}
Only parameters \(s\in G\cap(s_0-\rho,s_0+\rho)\) contribute to the ball on the left of
\eqref{eq:v047-integrated-front-frostman}.  For each such parameter its transverse section lies in
\(B(y_0,\rho)\), so \eqref{eq:v047-slice-frostman} bounds the conditional mass by
\(C_\varepsilon\rho^{3-\varepsilon}\).  Integration over an interval of length at most \(2\rho\)
proves \eqref{eq:v047-integrated-front-frostman}.  The mass distribution principle gives the final
claim.
\end{proof}

\begin{corollary}[Exact remaining contact--symplectic front endpoint]
\label{cor:v047-exact-spatial-endpoint}
In both the fixed-angle low-reuse branch and the angular Maslov branch, Remez persistence supplies a
set \(G_\infty\subset J\) of fixed positive measure on which front seeds occur at arbitrarily small
scales.  To obtain \(\dim_HK=4\), it is therefore enough to upgrade the persistent cluster measures,
on one positive-measure subset of \(G_\infty\), to the transverse kernels
\eqref{eq:v047-slice-frostman} for every \(\varepsilon>0\).  Once that purely spatial upgrade is
proved, \cref{lem:v047-slice-frostman-lift} performs the longitudinal and Hausdorff passage with no
further loss.
\end{corollary}

\begin{definition}[Weighted Maslov residual content]
\label{def:v048-maslov-residual-content}
For \(0<\rho<1\), define
\begin{equation}
\label{eq:v048-maslov-residual-content}
\boxed{
\mathfrak Z_\rho^{J^+}(\mu)
:=
\iint_{a\ne a'}
\frac{
\mathbf 1_{\{s_*(a,a')\in J^+\}}
\mathbf 1_{\{|r(a,a')|\le\rho\}}
}{|a-a'|}
\,d\sigma(a)d\sigma(a').
}
\end{equation}
\end{definition}

\begin{theorem}[Codimension-two Maslov residual criterion]
\label{thm:v048-residual-frostman-criterion}
Assume that for some \(0<\eta<2\),
\begin{equation}
\label{eq:v048-residual-codim-two}
\boxed{
\mathfrak Z_\rho^{J^+}(\mu)
\le
C_\eta\rho^{2-\eta}
\qquad(0<\rho<1).
}
\end{equation}
Then, for every
\[
1<t<3-\eta,
\]
the projected measures \(\nu_s=(\Pi_s)_\#\mu\), with
\(I_t(\nu_s):=\iint|y-y'|^{-t}\,d\nu_s(y)d\nu_s(y')\), satisfy
\begin{equation}
\label{eq:v048-average-projection-energy}
\boxed{
\int_J I_t(\nu_s)\,ds
=
\int_J\iint
|\Pi_s(x)-\Pi_s(x')|^{-t}
\,d\mu(x)d\mu(x')\,ds
<\infty.
}
\end{equation}
Consequently
\begin{equation}
\label{eq:v048-residual-dimension-output}
\dim_HK\ge4-\eta.
\end{equation}
If \eqref{eq:v048-residual-codim-two} holds with exponent \(2-\eta\) for every \(\eta>0\), then
\(\dim_HK=4\).
\end{theorem}

\begin{proof}
The Pythagorean collision identity gives
\begin{equation}
\label{eq:v048-pythagorean-energy}
|\beta+s\alpha|^2
=
|r|^2+|\alpha|^2|s-s_*|^2.
\end{equation}
Let \(d_0=\dist(J,\R\setminus J^+)>0\).  If \(s_*\notin J^+\), then for \(s\in J\),
\[
|\beta+s\alpha|
\ge
d_0|\alpha|,
\]
so this part of the averaged energy is bounded by a constant multiple of
\[
\iint|a-a'|^{-t}\,d\sigma(a)d\sigma(a'),
\]
which is finite for \(t<3\) because the direction marginal is bounded Lebesgue measure in
three dimensions.

If \(s_*\in J^+\), extending the \(s\)-integral to \(\R\) in
\eqref{eq:v048-pythagorean-energy} gives, for \(r\ne0\),
\begin{equation}
\label{eq:v048-one-pair-energy}
\int_J|\beta+s\alpha|^{-t}\,ds
\le
C_t|\alpha|^{-1}|r|^{1-t}.
\end{equation}
The hypothesis \eqref{eq:v048-residual-codim-two} forces the weighted pair measure of
\(\{s_*\in J^+,r=0\}\) to vanish.  On \(0<|r|\le1\), dyadic summation gives
\begin{align*}
&\iint_{\substack{s_*\in J^+\\0<|r|\le1}}
|\alpha|^{-1}|r|^{1-t}\,d\sigma d\sigma'
\\
&\qquad\lesssim_t
\sum_{k\ge0}
2^{(k+1)(t-1)}
\mathfrak Z_{2^{-k}}^{J^+}(\mu)
\\
&\qquad\lesssim_{t,\eta}
\sum_{k\ge0}
2^{-k(3-\eta-t)}
<\infty.
\end{align*}
The region \(|r|>1\) is harmless.  Tonelli's theorem and
\eqref{eq:v048-one-pair-energy} prove
\eqref{eq:v048-average-projection-energy}.

Hence \(I_t(\nu_s)<\infty\) for almost every \(s\in J\), so
\(\dim_HK_s\ge t\) for almost every such \(s\).  The standard slicing theorem for the analytic
Reeb front gives \(\dim_HK\ge t+1\).  Letting \(t\uparrow3-\eta\) proves
\eqref{eq:v048-residual-dimension-output}; then let \(\eta\downarrow0\) for the final assertion.
\end{proof}

\clearpage
\section{The codimension-two residual and polarized routing}
The energy, slicing, and Frostman principles are used in their standard geometric-measure form;
see \cite{Federer1969,Mattila1995}.


\subsection{Residual criteria and shell return}

\begin{proposition}[Collision energy is the discretized residual content]
\label{prop:v048-energy-residual-equivalence}
Let \(\mathfrak E_{\rho,\ge\tau}^{L}(\mu)\) denote
\cref{def:v006-energy} with the pair integral restricted to \(|a-a'|\ge\tau\) and time interval
\(L\).  Likewise append \(\ge\tau\) to \(\mathfrak Z\) when its pair integral is so restricted.
There is \(c=c(J^-,J,J^+)>0\) such that, whenever \(0<\rho\le c\tau\),
\begin{equation}
\label{eq:v048-energy-residual-equivalence}
\boxed{
c\rho^{-2}
\mathfrak Z_{\rho/2,\ge\tau}^{J^-}(\mu)
\le
\mathfrak E_{\rho,\ge\tau}^{J}(\mu)
\le
C\rho^{-2}
\mathfrak Z_{\rho,\ge\tau}^{J^+}(\mu).
}
\end{equation}
Thus the codimension-two target
\(\mathfrak Z_\rho\lesssim\rho^{2-o(1)}\) is exactly the subpower collision-energy target
\(\mathfrak E_\rho\lesssim\rho^{-o(1)}\) on every fixed-angle shell.
\end{proposition}

\begin{proof}
If \(|r|\le\rho/2\), \(s_*\in J^-\), and \(|\alpha|\ge\tau\), then
\(\rho\le c\tau\) and the separation of \(J^-\) from \(\partial J\) imply that the collision
window in \cref{prop:v006-window-formula} lies inside \(J\) and has length at least
\(c\rho/|\alpha|\).  Multiplication by the normalization \(\rho^{-3}\) proves the left inequality.

Conversely, a nonempty \(\rho\)-collision window forces \(|r|\le\rho\) and places \(s_*\) within
\(\rho/|\alpha|\le\rho/\tau\) of \(J\), hence in \(J^+\).  The universal window bound
\cref{lem:v007-window-upper} gives length at most \(2\rho/|\alpha|\), proving the right inequality.
\end{proof}

\begin{lemma}[Exact angular-shell ledger for the residual content]
\label{lem:v048-residual-shell-ledger}
For a dyadic shell write
\[
\mathfrak Z_{\rho,\tau}^{J^+}
:=
\iint_{\tau\le|a-a'|<2\tau}
\frac{
\mathbf 1_{\{s_*\in J^+\}}
\mathbf 1_{\{|r|\le\rho\}}
}{|a-a'|}
\,d\sigma(a)d\sigma(a').
\]
For every \(\rho\le\tau_0\le1\),
\begin{equation}
\label{eq:v048-residual-shell-ledger}
\boxed{
\mathfrak Z_\rho^{J^+}
\lesssim_U
\tau_0^2
+
\sum_{\substack{\tau\ \mathrm{dyadic}\\\tau_0\le\tau\lesssim1}}
\mathfrak Z_{\rho,\tau}^{J^+}.
}
\end{equation}
Consequently, for a fixed \(\eta>0\), choose
\[
\tau_0=\rho^{1-\eta/4}.
\]
Then the ultra-near term is \(O(\rho^{2-\eta/2})\), and
\eqref{eq:v048-final-residual-endpoint} follows if every remaining shell obeys
\begin{equation}
\label{eq:v048-dimensionless-shell-target}
\boxed{
\kappa_{\rho,\tau}
:=
\frac{\mathfrak Z_{\rho,\tau}^{J^+}}{\tau^2}
\lesssim_\eta
\frac{(\rho/\tau)^2\rho^{-\eta}}{\log(2/\rho)}.
}
\end{equation}
\end{lemma}

\begin{proof}
Discarding both indicators in the part \(|a-a'|<2\tau_0\) and applying the Newton estimate
\eqref{eq:v046-newton-near-diagonal} bounds that part by \(C_U\tau_0^2\).  The complementary pairs
belong to the displayed logarithmic family of dyadic shells, proving
\eqref{eq:v048-residual-shell-ledger}.  With the stated choice of \(\tau_0\), its first term is
smaller than \(\rho^{2-\eta}\).  Summing
\eqref{eq:v048-dimensionless-shell-target} over \(O(\log(2/\rho))\) shells proves the final claim.
\end{proof}

\begin{proposition}[Every sub-four-dimensional front forces a power-excess Maslov shell]
\label{prop:v048-counterexample-forces-shell}
Assume
\[
\dim_HK\le4-\delta
\]
for some \(\delta>0\), and fix \(0<\eta<\delta\).  Then there are scales
\(\rho_n\downarrow0\), dyadic angular radii \(\tau_n\), and numbers \(A_n\to\infty\) such that
\begin{equation}
\label{eq:v048-counterexample-shell-range}
\rho_n^{1-\eta/4}
\le
\tau_n
\lesssim1,
\qquad
\epsilon_n:=\rho_n/\tau_n\longrightarrow0,
\end{equation}
and
\begin{equation}
\label{eq:v048-counterexample-shell-excess}
\boxed{
\kappa_{\rho_n,\tau_n}
\ge
A_n
\frac{\epsilon_n^2\rho_n^{-\eta}}{\log(2/\rho_n)}.
}
\end{equation}
In particular a Hausdorff counterexample cannot hide in the ultra-near diagonal: it creates a
genuine Darboux-separated, power-excess Lagrangian-pencil collision shell.
\end{proposition}

\begin{proof}
If \(\mathfrak Z_\rho^{J^+}\lesssim\rho^{2-\eta}\) held at all sufficiently small scales, then
\cref{thm:v048-residual-frostman-criterion} would give
\[
\dim_HK\ge4-\eta>4-\delta,
\]
a contradiction.  Hence, after passing to a sequence, there are \(A_n\to\infty\) for which
\[
\mathfrak Z_{\rho_n}^{J^+}
\ge
A_n\rho_n^{2-\eta}.
\]
Apply \cref{lem:v048-residual-shell-ledger} with
\(\tau_{0,n}=\rho_n^{1-\eta/4}\).  Its ultra-near contribution is
\[
O(\rho_n^{2-\eta/2})
=
o(\rho_n^{2-\eta}),
\]
so, after decreasing \(A_n\) by an absolute factor, one of the
\(O(\log(2/\rho_n))\) remaining shells satisfies
\[
\mathfrak Z_{\rho_n,\tau_n}^{J^+}
\ge
A_n\frac{\rho_n^{2-\eta}}{\log(2/\rho_n)}.
\]
Divide by \(\tau_n^2\) and use
\(\epsilon_n=\rho_n/\tau_n\) to obtain
\eqref{eq:v048-counterexample-shell-excess}.  Finally
\(\tau_n\ge\rho_n^{1-\eta/4}\) gives
\(\epsilon_n\le\rho_n^{\eta/4}\to0\), proving
\eqref{eq:v048-counterexample-shell-range}.
\end{proof}

\begin{corollary}[Counterexample-driven return and the exact exponent gap]
\label{cor:v048-counterexample-trichotomy}
For every hypothetical Hausdorff deficit \(\delta>0\), the shells in
\cref{prop:v048-counterexample-forces-shell} may be passed through the bounded density and weight
decompositions of
\cref{prop:v046-critical-renormalization,lem:v046-weighted-code-edge-layer} and then tested by the
quantitative Maslov alternatives of the preceding stage.  The geometric outputs are:
\begin{enumerate}[label=\textnormal{(\roman*)}]
\item a threshold approximate Lagrangian bush with vanishing physical radius;
\item a further angular Darboux collapse with explicit scale separation;
\item a persistent nonplanar Reeb-front seed carrying the residual excess
\eqref{eq:v048-counterexample-shell-excess}.
\end{enumerate}
There is, however, one genuinely numerical alternative which must not be hidden.  If every normalized
patch falls into the present one-patch saving rather than a geometric output, then the shell may
satisfy simultaneously
\begin{equation}
\label{eq:v048-exact-exponent-gap}
\boxed{
A_n\frac{\epsilon_n^2\rho_n^{-\eta}}{\log(2/\rho_n)}
\lesssim
\kappa_{\rho_n,\tau_n}
\lesssim
\epsilon_n^{1/4}.
}
\end{equation}
The interval in \eqref{eq:v048-exact-exponent-gap} can be nonempty.  Therefore a residual-target
violation is not automatically above the fourth-root threshold.  A surviving counterexample must
either realize the persistent seed output (iii), or remain in precisely this diffuse exponent-gap
window.  Closing the proof requires a scale-persistent improvement from the upper exponent
\(1/4\) to the normal-slice exponent \(2-o(1)\); merely reapplying the same four-cycle estimate does
not supply it.
\end{corollary}

\begin{lemma}[The residual is exactly a two-coordinate symplectic normal map]
\label{lem:v049-two-normal-symplectic-coordinates}
Fix a dyadic shell \(\tau\le|\alpha|<2\tau\), where
\(\alpha=a-a'\), \(\beta=b(a)-b(a')\), and cover the sphere of directions
\(\widehat\alpha=\alpha/|\alpha|\) by finitely many conic charts.  On a chart choose a smooth
orthonormal frame \(e_1(\widehat\alpha),e_2(\widehat\alpha)\) of
\(\alpha^\perp\), put
\[
d=(\alpha,\beta),
\qquad
n_j=(e_j,0)\in\mathbb R_a^3\oplus\mathbb R_b^3,
\]
and define
\begin{equation}
\label{eq:v049-symplectic-normal-map}
\Phi_{\tau,j}(a,a')
:=
\tau^{-1}\omega(d,n_j),
\qquad
\Phi_\tau=(\Phi_{\tau,1},\Phi_{\tau,2}).
\end{equation}
Then, for the Maslov closest-time residual
\[
s_*=-\frac{\alpha\cdot\beta}{{|\alpha|}^2},
\qquad
r=\beta+s_*\alpha\in\alpha^\perp,
\]
one has the exact identity
\begin{equation}
\label{eq:v049-two-normal-identity}
\boxed{
\Phi_{\tau,j}=-\tau^{-1}e_j\cdot r,
\qquad
|\Phi_\tau|=\frac{|r|}{\tau}.
}
\end{equation}
If \((\chi_\ell)_\ell\) is a smooth partition of unity subordinate to the conic charts and
\begin{equation}
\label{eq:v049-normalized-pair-measure}
d\Lambda_{\tau,\ell}
:=
\tau^{-2}\chi_\ell(\widehat\alpha)
\mathbf 1_{\{\tau\le|\alpha|<2\tau\}}
\frac{d\sigma(a)d\sigma(a')}{|\alpha|},
\end{equation}
then \(\Lambda_{\tau,\ell}(\mathbb R^6)\lesssim_U1\) and, with
\(\epsilon=\rho/\tau\),
\begin{equation}
\label{eq:v049-kappa-as-symplectic-small-ball}
\boxed{
\kappa_{\rho,\tau}
=
\sum_\ell
\Lambda_{\tau,\ell}
\left\{
s_*\in J^+,\ |\Phi_\tau|\le\epsilon
\right\}.
}
\end{equation}
Thus the missing quadratic exponent is literally the two-dimensional small-ball exponent of the
symplectic normal map \(\Phi_\tau\); it is not an exponent imported from Euclidean endpoint
codes.
\end{lemma}

\begin{proof}
For the standard symplectic form
\[
\omega((x,y),(x',y'))=x\cdot y'-x'\cdot y
\]
one has
\[
\omega(d,n_j)=-e_j\cdot\beta=-e_j\cdot r,
\]
because \(e_j\perp\alpha\).  Since \(e_1,e_2\) are an orthonormal basis of
\(\alpha^\perp\), this proves \eqref{eq:v049-two-normal-identity}.  The Newton bound gives
\[
\iint_{\tau\le|a-a'|<2\tau}
\frac{d\sigma(a)d\sigma(a')}{|a-a'|}
\lesssim_U
\tau^{-1}\tau^3
=
\tau^2,
\]
which proves the uniform mass bound in \eqref{eq:v049-normalized-pair-measure}.  Finally sum the
partition of unity and use \eqref{eq:v049-two-normal-identity} in the definition
\(\kappa_{\rho,\tau}=\tau^{-2}\mathfrak Z_{\rho,\tau}^{J^+}\).
\end{proof}

\begin{proposition}[Conditional two-normal estimate closes one shell]
\label{prop:v049-conditional-two-normal-closure}
For a conic chart let
\[
\lambda_{\tau,\ell}
:=
\mathbf 1_{\{s_*\in J^+\}}\Lambda_{\tau,\ell},
\qquad
\Phi_\tau=(\phi_1,\phi_2).
\]
Suppose the first marginal has density at most \(L_1\),
\[
(\phi_1)_\#\lambda_{\tau,\ell}\le L_1\,du,
\]
and, in the disintegration
\[
\lambda_{\tau,\ell}
=
\int \lambda_u\,d((\phi_1)_\#\lambda_{\tau,\ell})(u),
\]
the conditional second marginals satisfy
\[
(\phi_2)_\#\lambda_u\le L_2\,dv
\]
for almost every \(u\).  Then
\begin{equation}
\label{eq:v049-conditional-two-normal-bound}
\boxed{
\lambda_{\tau,\ell}\{|\Phi_\tau|\le\epsilon\}
\lesssim
L_1L_2\epsilon^2.
}
\end{equation}
Consequently, if on every surviving chart
\begin{equation}
\label{eq:v049-two-normal-density-target}
L_1L_2
\lesssim_\eta
\frac{\rho^{-\eta}}{\log(2/\rho)},
\end{equation}
then \eqref{eq:v048-dimensionless-shell-target} holds.
\end{proposition}

\begin{proof}
The Euclidean ball \(\{|\Phi_\tau|\le\epsilon\}\) is contained in
\(\{|\phi_1|\le\epsilon,\ |\phi_2|\le\epsilon\}\).  Disintegrating first in \(\phi_1\), the
conditional estimate contributes at most \(2L_2\epsilon\); integrating over
\([-\epsilon,\epsilon]\) contributes at most \(2L_1\epsilon\).  This proves
\eqref{eq:v049-conditional-two-normal-bound}.  Sum over the fixed finite chart family and apply
\eqref{eq:v049-kappa-as-symplectic-small-ball}.
\end{proof}

\subsection{Anchored rank packets}

\begin{corollary}[Every residual excess creates a high-density symplectic normal flag]
\label{cor:v049-residual-excess-normal-flag}
Let
\[
\kappa_{\ell}
:=
\lambda_{\tau,\ell}\{|\Phi_\tau|\le\epsilon\},
\qquad
I_\epsilon=[-\epsilon,\epsilon],
\]
and use the disintegration of
\cref{prop:v049-conditional-two-normal-closure}.  Define the two successive normalized
concentrations
\begin{equation}
\label{eq:v049-two-stage-normal-densities}
D_1
:=
\frac{(\phi_1)_\#\lambda_{\tau,\ell}(I_\epsilon)}{2\epsilon},
\qquad
D_2
:=
\operatorname*{ess\,sup}_{u\in I_\epsilon}
\frac{(\phi_2)_\#\lambda_u(I_\epsilon)}{2\epsilon}.
\end{equation}
Then
\begin{equation}
\label{eq:v049-normal-flag-density-product}
\boxed{
D_1D_2
\ge
\frac{\kappa_\ell}{4\epsilon^2}.
}
\end{equation}
In particular, one chart of every shell in
\cref{prop:v048-counterexample-forces-shell} satisfies
\begin{equation}
\label{eq:v049-counterexample-normal-flag}
\boxed{
D_1D_2
\gtrsim
\frac{A_n\rho_n^{-\eta}}{\log(2/\rho_n)},
\qquad
\max\{D_1,D_2\}
\gtrsim
\left(
\frac{A_n\rho_n^{-\eta}}{\log(2/\rho_n)}
\right)^{1/2}.
}
\end{equation}
Thus the exponent-gap branch is not diffuse in both symplectic normal directions: it carries a
quantitatively divergent one-dimensional normal concentration after at most one conditioning.
\end{corollary}

\begin{proof}
The ball \(\{|\Phi_\tau|\le\epsilon\}\) lies inside the rectangle
\(\{\phi_1\in I_\epsilon,\phi_2\in I_\epsilon\}\).  Therefore
\[
\kappa_\ell
\le
\int_{I_\epsilon}
(\phi_2)_\#\lambda_u(I_\epsilon)
\,d((\phi_1)_\#\lambda_{\tau,\ell})(u)
\le
4\epsilon^2D_1D_2,
\]
which proves \eqref{eq:v049-normal-flag-density-product}.  There are only finitely many conic
charts, so \eqref{eq:v049-kappa-as-symplectic-small-ball} supplies one with
\(\kappa_\ell\gtrsim\kappa_{\rho_n,\tau_n}\).  Substitute
\eqref{eq:v048-counterexample-shell-excess} and use
\(\max\{D_1,D_2\}\ge(D_1D_2)^{1/2}\).
\end{proof}

\begin{proposition}[Arbitrary-threshold routing of a normalized Maslov graph]
\label{prop:v050-arbitrary-threshold-maslov-routing}
Use the normalized collision graph \(0\le h\le1\) and probability measure \(\nu\) from
\cref{thm:v046-quantitative-four-cycle}, and put
\[
M=e(h)=\iint h(x,y)\,d\nu(x)d\nu(y).
\]
For a four-cycle \(x_0,x_1,x_2,x_3,x_0\), write
\[
\mathcal D(x_0,x_1,x_2,x_3)
=
\left|
\det(a_1-a_0,a_2-a_0,a_3-a_0)
\right|.
\]
For every threshold \(0<\epsilon\le c\Delta<1\), at least one of the following holds:
\begin{enumerate}[label=\textnormal{(\roman*)}]
\item there are \(s_0\in J\) and \(c_0\in\mathbb R^3\) such that
\begin{equation}
\label{eq:v050-arbitrary-threshold-bush}
\boxed{
\nu\left\{(a,b):|b+s_0a-c_0|\le C\epsilon/\Delta\right\}
\gtrsim M^3;
}
\end{equation}
\item the determinant-small four-cycle content satisfies
\begin{equation}
\label{eq:v050-arbitrary-threshold-rank-loss}
\boxed{
\int
h_{01}h_{12}h_{23}h_{30}
\mathbf 1_{\{\mathcal D\le\Delta\}}
\,d\nu^4
\gtrsim M^4.
}
\end{equation}
\end{enumerate}
Thus the fourth-root estimate in
\eqref{eq:v046-common-neighbor-small} is only what results from discarding the determinant-small
cycles by their ambient sublevel measure.  If those cycles are retained, no lower threshold on
\(M\) is needed to obtain a bush-or-rank-loss alternative.
\end{proposition}

\begin{proof}
Twice applying Cauchy--Schwarz gives total four-cycle content at least \(M^4\).  Split this content
at \(\mathcal D=\Delta\).  If the determinant-small part carries at least half, then
\eqref{eq:v050-arbitrary-threshold-rank-loss} holds.

Otherwise let \(G(p,q)\) be the determinant-nondegenerate contribution of the two common neighbors
of the anchors \(p,q\), and let
\[
K(p,q)=\int h(p,y)h(q,y)\,d\nu(y).
\]
Then \(\iint G(p,q)\,d\nu(p)d\nu(q)\gtrsim M^4\), whereas
\[
\iint K(p,q)\,d\nu(p)d\nu(q)
=
\int\left(\int h(x,y)\,d\nu(x)\right)^2d\nu(y)
\le M.
\]
Hence some anchor pair satisfies \(G(p,q)/K(p,q)\gtrsim M^3\).  For every retained common-neighbor
pair, the Cramer-rule calculation in
\cref{thm:v046-quantitative-four-cycle} gives
\[
|s_{py}-s_{py'}|
\lesssim
\epsilon/\Delta.
\]
Partitioning \(J\) into intervals of this length and repeating the last pigeonhole in the proof of
\cref{thm:v046-common-neighbor-dichotomy} produces one interval carrying
\(\gtrsim M^3\) of the common-neighbor measure.  Since \(h\le1\), its underlying \(\nu\)-mass is at
least the same size, and the contact-edge equation places it in the bush
\eqref{eq:v050-arbitrary-threshold-bush}.
\end{proof}

\begin{proposition}[Bounded-density cost of every anchored Maslov rank packet]
\label{prop:v050-anchored-packet-density-cost}
Work in a normalized density-good Darboux patch, so the horizontal marginal of the selector measure
has density at most \(C_U\) on a fixed bounded region.  Use the hypotheses and notation of
\cref{prop:v046-anchored-rank-packet}, but denote its horizontal noncollinearity threshold by
\(0<\theta<1\).  If the de-multiplied output \(H_{Y_0}\) has direction mass
\[
h:=\sigma(H_{Y_0}),
\]
then the three alternatives of that proposition imply, respectively,
\begin{equation}
\label{eq:v050-three-packet-density-costs}
\boxed{
\begin{array}{rcl}
\textnormal{full-secant plane}&:&
h\lesssim_U R+\gamma/\theta,\\[1mm]
\textnormal{horizontal line}&:&
h\lesssim_U(R+\theta/\tau_1)^2,\\[1mm]
\textnormal{polarized plane}&:&
h\lesssim_U
R+\varepsilon_{\Delta,\gamma,\upsilon}/\theta^2.
\end{array}
}
\end{equation}
Consequently put
\[
\delta_{\Delta,\tau_1,\upsilon}
:=
\Delta^{2/15}\tau_1^{-14/15}\upsilon^{-2/5}.
\]
Provided the following thresholds lie in \((0,1)\), choose
\begin{equation}
\label{eq:v050-optimized-rank-thresholds}
\theta=\Delta^{1/15}\tau_1^{8/15}\upsilon^{-1/5},
\qquad
\gamma=\Delta^{1/5}\tau_1^{-2/5}\upsilon^{-3/5}.
\end{equation}
Then the alternative-independent estimate is
\begin{equation}
\label{eq:v050-optimized-packet-cost}
\boxed{
h
\lesssim_U
\delta_{\Delta,\tau_1,\upsilon}
+
R\Delta^{-2/15}\tau_1^{-16/15}\upsilon^{-3/5}.
}
\end{equation}
In particular, if
\[
R
\le
\Delta^{4/15}\tau_1^{2/15}\upsilon^{1/5},
\]
then
\begin{equation}
\label{eq:v050-rank-packet-mass-cap}
\boxed{
\lambda R^4m_\Gamma
\lesssim_U
h
\lesssim_U
\Delta^{2/15}\tau_1^{-14/15}\upsilon^{-2/5}.
}
\end{equation}
Thus a determinant-small anchored family cannot carry substantial de-multiplied direction mass at
arbitrarily small determinant threshold.
\end{proposition}

\begin{proof}
If \(F\subset\mathbb R^6\) is the full-secant plane in
\eqref{eq:v046-anchored-full-rank-collapse}, then \(\dim\pi_a(F)\le2\).  Hence the horizontal
projection of \(H_{Y_0}\) lies in an
\(O(R+\gamma/\theta)\)-neighborhood of an affine subspace of dimension at most two.  Its intersection
with the fixed bounded direction region has Lebesgue measure
\(O_U(R+\gamma/\theta)\).  The bounded horizontal density proves the first line of
\eqref{eq:v050-three-packet-density-costs}.

In the horizontal-line alternative, the corresponding tubular neighborhood has three-dimensional
Lebesgue measure \(O_U((R+\theta/\tau_1)^2)\), proving the second line.  In the polarized
alternative, \eqref{eq:v046-anchored-polarized-packet} places the directions within
\(O(R+\varepsilon_{\Delta,\gamma,\upsilon}/\theta^2)\) of the affine two-plane
\(a_p+E\), proving the third line.

Now use \eqref{eq:v046-anchored-packet-error}.  Under
\eqref{eq:v050-optimized-rank-thresholds},
\[
\frac{\gamma}{\theta}
=
\frac{\theta^2}{\tau_1^2}
=
\delta_{\Delta,\tau_1,\upsilon},
\qquad
\frac{(\Delta/\gamma)^{1/3}}{\upsilon\theta^2}
=
\delta_{\Delta,\tau_1,\upsilon}.
\]
The remaining \(R\)-term in the polarized alternative is
\[
O\!\left(
R\Delta^{-2/15}\tau_1^{-16/15}\upsilon^{-3/5}
\right),
\]
which proves \eqref{eq:v050-optimized-packet-cost}.  Under the displayed condition on \(R\), this
term is at most \(\delta_{\Delta,\tau_1,\upsilon}\).  Combine the result with the lower mass bound
\eqref{eq:v046-anchored-packet-mass} to obtain
\eqref{eq:v050-rank-packet-mass-cap}.
\end{proof}

\begin{lemma}[Weighted cycle-relative anchored time trichotomy]
\label{lem:v050-weighted-relative-time-trichotomy}
Fix anchors \(p,q\), let \(\lambda=\lambda_{p,q}\) be the common-neighbor measure of total mass
\(T>0\), and let \(\Gamma\subset Y\times Y\) be symmetric with
\[
Z:=(\lambda\times\lambda)(\Gamma)>0.
\]
Give every \(y\in Y\) its two contact labels \(s_y,t_y\), and fix
\(R\le\delta_s<1\) and \(0<\tau_1<1\).  At least one of the following holds:
\begin{enumerate}[label=\textnormal{(\roman*)}]
\item \(|a_p-a_q|<\tau_1\);
\item a subfamily \(\Gamma_{\mathrm{sep}}\subset\Gamma\), of product measure
\begin{equation}
\label{eq:v050-separated-time-cycle-mass}
(\lambda\times\lambda)(\Gamma_{\mathrm{sep}})
\gtrsim Z,
\end{equation}
has three pairwise \(c\delta_s\)-separated labels among
\(s_y,t_y,s_{y'},t_{y'}\) for every \((y,y')\in\Gamma_{\mathrm{sep}}\);
\item there are two intervals \(I_1,I_2\subset J\), each of length at most
\(C_{U,\tau_1}(R+\delta_s)\), and a set \(Y_0\subset Y\) such that
\begin{equation}
\label{eq:v050-two-time-neighbor-mass}
\boxed{
\lambda(Y_0)
\gtrsim
Z/T,
\qquad
s_y,t_y\in I_1\cup I_2
\quad(y\in Y_0).
}
\end{equation}
\end{enumerate}
In case \textnormal{(iii)}, if the intervals are not quantitatively separated, their union gives an
\(O_{U,\tau_1}(R+\delta_s)\)-bush.  Otherwise the pointwise argument of
\cref{lem:v046-two-time-common-neighbors} partitions \(Y_0\) into four sets, one of
\(\nu\)-mass \(\gtrsim Z/T\), which is a bush or an angular packet of radius
\begin{equation}
\label{eq:v050-weighted-two-time-angular-radius}
\boxed{
O_{U,\tau_1}\!\left(
\frac{R+\delta_s}{g^2}
\right),
\qquad
g:=\operatorname{dist}(I_1,I_2).
}
\end{equation}
\end{lemma}

\begin{proof}
Assume \textnormal{(i)} fails.  Call a pair \((y,y')\) incompatible when its four labels cannot be
covered by two intervals of length \(10\delta_s\).  The elementary one-dimensional packing argument
from \cref{lem:v046-anchored-time-packet-trichotomy} shows that every incompatible pair contains
three pairwise \(c\delta_s\)-separated labels.  If the incompatible part of \(\Gamma\) has measure at
least \(Z/4\), then \textnormal{(ii)} holds.

Otherwise the compatible-edge subgraph \(\Gamma_{\mathrm c}\) has measure at least \(3Z/4\).  Put
\[
Y_{\mathrm{far}}
=
\{y:|s_y-t_y|>100\delta_s\},
\qquad
d_{\mathrm c}(y)
=
\lambda\{y':(y,y')\in\Gamma_{\mathrm c}\}.
\]
If some \(y_0\in Y_{\mathrm{far}}\) has
\(d_{\mathrm c}(y_0)\ge Z/(8T)\), compatibility forces both labels of every compatible neighbor of
\(y_0\) into the two intervals of length \(10\delta_s\) about its two labels.  These intervals are
separated by more than eight times their length, and the neighbor set proves \textnormal{(iii)}.

Suppose no such far vertex exists.  Then
\[
\int_{Y_{\mathrm{far}}}d_{\mathrm c}(y)\,d\lambda(y)
<
Z/8.
\]
By symmetry, compatible edges with at least one far endpoint have measure below \(Z/4\).  Hence the
close--close compatible edges have measure at least \(Z/2\), and some close vertex has a set of close
compatible neighbors of \(\lambda\)-mass at least \(Z/(2T)\).  For every close vertex subtract the
two contact-edge equations and put \(r_y=(s_y+t_y)/2\).  Since
\(|a_p-a_q|\ge\tau_1\),
\[
|b_q-b_p+r_y(a_q-a_p)|
\lesssim
R+\delta_s
\]
places all \(r_y\)'s in one interval of length
\(C_{U,\tau_1}(R+\delta_s)\).  Hence both labels of the selected close-neighbor set lie in one
enlarged interval, proving \eqref{eq:v050-two-time-neighbor-mass}.

If the two intervals in \textnormal{(iii)} are not quantitatively separated, their union is one
enlarged time interval and the first contact-edge equation gives a bush.  If they are separated, the
four pointwise cases in \cref{lem:v046-two-time-common-neighbors} give a bush or an angular packet
with the stated bounds.  Since
\(d\lambda=h(p,y)h(q,y)d\nu(y)\le d\nu(y)\), passing to underlying direction mass loses nothing.
\end{proof}

\begin{proposition}[Weighted continuous anchored rank packet]
\label{prop:v050-continuous-anchored-rank-packet}
Fix anchors \(p,q\) with \(|a_p-a_q|\ge\tau_1\), and let
\begin{equation}
\label{eq:v050-common-neighbor-measure}
d\lambda_{p,q}(y)
:=
h(p,y)h(q,y)\,d\nu(y),
\qquad
T:=\lambda_{p,q}(Y)>0.
\end{equation}
Let \(\Gamma\subset Y\times Y\) be a symmetric measurable family of anchored cycles, and put
\[
Z:=(\lambda_{p,q}\times\lambda_{p,q})(\Gamma)>0,
\qquad
m_\Gamma:=Z/T.
\]
Assume every pair \((y,y')\in\Gamma\) has horizontal determinant at most \(\Delta\), and three of
its four contact labels have Vandermonde product at least \(\upsilon\).  Then, for every
\(0<\gamma,\theta<1\), there are \(y_0\in Y\) and a measurable set \(Y_0\subset Y\) such that
\begin{equation}
\label{eq:v050-continuous-packet-mass}
\boxed{
\nu(Y_0)
\ge
\lambda_{p,q}(Y_0)
\gtrsim
\frac ZT
=
m_\Gamma,
}
\end{equation}
and one of the three geometric conclusions in
\cref{prop:v046-anchored-rank-packet} holds on \(Y_0\), with its parameter \(\kappa\) there replaced
by \(\theta\) and with exactly the same errors
\cref{eq:v046-anchored-full-rank-collapse,eq:v046-anchored-line-packet,%
eq:v046-anchored-polarized-packet}.
\end{proposition}

\begin{proof}
The proof of \cref{prop:v046-anchored-rank-packet} uses only deletion and average degree, so it
extends to the finite measure \(\lambda_{p,q}\).  For completeness, split \(\Gamma\) according to
whether the full secant volume \(\mathfrak g(y,y')\) is below \(\gamma\), retain a part of product
measure at least \(Z/2\), and rename it \(\Gamma\).  If the retained part is the low-volume part,
define
\[
Y_{\mathrm{full-line}}
=
\left\{
y:
|(z_y-z_p)\wedge(z_q-z_y)|<\theta
\right\}.
\]
If \(\lambda_{p,q}(Y_{\mathrm{full-line}})\gtrsim Z/T\), this set gives the line alternative.
Otherwise the cycles incident to it have product measure at most
\(2T\lambda_{p,q}(Y_{\mathrm{full-line}})\), so deletion retains a fixed fraction of the relevant
cycle mass.  Weighted average degree then selects \(y_0\) whose retained neighbor set has
\(\lambda_{p,q}\)-mass \(\gtrsim Z/T\); division of the triple-volume bound by the base two-volume
gives the full-secant plane alternative.

If instead the retained part is the full-secant-volume part, repeat the same argument with
\[
Y_{\mathrm{line}}
=
\left\{
y:
|(a_y-a_p)\times(a_q-a_p)|<\theta
\right\}.
\]
Either this set has \(\lambda_{p,q}\)-mass \(\gtrsim Z/T\), giving the horizontal-line alternative,
or weighted average degree fixes \(y_0\) and a neighbor set of that mass on which the pointwise
three-time interpolation and shared-flag argument from
\cref{prop:v046-anchored-rank-packet} gives the polarized-plane alternative.  Finally
\(h(p,y)h(q,y)\le1\) implies \(\nu(Y_0)\ge\lambda_{p,q}(Y_0)\).
\end{proof}

\begin{corollary}[Determinant-small Maslov mass yields an actual direction packet]
\label{cor:v050-rank-loss-cycle-to-packet}
Assume the determinant-small output
\eqref{eq:v050-arbitrary-threshold-rank-loss} occurs, and that a fixed fraction of that cycle content
has anchors satisfying \(|a_p-a_q|\ge\tau_1\) and has three-label Vandermonde product at least
\(\upsilon\).  Then one anchored rank-loss packet has actual normalized direction mass
\begin{equation}
\label{eq:v050-rank-loss-packet-lower}
\boxed{h\gtrsim M^3.}
\end{equation}
Consequently the bounded-density estimate
\eqref{eq:v050-optimized-packet-cost} gives
\begin{equation}
\label{eq:v050-rank-loss-mass-comparison}
\boxed{
M^3
\lesssim_U
\Delta^{2/15}\tau_1^{-14/15}\upsilon^{-2/5}
+
R\Delta^{-2/15}\tau_1^{-16/15}\upsilon^{-3/5}.
}
\end{equation}
If
\(R\le\Delta^{4/15}\tau_1^{2/15}\upsilon^{1/5}\), this simplifies to
\begin{equation}
\label{eq:v050-rank-loss-saving}
\boxed{
M
\lesssim_U
\Delta^{2/45}\tau_1^{-14/45}\upsilon^{-2/15}.
}
\end{equation}
\end{corollary}

\begin{proof}
Let \(Z(p,q)\) be the determinant-small common-neighbor pair content for the anchors \(p,q\), and
let \(T(p,q)=\lambda_{p,q}(Y)\).  The retained fixed fraction gives
\[
\iint Z(p,q)\,d\nu(p)d\nu(q)
\gtrsim
M^4,
\]
whereas the calculation in
\cref{prop:v050-arbitrary-threshold-maslov-routing} gives
\[
\iint T(p,q)\,d\nu(p)d\nu(q)
\le
M.
\]
Therefore some retained anchor pair obeys
\[
\frac{Z(p,q)}{T(p,q)}
\gtrsim
M^3.
\]
Apply \cref{prop:v050-continuous-anchored-rank-packet} and then
\cref{prop:v050-anchored-packet-density-cost}.  Taking a cube root gives
\eqref{eq:v050-rank-loss-saving}.
\end{proof}

\begin{theorem}[Exact local routing below the fourth-root scale]
\label{thm:v050-complete-local-maslov-routing}
Fix
\[
0<\epsilon\le c\Delta<1,
\qquad
\epsilon\le\delta_s<1,
\qquad
0<\tau_1<1.
\]
For a normalized Maslov graph of edge mass \(M\), at least one of the following holds:
\begin{enumerate}[label=\textnormal{(\roman*)}]
\item an \(O(\epsilon/\Delta)\)-Lagrangian bush has direction mass \(\gtrsim M^3\);
\item a common-neighbor family of direction mass \(\gtrsim M^3\) has an anchor pair satisfying
\(|a_p-a_q|<\tau_1\), and hence enters the next angular Darboux refinement;
\item a two-time packet of direction mass \(\gtrsim M^3\) lies in intervals of width
\(O_{U,\tau_1}(\epsilon+\delta_s)\); if their gap is \(g\), it gives either one
\(O_{U,\tau_1}(\epsilon+\delta_s+g)\)-bush or an angular packet of radius
\(O_{U,\tau_1}((\epsilon+\delta_s)/g^2)\);
\item a separated-time anchored rank packet has direction mass \(h\gtrsim M^3\) and obeys
\eqref{eq:v050-rank-loss-mass-comparison}, with \(R\) there read as \(\epsilon\); in the
physical-error range of
\eqref{eq:v050-rank-loss-saving}, one has
\[
M
\lesssim_U
\Delta^{2/45}\tau_1^{-14/45}\upsilon^{-2/15},
\qquad
\upsilon\gtrsim\delta_s^3.
\]
\end{enumerate}
This alternative holds without assuming \(M\gtrsim\epsilon^{1/4}\).  Thus every local patch in the
former fourth-root exponent gap is now routed to a physical bush, a deeper angular packet, or a
quantitatively mass-capped Maslov rank packet.
\end{theorem}

\begin{proof}
Apply \cref{prop:v050-arbitrary-threshold-maslov-routing}.  Its first output is
\textnormal{(i)}.  In its determinant-small output, write \(Z(p,q)\) and \(T(p,q)\) as in the proof
of \cref{cor:v050-rank-loss-cycle-to-packet}.  Split the total \(Z\)-mass according to whether
\(|a_p-a_q|<\tau_1\).  One part carries a fixed fraction.  Since
\[
\iint T(p,q)\,d\nu(p)d\nu(q)
\le M,
\]
that part contains anchors with
\[
Z(p,q)/T(p,q)
\gtrsim M^3.
\]
If the close-anchor part was selected, then \(Z(p,q)\le T(p,q)^2\) implies
\(T(p,q)\gtrsim M^3\), giving \textnormal{(ii)}.

Otherwise apply \cref{lem:v050-weighted-relative-time-trichotomy}, with \(R=\epsilon\), to the
determinant-small cycle family for the selected separated anchors.  Its time-concentrated output is
\textnormal{(iii)} and
retains direction mass \(\gtrsim Z/T\gtrsim M^3\).  Its separated-time output has
\(\upsilon\gtrsim\delta_s^3\); apply
\cref{prop:v050-continuous-anchored-rank-packet,prop:v050-anchored-packet-density-cost} to obtain
\textnormal{(iv)}.
\end{proof}

\begin{corollary}[The weighted two-time output has no scale-persistent drift]
\label{cor:v050-two-time-scale-persistence}
Consider a sequence of outputs \textnormal{(iii)} in
\cref{thm:v050-complete-local-maslov-routing} with
\(\epsilon_n\to0\) and choose \(\delta_{s,n}\to0\).  After passing to a subsequence, the same
direction-mass scale \(\gtrsim M_n^3\) enters one of the following:
\begin{enumerate}[label=\textnormal{(\roman*)}]
\item the interval gap \(g_n\to0\), and their union is one Lagrangian bush of radius
\(O_{U,\tau_1}(\epsilon_n+\delta_{s,n}+g_n)\to0\);
\item \(g_n\ge g_0>0\), and
\eqref{eq:v050-weighted-two-time-angular-radius} is a shrinking angular packet.
\end{enumerate}
Thus the weighted time-packet passage preserves the \(Z/T\) mass scale and introduces no third
longitudinal-drift alternative.
\end{corollary}

\subsection{Polarization and affine-offset control}

\begin{lemma}[Transverse Maslov polarizations supply the second horizontal codimension]
\label{lem:v053-transverse-polarization-intersection}
Let \(\Pi_1,\Pi_2\subset\mathbb R^3\) be affine two-planes with unit normals \(n_1,n_2\), and assume
\[
|n_1\times n_2|\ge\vartheta>0.
\]
For every bounded direction region \(U\) and \(0<w<1\),
\begin{equation}
\label{eq:v053-two-plane-tube-volume}
\boxed{
\left|
U\cap N_w(\Pi_1)\cap N_w(\Pi_2)
\right|
\lesssim_U
\frac{w^2}{\vartheta}.
}
\end{equation}
Consequently, if a selector-direction subset \(H\) has density bounded by \(C_U\) and is carried by
two anchored Maslov packets of thickness \(w\) whose polarization planes have Grassmannian distance
at least \(c\vartheta\), then
\begin{equation}
\label{eq:v053-transverse-polarization-mass}
\boxed{
\sigma(H)
\lesssim_U
\frac{w^2}{\vartheta}.
}
\end{equation}
This is the missing quadratic horizontal factor on every transversely reused graph cell.
\end{lemma}

\begin{proof}
After a rigid motion, the two plane conditions are
\[
|n_1\cdot x-c_1|\le w,
\qquad
|n_2\cdot x-c_2|\le w.
\]
Use coordinates along \(n_1,n_2\), and along the intersection direction
\(n_1\times n_2\).  The Jacobian of this coordinate map is
\(|n_1\times n_2|\ge\vartheta\).  The first two coordinates range over intervals of length
\(O(w)\), while boundedness of \(U\) bounds the third.  This proves
\eqref{eq:v053-two-plane-tube-volume}; the horizontal density bound gives
\eqref{eq:v053-transverse-polarization-mass}.
\end{proof}

\begin{corollary}[Weighted summation of the transverse high-reuse cells]
\label{cor:v053-transverse-reuse-weighted-sum}
Use the reuse multiplicities \(r(\chi)\) of
\cref{prop:v046-anchor-cell-reuse}.  Let \(\mathcal C_{\mathrm{tr}}(K,\vartheta)\) be a family of
pairwise disjoint graph cells such that \(r(\chi)\ge K\) and two packets through every \(\chi\) have
polarization planes separated by at least \(c\vartheta\).  If all cells and packet tubes have common
direction thickness \(O(w)\), then their actual direction union satisfies
\begin{equation}
\label{eq:v053-transverse-reuse-sum}
\boxed{
\sigma(H_{\mathrm{tr}})
\lesssim_U
\frac{w^2}{K\vartheta}
\sum_{\chi\in\mathcal C_{\mathrm{tr}}(K,\vartheta)}r(\chi).
}
\end{equation}
Thus the transverse quadratic factor is summed with the reciprocal reuse gain \(K^{-1}\), not with
the total number of ambient graph cells.
\end{corollary}

\begin{proof}
The direction cube underlying a reused graph cell lies, after a fixed enlargement by its own
diameter, in the two transverse plane tubes.  Hence
\cref{lem:v053-transverse-polarization-intersection} bounds its actual selector-direction mass by
\(C_Uw^2/\vartheta\).  The graph cells are disjoint, and
\[
\#\mathcal C_{\mathrm{tr}}(K,\vartheta)
\le
K^{-1}
\sum_{\chi\in\mathcal C_{\mathrm{tr}}(K,\vartheta)}r(\chi).
\]
Multiplication gives \eqref{eq:v053-transverse-reuse-sum}.
\end{proof}

\begin{lemma}[A polarized cycle packet controls the vertical affine offset]
\label{lem:v053-polarized-packet-vertical-offset}
In alternative \textnormal{(iii)} of
\cref{prop:v046-anchored-rank-packet}, with its noncollinearity parameter denoted by \(\theta\), every
retained neighbor also satisfies
\begin{equation}
\label{eq:v053-polarized-vertical-plane}
\boxed{
\operatorname{dist}(b_y,b_p+E)
\lesssim_{U,I,J}
\varepsilon_{\Delta,\gamma,\upsilon}/\theta^2.
}
\end{equation}
Hence a coherent polarized packet is contained, up to the displayed error, in one affine phase-space
four-plane
\[
(a_p,b_p)+(E\oplus E).
\]
\end{lemma}

\begin{proof}
The cycle secant \(z_y-z_p\) belongs to \(V_{y_0,y}\).  By
\eqref{eq:v046-anchored-polarized-packet}, that three-plane lies within
\(C\varepsilon_{\Delta,\gamma,\upsilon}/\theta^2\) in the Grassmannian of a three-plane contained in
\(E\oplus E\).  All secants stay in a fixed bounded region, so
\[
\operatorname{dist}(z_y-z_p,E\oplus E)
\lesssim_{U,I,J}
\varepsilon_{\Delta,\gamma,\upsilon}/\theta^2.
\]
Take the vertical component.
\end{proof}

\begin{lemma}[A common direction cell aligns coherent horizontal plane offsets]
\label{lem:v053-common-cell-offset-alignment}
Write affine polarization planes as
\[
\Pi_i=\{a:n_i\cdot a=c_i\},
\qquad
|n_i|=1.
\]
Assume \(n_i\) all lie within \(\vartheta\) of one unit normal \(n_\infty\), and a direction cube
\(Q\subset U\) of diameter \(O(w)\) meets every tube \(N_w(\Pi_i)\).  Then
\begin{equation}
\label{eq:v053-common-cell-offset-bound}
\boxed{
|c_i-c_j|
\lesssim_U
w+\vartheta
\qquad(i,j).
}
\end{equation}
Thus reuse of one graph cell by coherently polarized packets automatically aligns their affine
horizontal offsets at the geometric resolution of that cell.
\end{lemma}

\begin{proof}
Choose \(a_i\in Q\) with \(|n_i\cdot a_i-c_i|\le w\).  Since \(Q\) has diameter \(O(w)\),
\[
\begin{aligned}
|c_i-c_j|
&\le
|c_i-n_i\cdot a_i|
+|n_i\cdot(a_i-a_j)|
+|(n_i-n_j)\cdot a_j|
+|n_j\cdot a_j-c_j|\\
&\lesssim_U
w+\vartheta.
\end{aligned}
\]
This is \eqref{eq:v053-common-cell-offset-bound}.
\end{proof}

\begin{corollary}[A common product cell aligns both phase offsets]
\label{cor:v053-common-cell-phase-offsets}
For a coherently polarized packet put
\[
c_i^a=n_i\cdot a_{p_i},
\qquad
c_i^b=n_i\cdot b_{p_i}.
\]
Suppose a graph product cell of diameter \(O(w)\) meets packet (i) inside the affine phase-plane
tube supplied by
\cref{lem:v053-polarized-packet-vertical-offset}, of additional thickness \(e_i\).  If
\(|n_i-n_\infty|\le\vartheta\), then all packets through that cell satisfy
\begin{equation}
\label{eq:v053-common-cell-phase-offset-bound}
\boxed{
|c_i^a-c_j^a|+|c_i^b-c_j^b|
\lesssim_U
w+\vartheta+e_i+e_j.
}
\end{equation}
\end{corollary}

\begin{proof}
Apply the calculation in \cref{lem:v053-common-cell-offset-alignment} first to the horizontal
coordinates of a point in the common product cell and then to its vertical coordinates.  The latter
application uses \eqref{eq:v053-polarized-vertical-plane}; the two packet errors contribute
\(e_i+e_j\).
\end{proof}

\begin{lemma}[One-dimensional Reeb-window sum for parallel phase offsets]
\label{lem:v055-parallel-offset-window-sum}
Let \(\zeta\) be a finite measure of total mass \(m\) on a fixed bounded subset of the phase-offset
plane \(\mathbb R_q\times\mathbb R_p\).  Assume its \(q\)-marginal satisfies
\begin{equation}
\label{eq:v055-offset-frostman}
\zeta_q(I)
\le
D|I|
\qquad
\text{for every interval \(I\) with \(|I|\ge W\)}.
\end{equation}
Then, for every fixed bounded Reeb interval \(J\),
\begin{equation}
\label{eq:v055-parallel-offset-window-bound}
\boxed{
\int_J
\iint
\mathbf 1_{\{|(p-p')+s(q-q')|\le W\}}
\,d\zeta(q,p)d\zeta(q',p')\,ds
\lesssim_J
DmW\log(2/W).
}
\end{equation}
This estimate is purely the contact/Reeb window calculation for the normal phase-offset pair; it
uses no Fourier or spectral argument.
\end{lemma}

\begin{proof}
For \(|q-q'|<W\), the pair mass is at most
\[
\int \zeta_q([q-W,q+W])\,d\zeta(q,p)
\lesssim
DmW,
\]
and the time interval has bounded length.  On the dyadic shell
\(r\le|q-q'|<2r\), \(r\ge W\), the collision-time window has length at most \(C_JW/r\).  On the other
hand \eqref{eq:v055-offset-frostman} bounds the pair mass of that shell by \(CDrm\).  Hence every
dyadic shell contributes \(O_J(DmW)\); there are \(O(\log(2/W))\) shells.
\end{proof}

\begin{lemma}[Pointwise Reeb degree of one phase offset]
\label{lem:v056-pointwise-offset-degree}
Under the hypotheses of \cref{lem:v055-parallel-offset-window-sum}, every fixed
\(\xi=(q,p)\) satisfies
\begin{equation}
\label{eq:v056-pointwise-offset-degree}
\boxed{
\int_J\int
\mathbf 1_{\{|(p-p')+s(q-q')|\le W\}}
\,d\zeta(q',p')ds
\lesssim_J
DW\log(2/W).
}
\end{equation}
Consequently the two coordinate marginals of the normal-window pair measure are each dominated by
\(C_JDW\log(2/W)\zeta\), not merely by the trivial measure \(|J|m\zeta\).
\end{lemma}

\begin{proof}
For \(|q-q'|<W\), the \(q'\)-mass is at most \(CDW\), and the time interval has bounded length.
On \(r\le|q-q'|<2r\), \(r\ge W\), the \(q'\)-mass is at most \(CDr\), while the admissible time
interval has length at most \(C_JW/r\).  Each dyadic shell therefore costs \(O_J(DW)\); summing the
shells proves \eqref{eq:v056-pointwise-offset-degree}.  The marginal statement follows by
integrating the pointwise estimate against the first, respectively second, copy of \(\zeta\).
\end{proof}

\begin{lemma}[Exact normal--tangential factorization in a Maslov polarization]
\label{lem:v056-polarized-collision-factorization}
Fix an oriented horizontal two-plane \(E\subset\mathbb R^3\), let \(n\) be its unit normal, and
write every phase point uniquely as
\begin{equation}
\label{eq:v056-polarized-coordinates}
a=x+qn,
\qquad
b=y+pn,
\qquad
x,y\in E,\quad q,p\in\mathbb R.
\end{equation}
In these coordinates the canonical symplectic form and the incidence Lagrangian split as
\[
\omega
=
\omega_E+dq\wedge dp,
\qquad
L_s
=
\{(u,-su):u\in E\}
\oplus
\{(rn,-srn):r\in\mathbb R\}.
\]
For two phase points \(z=(a,b)\), \(z'=(a',b')\), put
\[
T_s(z,z')=(y-y')+s(x-x')\in E,
\qquad
N_s(z,z')=(p-p')+s(q-q')\in\mathbb R.
\]
Then
\begin{equation}
\label{eq:v056-orthogonal-front-split}
\boxed{
|\Pi_s(z)-\Pi_s(z')|^2
=
|T_s(z,z')|^2+|N_s(z,z')|^2.
}
\end{equation}
Consequently, for every \(W>0\),
\begin{equation}
\label{eq:v056-kernel-sandwich}
\mathbf 1_{\{|T_s|\le W/\sqrt2\}}
\mathbf 1_{\{|N_s|\le W/\sqrt2\}}
\le
\mathbf 1_{\{|\Pi_s(z)-\Pi_s(z')|\le W\}}
\le
\mathbf 1_{\{|T_s|\le W\}}
\mathbf 1_{\{|N_s|\le W\}}.
\end{equation}
This is the intrinsic two-stage form of the coherent Maslov collision: the scalar factor is the
Reeb-normal window, while the remaining factor is the tangential front collision in
\(E\oplus E\).
\end{lemma}

\begin{proof}
The orthogonal splitting \(\mathbb R^3=E\oplus\mathbb Rn\) gives
\[
(b-b')+s(a-a')
=
T_s(z,z')+N_s(z,z')n.
\]
Pythagoras proves \eqref{eq:v056-orthogonal-front-split}; the two indicator inequalities follow
immediately.
\end{proof}

\subsection{Weighted slices and defect routing}

\begin{lemma}[Weighted finite-scale \(q\)-slice retention]
\label{lem:v056-weighted-q-slice-retention}
Let \(G\) be one quantitative Sticky packing piece with selector measure \(\mu\), bounded direction
density, and
\begin{equation}
\label{eq:v056-global-packing-piece}
N_r(G)\le C_\zeta r^{-3-\zeta}.
\end{equation}
Let \(\Gamma\) be a finite nonnegative measure on \(G\times G\times J\), of total mass
\(M>0\), such that both graph-coordinate marginals are dominated by \(L\mu\).  For a dyadic
\(q\)-interval \(I\) of length \(r\), let \(\mathcal N_I(r)\) be the number of full
\((x,q,y,p)\) product cells of side \(r\) which meet \(G\) over \(I\).  There is a collection
\(\mathcal I_{\rm good}\) for which
\begin{equation}
\label{eq:v056-weighted-good-slice-mass}
\boxed{
\Gamma\{q(z),q(z')\text{ both lie in good intervals}\}
\ge\frac12M
}
\end{equation}
and every \(I\in\mathcal I_{\rm good}\) obeys
\begin{equation}
\label{eq:v056-weighted-slice-count}
\boxed{
\mathcal N_I(r)
\lesssim_{\zeta,U}
\frac{L}{M}\,r^{-2-\zeta}.
}
\end{equation}
The number of tangential \((x,y)\)-cells occurring over such an interval is bounded by the same
right-hand side.
\end{lemma}

\begin{proof}
A fixed orthogonal change of coordinates changes covering numbers only by a dimensional constant,
so \eqref{eq:v056-global-packing-piece} gives
\[
\sum_I\mathcal N_I(r)
\lesssim_\zeta r^{-3-\zeta}.
\]
Call \(I\) bad when the right side of \eqref{eq:v056-weighted-slice-count}, with a sufficiently
large fixed implicit constant, is exceeded.  Then
\[
\#\{I:I\text{ bad}\}
\lesssim_U
\frac{M}{Lr}.
\]
Bounded direction density gives
\(
\mu\{q\text{ lies in a bad interval}\}
\lesssim_U r\#\{I:I\text{ bad}\}
\), and the choice of the implicit constant makes this at most \(M/(4L)\).  Domination of the first
\(\Gamma\)-marginal shows that pairs whose first coordinate is bad have mass at most \(M/4\); the
same argument applies to the second coordinate.  Removing both sets proves
\eqref{eq:v056-weighted-good-slice-mass}.  Projection of a full product cell to \((x,y)\) meets only
boundedly many tangential cells, which proves the last assertion.
\end{proof}

\begin{lemma}[Collision-preserving multiscale vertical-entropy pruning]
\label{lem:v057-collision-preserving-multiscale-pruning}
Retain the hypotheses and notation of
\cref{lem:v056-weighted-q-slice-retention}.  For a dyadic direction cube
\(A\subset\mathbb R_a^3\) of side \(r\), let
\[
v_r(A)
:=
\#\{B:\ (A\times B)\cap G\ne\varnothing\},
\]
where \(B\) ranges over the dyadic intercept cubes of side \(r\).  Fix
\(0<W<1\), let
\[
\mathcal D_W
=
\{r:\ r\text{ dyadic},\ W\le r\le1\},
\qquad
\ell_W=1+\#\mathcal D_W\lesssim\log(2/W),
\]
and let \(0\le\Theta\le\Gamma\) be any submeasure of total mass
\(H>0\).  There is a Borel endpoint set \(G_{\Theta,W}\subset G\) such that
\begin{equation}
\label{eq:v057-collision-pruned-mass}
\boxed{
\Theta(G_{\Theta,W}\times G_{\Theta,W}\times J)
\ge
\frac12H
}
\end{equation}
and, for every \(z=(a,b(a))\in G_{\Theta,W}\) and every
\(r\in\mathcal D_W\),
\begin{equation}
\label{eq:v057-collision-pruned-branch}
\boxed{
v_r(A_r(a))
\lesssim_{\zeta,U}
\frac{L}{H}\,\ell_W r^{-\zeta}.
}
\end{equation}
Here \(A_r(a)\) is the direction \(r\)-cube containing \(a\).

Moreover, fix a polarized normal coordinate interval \(I\) of length \(W\).  For every
\(r\in\mathcal D_W\) and every tangential direction \(r\)-cube \(X\subset E_x\), the part of
\(G_{\Theta,W}\) with \(q\in I\) and \(x\in X\) meets at most
\begin{equation}
\label{eq:v057-tangential-branch-all-scales}
\boxed{
\lesssim_{\zeta,U}
\frac{L}{H}\,\ell_W r^{-\zeta}
}
\end{equation}
tangential intercept \(r\)-cubes in \(E_y\).  Thus the same estimate controls the number of
tangential contact-front tube branches above \(X\) at every intermediate scale, not only the total
number of \(W\)-cells in one \(q\)-slice.
\end{lemma}

\begin{proof}
The product-grid form of \eqref{eq:v056-global-packing-piece} gives
\[
\sum_Av_r(A)
\lesssim_{\zeta,U}
r^{-3-\zeta}.
\]
For each \(r\in\mathcal D_W\), call \(A\) bad when
\[
v_r(A)>C_{\zeta,U}\frac{L}{H}\ell_Wr^{-\zeta},
\]
with the constant chosen below.  The number of bad cubes is at most
\[
\lesssim_{\zeta,U}
\frac{H}{L\ell_W}r^{-3}.
\]
Bounded direction density therefore gives
\[
\mu\{a\text{ lies in a bad }r\text{-cube}\}
\lesssim_{\zeta,U}
\frac{H}{L\ell_W}.
\]
After increasing \(C_{\zeta,U}\), domination of either coordinate marginal of \(\Gamma\) by
\(L\mu\), together with \(\Theta\le\Gamma\), makes the \(\Theta\)-mass for which a prescribed
endpoint is bad at a prescribed scale at most \(H/(4\ell_W)\).  Sum over the two endpoints and all
scales in \(\mathcal D_W\).  Removing their union costs at most \(H/2\), proving
\eqref{eq:v057-collision-pruned-mass} and \eqref{eq:v057-collision-pruned-branch}.

If \(q\in I\) and \(r\ge W\), then \(I\) meets only \(O(1)\) dyadic \(q\)-intervals of length
\(r\).  Likewise \(X\times I\) is contained in \(O(1)\) full direction \(r\)-cubes.  Above each
one, projection of a full intercept cube in \((y,p)\) meets only \(O(1)\) tangential
\(y\)-cubes.  Hence \eqref{eq:v057-collision-pruned-branch} implies
\eqref{eq:v057-tangential-branch-all-scales}.
\end{proof}

\begin{proposition}[Collision-preserving total Maslov-defect ledger]
\label{prop:v057-total-maslov-defect-ledger}
Normalize the coherent selector occurrence measure to total mass one, and absorb \(D\), reuse
caps, and logarithms into \(W^{-o(1)}\).  Pass to a scale subsequence on which
\begin{equation}
\label{eq:v057-two-defect-exponents}
\|\Lambda_W\|
=
W^{1+\beta+o(1)},
\qquad
\frac{\mathcal C_W}{\|\Lambda_W\|}
=
W^{\gamma+o(1)},
\qquad
\chi:=\beta+\gamma,
\end{equation}
where \(\beta,\gamma\ge0\).  Then
\begin{equation}
\label{eq:v057-full-collision-total-defect}
\boxed{
\mathcal C_W
=
W^{1+\chi+o(1)}.
}
\end{equation}
Consequently \(\chi\ge2-o(1)\) is already the desired full collision bound.  If
\(0\le\chi<2\), apply
\cref{lem:v057-collision-preserving-multiscale-pruning} to
\(\Theta_W\le\Lambda_W\).  At least half of the actual collision mass survives on an endpoint set
for which, simultaneously for every dyadic \(W\le r\le1\),
\begin{equation}
\label{eq:v057-total-defect-vertical-branch}
\boxed{
v_r(A_r(a))
\le
W^{-\chi-o(1)}r^{-\zeta}.
}
\end{equation}
Inside every normal \(W\)-slab, the same bound holds for tangential contact-front branches above
one tangential direction \(r\)-cube.  Moreover each endpoint projection of the retained collision
measure requires direction support of mass at least
\begin{equation}
\label{eq:v057-total-defect-support-mass}
\boxed{
\gtrsim
W^{\chi+o(1)}.
}
\end{equation}
Thus the collision-preserving branch tax and the endpoint-support mass are reciprocal at the power
scale.  The exponent \(\chi\), not \(\beta\) alone, is the invariant defect carried by the
self-selected Maslov diagonal.
\end{proposition}

\begin{proof}
The first identity is multiplication in
\eqref{eq:v057-two-defect-exponents}.  The pointwise Reeb-degree estimate gives both coordinate
marginals of \(\Lambda_W\), and hence of \(\Theta_W\), the domination
\[
(\Theta_W)_1,(\Theta_W)_2
\le
L\nu,
\qquad
L=W^{1-o(1)}.
\]
Use \(H=\mathcal C_W=W^{1+\chi+o(1)}\) in
\eqref{eq:v057-collision-pruned-branch}; this gives
\eqref{eq:v057-total-defect-vertical-branch}.  The tangential assertion is
\eqref{eq:v057-tangential-branch-all-scales}.  Finally, if \(S\) supports either endpoint marginal
of the retained measure, then
\[
\frac12\mathcal C_W
\le
(\Theta_W)_i(S)
\le
L\nu(S).
\]
Hence \(\nu(S)\gtrsim\mathcal C_W/L=W^{\chi+o(1)}\), proving
\eqref{eq:v057-total-defect-support-mass}.
\end{proof}

\begin{proposition}[The total-defect layer re-enters the Maslov four-cycle routing]
\label{prop:v057-total-defect-four-cycle-return}
Work on a normalized angular shell
\(c_0\le|a-a'|\le C_0\) and normalize \(\nu(G)=1\).  Choose a fixed
\(C_{c_0}\) larger than the maximum \(W\)-collision-window length divided by \(W\), and define
\begin{equation}
\label{eq:v057-normalized-maslov-graphon}
h_W(z,z')
:=
\frac1{C_{c_0}W}
\int_J
\mathbf 1_{\{|N_s(z,z')|\le W\}}
\mathbf 1_{\{|T_s(z,z')|\le W\}}
\,ds.
\end{equation}
Then \(h_W\) is a symmetric graphon with \(0\le h_W\le1\), and its edge mass is
\begin{equation}
\label{eq:v057-normalized-maslov-edge-mass}
\boxed{
M_W:=e(h_W)
=
\frac{\mathcal C_W}{C_{c_0}W}.
}
\end{equation}
Every edge with \(h_W(z,z')>0\) has a measurable label \(s_{zz'}\in J\) satisfying
\begin{equation}
\label{eq:v057-return-edge-contact-relation}
|(b-b')+s_{zz'}(a-a')|
\le
\sqrt2W.
\end{equation}
Consequently, for every threshold \(W\le c\Delta<1\), the arbitrary-threshold Maslov routing
\cref{prop:v050-arbitrary-threshold-maslov-routing} gives at least one of:
\begin{enumerate}[label=\textnormal{(\roman*)}]
\item an \(O(W/\Delta)\)-approximate Lagrangian point-probe bush of \(\nu\)-mass
\begin{equation}
\label{eq:v057-total-defect-bush-output}
\boxed{\gtrsim M_W^3;}
\end{equation}
\item determinant-small Maslov four-cycle content
\begin{equation}
\label{eq:v057-total-defect-rank-output}
\boxed{
\int
(h_W)_{01}(h_W)_{12}(h_W)_{23}(h_W)_{30}
\mathbf 1_{\{\mathcal D\le\Delta\}}\,d\nu^4
\gtrsim M_W^4.
}
\end{equation}
\end{enumerate}
On a total-defect level \(\mathcal C_W=W^{1+\chi+o(1)}\), these outputs are respectively
\begin{equation}
\label{eq:v057-total-defect-four-cycle-powers}
\boxed{
\text{bush mass }\gtrsim W^{3\chi+o(1)},
\qquad
\text{rank-loss four-cycle content }\gtrsim W^{4\chi+o(1)}.
}
\end{equation}
Thus the collision-preserving correction does not create a new analytic branch: it returns the
entire bad layer, with its exact defect exponent, to the manuscript's existing
Lagrangian-bush/Maslov-rank-loss alternative.
\end{proposition}

\begin{proof}
For a fixed pair on the normalized angular shell,
\(s\mapsto(b-b')+s(a-a')\) has speed at least \(c_0\).  The set on which both indicators in
\eqref{eq:v057-normalized-maslov-graphon} equal one lies inside a full-vector collision interval of
length at most \(C_{c_0}W\).  Hence \(0\le h_W\le1\); symmetry is immediate.  Tonelli gives
\eqref{eq:v057-normalized-maslov-edge-mass}.

The nonempty compact collision-time section of an edge admits a Borel choice of
\(s_{zz'}\); alternatively, choose the midpoint of its interval.  The orthogonal factorization
\eqref{eq:v056-orthogonal-front-split} gives
\eqref{eq:v057-return-edge-contact-relation}.  The horizontal marginal of \(\nu\) has the bounded
density required in \cref{prop:v050-arbitrary-threshold-maslov-routing}.  Apply that proposition
with collision error \(\sqrt2W\), absorbing the fixed factor into \(c\).  This proves
\eqref{eq:v057-total-defect-bush-output} and
\eqref{eq:v057-total-defect-rank-output}.  Finally
\(M_W=\mathcal C_W/(C_{c_0}W)=W^{\chi+o(1)}\), which yields
\eqref{eq:v057-total-defect-four-cycle-powers}.
\end{proof}

\begin{proposition}[Total defect extracts a quantitative common contact-star layer]
\label{prop:v057-total-defect-contact-star}
Work on a normalized angular shell
\(c_0\le|a-a'|\le C_0\), normalize \(\nu(G)=1\), and define the full polarized collision degree
\begin{equation}
\label{eq:v057-full-collision-degree}
d_W(z)
:=
\int_J\int
\mathbf 1_{\{|N_s(z,z')|\le W\}}
\mathbf 1_{\{|T_s(z,z')|\le W\}}
\,d\nu(z')ds.
\end{equation}
Then \(\int d_W\,d\nu=\mathcal C_W\) and
\(d_W(z)\le L\), where \(L\lesssim DW\log(2/W)\).  There are a dyadic
\(h\in[c\mathcal C_W,CL]\) and a Borel center set
\[
S_h=\{z:\ h\le d_W(z)<2h\}
\]
such that, with \(\ell_{W,\mathcal C}=1+\log(2L/\mathcal C_W)\),
\begin{equation}
\label{eq:v057-contact-star-center-mass}
\boxed{
\nu(S_h)
\gtrsim
\frac{\mathcal C_W}{h\ell_{W,\mathcal C}}.
}
\end{equation}
For every \(z\in S_h\), there is a Reeb time \(s_z\in J\) for which the simultaneous neighbor
set obeys
\begin{equation}
\label{eq:v057-simultaneous-contact-star}
\boxed{
\nu\{z':|N_{s_z}(z,z')|\le W,\ |T_{s_z}(z,z')|\le W\}
\gtrsim_J h.
}
\end{equation}
All selected lines in this set pass through the \(O(W)\)-ball about the same front point
\((\Pi_{s_z}(z),s_z)\); hence its geometric support is a \(W\)-scale common contact star carrying
\(\nu\)-weight \(\gtrsim h\).  If \(\nu\) is an occurrence measure rather than the unrepeated
selector measure, conversion of this weight to actual direction mass still requires the existing
reuse de-multiplication.  Its time-swept set of distinct \(\nu\)-leaves satisfies
\begin{equation}
\label{eq:v057-swept-contact-star-leaves}
\boxed{
\nu\{z':\exists s\in J, |N_s(z,z')|\le W, |T_s(z,z')|\le W\}
\gtrsim_{c_0,J}
\frac{h}{W}.
}
\end{equation}

Under the power normalization of
\cref{prop:v057-total-maslov-defect-ledger}, one may write
\begin{equation}
\label{eq:v057-contact-star-tradeoff-parameter}
h=W^{1+\theta+o(1)}
\qquad\text{for some}\qquad
0\le\theta\le\chi.
\end{equation}
The three retained quantities then have the exact power ledger
\begin{equation}
\label{eq:v057-contact-star-tradeoff}
\boxed{
\begin{array}{c|c|c}
\text{center \(\nu\)-mass}&
\text{one-time star \(\nu\)-weight}&
\text{swept leaf \(\nu\)-mass}\\ \hline
W^{\chi-\theta+o(1)}&
W^{1+\theta+o(1)}&
W^{\theta+o(1)}.
\end{array}
}
\end{equation}
If the occurrence measure is dominated by the unrepeated selector measure in the form
\(\nu\le K\sigma\), then the corresponding actual direction masses are bounded below by the three
entries in \eqref{eq:v057-contact-star-tradeoff} divided by \(K\).  In particular, when the
coherent reuse cap is subpower, the passage from the weighted star to a physical star loses no
power of \(W\).
Thus a total defect \(\chi<2\) is already concentrated on a quantitative family of weighted common
contact stars.  The final step below de-multiplies the \(\theta\)-layer through the coherent reuse ledger and either cross an existing Lagrangian-bush
removal threshold or sum the remaining weighted Maslov incidences.
\end{proposition}

\begin{proof}
The identity \(\int d_W\,d\nu=\mathcal C_W\) is Tonelli's theorem.  The pointwise bound follows
from \(d_W\le d_W^{\rm normal}\) and
\cref{lem:v056-pointwise-offset-degree}.  Degrees below \(\mathcal C_W/2\) contribute at most
\(\mathcal C_W/2\) because \(\nu(G)=1\).  Dyadically decompose the remaining interval
\([\mathcal C_W/2,L]\).  One level contributes at least
\(c\mathcal C_W/\ell_{W,\mathcal C}\) to \(\int d_W\,d\nu\); since \(d_W<2h\) on that level,
\eqref{eq:v057-contact-star-center-mass} follows.

For fixed \(z\in S_h\), put
\[
F_z(s)
=
\nu\{z':|N_s(z,z')|\le W,\ |T_s(z,z')|\le W\}.
\]
Then \(\int_JF_z(s)ds=d_W(z)\ge h\), so averaging in \(s\) proves
\eqref{eq:v057-simultaneous-contact-star}.  The two inequalities imply
\(|\Pi_{s_z}(z)-\Pi_{s_z}(z')|\le\sqrt2W\), giving the common front point.

For each fixed pair in the normalized angular shell, the full collision-time set has length at
most \(C_{c_0}W\): indeed
\(s\mapsto(b-b')+s(a-a')\) has speed at least \(c_0\).  If \(\mathcal N_W(z)\) denotes the swept
neighbor set in \eqref{eq:v057-swept-contact-star-leaves}, then
\[
h
\le
d_W(z)
\le
C_{c_0}W\nu(\mathcal N_W(z)),
\]
which proves that estimate.  Finally
\(\mathcal C_W=W^{1+\chi+o(1)}\le h\le L=W^{1-o(1)}\) gives
\eqref{eq:v057-contact-star-tradeoff-parameter}; substituting it into
\eqref{eq:v057-contact-star-center-mass},
\eqref{eq:v057-simultaneous-contact-star}, and
\eqref{eq:v057-swept-contact-star-leaves} yields
\eqref{eq:v057-contact-star-tradeoff}.
\end{proof}

\begin{lemma}[Polarized branch coloring into direction-separated tangential families]
\label{lem:v058-polarized-branch-coloring}
Discretize a collision-preserved endpoint set at scale \(W\), using full product cells in
\((x,q,y,p)\).  Fix a normal \(q\)-interval \(I\) of length \(W\).  Every occupied labeled cell
determines a tangential \(W\)-tube in \(E\times J\),
\[
\mathbb T_{\mathfrak c}
=
\{(u,s):|u-(y_\chi+s x_\chi)|\le C W\},
\]
while retaining its normal \((q,p)\)-label.  Suppose that, above every tangential direction
\(W\)-cell, there are at most \(B\) such labeled tubes with \(q\in I\).  Then this labeled tube
family is the union of at most \(C B\) subfamilies, each of which has essentially
\(W\)-separated tangential directions.  Summing over all normal \(q\)-intervals in a bounded
region gives at most
\begin{equation}
\label{eq:v058-global-tangential-color-count}
\boxed{J_W\lesssim B W^{-1}}
\end{equation}
direction-separated tangential color classes.
\end{lemma}

\begin{proof}
Split the tangential direction grid into a fixed number of congruence classes so that distinct
cells of one class have \(cW\)-separated centers.  Above one direction cell enumerate the at most
\(B\) labeled tubes and give the \(k\)-th tube color \(k\), together with the congruence color of
its direction cell.  A color class contains at most one tube above each direction cell and its
directions are essentially \(W\)-separated.  There are \(O(B)\) colors for one \(I\), and
\(O(W^{-1})\) normal intervals.
\end{proof}

\begin{lemma}[Dense-shading multiplicity bound in one tangential color]
\label{lem:v058-dense-shading-multiplicity}
Fix \(\varepsilon>0\).  Let \(\eta=\eta(\varepsilon)>0\) be supplied by the
three-dimensional shaded Kakeya theorem of Guth--Wang--Zahl
\cite{GuthWangZahl2026}.  Let \(\mathbb T\) be one essentially direction-separated family of
unit \(W\)-tubes in a bounded subset of \(E\times J\), and put
\[
m(u,s)=\sum_{T\in\mathbb T}\mathbf 1_T(u,s).
\]
For a dyadic multiplicity level
\(X_k=\{2^k\le m<2^{k+1}\}\), give each tube the induced shading
\(Y_k(T)=T\cap X_k\).  Write
\[
\mathbb T_{k,\mathrm{dense}}
=
\{T:|Y_k(T)|\ge W^\eta|T|\}.
\]
Then at least one of the following holds:
\begin{enumerate}[label=\textnormal{(\roman*)}]
\item the dense tubes carry at least half of the incidence on \(X_k\), and
\begin{equation}
\label{eq:v058-dense-level-multiplicity}
\boxed{2^k\lesssim_\varepsilon W^{-\varepsilon};}
\end{equation}
\item more than half of that incidence is carried by tube shadings satisfying
\begin{equation}
\label{eq:v058-sparse-tube-shading}
\boxed{|Y_k(T)|<W^\eta|T|.}
\end{equation}
\end{enumerate}
Consequently, if alternative \textnormal{(ii)} does not carry a fixed portion of any multiplicity
level, then
\begin{equation}
\label{eq:v058-one-color-L2}
\boxed{
\int m^2
\lesssim_\varepsilon
W^{-\varepsilon}\int m.
}
\end{equation}
\end{lemma}

\begin{proof}
An essentially direction-separated tube family has
\(\Delta_{\max}(\mathbb T)\lesssim1\).  In alternative \textnormal{(i)}, restrict to
\(\mathbb T_{k,\mathrm{dense}}\) with the shadings \(Y_k\).  The shaded three-dimensional Kakeya
theorem gives
\[
|X_k|
\ge
|U(\mathbb T_{k,\mathrm{dense}},Y_k)|
\gtrsim_\varepsilon
W^\varepsilon
\#\mathbb T_{k,\mathrm{dense}}\,|T|.
\]
On the other hand, the dense incidence is at least a fixed multiple of
\(2^k|X_k|\) and at most
\(\#\mathbb T_{k,\mathrm{dense}})|T|\).  Combining the inequalities proves
\eqref{eq:v058-dense-level-multiplicity}.  If alternative \textnormal{(ii)} never carries a fixed
portion, dyadic summation gives
\[
\int m^2
\asymp
\sum_k2^k\int_{X_k}m
\lesssim_\varepsilon
W^{-\varepsilon}\sum_k\int_{X_k}m,
\]
which is \eqref{eq:v058-one-color-L2}.
\end{proof}

\begin{proposition}[The dense tangential contribution closes every total defect below one half]
\label{prop:v058-dense-tangential-half-range}
Use the collision-preserved endpoint set in
\cref{prop:v057-total-maslov-defect-ledger}, and let \(\nu_*\) be the restriction of \(\nu\) to
that set.  At scale \(W\), replace \(\nu_*\) by its weighted occupied product cells.  After
discarding endpoint weight \(o(\mathcal C_W)\) and making a logarithmic dyadic refinement, the cell
weights are constant on each weight layer, every weight is at most \(C_UW^3\), and there are only
\(W^{-o(1)}\) relevant layers.  The discarded endpoint weight removes only
\(o(\mathcal C_W)\) from the collision measure.

Suppose that, in every relevant weight layer, color class, and multiplicity level, the dense
alternative \textnormal{(i)} of
\cref{lem:v058-dense-shading-multiplicity} holds.  If the vertical branch bound is \(B\), then
\begin{equation}
\label{eq:v058-dense-tangential-collision-bound}
\boxed{
\int_J\iint
\mathbf 1_{\{|T_s(z,z')|\le W\}}
\,d\nu_*(z)d\nu_*(z')ds
\lesssim_\varepsilon
B W^{2-\varepsilon-o(1)}\nu_*(G).
}
\end{equation}
For the collision-preserving branch
\(B=W^{-\chi-\zeta-o(1)}\), this is
\begin{equation}
\label{eq:v058-dense-total-defect-bound}
\boxed{
\lesssim
W^{2-\chi-\varepsilon-\zeta-o(1)}.
}
\end{equation}
Since the retained full collision mass is at least
\(\frac12\mathcal C_W=W^{1+\chi+o(1)}\), every level with
\begin{equation}
\label{eq:v058-half-range-condition}
\boxed{\chi<\frac12}
\end{equation}
is impossible in the dense-shading branch: choose
\(\varepsilon,\zeta>0\) with
\(\varepsilon+\zeta<1-2\chi\).  Therefore every hypothetical bad level with
\(\chi<1/2\) is forced into the short-Reeb-shading alternative
\eqref{eq:v058-sparse-tube-shading}.
\end{proposition}

\begin{proof}
Let \(\mathscr Q_W\) be the occupied full product cells and put
\(w_{\mathfrak c}=\nu_*(\mathfrak c)\).  Bounded horizontal density gives
\[
\sum_{\mathfrak c\text{ above }A}w_{\mathfrak c}
\lesssim_U W^3
\]
for every full direction \(W\)-cell \(A\), so in particular
\(w_{\mathfrak c}\lesssim_UW^3\).  Since
\(\#\mathscr Q_W\lesssim_\zeta W^{-3-\zeta}\), weights below a sufficiently large fixed power of
\(W\) have total mass \(o(\mathcal C_W)\); choose that power after fixing the total-defect level.
Because the collision kernel is bounded by one and \(\nu(G)=1\), deleting endpoint mass \(e\)
deletes at most \(O_J(e)\) collision mass.  The remaining weights have logarithmically many dyadic
levels.

For one weight level of size \(w\), let
\[
F(u,s)=\sum_{\mathfrak c}w\mathbf 1_{\mathbb T_{\mathfrak c}}(u,s).
\]
By \cref{lem:v058-polarized-branch-coloring}, write its labeled tubes as the union of
\(J_W\lesssim BW^{-1}\) direction-separated color classes and write
\(F=\sum_{j=1}^{J_W}F_j\).  On the dense-shading branch,
\cref{lem:v058-dense-shading-multiplicity} gives
\[
\int F_j^2
\lesssim_\varepsilon
wW^{-\varepsilon}\int F_j.
\]
Pointwise Cauchy--Schwarz across the colors, the weight bound \(w\lesssim W^3\), and
\(\sum_j\int F_j\lesssim W^2\nu_*(G)\) give
\[
\int F^2
\le
J_W\sum_j\int F_j^2
\lesssim_\varepsilon
(BW^{-1})W^3W^{-\varepsilon}W^2\nu_*(G)
=
BW^{4-\varepsilon}\nu_*(G).
\]
The logarithmic sum over weight layers is absorbed by \(W^{-o(1)}\).

If \(|T_s(z,z')|\le W\), the two enlarged tangential tube cross-sections at time \(s\) have
intersection area \(\gtrsim W^2\).  After replacing points by their product-cell representatives,
which changes all radii by a fixed factor,
\[
\int_J\iint
\mathbf 1_{\{|T_s|\le W\}}d\nu_*d\nu_*'ds
\lesssim
W^{-2}\int_{E\times J}F^2.
\]
This proves \eqref{eq:v058-dense-tangential-collision-bound}.  Substitute the branch bound from
\eqref{eq:v057-total-defect-vertical-branch} to obtain
\eqref{eq:v058-dense-total-defect-bound}.  Comparing it with the retained mass
\eqref{eq:v057-collision-pruned-mass} gives the last assertion.
\end{proof}

\begin{corollary}[Mass-sensitive dense cross-collision bound]
\label{cor:v059-mass-sensitive-dense-cross-bound}
Let \(\nu_1,\nu_2\le\nu\) be two endpoint restrictions of masses
\(m_1,m_2\).  Assume their scale-\(W\) tube families obey the same vertical branch cap \(B\), and
that all relevant weight, color, and multiplicity levels lie in the dense alternative of
\cref{lem:v058-dense-shading-multiplicity}.  Then
\begin{equation}
\label{eq:v059-mass-sensitive-dense-cross-bound}
\boxed{
\int_J\iint
\mathbf 1_{\{|T_s(z,z')|\le W\}}
\,d\nu_1(z)d\nu_2(z')ds
\lesssim_\varepsilon
B W^{2-\varepsilon-o(1)}(m_1m_2)^{1/2}.
}
\end{equation}
\end{corollary}

\begin{proof}
Let \(F_1,F_2\) be the weighted tangential tube functions used in the proof of
\cref{prop:v058-dense-tangential-half-range}.  The tube-intersection comparison and
Cauchy--Schwarz give
\[
\text{left side of \eqref{eq:v059-mass-sensitive-dense-cross-bound}}
\lesssim
W^{-2}\int F_1F_2
\le
W^{-2}\|F_1\|_2\|F_2\|_2.
\]
The mass-sensitive form of the preceding proof gives
\(\|F_i\|_2^2\lesssim BW^{4-\varepsilon-o(1)}m_i\).  After replacing
\(\varepsilon\) by \(\varepsilon/2\), the displayed estimate follows.
\end{proof}

\begin{lemma}[Normal-compatible collision shading regularization]
\label{lem:v060-normal-compatible-shading-regularization}
Discretize the retained measure \(\Theta_W^*\) at scale \(W\), put
\(\nu^*=\nu|_{G_{\Theta,W}}\), and let
\(w_{\mathfrak c}=\nu^*(\mathfrak c)\) be the endpoint-cell weights.  For two labeled cells define
the measurable edge-time density
\begin{equation}
\label{eq:v060-cell-edge-time-density}
g_{\mathfrak c\mathfrak d}(s)
:=
\iint_{\mathfrak c\times\mathfrak d}
\mathbf 1_{\{|N_s(z,z')|\le W\}}
\mathbf 1_{\{|T_s(z,z')|\le W\}}
\,d\nu^*(z)d\nu^*(z').
\end{equation}
Thus
\begin{equation}
\label{eq:v060-cell-edge-time-sum}
\boxed{
\mathcal C_W^*
=
\sum_{\mathfrak c,\mathfrak d}\int_J
g_{\mathfrak c\mathfrak d}(s)\,ds,
\qquad
0\le g_{\mathfrak c\mathfrak d}(s)
\le w_{\mathfrak c}w_{\mathfrak d}.
}
\end{equation}
After deleting \(o(\mathcal C_W^*)\), there is a decomposition into
\(W^{-o(1)}\) nonnegative symmetric pieces.  One piece, of mass
\(\mathcal C_\alpha\gtrsim W^{o(1)}\mathcal C_W^*\), has two endpoint families
\(\mathscr V_{\alpha,1},\mathscr V_{\alpha,2}\) and numbers
\(m_i,h_i,\lambda_i>0\), \(i=1,2\), with the following properties:
throughout this lemma, \(\asymp_{W^{o(1)}}\) denotes two-sided comparison with losses bounded by
\(W^{-o(1)}\).
\begin{enumerate}[label=\textnormal{(\roman*)}]
\item the endpoint weight and integrated degree are biregular at the power scale,
\begin{equation}
\label{eq:v060-biregular-mass-degree}
\boxed{
\sum_{\mathfrak c\in\mathscr V_{\alpha,i}}w_{\mathfrak c}
\asymp_{W^{o(1)}}m_i,
\qquad
\int_JD_{\alpha,i,\mathfrak c}(s)ds
\asymp_{W^{o(1)}}h_i,
\qquad
m_ih_i\asymp_{W^{o(1)}}\mathcal C_\alpha;
}
\end{equation}
\item for every retained endpoint cell \(\mathfrak c\), its active time set
\[
S_{\alpha,i,\mathfrak c}
=
\{s:D_{\alpha,i,\mathfrak c}(s)>0\}
\]
has \(|S_{\alpha,i,\mathfrak c}|\asymp_{W^{o(1)}}\lambda_i\), and on a subset of comparable
measure
\begin{equation}
\label{eq:v060-instantaneous-degree}
\boxed{
D_{\alpha,i,\mathfrak c}(s)
\asymp_{W^{o(1)}}
\frac{h_i}{\lambda_i};
}
\end{equation}
\item the time-saturated tangential shading
\begin{equation}
\label{eq:v060-time-saturated-shading}
Y_{\alpha,i}(\mathbb T_{\mathfrak c})
:=
\{(u,s)\in\mathbb T_{\mathfrak c}:s\in S_{\alpha,i,\mathfrak c}\}
\end{equation}
has density \(\asymp_{W^{o(1)}}\lambda_i\) inside its tube.
\item the decomposition is simultaneous with the tangential color and multiplicity levels used in
\cref{lem:v058-dense-shading-multiplicity}: for every selected edge-time, the midpoint of the two
enlarged tangential tube cross-sections lies in fixed dyadic multiplicity levels for the two
endpoint colors.  Thus, when \(\lambda_i\ge W^\eta\), every retained tube on side \(i\) belongs to
the dense class for that selected multiplicity level.
\end{enumerate}
Here \(D_{\alpha,i,\mathfrak c}(s)\) is the total opposite-endpoint weight joined to
\(\mathfrak c\) by the selected piece at time \(s\), including the selected dyadic edge-density
factor.  In particular the dichotomy
\begin{equation}
\label{eq:v060-dense-sparse-actual-collision}
\boxed{
\lambda_1,\lambda_2\ge W^\eta
\quad\text{or}\quad
\min\{\lambda_1,\lambda_2\}<W^\eta
}
\end{equation}
is a decomposition of the actual normal-compatible collision mass, not of an enlarged unweighted
tangential overlap.
\end{lemma}

\begin{proof}
Identity \eqref{eq:v060-cell-edge-time-sum} is Tonelli's theorem, and the upper bound follows by
deleting the two indicators.  Discard endpoint cells of sufficiently small weight exactly as in
the proof of \cref{prop:v058-dense-tangential-half-range}; their total collision contribution is
\(o(\mathcal C_W^*)\).  On the remaining finite grid write
\[
r_{\mathfrak c\mathfrak d}(s)
=
\frac{g_{\mathfrak c\mathfrak d}(s)}
{w_{\mathfrak c}w_{\mathfrak d}}
\in[0,1].
\]
The part where \(r_{\mathfrak c\mathfrak d}<W^A\), with \(A\) chosen after the total-defect
level, has mass at most \(W^A|J|\) and is negligible.  Dyadically partition the two endpoint
weights and the remaining values of \(r_{\mathfrak c\mathfrak d}\).  This costs only a power of
\(\log(2/W)\).

For every nonzero edge-time, choose the midpoint of the two tangential front centers; after a fixed
enlargement it lies in both tangential tubes.  Record the dyadic multiplicities of the two endpoint
color classes at that point.  This adds only logarithmically many choices, without selecting one
of the power-many colors.  For each resulting piece define its weighted time degree at either endpoint by summing the
opposite endpoint weights times the selected \(r\)-level.  Dyadically partition the positive
degrees, then the integrals of those degrees along each endpoint, and finally the measures of their
time supports after the selected tangential multiplicity restriction.  Remove endpoint classes carrying less than a fixed logarithmic fraction of the
piece and repeat once from the opposite side.  The usual two-sided degree pruning leaves a
subpiece with a logarithmic fraction of the mass on which both endpoint degrees have the stated
lower and upper bounds.  Summing an endpoint's integrated degree against its cell weight gives the
piece mass; this proves \eqref{eq:v060-biregular-mass-degree}.  The degree and support-size
pigeonholes give \eqref{eq:v060-instantaneous-degree}.  A unit tangential \(W\)-tube has
cross-sectional area \(\asymp W^2\), so the time-saturated set
\eqref{eq:v060-time-saturated-shading} has volume
\(\asymp W^2\lambda_i\), proving its density statement.  There are only logarithmically many
choices at each stage, which is \(W^{-o(1)}\); one piece carries the asserted mass.
\end{proof}

\begin{proposition}[Reeb-time support converts collision degree into physical bush volume]
\label{prop:v062-time-support-bush-volume}
Let one piece of
\cref{lem:v060-normal-compatible-shading-regularization} have endpoint data
\((m_i,h_i,\lambda_i)\), \(i=1,2\), and suppose
\(\nu\le K\sigma\).  For either endpoint orientation and every retained endpoint cell
\(\mathfrak c\), there is a time \(s_{i,\mathfrak c}\in J\) and an actual selector-direction set
\(H_{i,\mathfrak c}\) which is a physical \(CW\)-bush at the corresponding front point and obeys
\begin{equation}
\label{eq:v062-time-support-bush-mass}
\boxed{
\sigma(H_{i,\mathfrak c})
\gtrsim
\frac{h_i}{K\lambda_i}
W^{o(1)}.
}
\end{equation}
Consequently
\begin{equation}
\label{eq:v062-time-support-volume}
\boxed{
\mathcal L^4\!\left(N_{CW}(K_b)\right)
\gtrsim
\frac{h_i}{K\lambda_i}
W^{o(1)}
\qquad(i=1,2).
}
\end{equation}
This holds for dense and sparse pieces alike; no summation over center cells and no escaped-bush
overlap estimate are required.
\end{proposition}

\begin{proof}
For every retained endpoint cell, the biregularization gives a subset of its active time set on
which
\[
D_{\alpha,i,\mathfrak c}(s)
\asymp_{W^{o(1)}}
\frac{h_i}{\lambda_i}.
\]
Choose one such time.  Since the selected edge density is at most one, the actual opposite-endpoint
occurrence weight of the corresponding neighbor set is at least the left side.  Division by the
reuse cap gives \eqref{eq:v062-time-support-bush-mass}.  The two polarized inequalities in
\eqref{eq:v056-orthogonal-front-split} place every leaf front in the same \(O(W)\)-ball; the diameter
of a phase-space cell changes the constant only.  Apply the arbitrary-mass Reeb escape formula
\cref{prop:v010-bush-volume} at radius \(CW\) to obtain
\eqref{eq:v062-time-support-volume}.
\end{proof}

\subsection{Normal shear and direct packet bounds}

\begin{proposition}[A \(q\)-local dense collision gains the missing color factor]
\label{prop:v063-q-local-dense-bound}
Let a dense piece of
\cref{lem:v060-normal-compatible-shading-regularization} have endpoint masses \(m_1,m_2\), and
suppose its collision-preserving vertical branch cap is \(B\).  Partition the normal direction
coordinate \(q\) into half-open intervals of length \(W\), and let
\(\mathcal C_{\rm loc}\) be the part of the actual normal-compatible collision mass joining cells
whose \(q\)-intervals have distance \(O(W)\).  Then
\begin{equation}
\label{eq:v063-q-local-dense-bound}
\boxed{
\mathcal C_{\rm loc}
\lesssim_\varepsilon
B\,W^{3-\varepsilon-o(1)}
(m_1m_2)^{1/2}.
}
\end{equation}
The extra factor \(W\), compared with
\eqref{eq:v059-mass-sensitive-dense-cross-bound}, comes from retaining one normal-offset slab before
coloring the tangential tubes.
\end{proposition}

\begin{proof}
For one \(q\)-interval \(I\), let \(F_{i,I}\) be the weighted tangential tube function of endpoint
side \(i\), restricted to the selected cell-weight and multiplicity layer.  Inside \(I\),
\cref{lem:v058-polarized-branch-coloring} uses \(O(B)\) colors, not the global
\(O(BW^{-1})\) colors.  On one color the dense-shading estimate and the cell-weight bound
\(w_{\mathfrak c}\lesssim W^3\) give
\[
\|F_{i,I,j}\|_2^2
\lesssim_\varepsilon
W^{3-\varepsilon-o(1)}
\int F_{i,I,j}.
\]
Since a tangential unit \(W\)-tube has volume \(\asymp W^2\),
\[
\sum_j\int F_{i,I,j}
\lesssim
W^2m_{i,I},
\]
where \(m_{i,I}\) is the endpoint mass over \(I\).  Pointwise Cauchy--Schwarz over the \(O(B)\)
colors therefore yields
\[
\|F_{i,I}\|_2^2
\lesssim_\varepsilon
B\,W^{5-\varepsilon-o(1)}m_{i,I}.
\]
For a pair \(I,I'\) at distance \(O(W)\), Cauchy--Schwarz between the two endpoint sides and the
standard tangential tube-intersection comparison give
\[
\mathcal C_{\rm loc}(I,I')
\lesssim
W^{-2}\int F_{1,I}F_{2,I'}
\lesssim_\varepsilon
B\,W^{3-\varepsilon-o(1)}
(m_{1,I}m_{2,I'})^{1/2}.
\]
Every \(I\) has only \(O(1)\) admissible \(I'\).  Sum and apply discrete Cauchy--Schwarz, using
\(\sum_I m_{i,I}=m_i\), to prove
\eqref{eq:v063-q-local-dense-bound}.  Dropping the normal indicator and the selected edge-density
factor only enlarges the expression, so the estimate applies to the actual cellwise collision
density \eqref{eq:v060-cell-edge-time-density}.
\end{proof}

\begin{corollary}[Below unit defect, every dense survivor has genuine normal shear]
\label{cor:v063-dense-survivor-normal-shear}
Let \(1/2\le\chi<1\), take
\[
B=W^{-\chi-\zeta-o(1)},
\qquad
m_i=W^{\chi-\theta_i+o(1)},
\qquad
0\le\theta_i\le\chi,
\]
and choose \(\varepsilon,\zeta>0\) with
\(\varepsilon+\zeta<2(1-\chi)\).  Then
\begin{equation}
\label{eq:v063-local-layer-too-small}
\boxed{
\mathcal C_{\rm loc}
\lesssim
W^{3-(\theta_1+\theta_2)/2-\varepsilon-\zeta-o(1)}
=
o\!\left(W^{1+\chi+o(1)}\right).
}
\end{equation}
Hence a dense retained collision layer of mass \(W^{1+\chi+o(1)}\) cannot remain in
\(|q-q'|=O(W)\).  After a logarithmic dyadic refinement, a subpower fraction of that mass lies on
\begin{equation}
\label{eq:v063-normal-offset-annulus}
\boxed{
r\le|q-q'|<2r,
\qquad
W\lesssim r\le1.
}
\end{equation}
\end{corollary}

\begin{proof}
Substitute the displayed powers in
\eqref{eq:v063-q-local-dense-bound}.  Since
\((\theta_1+\theta_2)/2\le\chi\), the exponent on the right of
\eqref{eq:v063-local-layer-too-small} is at least
\(3-\chi-\varepsilon-\zeta>1+\chi\).  This proves the strict power saving.  The complement of the
local part is covered by \(O(\log(2/W))=W^{-o(1)}\) dyadic \(q\)-separation annuli, so one annulus
retains a subpower fraction.
\end{proof}

\begin{proposition}[A normal-shear annulus labels the full collision windows]
\label{prop:v063-normal-shear-output}
Apply two-sided degree regularization once more inside the annulus
\eqref{eq:v063-normal-offset-annulus}, absorbing logarithms into \(W^{-o(1)}\).  For a retained
endpoint cell on side \(i\), let \(h_i\) be its integrated annular collision degree and
\(\lambda_i\) its active-time support length.  Its swept opposite-endpoint selector-direction set
\(L_{i,\mathfrak c}\) satisfies the fixed-angle bound
\begin{equation}
\label{eq:v063-swept-leaf-mass}
\boxed{
\sigma(L_{i,\mathfrak c})
\gtrsim
\frac{h_i}{K W}
W^{o(1)}.
}
\end{equation}
The active time set requires at least
\begin{equation}
\label{eq:v063-distinct-normal-windows}
\boxed{
N_{{\rm edge},i}
\gtrsim
\max\left\{1,\frac{\lambda_i}{W}\right\}
W^{o(1)}
}
\end{equation}
distinct full collision windows of length \(O(W)\).  In addition, every such edge between phase
offsets \((q,p)\) and \((q',p')\) has a normal phase
\begin{equation}
\label{eq:v063-normal-phase-label}
\boxed{
s_N
:=
-\frac{p-p'}{q-q'},
\qquad
|s_e-s_N|
\lesssim
\frac Wr,
}
\end{equation}
where \(s_e\) is any selected full collision time for that edge.

If
\[
h_i=W^{1+\theta_i+o(1)},
\qquad
\lambda_i=W^{\ell_i+o(1)},
\qquad
r=W^{\rho+o(1)},
\quad 0\le\rho\le1,
\]
then
\begin{equation}
\label{eq:v063-normal-shear-power-ledger}
\boxed{
\sigma(L_{i,\mathfrak c})
\gtrsim
W^{\theta_i+o(1)},
\qquad
N_{{\rm edge},i}
\gtrsim
\max\{1,W^{\ell_i-1+o(1)}\},
\qquad
|s_e-s_N|
\lesssim
W^{1-\rho+o(1)}.
}
\end{equation}
Thus the annular survivor is a family of fixed-angle contact edges with power-counted Reeb windows
and an additional normal-symplectic phase label; the latter, rather than a weaker window count, is
the new information supplied by the \(q\)-shear.
\end{proposition}

\begin{proof}
On the normalized angular shell, every full-vector collision section of one edge has length
\(O(W)\).  Fubini therefore gives
\[
h_i
\lesssim
W\,\nu(L_{i,\mathfrak c})W^{-o(1)},
\]
and division by the reuse cap proves \eqref{eq:v063-swept-leaf-mass}.  The same \(O(W)\) length bound
covers an active set of measure \(\lambda_iW^{o(1)}\); a greedy disjoint-subinterval selection proves
\eqref{eq:v063-distinct-normal-windows}.  On the annulus \(|q-q'|\asymp r\), the normal component of
the polarized contact equation is
\[
|(p-p')+s_e(q-q')|\le CW.
\]
Division by \(|q-q'|\) proves \eqref{eq:v063-normal-phase-label}.  Substitution of the three power
parameters and \(K=W^{-o(1)}\) gives \eqref{eq:v063-normal-shear-power-ledger}.
\end{proof}

\begin{proposition}[Mesoscopic \(q\)-near collisions retain the color gain]
\label{prop:v064-mesoscopic-q-near}
Under the hypotheses of \cref{prop:v063-q-local-dense-bound}, fix
\(W\le R_q\le1\), and let \(\mathcal C_{\le R_q}\) be the actual normal-compatible collision mass
on \(|q-q'|\le C R_q\).  Then
\begin{equation}
\label{eq:v064-mesoscopic-q-near}
\boxed{
\mathcal C_{\le R_q}
\lesssim_\varepsilon
B\,R_q\,W^{2-\varepsilon-o(1)}
(m_1m_2)^{1/2}.
}
\end{equation}
For \(R_q=W\) this is \eqref{eq:v063-q-local-dense-bound}; for \(R_q<1\) it retains the precise
fraction \(R_q\) of the global \(q\)-color count.
\end{proposition}

\begin{proof}
Partition the bounded \(q\)-range into half-open intervals \(I\) of length \(R_q\).  Such an
interval contains \(O(R_q/W)\) scale-\(W\) normal direction slabs.  Apply
\cref{lem:v058-polarized-branch-coloring} in each of them: the tangential tubes over \(I\) use at
most \(O(BR_q/W)\) direction-separated colors.  The one-color calculation in the proof of
\cref{prop:v063-q-local-dense-bound} is unchanged,
\[
\|F_{i,I,j}\|_2^2
\lesssim_\varepsilon
W^{5-\varepsilon-o(1)}m_{i,I,j}.
\]
Pointwise Cauchy--Schwarz over \(O(BR_q/W)\) colors gives
\[
\|F_{i,I}\|_2^2
\lesssim_\varepsilon
B R_q W^{4-\varepsilon-o(1)}m_{i,I}.
\]
Only \(O(1)\) pairs of coarse \(q\)-intervals can satisfy \(|q-q'|\le C R_q\).  Apply
Cauchy--Schwarz between the endpoint functions, multiply by the tangential intersection factor
\(W^{-2}\), and sum the coarse interval pairs.  This proves
\eqref{eq:v064-mesoscopic-q-near}.
\end{proof}

\begin{lemma}[A marked edge costs the two regularized graphon degrees]
\label{lem:v064-marked-edge-four-cycle}
Let \(0\le h\le1\) be a bipartite graphon on \(X_1\times X_2\), with row and column degrees bounded
by \(D_1,D_2\).  If \(0\le h_{\rm m}\le h\) has edge mass
\[
M_{\rm m}
:=
\iint h_{\rm m}(x,y)\,d\nu_1(x)d\nu_2(y),
\]
then the alternating four-cycle mass whose first edge is marked obeys
\begin{equation}
\label{eq:v064-marked-edge-four-cycle}
\boxed{
\int
h_{\rm m}(x_0,y_1)
h(x_2,y_1)
h(x_2,y_3)
h(x_0,y_3)
\,d\nu_1(x_0)d\nu_2(y_1)d\nu_1(x_2)d\nu_2(y_3)
\le
M_{\rm m}D_1D_2.
}
\end{equation}
Consequently cycles having at least one marked edge have mass at most
\(4M_{\rm m}D_1D_2\).
\end{lemma}

\begin{proof}
For fixed \(x_0,y_1,x_2\),
\[
\int h(x_2,y_3)h(x_0,y_3)\,d\nu_2(y_3)
\le
D_1.
\]
Integrating \(h(x_2,y_1)\) in \(x_2\) contributes at most \(D_2\).  The remaining integral is
\(M_{\rm m}\), proving \eqref{eq:v064-marked-edge-four-cycle}.  Sum over the four possible marked
positions.
\end{proof}

\begin{proposition}[Mesoscopic normal phases survive on every rank-loss cycle below \(3/5\)]
\label{prop:v064-phase-labeled-rank-loss-cycles}
Apply the arbitrary-threshold Maslov routing after passing to the two-sided degree block; equivalently
symmetrize that bipartite block on two disjoint copies.  Put
\[
\bar\theta=\frac{\theta_1+\theta_2}{2},
\qquad
M_W=W^{\chi+o(1)},
\qquad
D_i=W^{\theta_i+o(1)}.
\]
Suppose the determinant-small output has four-cycle mass \(\gtrsim M_W^4\).  Let
\(R_q=W^\alpha\), where
\begin{equation}
\label{eq:v064-alpha-gap}
\boxed{
0<\alpha<1,
\qquad
1+\alpha+\bar\theta>4\chi.
}
\end{equation}
Then, after deleting \(o(M_W^4)\) and making four logarithmic dyadic refinements, every edge of the
retained rank-loss cycles satisfies
\begin{equation}
\label{eq:v064-four-shear-annuli}
\boxed{
r_e\le|q-q'|<2r_e,
\qquad
R_q\lesssim r_e\le1,
}
\end{equation}
and carries a normal phase label
\begin{equation}
\label{eq:v064-four-normal-phases}
\boxed{
s_{N,e}
:=
-\frac{p-p'}{q-q'},
\qquad
|s_e-s_{N,e}|
\lesssim
\frac{W}{r_e}
\le
C W^{1-\alpha}.
}
\end{equation}

The condition \eqref{eq:v064-alpha-gap} can always be met when
\begin{equation}
\label{eq:v064-automatic-three-fifths}
\boxed{\frac12\le\chi<\frac35.}
\end{equation}
Indeed the dense wedge gives \(\bar\theta\ge1-\chi-o(1)\), so it is enough to choose
\[
5\chi-2<\alpha<1.
\]
\end{proposition}

\begin{proof}
Normalize \eqref{eq:v064-mesoscopic-q-near} by the universal collision-window factor \(W\).  With
\(B=W^{-\chi-\zeta-o(1)}\), \(R_q=W^\alpha\), and
\((m_1m_2)^{1/2}=W^{\chi-\bar\theta+o(1)}\), the marked \(q\)-near graphon edge mass satisfies
\[
M_{\le R_q}
\lesssim
W^{1+\alpha-\bar\theta-\varepsilon-\zeta-o(1)}.
\]
Apply \cref{lem:v064-marked-edge-four-cycle}.  Since
\(D_1D_2=W^{2\bar\theta+o(1)}\), all rank-loss cycles containing a \(q\)-near edge have mass at most
\begin{equation}
\label{eq:v064-near-edge-cycle-bound}
\boxed{
\lesssim
W^{1+\alpha+\bar\theta-\varepsilon-\zeta-o(1)}.
}
\end{equation}
Choose \(\varepsilon,\zeta>0\) smaller than the strict gap in
\eqref{eq:v064-alpha-gap}.  Then \eqref{eq:v064-near-edge-cycle-bound} is
\(o(W^{4\chi+o(1)})\), so those cycles may be deleted.  Dyadically partition the four remaining
\(q\)-separations; one four-tuple of annuli retains a subpower fraction.  Equation
\eqref{eq:v063-normal-phase-label} gives \eqref{eq:v064-four-normal-phases}.  Finally
\(\bar\theta\ge1-\chi-o(1)\) reduces the strict inequality in
\eqref{eq:v064-alpha-gap} to \(\alpha>5\chi-2+o(1)\), which admits an \(\alpha<1\) precisely below
\(\chi=3/5\).
\end{proof}

\begin{lemma}[The anchored horizontal determinant gives the packet thickness directly]
\label{lem:v066-direct-horizontal-packet-thickness}
Use the polarized-plane alternative of
\cref{prop:v046-anchored-rank-packet,prop:v050-continuous-anchored-rank-packet}, and denote its
horizontal noncollinearity threshold by \(\theta_{\rm hr}\).  Put
\[
u=a_{y_0}-a_p,\qquad
v=a_q-a_{y_0},\qquad
E=\operatorname{Span}\{u,v\}.
\]
Then every retained \(y\in Y_0\) satisfies
\begin{equation}
\label{eq:v066-direct-horizontal-packet-thickness}
\boxed{
\operatorname{dist}(a_y,a_p+E)
\le
\frac{\Delta}{|u\times v|}
\lesssim
\frac{\Delta}{\theta_{\rm hr}}.
}
\end{equation}
Consequently the actual selector directions obtained after graph-cell de-multiplication lie in the
\begin{equation}
\label{eq:v066-actual-horizontal-packet-thickness}
\boxed{
t_{\rm dir}
\lesssim
W+\frac{\Delta}{\theta_{\rm hr}}
}
\end{equation}
neighborhood of \(a_p+E\).  Unlike the full phase-space conclusion
\eqref{eq:v046-anchored-polarized-packet}, this horizontal estimate contains neither
\(\gamma^{-1}\) nor \(\upsilon^{-1}\).
\end{lemma}

\begin{proof}
For the anchored cycle \(p,y_0,q,y,p\), the determinant-small hypothesis is
\[
|\det(u,v,a_y-a_q)|\le\Delta.
\]
Moreover
\[
u\times(a_q-a_p)=u\times(u+v)=u\times v,
\]
so the noncollinearity condition in the polarized alternative gives
\(|u\times v|\ge\theta_{\rm hr}\).  The point-to-plane formula therefore yields
\[
\operatorname{dist}(a_y-a_q,E)
=
\frac{|\det(u,v,a_y-a_q)|}{|u\times v|}
\le
\frac{\Delta}{\theta_{\rm hr}}.
\]
Since \(a_q-a_p=u+v\in E\), this is exactly
\eqref{eq:v066-direct-horizontal-packet-thickness}.  The final \(O(W)\) is the graph-cell
diameter already present in the de-multiplication conclusion of
\cref{prop:v046-anchored-rank-packet}.
\end{proof}

\begin{lemma}[The contact equations give the same direct vertical thickness]
\label{lem:v067-direct-contact-packet-thickness}
Under \cref{lem:v066-direct-horizontal-packet-thickness}, suppose every retained neighbor \(y\)
is joined to the anchor \(p\) by a selected contact edge
\[
|(b_y-b_p)+s_{py}(a_y-a_p)|\le C W,
\qquad s_{py}\in J.
\]
Then
\begin{equation}
\label{eq:v067-direct-contact-packet-thickness}
\boxed{
\operatorname{dist}\!\left(
(a_y,b_y),\,
(a_p,b_p)+(E\oplus E)
\right)
\lesssim_J
W+\frac{\Delta}{\theta_{\rm hr}}.
}
\end{equation}
This full phase-space packet estimate uses only the horizontal determinant and the original contact
edge equation; it is independent of \(\gamma\), the number of separated labels, and their
Vandermonde product.
\end{lemma}

\begin{proof}
By \eqref{eq:v066-direct-horizontal-packet-thickness}, choose \(e_y\in E\) such that
\[
|a_y-a_p-e_y|
\lesssim
\Delta/\theta_{\rm hr}.
\]
Because \(E\) is linear and \(s_{py}\) stays in the fixed interval \(J\),
\[
\begin{aligned}
\operatorname{dist}(b_y-b_p,E)
&\le
|(b_y-b_p)+s_{py}(a_y-a_p)|
+|s_{py}|\,\operatorname{dist}(a_y-a_p,E)\\
&\lesssim_J
W+\Delta/\theta_{\rm hr}.
\end{aligned}
\]
Combine the horizontal and vertical bounds.
\end{proof}

\begin{proposition}[Anchored determinant-small cycles form a contact-polarized packet without a
time-separation hypothesis]
\label{prop:v067-no-time-anchored-packet}
Fix anchors \(p,q\) with \(|a_p-a_q|\ge\tau_1\), let
\[
d\lambda_{p,q}(y)=h(p,y)h(q,y)\,d\nu(y),
\qquad
T=\lambda_{p,q}(Y)>0,
\]
and let \(\Gamma\subset Y\times Y\) be a symmetric measurable family of anchored cycles of
\((\lambda_{p,q}\times\lambda_{p,q})\)-mass \(Z>0\).  Assume only that every cycle in \(\Gamma\)
has horizontal determinant at most \(\Delta\); impose no condition on its four contact labels.
For every \(0<\theta_{\rm hr}<1\), there are \(y_0\) and a measurable \(Y_0\subset Y\) with
\begin{equation}
\label{eq:v067-no-time-packet-mass}
\boxed{
\nu(Y_0)\ge\lambda_{p,q}(Y_0)
\gtrsim Z/T
}
\end{equation}
for which one of the following holds:
\begin{enumerate}[label=\textnormal{(\roman*)}]
\item its directions lie in an
  \(O(\theta_{\rm hr}/\tau_1)\)-tube about the horizontal line through \(a_p,a_q\);
\item for
  \(E=\operatorname{Span}\{a_{y_0}-a_p,a_q-a_{y_0}\}\), all graph points in \(Y_0\) lie in the
  contact-polarized phase-space packet
  \begin{equation}
  \label{eq:v067-no-time-contact-polarized-output}
  \boxed{
  N_{C(W+\Delta/\theta_{\rm hr})}
  \bigl((a_p,b_p)+(E\oplus E)\bigr).
  }
  \end{equation}
\end{enumerate}
\end{proposition}

\begin{proof}
Split \(Y\) by
\[
|(a_y-a_p)\times(a_q-a_p)|<\theta_{\rm hr}.
\]
If this set has \(\lambda_{p,q}\)-mass \(\gtrsim Z/T\), the point-to-line formula and
\(|a_q-a_p|\ge\tau_1\) give alternative \textnormal{(i)}.  Otherwise the \(\Gamma\)-edges incident
to it have product measure at most \(2T\) times its measure.  Choosing the implicit constant small
enough, deletion retains a fixed fraction of \(Z\).  Weighted average degree then fixes
\(y_0\) and a neighbor set \(Y_0\) of \(\lambda_{p,q}\)-mass \(\gtrsim Z/T\), with
\[
|(a_{y_0}-a_p)\times(a_q-a_p)|
\ge\theta_{\rm hr}.
\]
The determinant hypothesis and
\cref{lem:v066-direct-horizontal-packet-thickness} give the horizontal \(E\)-packet, while
\cref{lem:v067-direct-contact-packet-thickness} gives the vertical \(E\)-packet.  This proves
\textnormal{(ii)}.  Finally \(h(p,y)h(q,y)\le1\) gives
\eqref{eq:v067-no-time-packet-mass}.
\end{proof}

\begin{lemma}[The anchor-line azimuth has bounded marginal density]
\label{lem:v070-anchor-azimuth-density}
Fix distinct anchor directions \(a_p,a_q\in U\), put \(w=a_q-a_p\), and let
\[
\pi_{w^\perp}:\mathbb R^3\longrightarrow w^\perp
\]
be orthogonal projection.  Away from the affine line
\(L_{pq}=a_p+\mathbb Rw\), define
\[
\varphi_{pq}(a)
:=
\mathbb P\!\left(\pi_{w^\perp}(a-a_p)\right)
\in\mathbb{RP}^1.
\]
If the direction marginal has density at most \(C_U\), then every projective interval
\(\Omega\subset\mathbb{RP}^1\) of length \(\delta\le1\) satisfies
\begin{equation}
\label{eq:v070-azimuth-marginal-density}
\boxed{
\nu\{a\in U\setminus L_{pq}:\varphi_{pq}(a)\in\Omega\}
\lesssim_U \delta.
}
\end{equation}
The constant is uniform in the anchor line.
\end{lemma}

\begin{proof}
Use cylindrical coordinates about \(L_{pq}\).  Inside the fixed bounded region \(U\), the
longitudinal coordinate and radius range over bounded intervals, while the angular coordinate
belongs to a set of length \(O(\delta)\).  The cylindrical Jacobian is the radius, whose integral
is uniformly bounded.  This gives the Lebesgue-volume estimate \(O_U(\delta)\); apply the bounded
horizontal density.
\end{proof}

\begin{lemma}[The horizontal four-cycle determinant is an anchor-azimuth area]
\label{lem:v070-determinant-azimuth-identity}
For fixed anchors \(p,q\), put
\[
w=a_q-a_p,\qquad
u_y=\pi_{w^\perp}(a_y-a_p).
\]
Then every anchored cycle \(p,y,q,y',p\) satisfies the exact identity
\begin{equation}
\label{eq:v070-determinant-azimuth-identity}
\boxed{
\det(a_y-a_p,a_q-a_y,a_{y'}-a_q)
=
\det(a_y-a_p,w,a_{y'}-a_p),
}
\end{equation}
and hence
\begin{equation}
\label{eq:v070-determinant-azimuth-area}
\boxed{
\left|
\det(a_y-a_p,a_q-a_y,a_{y'}-a_q)
\right|
=
|w|\,|u_y|\,|u_{y'}|\,
\left|\sin(\varphi_{pq}(a_y)-\varphi_{pq}(a_{y'}))\right|.
}
\end{equation}
If
\[
|(a_y-a_p)\times w|,
\ |(a_{y'}-a_p)\times w|
\ge\theta_{\rm az},
\]
then determinant at most \(\Delta\) implies
\begin{equation}
\label{eq:v070-azimuth-closeness}
\boxed{
d_{\mathbb{RP}^1}
\bigl(\varphi_{pq}(a_y),\varphi_{pq}(a_{y'})\bigr)
\lesssim_U
\frac{\Delta}{\theta_{\rm az}^2}.
}
\end{equation}
\end{lemma}

\begin{proof}
Write \(a_y-a_p=u\) and \(a_{y'}-a_p=v\).  Multilinearity gives
\[
\det(u,w-u,v-w)=\det(u,w,v),
\]
which is \eqref{eq:v070-determinant-azimuth-identity}.  Components of \(u,v\) parallel to \(w\)
do not contribute, proving \eqref{eq:v070-determinant-azimuth-area}.  Moreover
\[
|u_y|=\frac{|(a_y-a_p)\times w|}{|w|},
\qquad
|u_{y'}|=\frac{|(a_{y'}-a_p)\times w|}{|w|}.
\]
The anchor directions remain in a fixed bounded chart, so
\[
|w|\,|u_y|\,|u_{y'}|
\gtrsim_U
\theta_{\rm az}^2.
\]
For a sufficiently small right side, sine is comparable to projective angular distance; for a
larger right side \eqref{eq:v070-azimuth-closeness} is trivial after changing the constant.
\end{proof}

\begin{proposition}[Direct bound for the complete anchored determinant-small graph]
\label{prop:v070-complete-anchored-determinant-bound}
Use the anchored graphon notation
\[
d\lambda_{p,q}(y)=h(p,y)h(q,y)\,d\nu(y),
\qquad
T=\lambda_{p,q}(Y),
\]
and let \(\Gamma\subset Y\times Y\) be symmetric, of
\((\lambda_{p,q}\times\lambda_{p,q})\)-mass \(Z\), with horizontal determinant at most
\(\Delta\) on every pair.  Assume \(|a_p-a_q|\ge\tau_1\).  Then, for every
\(0<\theta_{\rm az}<1\),
\begin{equation}
\label{eq:v070-complete-anchored-determinant-bound}
\boxed{
\frac ZT
\lesssim_U
\left(\frac{\theta_{\rm az}}{\tau_1}\right)^2
+\frac{\Delta}{\theta_{\rm az}^2}.
}
\end{equation}
With
\[
\theta_{\rm az}=(\Delta\tau_1^2)^{1/4},
\]
provided this lies in \((0,1)\), one obtains
\begin{equation}
\label{eq:v070-optimized-anchored-determinant-bound}
\boxed{
\frac ZT
\lesssim_U
\Delta^{1/2}\tau_1^{-1}.
}
\end{equation}
This estimate controls the entire anchored cycle graph and does not first select one
average-degree packet.
\end{proposition}

\begin{proof}
Let
\[
Y_{\rm line}
=
\{y:|(a_y-a_p)\times(a_q-a_p)|<\theta_{\rm az}\}.
\]
The point-to-line formula and \(|a_q-a_p|\ge\tau_1\) give
\[
\lambda_{p,q}(Y_{\rm line})
\le\nu(Y_{\rm line})
\lesssim_U
\left(\theta_{\rm az}/\tau_1\right)^2.
\]
Pairs with at least one endpoint in this set have product measure at most
\[
2T\,\lambda_{p,q}(Y_{\rm line}).
\]

On the complementary pair set, \cref{lem:v070-determinant-azimuth-identity} places \(y'\), for
each fixed \(y\), in a projective angular interval of length
\[
O_U(\Delta/\theta_{\rm az}^2).
\]
Since \(\lambda_{p,q}\le\nu\), \cref{lem:v070-anchor-azimuth-density} bounds the
\(\lambda_{p,q}\)-mass of this interval by the same quantity.  Integration in \(y\) therefore
bounds the far--far pair mass by
\[
C_U T\Delta/\theta_{\rm az}^2.
\]
Add the two contributions and divide by \(T\), proving
\eqref{eq:v070-complete-anchored-determinant-bound}.  The displayed parameter choice balances
\(\theta_{\rm az}^2/\tau_1^2\) with
\(\Delta/\theta_{\rm az}^2\).
\end{proof}

\subsection{Azimuth sectors and rooted contact stars}

\begin{lemma}[Sector-square form of the anchored azimuth estimate]
\label{lem:v071-azimuth-sector-square}
Retain the hypotheses and notation of
\cref{prop:v070-complete-anchored-determinant-bound}.  Put
\[
Y_{\rm line}(p,q;\theta_{\rm az})
=
\{y:|(a_y-a_p)\times(a_q-a_p)|<\theta_{\rm az}\},
\qquad
L_{p,q}=\lambda_{p,q}(Y_{\rm line}),
\]
and suppose
\[
\delta_{\rm az}:=C_U\frac{\Delta}{\theta_{\rm az}^2}\le c_U.
\]
There is a bounded-overlap family \(\mathscr I_{p,q}\) of projective intervals of length
\(O_U(\delta_{\rm az})\) such that, on writing
\[
A_{p,q,\omega}
=
\{y\notin Y_{\rm line}:\varphi_{pq}(a_y)\in\omega\},
\qquad
m_{p,q,\omega}=\lambda_{p,q}(A_{p,q,\omega}),
\]
one has
\begin{equation}
\label{eq:v071-sector-square-anchored}
\boxed{
Z(p,q)
\lesssim
T(p,q)L_{p,q}
+
\sum_{\omega\in\mathscr I_{p,q}}m_{p,q,\omega}^{2}.
}
\end{equation}
In particular, without replacing either term by a density cap, one has the following exact
alternative:
\begin{equation}
\label{eq:v071-line-or-square-alternative}
\boxed{
L_{p,q}\gtrsim Z(p,q)/T(p,q)
\quad\hbox{or}\quad
\sum_{\omega\in\mathscr I_{p,q}}m_{p,q,\omega}^{2}
\gtrsim Z(p,q).
}
\end{equation}
The first output is the horizontal anchor-line packet itself; the second retains all azimuth
packet occurrences and their quadratic masses.
\end{lemma}

\begin{proof}
Partition \(\mathbb{RP}^1\) into half-open intervals of length comparable to
\(\delta_{\rm az}\), and enlarge each interval by a fixed number of neighbors.  The enlarged
family has bounded overlap.  By \eqref{eq:v070-azimuth-closeness}, every determinant-small pair
with both endpoints outside \(Y_{\rm line}\) belongs to one enlarged interval.  Its product mass
is therefore at most a constant multiple of
\(\sum_\omega m_{p,q,\omega}^2\).  Pairs with at least one endpoint in
\(Y_{\rm line}\) have mass at most \(2T(p,q)L_{p,q}\).  This proves
\eqref{eq:v071-sector-square-anchored}; splitting according to which displayed contribution is at
least a fixed fraction of \(Z(p,q)\) proves
\eqref{eq:v071-line-or-square-alternative}.
\end{proof}

\begin{lemma}[Every azimuth sector is a contact-polarized packet]
\label{lem:v071-sector-contact-packet}
For \(\omega\in\mathscr I_{p,q}\), choose a representative projective ray
\(e_\omega\subset(a_q-a_p)^\perp\), and set
\[
E_{p,q,\omega}
=
\Span\{a_q-a_p,e_\omega\}\subset\mathbb R^3.
\]
Then
\begin{equation}
\label{eq:v071-sector-horizontal-packet}
\operatorname{dist}
\bigl(a_y,a_p+E_{p,q,\omega}\bigr)
\lesssim_U\delta_{\rm az}
\qquad(y\in A_{p,q,\omega}).
\end{equation}
If, as supplied by the normalized contact graphon, every retained edge \((p,y)\) has a selected
Reeb label \(s_{py}\in J\) satisfying
\[
|(b_y-b_p)+s_{py}(a_y-a_p)|\le C W,
\]
then the same sector obeys the full phase-space estimate
\begin{equation}
\label{eq:v071-sector-contact-packet}
\boxed{
\operatorname{dist}\!\left(
(a_y,b_y),
(a_p,b_p)+(E_{p,q,\omega}\oplus E_{p,q,\omega})
\right)
\lesssim_{U,J}W+\delta_{\rm az}.
}
\end{equation}
Thus every summand \(m_{p,q,\omega}^2\) in
\eqref{eq:v071-sector-square-anchored} is the pair mass of one explicitly identified
contact--symplectic polarization occurrence.
\end{lemma}

\begin{proof}
The component of \(a_y-a_p\) parallel to \(a_q-a_p\) already lies in
\(E_{p,q,\omega}\).  Its orthogonal component has projective direction within
\(O_U(\delta_{\rm az})\) of \(e_\omega\), and its length is uniformly bounded because the
direction chart is bounded.  This proves \eqref{eq:v071-sector-horizontal-packet}.  Since
\(E_{p,q,\omega}\) is linear and \(s_{py}\) remains in the fixed interval \(J\),
\[
\begin{aligned}
\operatorname{dist}(b_y-b_p,E_{p,q,\omega})
&\le |(b_y-b_p)+s_{py}(a_y-a_p)|\\
&\quad+|s_{py}|\,
\operatorname{dist}(a_y-a_p,E_{p,q,\omega})\\
&\lesssim_{U,J}W+\delta_{\rm az}.
\end{aligned}
\]
Combine the horizontal and vertical estimates.  This is precisely the contact lift used in
\cref{lem:v067-direct-contact-packet-thickness}, now applied sector by sector rather than after an
average-degree selection.
\end{proof}

\begin{lemma}[The sector-square ledger has an exact rooted contact load]
\label{lem:v072-rooted-sector-load}
Let \(0\le h\le1\) be the biregular bipartite Maslov graphon on
\(X_1\times X_2\), with row and column degrees bounded by \(D_1,D_2\).  Anchor on
\(p,q\in X_1\), and use the sector sets of
\cref{lem:v071-azimuth-sector-square}.  Thus
\[
d\lambda_{p,q}(y)=h(p,y)h(q,y)\,d\nu_2(y),
\qquad
m_{p,q,\omega}=\lambda_{p,q}(A_{p,q,\omega}).
\]
Define the rooted sector load on \(X_2\) by
\begin{equation}
\label{eq:v072-rooted-sector-load}
\mathcal R_2(y)
:=
\iint_{X_1\times X_1}
\sum_{\omega\in\mathscr I_{p,q}}
\mathbf 1_{A_{p,q,\omega}}(y)
h(p,y)h(q,y)m_{p,q,\omega}
\,d\nu_1(p)d\nu_1(q).
\end{equation}
Then the sector-square mass is rooted exactly:
\begin{equation}
\label{eq:v072-rooted-square-identity}
\boxed{
\int_{X_2}\mathcal R_2(y)\,d\nu_2(y)
=
\iint_{X_1\times X_1}
\sum_{\omega\in\mathscr I_{p,q}}m_{p,q,\omega}^{2}
\,d\nu_1(p)d\nu_1(q)
=
\mathfrak S_{\rm az}.
}
\end{equation}
Moreover, if
\[
d_2(y):=\int_{X_1}h(x,y)\,d\nu_1(x),
\]
then one has the pointwise weighted-reuse estimate
\begin{equation}
\label{eq:v072-root-load-degree-square}
\boxed{
\mathcal R_2(y)
\lesssim
D_1d_2(y)^2
\le
D_1D_2^2.
}
\end{equation}
The symmetric assertion holds after interchanging the two endpoint sides.
\end{lemma}

\begin{proof}
The first identity is Tonelli:
\[
\int_{X_2}
\mathbf 1_{A_{p,q,\omega}}(y)h(p,y)h(q,y)\,d\nu_2(y)
=m_{p,q,\omega}.
\]
For the pointwise estimate, put
\[
c(p,q):=\int_{X_2}h(p,y')h(q,y')\,d\nu_2(y').
\]
For fixed \(p,q,y\), bounded overlap of the enlarged azimuth intervals gives
\[
\sum_\omega
\mathbf 1_{A_{p,q,\omega}}(y)m_{p,q,\omega}
\lesssim c(p,q).
\]
Also \(c(p,q)\le\min\{\deg(p),\deg(q)\}\le D_1\).  Hence
\[
\mathcal R_2(y)
\lesssim
D_1
\iint h(p,y)h(q,y)\,d\nu_1(p)d\nu_1(q)
=D_1d_2(y)^2,
\]
which proves \eqref{eq:v072-root-load-degree-square}.
\end{proof}

\begin{proposition}[Occurrence-dependent lightness is low rooted load or a common contact star]
\label{prop:v072-conditional-light-to-star}
For every occurrence \(Q=(p,q,\omega)\), let
\(B_Q\subset A_{p,q,\omega}\) be an arbitrary measurable subset; in particular \(B_Q\) may be the
part of the packet lying in a \(q\)-slab declared light only after conditioning on \(Q\).  Define its
weighted paired content by
\begin{equation}
\label{eq:v072-conditioned-occurrence-content}
\mathcal L_2(\mathscr B)
:=
\iint_{X_1\times X_1}
\sum_{\omega\in\mathscr I_{p,q}}
m_{p,q,\omega}\lambda_{p,q}(B_Q)
\,d\nu_1(p)d\nu_1(q).
\end{equation}
Let \(m_2=\nu_2(\supp d_2)\).  For every \(K>0\),
\begin{equation}
\label{eq:v072-rooted-low-high-split}
\boxed{
\mathcal L_2(\mathscr B)
\le
K m_2
+
\int_{\{\mathcal R_2>K\}}\mathcal R_2(y)\,d\nu_2(y).
}
\end{equation}
Every root in the second term has normalized graphon degree
\begin{equation}
\label{eq:v072-high-root-degree}
\boxed{
d_2(y)
\gtrsim
\left(\frac K{D_1}\right)^{1/2}.
}
\end{equation}
More importantly, the sector label is not lost.  For every such root there are anchors \(p,q\)
and a sector \(\omega\) with \(y\in A_{p,q,\omega}\),
\(h(p,y)h(q,y)>0\), and
\begin{equation}
\label{eq:v072-high-root-flagged-packet}
\boxed{
m_{p,q,\omega}
\gtrsim
\frac{\mathcal R_2(y)}{d_2(y)^2}
\gtrsim
\frac K{D_2^2}.
}
\end{equation}
Thus a high root is a common contact star carrying one actual anchor pair together with a
same-azimuth \(E_{p,q,\omega}\oplus E_{p,q,\omega}\) packet; it is not merely a vertex of large
unlabeled degree.
For the normalized contact graphon \(h_W\), such a root therefore produces, by the original
Reeb-label averaging, a one-time common contact star of occurrence weight
\begin{equation}
\label{eq:v072-high-root-one-time-star}
\boxed{
\gtrsim_J
W\left(\frac K{D_1}\right)^{1/2},
}
\end{equation}
and a swept neighbor set of occurrence mass
\begin{equation}
\label{eq:v072-high-root-swept-star}
\boxed{
\gtrsim
\left(\frac K{D_1}\right)^{1/2}.
}
\end{equation}
After the existing occurrence de-multiplication, these are precisely the common-contact-star inputs
to the Reeb-time/Lagrangian-bush alternative of
\cref{prop:v057-total-defect-contact-star,prop:v062-time-support-bush-volume}.
\end{proposition}

\begin{proof}
Root \eqref{eq:v072-conditioned-occurrence-content} at its restricted endpoint \(y\).  The resulting
density \(\mathcal R_{2,\mathscr B}(y)\) satisfies
\[
0\le\mathcal R_{2,\mathscr B}(y)\le\mathcal R_2(y)
\]
because \(B_Q\subset A_{p,q,\omega}\).  On
\(\{\mathcal R_2\le K\}\), its integral is at most \(Km_2\); on the complement it is at most the
second term of \eqref{eq:v072-rooted-low-high-split}.  This proves that formula.
Equation \eqref{eq:v072-high-root-degree} follows immediately from
\eqref{eq:v072-root-load-degree-square}.
For the flagged assertion, the total anchor-sector weight at a fixed root is bounded by
\[
\iint
\sum_\omega
\mathbf 1_{A_{p,q,\omega}}(y)h(p,y)h(q,y)
\,d\nu_1(p)d\nu_1(q)
\lesssim d_2(y)^2
\]
because the enlarged sector family has bounded overlap.  The definition
\eqref{eq:v072-rooted-sector-load} is the weighted average of
\(m_{p,q,\omega}\) against this measure.  Some contributing anchor-sector flag therefore satisfies
\[
m_{p,q,\omega}\gtrsim\mathcal R_2(y)/d_2(y)^2.
\]
Use \(\mathcal R_2(y)>K\) and \(d_2(y)\le D_2\).

For \(h_W\), undoing the normalization by the universal collision-window length \(W\) makes the
full collision-time degree comparable to \(Wd_2(y)\).  Averaging that degree over the fixed Reeb
interval gives a time at which the simultaneous contact neighbor mass is
\(\gtrsim_JWd_2(y)\).  Conversely, each swept neighbor contributes to the collision-time integral
for at most \(O(W)\), so its union has mass \(\gtrsim d_2(y)\).  Substitute
\eqref{eq:v072-high-root-degree}.  This is the pointwise form of the argument in
\cref{prop:v057-total-defect-contact-star}; the cited de-multiplication and bush map supply the
final routing statement.
\end{proof}

\begin{corollary}[Rank-loss mass forces high rooted contact stars in every non-line orientation]
\label{cor:v072-rank-loss-high-root-stars}
On a total-defect biregular layer, write
\[
M=W^{\chi+o(1)},\qquad
m_i=W^{\chi-\theta_i+o(1)},\qquad
D_i=W^{\theta_i+o(1)}.
\]
Suppose that, after anchoring on the side opposite \(X_i\), an occurrence-dependent subfamily of
the non-line sector ledger retains content \(\gtrsim M^4\).  Fix \(\eta>0\) and set
\[
K_i=W^{3\chi+\theta_i+\eta+o(1)}.
\]
Then its low-root part on side \(i\) is
\begin{equation}
\label{eq:v072-low-root-negligible}
\boxed{
K_im_i
=W^{4\chi+\eta+o(1)}
=o(M^4).
}
\end{equation}
Consequently the surviving high-root part returns, in every endpoint orientation satisfying this
non-line hypothesis, to a normalized common-contact star with the corresponding degree bound
\begin{equation}
\label{eq:v072-high-root-star-powers}
\boxed{
\begin{aligned}
d_1&\gtrsim
W^{(3\chi+\theta_1-\theta_2+\eta)/2+o(1)},\\
d_2&\gtrsim
W^{(3\chi+\theta_2-\theta_1+\eta)/2+o(1)}.
\end{aligned}
}
\end{equation}
The associated one-time star weights are the displayed quantities multiplied by \(W\); their
swept leaf masses have no additional \(W\)-loss.  In the same orientations, the retained flagged
contact packets satisfy
\begin{equation}
\label{eq:v072-flagged-packet-powers}
\boxed{
\begin{aligned}
m_{\mathrm{flag},1}&\gtrsim
W^{3\chi-\theta_1+\eta+o(1)},\\
m_{\mathrm{flag},2}&\gtrsim
W^{3\chi-\theta_2+\eta+o(1)}.
\end{aligned}
}
\end{equation}
\end{corollary}

\begin{proof}
Apply \cref{prop:v072-conditional-light-to-star} with \(X_2\) rooted and \(K=K_2\).
Equation \eqref{eq:v072-low-root-negligible} shows that the low-root term cannot carry a fixed
portion of \(M^4\), so \eqref{eq:v072-high-root-degree} gives the second line of
\eqref{eq:v072-high-root-star-powers}.  If the reversed anchoring also remains in its non-line
sector branch, repeat with \(K=K_1\) to obtain the first line.  If it does not, that portion has
already entered the anchor-line angular-collapse branch.  The time and swept-mass statements are
\eqref{eq:v072-high-root-one-time-star}--\eqref{eq:v072-high-root-swept-star}.
Finally \eqref{eq:v072-high-root-flagged-packet} gives
\(m_{\mathrm{flag},i}\gtrsim K_i/D_i^2\), which is
\eqref{eq:v072-flagged-packet-powers}.
\end{proof}

\begin{lemma}[A high root retains a weighted family of Maslov flags]
\label{lem:v073-high-root-flag-layer}
Use the \(X_2\)-rooted notation of \cref{lem:v072-rooted-sector-load}.  For fixed \(y\), put
\[
d\eta_y(p,q,\omega)
:=
\mathbf 1_{A_{p,q,\omega}}(y)
h(p,y)h(q,y)\,d\nu_1(p)d\nu_1(q),
\]
where the bounded-overlap sector index is counted as part of the measure space.  Then
\begin{equation}
\label{eq:v073-root-flag-first-moment}
\int m_{p,q,\omega}\,d\eta_y(p,q,\omega)
=\mathcal R_2(y),
\qquad
\|\eta_y\|\lesssim d_2(y)^2.
\end{equation}
If \(\mathcal R_2(y)>0\), there is a dyadic number \(H_y\) satisfying
\begin{equation}
\label{eq:v073-flag-layer-scale}
\boxed{
\frac{\mathcal R_2(y)}{C d_2(y)^2}
\lesssim H_y\le D_1
}
\end{equation}
and a flag family
\[
\mathscr F_y(H_y)
:=
\{(p,q,\omega):H_y\le m_{p,q,\omega}<2H_y\}
\]
such that
\begin{equation}
\label{eq:v073-flag-layer-product}
\boxed{
H_y\,\eta_y(\mathscr F_y(H_y))
\gtrsim
\frac{\mathcal R_2(y)}
{1+\log\!\left(2+D_1d_2(y)^2/\mathcal R_2(y)\right)}.
}
\end{equation}
Every flag in this family consists of the same root \(y\), two actual contact neighbors \(p,q\),
and an \(H_y\)-mass common-neighbor set lying in the corresponding
\(E_{p,q,\omega}\oplus E_{p,q,\omega}\) contact-polarized packet.  The symmetric statement holds
with the endpoint sides reversed.
\end{lemma}

\begin{proof}
The first equality is the definition \eqref{eq:v072-rooted-sector-load}; bounded overlap of
\(\mathscr I_{p,q}\) gives the estimate for \(\|\eta_y\|\).  Since
\(m_{p,q,\omega}\le c(p,q)\le D_1\), all packet masses lie below \(D_1\).  Set
\[
H_0=\frac{\mathcal R_2(y)}{2C d_2(y)^2},
\]
with \(C\) larger than the implicit constant in
\eqref{eq:v073-root-flag-first-moment}.  The flags with packet mass below \(H_0\) contribute at most
\(\mathcal R_2(y)/2\).  The interval \([H_0,D_1]\) has
\[
O\!\left(1+\log(2+D_1/H_0)\right)
\]
dyadic layers.  One layer retains the reciprocal of this logarithm of the remaining first moment.
On that layer \(m_{p,q,\omega}<2H_y\), so its contribution is at most
\(2H_y\eta_y(\mathscr F_y(H_y))\).  This proves
\eqref{eq:v073-flag-layer-scale}--\eqref{eq:v073-flag-layer-product}.
\end{proof}

\begin{corollary}[The rooted \(M^4\) output enters the original graph-cell reuse ledger]
\label{cor:v073-rooted-to-anchor-reuse}
Suppose a high-root portion of the non-line sector ledger carries a fixed part of \(M^4\).  After
dyadic decomposition of the root load, the flag mass \(H_y\), and the original graphon edge
weights, there is a layer losing only \(W^{-o(1)}\) on which
\begin{equation}
\label{eq:v073-weighted-anchor-occurrence}
\boxed{
\int H_y\,\eta_y(\mathscr F_y(H_y))\,d\nu_i(y)
\gtrsim
M^4W^{o(1)}.
}
\end{equation}
After the manuscript's scale-\(W\) graph-cell de-multiplication, the measure
\(\eta_y(\mathscr F_y(H_y))\) is the anchor-pair occurrence weight and \(H_y\) is the retained
endpoint-packet mass.  Hence \eqref{eq:v073-weighted-anchor-occurrence} is the weighted continuous
counterpart of the input \(S\) in
\cref{prop:v046-anchor-cell-reuse}; it retains the root, both contact edges, the azimuth sector,
and the polarization plane.
\end{corollary}

\begin{proof}
Integrate \eqref{eq:v073-flag-layer-product} over the high-root set.  All logarithms are
\(W^{-o(1)}\) on a fixed power layer.  The root loads and \(H_y\) occupy only
\(W^{-o(1)}\) dyadic classes, and the graphon construction already includes the edge-weight
regularization.  Pigeonhole these classes.  Scale-\(W\) de-multiplication changes the resulting
occurrence ledger only by the cell factors already recorded in
\cref{prop:v046-anchor-cell-reuse}; no new normalization of a packet or root is performed.
\end{proof}

\begin{proposition}[Weighted continuous high-reuse contact-star trichotomy]
\label{prop:v074-continuous-high-reuse-trichotomy}
Fix a root \(y\in X_2\), and let
\(\mathscr F\subset\mathscr F_y(H)\) be a symmetric measurable flag family from
\cref{lem:v073-high-root-flag-layer}.  Put
\[
d\mu_y(p):=h(p,y)\,d\nu_1(p),
\qquad
g_y(p,q):=
\sum_\omega\mathbf 1_{\mathscr F}(p,q,\omega),
\qquad
P:=\iint g_y(p,q)\,d\mu_y(p)d\mu_y(q).
\]
The bounded-overlap sector construction gives \(0\le g_y\lesssim1\).  Assume the normalized
angular shell supplies
\[
\tau_0\le |a_p-a_y|\le C_U
\]
for all participating contact neighbors.  For every \(B\ge1\) at least one of the following holds.
\begin{enumerate}[label=\textnormal{(\roman*)}]
\item The \(\mu_y\)-mass \(S\) of the neighbor support used by \(g_y\) satisfies
  \begin{equation}
  \label{eq:v074-amplified-neighbor-support}
  \boxed{S\gtrsim BP^{1/2}.}
  \end{equation}
\item There are one contact neighbor \(p_0\) and a measurable fiber \(Q\) of contact neighbors with
  \begin{equation}
  \label{eq:v074-large-anchor-fiber}
  \boxed{\mu_y(Q)\gtrsim B^{-1}P^{1/2}.}
  \end{equation}
  For every \(W\ll\zeta\le\kappa\ll1\), that fiber has at least one of:
  \begin{enumerate}[label=\textnormal{(\alph*)}]
  \item an angular subset \(Q_{\rm ang}\subset Q\) of mass
    \(\mu_y(Q_{\rm ang})\gtrsim\mu_y(Q)\), whose projective directions
    \(a_q-a_y\) lie in one \(O_{\tau_0}(\kappa)\)-cap;
  \item a subset \(Q_{\rm nd}\subset Q\) of comparable mass for which
    \begin{equation}
    \label{eq:v074-nondegenerate-substar-pairs}
    \boxed{
    (\mu_y\times\mu_y)
    \left\{
    (q_0,q_1)\in Q_{\rm nd}^2:
    |\det(a_{p_0}-a_y,a_{q_0}-a_y,a_{q_1}-a_y)|
    \ge\zeta
    \right\}
    \gtrsim\mu_y(Q)^2;
    }
    \end{equation}
    this is the positive pair-mass tetrahedral input to the four-cluster Reeb-front seed;
  \item a subset \(Q_{\rm pol}\subset Q\) of mass
    \(\mu_y(Q_{\rm pol})\gtrsim\mu_y(Q)\) and a horizontal plane
    \(E\subset\mathbb R^3\) such that
    \begin{equation}
    \label{eq:v074-continuous-contact-plane}
    \boxed{
    \begin{aligned}
    \operatorname{dist}(a_q-a_y,E)&\lesssim\zeta/\kappa,\\
    \operatorname{dist}(b_q-b_y,E)&\lesssim_J W+\zeta/\kappa
    \end{aligned}
    \qquad(q\in Q_{\rm pol}).
    }
    \end{equation}
    Thus this entire substar lies in one
    \(O_J(W+\zeta/\kappa)\)-neighborhood of
    \(z_y+(E\oplus E)\).
  \end{enumerate}
\end{enumerate}
Every retained pair \((p,q)\) still carries an azimuth-sector common-neighbor packet of mass
\(\asymp H\); none of the alternatives forgets that Maslov flag.
\end{proposition}

\begin{proof}
Let
\[
N=\left\{p:\int g_y(p,q)\,d\mu_y(q)>0\right\},
\qquad S=\mu_y(N).
\]
If \eqref{eq:v074-amplified-neighbor-support} fails, averaging the fiber degree over \(N\) gives
some \(p_0\in N\) for which
\[
\int g_y(p_0,q)\,d\mu_y(q)
\gtrsim P/S
\gtrsim B^{-1}P^{1/2}.
\]
Because \(g_y\lesssim1\), its underlying neighbor fiber has the mass in
\eqref{eq:v074-large-anchor-fiber}.

Put \(\alpha_q=a_q-a_y\).  If a fixed fraction of this fiber satisfies
\[
|\alpha_{p_0}\times\alpha_q|<\kappa,
\]
the fixed lower radial bounds place that part in one projective
\(O_{\tau_0}(\kappa)\)-cap, giving (a).  Otherwise retain
\(Q_{\rm far}\) of comparable mass on which
\(|\alpha_{p_0}\times\alpha_q|\ge\kappa\).
If the determinant-large pair set in \(Q_{\rm far}^2\) has a fixed fraction of the product mass,
take \(Q_{\rm nd}=Q_{\rm far}\) and obtain
\eqref{eq:v074-nondegenerate-substar-pairs}.

In the remaining case, Fubini selects \(q_0\in Q_{\rm far}\) such that a comparable-mass set
\(Q_{\rm pol}\subset Q_{\rm far}\) satisfies
\[
|\det(\alpha_{p_0},\alpha_{q_0},\alpha_q)|<\zeta.
\]
For \(E=\Span\{\alpha_{p_0},\alpha_{q_0}\}\), the point-to-plane formula and
\(|\alpha_{p_0}\times\alpha_{q_0}|\ge\kappa\) give the first line of
\eqref{eq:v074-continuous-contact-plane}.  Each edge \(yq\) has a selected Reeb label
\(s_{yq}\in J\) with
\[
|(b_q-b_y)+s_{yq}(a_q-a_y)|\lesssim W.
\]
Since \(E\) is linear and \(J\) is bounded, the vertical estimate follows.  The packet flag of
each pair was part of \(\mathscr F\) throughout.
\end{proof}

\begin{corollary}[Power ledger for the continuous flagged trichotomy]
\label{cor:v074-continuous-flag-power-ledger}
Root on side \(i\), with the anchors on side \(j\), and use the high-root threshold
\[
K_i=W^{3\chi+\theta_i+\eta+o(1)}
\]
from \cref{cor:v072-rank-loss-high-root-stars}.  On a dyadic flag layer write
\[
H_i=W^{\phi_i+o(1)}.
\]
Then
\begin{equation}
\label{eq:v074-flag-exponent-range}
\boxed{
\theta_j
\le\phi_i
\le3\chi-\theta_i+\eta,
}
\end{equation}
and the corresponding anchor-pair mass \(P_i\) obeys
\begin{equation}
\label{eq:v074-flag-product-power}
\boxed{
H_iP_i
\gtrsim W^{3\chi+\theta_i+\eta+o(1)},
\qquad
P_i
\gtrsim W^{3\chi+\theta_i-\phi_i+\eta+o(1)}.
}
\end{equation}
Consequently \cref{prop:v074-continuous-high-reuse-trichotomy} gives either a neighbor-support
star of mass \(\gtrsim BP_i^{1/2}\), or a fixed-anchor substar of mass
\begin{equation}
\label{eq:v074-fiber-star-power}
\boxed{
q_i
\gtrsim
B^{-1}
W^{(3\chi+\theta_i-\phi_i+\eta)/2+o(1)}
}
\end{equation}
which is angular, tetrahedrally nondegenerate, or contact-polarized as in
\eqref{eq:v074-continuous-contact-plane}.
\end{corollary}

\begin{proof}
The lower bound in \eqref{eq:v073-flag-layer-scale}, with
\(\mathcal R_i\ge K_i\) and \(d_i\le D_i\), gives
\(H_i\gtrsim K_i/D_i^2=W^{3\chi-\theta_i+\eta+o(1)}\), which is the upper bound for
\(\phi_i\).  The common-neighbor mass is at most the degree \(D_j\) of either anchor, giving
\(\phi_i\ge\theta_j\).  Equation \eqref{eq:v073-flag-layer-product} gives
\eqref{eq:v074-flag-product-power} after absorbing its logarithm into \(W^{o(1)}\).
Apply \eqref{eq:v074-large-anchor-fiber}.
\end{proof}

\begin{lemma}[Continuous Reeb-time synchronization of a contact-neighbor star]
\label{lem:v075-continuous-star-time-bush}
Let \(N\subset X_j\) be a contact-neighbor set of a root \(y\in X_i\), measured by
\[
d\mu_y(x)=h_W(x,y)\,d\nu_j(x),
\qquad
\mu_y(N)=S.
\]
For every \(W\le\delta_s\le1\), there are a time \(s_0\in J\) and a subset
\(N_{s_0}\subset N\) such that
\begin{equation}
\label{eq:v075-time-bush-mass}
\boxed{
\nu_j(N_{s_0})
\ge\mu_y(N_{s_0})
\gtrsim_J\delta_sS,
}
\end{equation}
and all its selected lines form an
\(O_J(W+\delta_s)\)-Lagrangian point-probe bush centered at the front point
\((b_y+s_0a_y,s_0)\).
\end{lemma}

\begin{proof}
Choose measurably one contact label \(s_{xy}\in J\) for every edge with \(h_W(x,y)>0\).
Partition \(J\) into \(O_J(\delta_s^{-1})\) half-open intervals and retain one containing
\(\mu_y\)-mass \(\gtrsim_J\delta_sS\).  If \(s_0\) is its midpoint, then
\[
|(b_x-b_y)+s_0(a_x-a_y)|
\lesssim_J W+\delta_s
\]
on the retained set.  This is exactly the asserted point-probe bush, and
\(\mu_y\le\nu_j\) proves \eqref{eq:v075-time-bush-mass}.
\end{proof}

\begin{lemma}[A continuous contact-plane substar returns to one angular cap]
\label{lem:v075-continuous-plane-angular}
Let \(Q\) be the polarized substar in
\eqref{eq:v074-continuous-contact-plane}, of \(\mu_y\)-mass \(q\), and assume
\[
W+\zeta/\kappa\le\varpi\ll1.
\]
Then one projective cap of aperture \(O_{U,\tau_0}(\varpi)\) contains a contact-neighbor subset of
mass
\begin{equation}
\label{eq:v075-plane-angular-mass}
\boxed{
\mu_y(Q_{\rm cap})
\gtrsim_{U,\tau_0}
\varpi q.
}
\end{equation}
\end{lemma}

\begin{proof}
After radial normalization, the first line of
\eqref{eq:v074-continuous-contact-plane} places all directions in an
\(O_{U,\tau_0}(\varpi)\)-neighborhood of
\(\mathbb P(E)\subset\mathbb{RP}^2\).  Cover this projective circle by
\(O_{U,\tau_0}(\varpi^{-1})\) caps of aperture \(O_{U,\tau_0}(\varpi)\), using half-open cells to
avoid overcounting, and pigeonhole with the measure \(\mu_y\).
\end{proof}

\begin{proposition}[Adaptive continuous routing of every high-reuse root]
\label{prop:v075-adaptive-continuous-routing}
In \cref{prop:v074-continuous-high-reuse-trichotomy}, fix a desired bush mass
\(0<\mathfrak b\le1\), a time width \(W\le\delta_s\ll1\), and parameters
\[
W\ll\zeta\le\kappa\ll1,
\qquad
W+\zeta/\kappa\le\varpi\ll1.
\]
Set
\begin{equation}
\label{eq:v075-adaptive-time-amplification}
\boxed{
B_{\rm time}
:=
\max\left\{
1,\frac{\mathfrak b}{\delta_sP^{1/2}}
\right\}.
}
\end{equation}
Then every high-reuse root has at least one of:
\begin{enumerate}[label=\textnormal{(\roman*)}]
\item an \(O_J(W+\delta_s)\)-Lagrangian bush of actual direction mass
  \begin{equation}
  \label{eq:v075-target-bush}
  \boxed{\gtrsim_J\mathfrak b;}
  \end{equation}
\item a projective angular cap of aperture
  \(O_{U,\tau_0}(\max\{\kappa,\varpi\})\) carrying contact-neighbor mass
  \begin{equation}
  \label{eq:v075-residual-angular-mass}
  \boxed{
  \gtrsim
  \varpi B_{\rm time}^{-1}P^{1/2};
  }
  \end{equation}
\item a tetrahedrally nondegenerate substar whose pair mass is
  \begin{equation}
  \label{eq:v075-residual-seed-mass}
  \boxed{
  \gtrsim
  B_{\rm time}^{-2}P.
  }
  \end{equation}
\end{enumerate}
Thus the continuous high-reuse branch has exactly the three exits prescribed by the original
contact--symplectic proof route; the contact-plane output is not a fourth remainder.
\end{proposition}

\begin{proof}
Apply \cref{prop:v074-continuous-high-reuse-trichotomy} with \(B=B_{\rm time}\).
In its large-support alternative,
\cref{lem:v075-continuous-star-time-bush} gives bush mass
\[
\gtrsim_J
\delta_sB_{\rm time}P^{1/2}
\ge\mathfrak b.
\]
In the fixed-anchor alternative, the fiber mass is
\(q\gtrsim B_{\rm time}^{-1}P^{1/2}\).  Its direct angular branch gives at least this mass in an
\(O(\kappa)\)-cap.  Its contact-plane branch gives
\(\gtrsim\varpi q\) in an \(O(\varpi)\)-cap by
\cref{lem:v075-continuous-plane-angular}.  The weaker common bound is
\eqref{eq:v075-residual-angular-mass}.  Finally
\eqref{eq:v074-nondegenerate-substar-pairs} gives pair mass
\(\gtrsim q^2\gtrsim B_{\rm time}^{-2}P\), proving
\eqref{eq:v075-residual-seed-mass}.
\end{proof}

\begin{corollary}[Power form at the rank-loss bush scale]
\label{cor:v075-rank-loss-routing-powers}
In the setting of \cref{cor:v074-continuous-flag-power-ledger}, take
\(\mathfrak b=M^3=W^{3\chi+o(1)}\), write
\(\delta_s=W^{u+o(1)}\), and suppress subpower angular factors.  If
\(B_{\rm time}>1\), the two non-bush outputs in
\cref{prop:v075-adaptive-continuous-routing} have respectively
\begin{equation}
\label{eq:v075-angular-seed-power-ledger}
\boxed{
\begin{aligned}
\text{angular-star mass}
&\gtrsim
W^{u+\theta_i-\phi_i+o(1)},\\
\text{tetrahedral pair mass}
&\gtrsim
W^{2u+2\theta_i-2\phi_i+o(1)}.
\end{aligned}
}
\end{equation}
If \(B_{\rm time}=1\), both lower bounds improve to the corresponding powers of
\(P_i^{1/2}\) and \(P_i\) from \eqref{eq:v074-flag-product-power}.
\end{corollary}

\begin{proof}
When \(B_{\rm time}>1\),
\[
B_{\rm time}^{-1}P_i^{1/2}
=
\frac{\delta_sP_i}{M^3}.
\]
Substitute
\(P_i\gtrsim W^{3\chi+\theta_i-\phi_i+o(1)}\) from
\eqref{eq:v074-flag-product-power}; square for the seed term.  If \(B_{\rm time}=1\), use
\eqref{eq:v075-residual-angular-mass}--\eqref{eq:v075-residual-seed-mass} directly.
\end{proof}

\begin{lemma}[A surviving rooted load occupies a quantitative root set]
\label{lem:v076-high-root-support}
Root the sector-square ledger on side \(i\), with anchors on side \(j\), and let
\(\mathcal H_i\) be any root set carrying load
\[
\int_{\mathcal H_i}\mathcal R_i(y)\,d\nu_i(y)\gtrsim M^4.
\]
Then
\begin{equation}
\label{eq:v076-high-root-support}
\boxed{
\nu_i(\mathcal H_i)
\gtrsim
\frac{M^4}{D_jD_i^2}.
}
\end{equation}
On a total-defect biregular layer this becomes
\begin{equation}
\label{eq:v076-high-root-support-power}
\boxed{
\nu_i(\mathcal H_i)
\gtrsim
W^{4\chi-\theta_j-2\theta_i+o(1)}.
}
\end{equation}
\end{lemma}

\begin{proof}
The symmetric form of \eqref{eq:v072-root-load-degree-square} gives
\(\mathcal R_i(y)\lesssim D_jD_i^2\) pointwise.  Integrate this bound over
\(\mathcal H_i\), compare with \(M^4\), and substitute the degree powers.
\end{proof}

\begin{proposition}[Every non-bush high-reuse output is an actual persistent front seed]
\label{prop:v076-high-reuse-seed-promotion}
Suppose a root enters alternative \textnormal{(ii)} or \textnormal{(iii)} of
\cref{prop:v075-adaptive-continuous-routing}.  Then its fiber contains an unrepeated
selector-direction set \(H_y\) of mass
\begin{equation}
\label{eq:v076-seed-direction-mass}
\boxed{
\mathfrak m_y:=\sigma(H_y)
\gtrsim
\varpi B_{\rm time}^{-1}P^{1/2}.
}
\end{equation}
For every \(0<r\le c_U\mathfrak m_y\), the selector graph over \(H_y\) contains four graph
\(r\)-clusters satisfying
\begin{equation}
\label{eq:v076-seed-cluster-ledger}
\boxed{
|\det(a_1-a_0,a_2-a_0,a_3-a_0)|
\gtrsim_U\mathfrak m_y,
\qquad
\mu_{H_y}(\chi_k)
\gtrsim_{\varepsilon,U}
\mathfrak m_y r^{3+\varepsilon}.
}
\end{equation}
Moreover a set \(G_y\subset J\), of measure bounded below independently of
\(\mathfrak m_y,r\), sees the corresponding four separated Reeb-front clusters.  Along any scale
sequence, a fixed positive-measure limsup set of Reeb times sees such seeds infinitely often.
\end{proposition}

\begin{proof}
In the angular alternative,
\eqref{eq:v075-residual-angular-mass} already gives a contact-neighbor set of
\(\mu_y\)-mass at least the right side of \eqref{eq:v076-seed-direction-mass}.  Since
\(\mu_y=h_W(\,\cdot\,,y)\nu_j\le\nu_j\), its underlying unrepeated selector-direction mass is at
least the same; trim it by nonatomicity.  In the tetrahedral alternative the fixed-anchor fiber
itself has mass \(\gtrsim B_{\rm time}^{-1}P^{1/2}\), which is stronger because
\(\varpi\le1\).  Apply
\cref{lem:v046-sticky-four-cluster-extraction} to \(H_y\), obtaining
\eqref{eq:v076-seed-cluster-ledger}.  Uniform Remez persistence is
\eqref{eq:v047-uniform-good-time}; the scale-limsup assertion is
\cref{cor:v046-longitudinal-limsup}.
\end{proof}

\begin{lemma}[Exact two-root identity for a reused contact leaf]
\label{lem:v077-two-root-contact-identity}
Let \(x\) be a selector graph point joined by normalized contact edges to two roots \(y,y'\).
Choose edge labels \(s,t\in J\) so that
\[
|(b_x-b_y)+s(a_x-a_y)|\lesssim W,
\qquad
|(b_x-b_{y'})+t(a_x-a_{y'})|\lesssim W.
\]
Then
\begin{equation}
\label{eq:v077-two-root-contact-identity}
\boxed{
(b_y+s a_y)-(b_{y'}+t a_{y'})
=(s-t)a_x+O(W).
}
\end{equation}
Consequently, for every \(W\le\delta_s\le1\), exactly one of the following estimates applies:
\begin{enumerate}[label=\textnormal{(\roman*)}]
\item if \(|s-t|\le\delta_s\), then the two roots have a secondary contact collision,
  \begin{equation}
  \label{eq:v077-close-label-root-recollision}
  \boxed{
  |(b_y-b_{y'})+s(a_y-a_{y'})|
  \lesssim_J W+\delta_s;
  }
  \end{equation}
\item if \(|s-t|>\delta_s\), then the shared leaf direction lies in the radial Darboux fiber
  \begin{equation}
  \label{eq:v077-separated-label-radial-fiber}
  \boxed{
  \left|
  a_x-
  \frac{(b_y+s a_y)-(b_{y'}+t a_{y'})}{s-t}
  \right|
  \lesssim
  \frac W{\delta_s}.
  }
  \end{equation}
\end{enumerate}
\end{lemma}

\begin{proof}
The two edge equations say
\[
b_x+s a_x=b_y+s a_y+O(W),
\qquad
b_x+t a_x=b_{y'}+t a_{y'}+O(W).
\]
Subtract them to obtain \eqref{eq:v077-two-root-contact-identity}.  In the close-label case, add
\((t-s)a_{y'}\) and use boundedness of the direction chart.  In the separated-label case, divide
by \(s-t\).
\end{proof}

\begin{proposition}[The seed-overlap second moment has only collision and angular sources]
\label{prop:v077-seed-overlap-routing}
Let \((\mathcal H,\kappa)\) be a finite root measure, and let every root \(y\in\mathcal H\) carry a
measurable contact-neighbor seed set \(H_y\) with selected edge labels \(s_y(x)\in J\).  Define the
leaf multiplicity
\[
\mathcal M(x)
:=
\int_{\mathcal H}\mathbf 1_{H_y}(x)\,d\kappa(y).
\]
Then
\begin{equation}
\label{eq:v077-seed-multiplicity-L1}
\int\mathcal M(x)\,d\nu(x)
=
\int_{\mathcal H}\nu(H_y)\,d\kappa(y),
\end{equation}
and, for every \(W\le\delta_s\le1\),
\begin{equation}
\label{eq:v077-overlap-close-far-split}
\boxed{
\int\mathcal M(x)^2\,d\nu(x)
=
\mathcal O_{\rm rec}(\delta_s)
+
\mathcal O_{\rm rad}(\delta_s),
}
\end{equation}
where \(\mathcal O_{\rm rec}\) is supported on root pairs satisfying
\eqref{eq:v077-close-label-root-recollision}, and \(\mathcal O_{\rm rad}\) is supported on shared
leaves satisfying \eqref{eq:v077-separated-label-radial-fiber}.
In particular,
\begin{equation}
\label{eq:v077-seed-union-cauchy}
\boxed{
\nu\!\left(\bigcup_{y\in\mathcal H}H_y\right)
\ge
\frac{\left(\int_{\mathcal H}\nu(H_y)\,d\kappa(y)\right)^2}
{\mathcal O_{\rm rec}(\delta_s)+\mathcal O_{\rm rad}(\delta_s)}.
}
\end{equation}
Thus failure of transverse de-multiplication is not an arbitrary overlap phenomenon: it is either
another collision layer of the original contact graph or a radial/angular Darboux concentration.
\end{proposition}

\begin{proof}
Equation \eqref{eq:v077-seed-multiplicity-L1} is Tonelli.  Expand
\[
\int\mathcal M^2\,d\nu
=
\iint_{\mathcal H\times\mathcal H}
\nu(H_y\cap H_{y'})\,d\kappa(y)d\kappa(y').
\]
On every shared leaf split the two selected labels at
\(|s_y(x)-s_{y'}(x)|=\delta_s\), and apply
\cref{lem:v077-two-root-contact-identity}; this defines the two nonnegative terms in
\eqref{eq:v077-overlap-close-far-split}.  Cauchy--Schwarz on
\(\supp\mathcal M\) gives \eqref{eq:v077-seed-union-cauchy}.
\end{proof}

\begin{lemma}[Existence of a fixed-angle root re-collision is paid by the normalized contact graph]
\label{lem:v078-recollision-existence-to-graph}
Fix \(\tau_0>0\).  Let \(J\Subset J^+\), \(0<R_s\ll\tau_0\), and let \(y,y'\) satisfy
\[
\tau_0\le|a_y-a_{y'}|\le C_U.
\]
If for some \(s_0\in J\),
\[
|(b_y-b_{y'})+s_0(a_y-a_{y'})|\le R_s,
\]
then the enlarged collision window obeys
\begin{equation}
\label{eq:v078-recollision-window-lower}
\left|
\left\{
s\in J^+:
|(b_y-b_{y'})+s(a_y-a_{y'})|\le2R_s
\right\}
\right|
\gtrsim_U R_s.
\end{equation}
Consequently, if
\[
h_{2R_s}^{J^+}(y,y')
:=
\mathbf 1_{\{\tau_0\le|a_y-a_{y'}|\le C_U\}}
\frac1{C_{U,\tau_0}R_s}
\int_{J^+}
\mathbf 1_{\{|(b_y-b_{y'})+s(a_y-a_{y'})|\le2R_s\}}\,ds,
\]
then
\begin{equation}
\label{eq:v078-existence-indicator-domination}
\boxed{
\mathbf 1_{\{\exists s_0\in J:\,
|(b_y-b_{y'})+s_0(a_y-a_{y'})|\le R_s\}}
\lesssim_{U,\tau_0}
h_{2R_s}^{J^+}(y,y').
}
\end{equation}
\end{lemma}

\begin{proof}
For
\[
|s-s_0|\le c_UR_s,
\]
the change in the collision vector is at most
\(|s-s_0||a_y-a_{y'}|\le R_s\) after choosing \(c_U\) small.  Since \(J\Subset J^+\), this interval
lies in \(J^+\) for all sufficiently small \(R_s\).  It has length comparable to \(R_s\), proving
\eqref{eq:v078-recollision-window-lower}.  The fixed-shell lower cutoff \(\tau_0\) permits a
normalizing constant \(C_{U,\tau_0}\) for which the graph is bounded by one; division by that
normalization gives \eqref{eq:v078-existence-indicator-domination}.
\end{proof}

\begin{proposition}[Quantitative fixed-angle charge for the close-label seed overlap]
\label{prop:v078-close-overlap-charge}
Use \cref{prop:v077-seed-overlap-routing}, put
\[
R_s=C_J(W+\delta_s),
\]
and split its close-label root pairs at the fixed normalized cutoff
\(|a_y-a_{y'}|=\tau_0\), where \(R_s\ll\tau_0\).  If the root measure
\(\kappa\) is a submeasure of the retained selector measure, then
\begin{equation}
\label{eq:v078-close-overlap-fixed-angular}
\boxed{
\mathcal O_{\rm rec}(\delta_s)
\lesssim
\iint_{|a_y-a_{y'}|<\tau_0}d\kappa(y)d\kappa(y')
+
\iint h_{2R_s}^{J^+}(y,y')\,d\kappa(y)d\kappa(y').
}
\end{equation}
The first term is passed wholly to the next angular Darboux patch; the second is the edge mass of
the same normalized Maslov collision graph at the coarser physical scale \(R_s\).  In particular
the close-label overlap introduces no independent occurrence coefficient.
\end{proposition}

\begin{proof}
On the near-angle set discard the shared-leaf indicator and bound it by one.  On the complementary
set, \eqref{eq:v077-close-label-root-recollision} gives existence of an \(R_s\)-collision at a time
in \(J\).  Discarding the shared-leaf indicator and applying
\cref{lem:v078-recollision-existence-to-graph} gives the second term.
\end{proof}

\begin{lemma}[Separated Reeb labels lie in one affine radial Darboux plane]
\label{lem:v078-radial-darboux-surface}
Let \(U\Subset\mathbb R^3\) and \(J\Subset\mathbb R\) be bounded, and keep the root coordinates in
one bounded contact chart.  For fixed roots \(y,y'\), put
\[
\mathcal D_{\delta_s}
:=
\{(s,t)\in J^2:|s-t|\ge\delta_s\},
\qquad
\Phi_{y,y'}(s,t)
:=
\frac{(b_y+s a_y)-(b_{y'}+t a_{y'})}{s-t}.
\]
Then the full image of \(\mathcal D_{\delta_s}\) is contained in the fixed affine plane
\begin{equation}
\label{eq:v078-radial-darboux-plane}
\boxed{
\Pi_{y,y'}
:=
a_y+\operatorname{span}\{b_y-b_{y'},\,a_y-a_{y'}\}.
}
\end{equation}
If \(0<W\le\delta_s\le1\) and \(\rho=W/\delta_s\), then
\begin{equation}
\label{eq:v078-radial-darboux-tube-volume}
\boxed{
\left|
U\cap N_{C\rho}\!\left(\Phi_{y,y'}(\mathcal D_{\delta_s})\right)
\right|
\lesssim_{U,J,C}
\rho
=
W\delta_s^{-1}.
}
\end{equation}
The estimate is uniform in the two roots.  Thus the apparent four-variable radial fiber in
\eqref{eq:v077-separated-label-radial-fiber} is one contact-polarized plane and pays one full
direction codimension with no parameter-covering loss.
\end{lemma}

\begin{proof}
Write \(u=s-t\).  The contact numerator gives the exact homothety identity
\[
\Phi_{y,y'}(s,t)
=
a_y+\frac{b_y-b_{y'}}u+\frac t u(a_y-a_{y'}).
\]
This proves \eqref{eq:v078-radial-darboux-plane}.  The intersection of a bounded subset of
\(\mathbb R^3\) with the \(C\rho\)-neighborhood of an affine plane has volume
\(O_{U,C}(\rho)\).  This remains true, with a stronger bound, if the displayed span has dimension
zero or one.  Hence \eqref{eq:v078-radial-darboux-tube-volume} follows.
\end{proof}

\begin{proposition}[The separated-label seed overlap pays the radial Darboux codimension]
\label{prop:v078-radial-overlap-bound}
Under the hypotheses of \cref{prop:v077-seed-overlap-routing}, assume that the horizontal marginal
of \(\nu\) has density at most \(C_U\).  Then
\begin{equation}
\label{eq:v078-radial-overlap-bound}
\boxed{
\mathcal O_{\rm rad}(\delta_s)
\lesssim_{U,J}
\|\kappa\|^2W\delta_s^{-1}.
}
\end{equation}
In particular, for \(\delta_s=W^\alpha\) with \(0<\alpha<1\),
\begin{equation}
\label{eq:v078-radial-overlap-power}
\boxed{
R_s\asymp W^\alpha,
\qquad
\frac W{\delta_s}=W^{1-\alpha},
\qquad
\mathcal O_{\rm rad}(\delta_s)
\lesssim
\|\kappa\|^2W^{1-\alpha}.
}
\end{equation}
Both geometric scales vanish and the radial overlap has a strict power gain.
\end{proposition}

\begin{proof}
For a fixed root pair \((y,y')\), every shared leaf counted by
\(\mathcal O_{\rm rad}(\delta_s)\) has selected labels in
\(\mathcal D_{\delta_s}\), and \eqref{eq:v077-separated-label-radial-fiber} places its direction
in the \(C\rho\)-neighborhood of
\(\Phi_{y,y'}(\mathcal D_{\delta_s})\).  The bounded horizontal density and
\cref{lem:v078-radial-darboux-surface} bound its \(\nu\)-mass by
\(O_{U,J}(W\delta_s^{-1})\).  Integrate this uniform bound against
\(d\kappa(y)d\kappa(y')\).  Substituting \(\delta_s=W^\alpha\) gives
\eqref{eq:v078-radial-overlap-power}.
\end{proof}

\begin{proposition}[One-scale contact--symplectic closure of the seed second moment]
\label{prop:v078-one-scale-seed-closure}
Keep the hypotheses of
\cref{prop:v077-seed-overlap-routing,prop:v078-close-overlap-charge,%
prop:v078-radial-overlap-bound}.  Write
\[
K:=\|\kappa\|>0,
\qquad
A:=\int_{\mathcal H}\nu(H_y)\,d\kappa(y),
\qquad
\overline m:=A/K,
\]
and let \(0<Q\le1\).  There is a constant \(c_{U,J}>0\) such that, whenever
\begin{equation}
\label{eq:v078-radial-absorption-condition}
W\delta_s^{-1}
\le
c_{U,J}\frac{\overline m^2}{Q},
\end{equation}
at least one of the following holds:
\begin{enumerate}[label=\textnormal{(\roman*)}]
\item the seed directions de-multiply at the requested mass,
\begin{equation}
\label{eq:v078-seed-union-output}
\boxed{
\nu\!\left(\bigcup_{y\in\mathcal H}H_y\right)\ge Q;
}
\end{equation}
\item one root-centered angular Darboux patch carries mass
\begin{equation}
\label{eq:v078-angular-root-output}
\boxed{
\kappa\{y:|a_y-a_{y_0}|<\tau_0\}
\gtrsim_{U,J}
K\frac{\overline m^2}{Q}
}
\end{equation}
for some \(y_0\in\mathcal H\);
\item for \(\widehat\kappa:=K^{-1}\kappa\), the fixed-angle coarser-scale Maslov graph has edge
mass
\begin{equation}
\label{eq:v078-coarse-maslov-edge-output}
\boxed{
M_{R_s}
:=
\iint h_{2R_s}^{J^+}(y,y')\,d\widehat\kappa(y)d\widehat\kappa(y')
\gtrsim_{U,J}
\frac{\overline m^2}{Q}.
}
\end{equation}
For every \(2R_s\le c\Delta<1\), this last graph enters
\cref{prop:v050-arbitrary-threshold-maslov-routing}: it produces either an
\(O(R_s/\Delta)\)-Lagrangian point-probe bush of \(\widehat\kappa\)-mass
\(\gtrsim M_{R_s}^3\), or determinant-small Maslov four-cycle content
\(\gtrsim M_{R_s}^4\).
\end{enumerate}
Thus a failed seed aggregation has no set-theoretic remainder: after the radial Darboux
codimension is absorbed, it is exactly an angular change of chart or a quantified return to the
manuscript's original bush/rank-loss route.
\end{proposition}

\begin{proof}
Put
\[
\mathcal A_{\tau_0}
:=
\iint_{|a_y-a_{y'}|<\tau_0}d\kappa(y)d\kappa(y'),
\qquad
\mathcal E_{R_s}
:=
\iint h_{2R_s}^{J^+}(y,y')\,d\kappa(y)d\kappa(y').
\]
If \eqref{eq:v078-seed-union-output} fails, then
\eqref{eq:v077-seed-union-cauchy} gives
\[
\mathcal O_{\rm rec}(\delta_s)+\mathcal O_{\rm rad}(\delta_s)
>
\frac{A^2}{Q}.
\]
Use \cref{prop:v078-close-overlap-charge,prop:v078-radial-overlap-bound}.  After choosing the
constant in \eqref{eq:v078-radial-absorption-condition} below their fixed implicit constants, the
radial term accounts for at most a fixed small portion of \(A^2/Q\).  Hence
\begin{equation}
\label{eq:v078-angular-or-coarse-edge}
\mathcal A_{\tau_0}+\mathcal E_{R_s}
\gtrsim_{U,J}
\frac{A^2}{Q}.
\end{equation}
If the first summand carries a fixed portion, Fubini gives a root \(y_0\) for which
\[
\kappa\{y:|a_y-a_{y_0}|<\tau_0\}
\gtrsim
\mathcal A_{\tau_0}/K
\gtrsim
K\overline m^2/Q,
\]
which is \eqref{eq:v078-angular-root-output}.  Otherwise divide the second summand by \(K^2\) to
obtain \eqref{eq:v078-coarse-maslov-edge-output}.  The angular cutoff built into
\(h_{2R_s}^{J^+}\) makes it a symmetric graphon with values in \([0,1]\).  Its positive edges have
selected contact labels with error \(O(R_s)\), so
\cref{prop:v050-arbitrary-threshold-maslov-routing} applies with collision scale \(O(R_s)\) and
gives the stated bush/rank-loss alternative.
\end{proof}

\subsection{Two-generation flags and relative residual return}

\begin{corollary}[The input-overlap Maslov--radial balance occurs at the one-third Reeb scale]
\label{cor:v078-one-third-maslov-radial-balance}
Take
\begin{equation}
\label{eq:v078-one-third-choice}
\boxed{
\delta_s=W^{1/3},
\qquad
R_s\asymp W^{1/3},
\qquad
\frac W{\delta_s}=W^{2/3}.
}
\end{equation}
Suppose the angular output in
\eqref{eq:v078-angular-root-output} is passed to the angular Darboux branch, and decompose the
coarse graph dyadically by its actual edge mass.  Every coarse layer with
\(M_{R_s}>R_s^{2-o(1)}\) is a total-defect layer \(\chi_s<2\) and re-enters the existing
Lagrangian-bush/Maslov-rank-loss route.  On the complementary layers,
\cref{prop:v078-close-overlap-charge,prop:v078-radial-overlap-bound} and
\eqref{eq:v077-seed-union-cauchy} give
\begin{equation}
\label{eq:v078-one-third-seed-union}
\boxed{
\nu\!\left(\bigcup_{y\in\mathcal H}H_y\right)
\gtrsim
\min\left\{1,\overline m^2W^{-2/3+o(1)}\right\},
}
\end{equation}
up to the angular piece already routed away.  In particular, every seed layer with
\begin{equation}
\label{eq:v078-one-third-seed-threshold}
\boxed{
\overline m^2\gtrsim W^{2/3+o(1)}
}
\end{equation}
passes the input-mass gate in \eqref{eq:v078-relative-residual-sufficient}; the relative residual
gate is addressed separately by \cref{prop:v078-input-output-frostman-firewall}.
\end{corollary}

\begin{proof}
At \(\delta_s=W^{1/3}\), \Cref{prop:v078-radial-overlap-bound} gives
\(\mathcal O_{\rm rad}\lesssim K^2W^{2/3}\).  A coarse graph outside the bad total-defect layers
has edge mass at most \(R_s^{2-o(1)}=W^{2/3-o(1)}\), and hence contributes at most
\(K^2W^{2/3-o(1)}\).  Insert these two bounds into
\eqref{eq:v077-seed-union-cauchy}, using \(A=K\overline m\); the factor \(K^2\) cancels and yields
\eqref{eq:v078-one-third-seed-union}.  Equation
\eqref{eq:v078-one-third-seed-threshold} makes its right side subpower.
\end{proof}

\begin{corollary}[Power audit of the high-reuse flag layer]
\label{cor:v078-high-reuse-seed-power-audit}
Root on the side \(i\) for which \(\theta_i\le\theta_j\).  Then
\eqref{eq:v074-flag-exponent-range} gives \(\phi_i\ge\theta_j\ge\theta_i\).  At the choice
\(\delta_s=W^{1/3}\):
\begin{enumerate}[label=\textnormal{(\roman*)}]
\item if \(B_{\rm time}>1\), then
\begin{equation}
\label{eq:v078-amplified-seed-closes}
\overline m^2
\gtrsim
W^{2/3+2\theta_i-2\phi_i+o(1)}
\ge
W^{2/3+o(1)},
\end{equation}
so \cref{cor:v078-one-third-maslov-radial-balance} closes the input-side seed overlap;
\item if \(B_{\rm time}=1\), put
\begin{equation}
\label{eq:v078-unamplified-pair-exponent}
p_i:=3\chi+\theta_i-\phi_i+\eta.
\end{equation}
Then \(\overline m^2\gtrsim P_i\gtrsim W^{p_i+o(1)}\), so every layer with
\(p_i\le2/3\) also passes the input-overlap gate.  The only flag-power region not settled by the one-third balance is
\begin{equation}
\label{eq:v078-final-unamplified-wedge}
\boxed{
B_{\rm time}=1,
\qquad
3\chi+\theta_i-\phi_i+\eta>\frac23,
\qquad
\theta_i\le\theta_j\le\phi_i.
}
\end{equation}
\end{enumerate}
Thus neither the amplified high-reuse branch nor any flag layer of pair mass at least
\(W^{2/3+o(1)}\) remains in the global frontier.
\end{corollary}

\begin{proof}
When \(B_{\rm time}>1\), the first line of
\eqref{eq:v075-angular-seed-power-ledger}, together with
\eqref{eq:v076-seed-direction-mass}, gives
\(\overline m\gtrsim W^{1/3+\theta_i-\phi_i+o(1)}\); square it and use
\(\phi_i\ge\theta_i\).  When \(B_{\rm time}=1\), \Cref{prop:v075-adaptive-continuous-routing} and
\eqref{eq:v074-flag-product-power} give
\(\overline m^2\gtrsim P_i\gtrsim W^{p_i+o(1)}\).  Apply
\eqref{eq:v078-one-third-seed-threshold}.
\end{proof}

\begin{proposition}[The residual flag wedge retains its full packet-occurrence measure]
\label{prop:v078-flag-packet-incidence-lift}
Work on a dyadically regularized residual layer
\eqref{eq:v078-final-unamplified-wedge}.  Let \(\mathcal H_i\) have root mass \(K_i\), assume
\[
\eta_y(\mathscr F_y(H_i))\asymp P_i
\qquad(y\in\mathcal H_i),
\]
For \(f=(p,q,\omega)\), retain its actual common-neighbor density
\[
w_f(z)
:=
\mathbf 1_{A_{p,q,\omega}}(z)h(p,z)h(q,z),
\qquad
H_i\le\int w_f\,d\nu_i<2H_i.
\]
On the flag space put
\[
d\Lambda(y,f)
:=
d\nu_i(y)d\eta_y(f),
\qquad
\|\Lambda\|\asymp K_iP_i,
\]
and define the packet-occurrence density on the selector graph by
\[
\mathcal N_{\rm flag}(z)
:=
\int w_f(z)\,d\Lambda(y,f).
\]
Then
\begin{equation}
\label{eq:v078-flag-lift-first-moment}
\boxed{
\int\mathcal N_{\rm flag}\,d\nu_i
\asymp
K_iH_iP_i,
}
\end{equation}
and
\begin{equation}
\label{eq:v078-flag-lift-second-moment}
\boxed{
\int\mathcal N_{\rm flag}^2\,d\nu_i
=
\iint
\left(\int w_f(z)w_{f'}(z)\,d\nu_i(z)\right)
\,d\Lambda(y,f)d\Lambda(y',f').
}
\end{equation}
Consequently
\begin{equation}
\label{eq:v078-flag-packet-union}
\boxed{
\nu_i\!\left(\{z:\mathcal N_{\rm flag}(z)>0\}\right)
\gtrsim
\frac{(K_iH_iP_i)^2}
{\int\mathcal N_{\rm flag}^2\,d\nu_i}.
}
\end{equation}
In particular, the factor \(H_iP_i\) from
\eqref{eq:v073-weighted-anchor-occurrence} is present before any root normalization and may not be
replaced by the smaller single-fiber mass \(P_i^{1/2}\).

Moreover \(\mathcal N_{\rm flag}\) is exactly the root-restricted square load
\begin{equation}
\label{eq:v078-flag-load-self-duality}
\boxed{
\mathcal N_{\rm flag}(z)
=
\iint\sum_\omega
w_{p,q,\omega}(z)
\left(\int_{\mathcal H_i}w_{p,q,\omega}(y)\,d\nu_i(y)\right)
\,d\nu_j(p)d\nu_j(q)
\le
\mathcal R_i(z).
}
\end{equation}
Equality with the full rooted load holds when the root restriction is removed; on the retained
layer the displayed domination is all that is used.  Hence
\begin{equation}
\label{eq:v078-flag-load-pointwise-cap}
\boxed{
\mathcal N_{\rm flag}(z)
\lesssim
D_jd_i(z)^2
\lesssim
D_jD_i^2.
}
\end{equation}
\end{proposition}

\begin{proof}
Tonelli and \(H_i\le\int w_f\,d\nu_i<2H_i\) give
\eqref{eq:v078-flag-lift-first-moment}.  Expanding the square and applying Tonelli again gives
\eqref{eq:v078-flag-lift-second-moment}.  Cauchy--Schwarz on the support of
\(\mathcal N_{\rm flag}\) proves \eqref{eq:v078-flag-packet-union}.  Finally substitute the
definition
\(d\eta_y(p,q,\omega)=w_{p,q,\omega}(y)d\nu_j(p)d\nu_j(q)\), interchange the \(y\) and
\((p,q,\omega)\) integrals, and obtain \eqref{eq:v078-flag-load-self-duality}.  The pointwise rooted
load bound \eqref{eq:v072-root-load-degree-square} gives
\eqref{eq:v078-flag-load-pointwise-cap}.
\end{proof}

\begin{lemma}[Coherent flag reuse re-roots without losing occurrence mass]
\label{lem:v078-coherent-flag-rerooting}
For a selector point \(z\) on side \(i\), define the full pair mass of the same flag-packet layer by
\begin{equation}
\label{eq:v078-rerooted-flag-pair-mass}
P_i^{(H)}(z)
:=
\iint\sum_{\omega:\,H_i\le m_{p,q,\omega}<2H_i}
w_{p,q,\omega}(z)
\,d\nu_j(p)d\nu_j(q).
\end{equation}
Then the lifted occurrence density satisfies the pointwise bound
\begin{equation}
\label{eq:v078-flag-density-by-rerooted-pair-mass}
\boxed{
\mathcal N_{\rm flag}(z)
\le
2H_iP_i^{(H)}(z).
}
\end{equation}
For every threshold \(0<L\le1\), at least one of the following holds:
\begin{enumerate}[label=\textnormal{(\roman*)}]
\item the low-reuse packet union has actual selector-direction mass
\begin{equation}
\label{eq:v078-low-flag-reuse-union}
\boxed{
\nu_i\{0<\mathcal N_{\rm flag},\ P_i^{(H)}\le L\}
\gtrsim
\frac{K_iP_i}{L};
}
\end{equation}
\item the re-rooted high-pair set retains lifted load
\begin{equation}
\label{eq:v078-high-flag-rerooted-load}
\boxed{
\int_{\{P_i^{(H)}>L\}}
\mathcal R_i(z)\,d\nu_i(z)
\gtrsim
K_iH_iP_i.
}
\end{equation}
\end{enumerate}
Thus high coherent reuse is not a new multiplicity: it is the original rooted load on a root set
whose flag-pair mass has increased to at least \(L\).
\end{lemma}

\begin{proof}
Using \eqref{eq:v078-flag-load-self-duality} before the final domination, the inner restricted root
mass satisfies
\[
\int_{\mathcal H_i}w_{p,q,\omega}(y)\,d\nu_i(y)
\le
m_{p,q,\omega}<2H_i.
\]
Integration in \((p,q,\omega)\) proves
\eqref{eq:v078-flag-density-by-rerooted-pair-mass}.  Split the first moment
\eqref{eq:v078-flag-lift-first-moment} over
\(P_i^{(H)}\le L\) and \(P_i^{(H)}>L\).  If the first part carries at least half, then
\(\mathcal N_{\rm flag}\le2H_iL\) there, so its support has mass at least
\[
\frac{cK_iH_iP_i}{2H_iL}
\gtrsim
\frac{K_iP_i}{L},
\]
which is \eqref{eq:v078-low-flag-reuse-union}.  Otherwise the second part carries half.  Since
\(\mathcal N_{\rm flag}\le\mathcal R_i\), this gives
\eqref{eq:v078-high-flag-rerooted-load}.
\end{proof}

\begin{corollary}[The high coherent-reuse half reaches the one-third input threshold]
\label{cor:v078-coherent-high-reuse-closes}
Take \(L=W^{2/3+o(1)}\).  In alternative 
\textnormal{(ii)} of \cref{lem:v078-coherent-flag-rerooting}, dyadic regularization on the set
\(P_i^{(H)}>L\) produces a rooted flag layer with pair mass at least \(W^{2/3+o(1)}\), while
retaining a fixed portion of \(K_iH_iP_i\).  Its non-bush seed mass is at least
\((P_i^{(H)})^{1/2}\ge W^{1/3+o(1)}\), so
\cref{cor:v078-one-third-maslov-radial-balance} closes its input-side overlap.  Consequently the coherent survivor may
be assumed to satisfy the genuine low-reuse output
\begin{equation}
\label{eq:v078-final-low-reuse-direction-mass}
\boxed{
\nu_i\{0<\mathcal N_{\rm flag},\ P_i^{(H)}\le W^{2/3+o(1)}\}
\gtrsim
K_iP_iW^{-2/3-o(1)}.
}
\end{equation}
On the total-defect power layer, this lower bound is
\begin{equation}
\label{eq:v078-final-low-reuse-direction-power}
\boxed{
\gtrsim
W^{7\chi-\theta_j-\theta_i-\phi_i+\eta-2/3+o(1)}.
}
\end{equation}
\end{corollary}

\begin{proof}
Only the last displayed power requires verification.  Multiply the root-support power
\eqref{eq:v076-high-root-support-power} by
\(P_iW^{-2/3}=W^{3\chi+\theta_i-\phi_i+\eta-2/3+o(1)}\).
\end{proof}

\begin{proposition}[Input de-multiplication reaches the Frostman target only through relative Maslov energy]
\label{prop:v078-input-output-frostman-firewall}
Let \(H\) be a measurable subset of the selector graph, let
\(\mu_H:=\mathbf 1_H\mu\), and put \(m:=\mu_H(H)\).  For a half-open grid
\(\mathcal Q_r\) of physical \(r\)-cubes, define
\[
S_r(s)
:=
\sum_{Q\in\mathcal Q_r}
\mu_H(\Pi_s^{-1}(Q))^2.
\]
Then
\begin{equation}
\label{eq:v078-output-entropy-second-moment}
\boxed{
\int_JN_r(\Pi_s(H))\,ds
\gtrsim_J
\frac{m^2}
{mr^3+r\,\mathfrak Z_{Cr,\ge r}^{J^+}(\mu_H)}.
}
\end{equation}
Consequently the two conditions
\begin{equation}
\label{eq:v078-relative-residual-sufficient}
\boxed{
m=r^{o(1)},
\qquad
\mathfrak Z_{Cr,\ge r}^{J^+}(\mu_H)
\lesssim
m^2r^{2-o(1)}
}
\end{equation}
imply
\begin{equation}
\label{eq:v078-relative-residual-entropy-output}
\boxed{
\int_JN_r(\Pi_s(H))\,ds
\gtrsim
r^{-3+o(1)}.
}
\end{equation}
Thus a subpower input-direction union is necessary to pay the near-direction diagonal, but it is
not by itself a projection theorem: the separated pairs still require the quadratic relative
Maslov residual bound in \eqref{eq:v078-relative-residual-sufficient}.
\end{proposition}

\begin{proof}
For every \(s\), Cauchy--Schwarz over the occupied output cubes gives
\[
m^2
\le
N_r(\Pi_s(H))S_r(s).
\]
Cauchy--Schwarz in \(s\), or the harmonic--arithmetic mean inequality, therefore yields
\[
\int_JN_r(\Pi_s(H))\,ds
\gtrsim_J
\frac{m^2}{\int_JS_r(s)\,ds}.
\]
Two points counted in the same output cube satisfy
\(|\Pi_s(z)-\Pi_s(z')|\le Cr\).  On \(|a-a'|<r\), bounded horizontal density gives pair mass
\(O_U(mr^3)\).  On \(|a-a'|\ge r\), a nonempty collision window forces
\(s_*\in J^+\), \(|r(z,z')|\le Cr\), and has length at most
\(C_Jr/|a-a'|\).  Hence
\[
\int_JS_r(s)\,ds
\lesssim_{U,J}
mr^3+r\,\mathfrak Z_{Cr,\ge r}^{J^+}(\mu_H).
\]
This proves \eqref{eq:v078-output-entropy-second-moment}; substitution of
\eqref{eq:v078-relative-residual-sufficient} gives
\eqref{eq:v078-relative-residual-entropy-output}.
\end{proof}

\begin{lemma}[Coarse old flags and fine new collisions freeze the input mass]
\label{lem:v078-coarse-fine-mass-freezing}
First fix one coarse scale \(W>0\) and one coarse flag output \(H\) of mass \(m>0\).  As the new
collision scale \(r\downarrow0\), every finite constant depending only on this frozen coarse layer,
including \(m^{\pm1}\) and any fixed \(W\)-power, is \(r^{-o(1)}\).  Thus a residual estimate proved
uniformly for all sufficiently small \(r\), after freezing \(W,H\), may absorb all such coarse
constants and still feed the Hausdorff Frostman criterion.

For the scale-sequence version, let \(W_n\downarrow0\), and for each \(n\) let \(H_n\) be a
measurable selector subset produced at
the coarse flag scale \(W_n\), with mass
\[
0<m_n:=\mu(H_n)\le1.
\]
There is a fine collision scale \(r_n\downarrow0\) such that
\begin{equation}
\label{eq:v078-coarse-fine-scale-choice}
\boxed{
0<r_n
\le
\min\{W_n^n,m_n^n\}.
}
\end{equation}
Consequently
\begin{equation}
\label{eq:v078-frozen-mass-is-subpower}
\boxed{
m_n=r_n^{o(1)},
\qquad
r_n=o(W_n^A)
\quad\text{for every fixed }A>0.
}
\end{equation}
Thus the first condition in \eqref{eq:v078-relative-residual-sufficient} does not require the
coarse lower bound \eqref{eq:v078-final-low-reuse-direction-power} to have exponent zero.  It is
paid by separating the new Maslov-collision scale from the old flag scale.  This is a diagonal
choice of two scales, not an intrinsic recursive construction.  By itself the diagonal
scale-sequence conclusion gives only a scale-wise/Minkowski output; a Hausdorff conclusion still
requires the fixed-coarse-layer, all-fine-scales form above or the scale-coherent kernels in
\cref{lem:v047-slice-frostman-lift,cor:v047-exact-spatial-endpoint}.
\end{lemma}

\begin{proof}
The fixed-coarse-layer assertion follows because the logarithm of every such constant is fixed
while \(|\log r|\to\infty\).  For the sequence assertion, choose any positive \(r_n\) satisfying
\eqref{eq:v078-coarse-fine-scale-choice}.  Then
\[
\frac{|\log m_n|}{|\log r_n|}
\le
\frac1n,
\]
with the ratio interpreted as zero when \(m_n=1\).  This proves
\(m_n=r_n^{o(1)}\).  Likewise
\[
\frac{r_n}{W_n^A}
\le
W_n^{\,n-A}\longrightarrow0
\]
for each fixed \(A\), proving the second assertion.
\end{proof}

\begin{principle}[The final low-reuse task is a relative residual estimate, not a direction-union estimate]
\label{pr:v078-final-relative-residual-task}
For the Hausdorff endpoint, first freeze one sufficiently fine coarse low-reuse output \(H\) from
\eqref{eq:v078-final-low-reuse-direction-mass}, and then apply
\cref{prop:v078-input-output-frostman-firewall} uniformly over all sufficiently small new collision
scales \(r\).  The mass \(m=\mu(H)>0\) is then a coarse constant and hence \(m=r^{o(1)}\).  For a
scale-sequence/Minkowski statement one may instead use the diagonal scales from
\cref{lem:v078-coarse-fine-mass-freezing}.  In either case the preceding flag arguments remove
radial, transverse-plane, and high coherent occurrence losses, and reduce the proof to the second condition in \eqref{eq:v078-relative-residual-sufficient}.  The exact remaining estimate is
therefore
\begin{equation}
\label{eq:v078-final-low-reuse-relative-residual}
\boxed{
\mathfrak Z_{Cr,\ge r}^{J^+}(\mu_H)
\lesssim
\mu(H)^2r^{2-o(1)}.
}
\end{equation}
This is the relative form of the manuscript's original codimension-two Maslov target
\eqref{eq:v048-final-residual-endpoint}.  No Euclidean dimension credit is taken from the flag
occurrence bounds before \eqref{eq:v078-final-low-reuse-relative-residual} is established on the
fixed-coarse/all-fine scale family required by the Hausdorff endpoint.
\end{principle}

\begin{proposition}[Failure of the relative residual bound re-enters the normalized Maslov graph]
\label{prop:v078-relative-residual-maslov-return}
Let \(J^-\Subset J\Subset J^+\), work on a normalized fixed-angle Darboux shell, and let
\(H_{\rm low}\) be the low-reuse output of
\eqref{eq:v078-final-low-reuse-direction-mass}.  Put
\[
m:=\mu(H_{\rm low}),
\qquad
\widehat\mu:=m^{-1}\mu_{H_{\rm low}}.
\]
If, along a scale sequence \(r\downarrow0\),
\begin{equation}
\label{eq:v078-relative-residual-failure}
\mathfrak Z_{r/2,\ge\tau_0}^{J^-}(\mu_{H_{\rm low}})
\ge
A_r m^2r^{2-\eta},
\qquad
A_r\longrightarrow\infty,
\end{equation}
then the normalized collision graph
\begin{equation}
\label{eq:v078-low-output-normalized-maslov-graph}
h_r(z,z')
:=
\mathbf 1_{\{\tau_0\le|a-a'|\le C_U\}}
\frac1{C_{\tau_0}r}
\int_J
\mathbf 1_{\{|\Pi_s(z)-\Pi_s(z')|\le r\}}\,ds
\end{equation}
on \((H_{\rm low},\widehat\mu)\) has edge mass
\begin{equation}
\label{eq:v078-low-output-maslov-edge-mass}
\boxed{
M_r:=e(h_r)
\gtrsim
A_r r^{2-\eta}.
}
\end{equation}
For every \(r\le c\Delta<1\), \Cref{prop:v050-arbitrary-threshold-maslov-routing} therefore yields either
\begin{enumerate}[label=\textnormal{(\roman*)}]
\item an \(O(r/\Delta)\)-Lagrangian point-probe bush of normalized
\(\widehat\mu\)-mass \(\gtrsim M_r^3\); or
\item determinant-small Maslov four-cycle content \(\gtrsim M_r^4\) inside the low old-flag-reuse
set.
\end{enumerate}
No factor of \(m^{-1}\) remains in \eqref{eq:v078-low-output-maslov-edge-mass}; the required
quadratic factor \(m^2\) in the relative residual target is exactly canceled by probability
normalization.
\end{proposition}

\begin{proof}
The collision-window lower bound in
\eqref{eq:v048-energy-residual-equivalence}, applied to \(\widehat\mu\), gives
\[
e(h_r)
\gtrsim
\mathfrak Z_{r/2,\ge\tau_0}^{J^-}(\widehat\mu)
=
m^{-2}\mathfrak Z_{r/2,\ge\tau_0}^{J^-}(\mu_{H_{\rm low}}).
\]
Insert \eqref{eq:v078-relative-residual-failure}.  The fixed-angle cutoff makes
\(0\le h_r\le1\) after choosing \(C_{\tau_0}\), so
\cref{prop:v050-arbitrary-threshold-maslov-routing} applies and gives the two stated outputs.
\end{proof}

\begin{principle}[The last unclosed incidence is an old-flag/new-cycle compatibility]
\label{pr:v078-final-six-point-compatibility}
The bush output of \cref{prop:v078-relative-residual-maslov-return} is already a physical exit.
Hence a failure of \eqref{eq:v078-final-low-reuse-relative-residual} may be assumed to produce a
determinant-small Maslov four-cycle entirely inside \(H_{\rm low}\).  Every vertex of that cycle
still carries the old azimuth--Maslov packet flags from
\cref{prop:v078-flag-packet-incidence-lift}, while the re-rooted old-flag pair mass is bounded by
\(W^{2/3+o(1)}\) on the low-reuse set.  The final compatibility estimate must therefore bound a
new Maslov four-cycle decorated by an old common-neighbor flag, with the old flag occurrence kept
below the one-third threshold.  This is a concrete six-point contact incidence; it is not a
generic projection or set-overlap problem.
\end{principle}

\begin{lemma}[Two-generation polarization comparison at the decorated cycle vertex]
\label{lem:v078-old-flag-new-cycle-polarization}
Let an old flag packet lie in a phase-plane tube
\[
(\Pi_{\rm old}^a\times\Pi_{\rm old}^b)_{t_{\rm old}}
=
(z_{\rm old}+E_{\rm old}\oplus E_{\rm old})_{t_{\rm old}},
\]
and let the determinant-small new Maslov cycle routing place its decorated vertex in
\[
(\Pi_{\rm new}^a\times\Pi_{\rm new}^b)_{t_{\rm new}}
=
(z_{\rm new}+E_{\rm new}\oplus E_{\rm new})_{t_{\rm new}}.
\]
Suppose both packets meet one graph product cell of diameter \(O(r)\), and let
\(n_{\rm old},n_{\rm new}\) be the corresponding horizontal unit normals.  Put
\[
T:=t_{\rm old}+t_{\rm new}+r.
\]
For every \(T\le\vartheta\le1\), exactly one of the following geometric estimates applies:
\begin{enumerate}[label=\textnormal{(\roman*)}]
\item if \(|n_{\rm old}\times n_{\rm new}|\ge\vartheta\), then the possible decorated selector
directions occupy volume and selector mass at most
\begin{equation}
\label{eq:v078-two-generation-transverse-cost}
\boxed{
\lesssim_U
\frac{(t_{\rm old}+r)(t_{\rm new}+r)}{\vartheta};
}
\end{equation}
\item if \(|n_{\rm old}-n_{\rm new}|<\vartheta\), then their horizontal and vertical affine phase
offsets satisfy
\begin{equation}
\label{eq:v078-two-generation-phase-alignment}
\boxed{
|c_{\rm old}^a-c_{\rm new}^a|
+
|c_{\rm old}^b-c_{\rm new}^b|
\lesssim_U
T+\vartheta.
}
\end{equation}
Consequently both generations lie, on the common cell and up to thickness
\(O_U(T+\vartheta)\), in one coherent polarization
\(z_*+(E\oplus E)\).
\end{enumerate}
\end{lemma}

\begin{proof}
The horizontal projection of the common graph cell meets both affine plane tubes.  For transverse
normals, the two defining normal coordinates have Jacobian
\(|n_{\rm old}\times n_{\rm new}|\), while the remaining direction coordinate ranges in a fixed
bounded interval.  This gives \eqref{eq:v078-two-generation-transverse-cost}; bounded horizontal
density gives the selector-mass version.  In the coherent case, apply the common-cell offset
calculation of \cref{lem:v053-common-cell-offset-alignment} to the horizontal coordinates and
\cref{cor:v053-common-cell-phase-offsets} to the vertical coordinates, adding the two packet
thicknesses and the cell diameter.  This proves
\eqref{eq:v078-two-generation-phase-alignment}.
\end{proof}

\begin{lemma}[A coherent new Maslov plane makes its opposite anchors \(q\)-local]
\label{lem:v079-coherent-new-plane-q-local}
Fix the old coherent polarization \(E_*\subset\mathbb R^3\), let \(n_*\) be its unit normal, and
write
\[
\mathfrak q_*(a):=n_*\cdot a.
\]
For an anchored new rank-loss cycle \(p,y_0,q,y,p\), let
\[
E_{\rm new}
=
\Span\{a_{y_0}-a_p,a_q-a_{y_0}\},
\]
and choose the sign of its unit normal \(n_{\rm new}\) so that
\(n_{\rm new}\cdot n_*\ge0\).  If
\[
|n_{\rm new}-n_*|\le\vartheta,
\]
then
\begin{equation}
\label{eq:v079-coherent-anchor-q-local}
\boxed{
|\mathfrak q_*(a_q)-\mathfrak q_*(a_p)|
\lesssim_U
\vartheta.
}
\end{equation}
Conversely, if the left side is at least \(R_q\), then
\begin{equation}
\label{eq:v079-q-separated-anchor-transverse}
\boxed{
|n_{\rm new}\times n_*|
\gtrsim_U
R_q.
}
\end{equation}
Thus coherence of the old and new Maslov generations is not merely a statement about two plane
parameters: it forces the \emph{opposite anchor pair} of the new four-cycle onto the old
Reeb-normal diagonal.
\end{lemma}

\begin{proof}
Put
\[
u=a_{y_0}-a_p,
\qquad
v=a_q-a_{y_0}.
\]
Then \(u,v\in E_{\rm new}\) and \(a_q-a_p=u+v\in E_{\rm new}\).  Therefore
\[
\begin{aligned}
|\mathfrak q_*(a_q)-\mathfrak q_*(a_p)|
&=
|n_*\cdot(a_q-a_p)|\\
&=
|(n_*-n_{\rm new})\cdot(a_q-a_p)|
\lesssim_U
\vartheta,
\end{aligned}
\]
which proves \eqref{eq:v079-coherent-anchor-q-local}.  For the converse, the norm of the
restriction of \(a\mapsto n_*\cdot a\) to \(E_{\rm new}\) is
\(|n_{\rm new}\times n_*|\).  Test it on
\((a_q-a_p)/|a_q-a_p|\) and use the bounded direction chart.  This is also the opposite-anchor
form of \cref{lem:v066-anchor-plane-q-transverse}.
\end{proof}

\begin{lemma}[Weighted \(q\)-local opposite-anchor four-cycle bound]
\label{lem:v079-q-local-anchor-codegree}
Let \(0\le h\le1\) be a bipartite graphon on probability spaces
\((X,\mu_X)\) and \((Y,\mu_Y)\), and write
\[
M
:=
\iint h(x,y)\,d\mu_X(x)d\mu_Y(y).
\]
Assume the \(X\)-degrees obey
\[
d_X(x):=\int h(x,y)\,d\mu_Y(y)\le D_X
\]
and that a scalar coordinate \(\mathfrak q_*:X\to\mathbb R\) has the resolution-\(r_0\)
Frostman bound
\begin{equation}
\label{eq:v079-q-marginal-hypothesis}
(\mathfrak q_*)_\#\mu_X(I)
\le
L_q(|I|+r_0)
\qquad\text{for every interval }I.
\end{equation}
Define the common-neighbor codegree
\[
\mathsf T_X(x,x')
:=
\int h(x,y)h(x',y)\,d\mu_Y(y).
\]
Then, for every \(R\ge0\),
\begin{equation}
\label{eq:v079-q-local-codegree-square}
\boxed{
\iint_{\{|\mathfrak q_*(x)-\mathfrak q_*(x')|\le R\}}
\mathsf T_X(x,x')^2
\,d\mu_X(x)d\mu_X(x')
\lesssim
D_XL_q(R+r_0)M.
}
\end{equation}
The left side is exactly the weighted four-cycle mass whose opposite \(X\)-anchors are
\(q\)-local.  In particular the estimate keeps the full graphon edge weights; it does not replace
the collision graph by its support.
\end{lemma}

\begin{proof}
Since
\[
\mathsf T_X(x,x')
\le
\min\{d_X(x),d_X(x')\}
\le D_X,
\]
one has \(\mathsf T_X^2\le D_X\mathsf T_X\).  By Tonelli and
\eqref{eq:v079-q-marginal-hypothesis},
\[
\begin{aligned}
&\iint_{\{|\mathfrak q_*(x)-\mathfrak q_*(x')|\le R\}}
\mathsf T_X(x,x')
\,d\mu_X(x)d\mu_X(x')
\\
&\quad=
\int_Y
\iint_{\{|\mathfrak q_*(x)-\mathfrak q_*(x')|\le R\}}
h(x,y)h(x',y)
\,d\mu_X(x)d\mu_X(x')d\mu_Y(y)
\\
&\quad\le
C L_q(R+r_0)
\int_Y\int_X h(x,y)\,d\mu_X(x)d\mu_Y(y)
\\
&\quad=
C L_q(R+r_0)M.
\end{aligned}
\]
Multiply by \(D_X\).  Expanding \(\mathsf T_X^2\) gives the alternating product of the four
edge weights, proving the four-cycle interpretation.
\end{proof}

\begin{proposition}[The coherent two-generation four-cycle pays an anchor-diagonal debt]
\label{prop:v079-coherent-four-cycle-anchor-debt}
Apply the normalized Maslov graph construction on the final occurrence layer
\(G_*\), but keep the occurrence variables \((\mathfrak f,z)\) and their weights from
\cref{prop:v078-double-occurrence-collision-lift}.  After the collision-preserving two-sided degree
regularization, view the graph as bipartite on two copies and write its edge mass and first-side
degree cap as \(M\) and \(D_X\).  Suppose the normalized \(q\)-marginal satisfies
\eqref{eq:v079-q-marginal-hypothesis}.

First remove the physical-bush output of the Maslov graph routing, then route its
determinant-small four-cycles by
\cref{prop:v067-no-time-anchored-packet}.  Apart from the close-anchor and horizontal-line
destinations, compare the resulting \(E_{\rm new}\) with the old flag plane \(E_*\).  For every
angular threshold \(\vartheta\), the total four-cycle mass routed to
\[
|n_{\rm new}-n_*|\le\vartheta
\]
is bounded by
\begin{equation}
\label{eq:v079-coherent-four-cycle-bound}
\boxed{
\mathcal Z_{\rm coh}(\vartheta)
\lesssim
D_XL_q(\vartheta+r_0)M.
}
\end{equation}
Every complementary packet has
\(|n_{\rm new}\times n_*|\gtrsim\vartheta\) and therefore returns to the transverse
old/new-polarization intersection ledger of
\cref{lem:v078-old-flag-new-cycle-polarization}.

Consequently, if the rank-loss alternative carries four-cycle mass
\(\mathcal Z_{\rm rk}\gtrsim M^4\) and the coherent class carries a fixed portion of it, then
\begin{equation}
\label{eq:v079-coherent-survivor-power-debt}
\boxed{
M^3
\lesssim
D_XL_q(\vartheta+r_0).
}
\end{equation}
On a fine power layer
\[
M=r^{\chi+o(1)},
\qquad
D_X=r^{\theta+o(1)},
\qquad
L_q=r^{-o(1)},
\qquad
\vartheta=r^{\alpha+o(1)},
\qquad
r_0\lesssim\vartheta,
\]
a coherent survivor must therefore obey
\begin{equation}
\label{eq:v079-coherent-anchor-exponent-region}
\boxed{\alpha+\theta\le3\chi+o(1).}
\end{equation}
This is the exact anchor-diagonal debt left by the coherent old-flag/new-cycle geometry.  It is the necessary region used in the later conditional tangential estimate
\eqref{eq:v078-final-fine-conditional-tangential}.
\end{proposition}

\begin{proof}
For every coherent new packet,
\cref{lem:v079-coherent-new-plane-q-local} places its opposite anchors in
\[
|\mathfrak q_*(a_p)-\mathfrak q_*(a_q)|
\lesssim_U\vartheta.
\]
The rank-loss cycles routed to that packet form a subfamily of all four-cycles with such
\(q\)-local opposite anchors.  Apply
\cref{lem:v079-q-local-anchor-codegree} with \(R=C_U\vartheta\), proving
\eqref{eq:v079-coherent-four-cycle-bound}.  If the normals are not coherent, the second part of
\cref{lem:v079-coherent-new-plane-q-local} and the transverse tube-intersection calculation in
\cref{lem:v078-old-flag-new-cycle-polarization} give the stated return.

If the coherent part is \(\gtrsim M^4\), compare that lower bound with
\eqref{eq:v079-coherent-four-cycle-bound} and divide by \(M>0\), obtaining
\eqref{eq:v079-coherent-survivor-power-debt}.  Substitution of the displayed power
parameterization gives \eqref{eq:v079-coherent-anchor-exponent-region}.
\end{proof}

\begin{lemma}[A prescribed old-flag polarization cap costs one angular factor]
\label{lem:v080-old-flag-normal-cap}
Keep the flag layer of
\cref{prop:v078-flag-packet-incidence-lift}.  For
\(f=(p,q,\omega)\), let \(E_f=E_{p,q,\omega}\) be the contact-polarization plane from
\cref{lem:v071-sector-contact-packet}, and let
\([n_f]\in\mathbb {RP}^2\) be its unoriented normal.  For a prescribed
\([n]\in\mathbb {RP}^2\) and \(0<\vartheta\le1\), define
\[
\mathcal N_{\rm flag}^{[n],\vartheta}(z)
:=
\int
\mathbf 1_{\{d_{\mathbb {RP}^2}([n_f],[n])\le\vartheta\}}
w_f(z)\,d\Lambda(y,f).
\]
If the horizontal marginal of \(\nu_j\) has density bounded by \(C_U\), then
\begin{equation}
\label{eq:v080-old-flag-normal-cap}
\boxed{
\mathcal N_{\rm flag}^{[n],\vartheta}(z)
\lesssim_U
H_i\vartheta
\qquad\text{for every }z.
}
\end{equation}
The estimate is uniform in the prescribed plane and retains both old contact-edge weights and the
root occurrence weight.
\end{lemma}

\begin{proof}
Use the definition of the occurrence measure and integrate the root variable first:
\[
\begin{aligned}
\mathcal N_{\rm flag}^{[n],\vartheta}(z)
&=
\iint\sum_\omega
\mathbf 1_{\{d_{\mathbb {RP}^2}([n_{p,q,\omega}],[n])\le\vartheta\}}
w_{p,q,\omega}(z)
\\[-1mm]
&\hspace{29mm}\times
\left(
\int_{\mathcal H_i}w_{p,q,\omega}(y)\,d\nu_i(y)
\right)
d\nu_j(p)d\nu_j(q).
\end{aligned}
\]
On the retained flag layer the parenthesized factor is at most
\(m_{p,q,\omega}<2H_i\).  Moreover
\[
a_q-a_p\in E_{p,q,\omega}.
\]
Thus the angular-cap condition implies
\[
|n\cdot(a_q-a_p)|
\lesssim_U
\vartheta.
\]
For fixed \(p\), the allowed \(q\)'s lie in a direction slab of width
\(O_U(\vartheta)\), whose \(\nu_j\)-mass is \(O_U(\vartheta)\).  For fixed \(p,q,z\), bounded
overlap of the azimuth intervals leaves only \(O_U(1)\) values of \(\omega\) for which
\(w_{p,q,\omega}(z)>0\).  Delete the remaining factors, which lie in \([0,1]\), integrate first
in \(q\) and then in \(p\), and multiply by \(2H_i\).  This proves
\eqref{eq:v080-old-flag-normal-cap}.
\end{proof}

\clearpage
\section{Bush families, compactness, and Frostman stopping}
The terms bush and hairbrush refer to the classical Kakeya mechanisms of
\cite{Bourgain1991,Wolff1995}; the sticky implementation follows the geometric line developed in
\cite{KatzLabaTao2000,WangZahlSticky}.


\subsection{Flag amplification and reduction to bushes}

\begin{corollary}[Conditional flag normals remain angularly diffuse on the occurrence layer]
\label{cor:v080-conditional-flag-normal-cap}
Disintegrate the normalized occurrence measure over its physical marginal:
\[
d\widehat\rho_*(\mathfrak f,z)
=
d\kappa_z(\mathfrak f)\,d\widehat\mu_*(z),
\qquad
\widehat\mu_*
:=
\frac{\mathbf 1_{G_*}\mathcal N_{\rm flag}\,d\nu_i}
{\|\rho_*\|}.
\]
Then, for \(z\in G_*\),
\begin{equation}
\label{eq:v080-conditional-normal-cap}
\boxed{
\kappa_z
\left\{
\mathfrak f:
d_{\mathbb {RP}^2}([n_f],[n])\le\vartheta
\right\}
\lesssim_U
\min\left\{1,\frac{H_i}{\lambda_*}\vartheta\right\}.
}
\end{equation}
On the fixed-coarse/all-fine scale family,
\begin{equation}
\label{eq:v080-coarse-angular-tax-subpower}
\boxed{
\frac{H_i}{\lambda_*}
=
r^{-o(1)}.
}
\end{equation}
The same conclusion holds on the diagonal pairs of
\cref{lem:v078-coarse-fine-mass-freezing}.
\end{corollary}

\begin{proof}
The conditional measure is
\[
d\kappa_z(\mathfrak f)
=
\mathcal N_{\rm flag}(z)^{-1}
w_f(z)\,d\Lambda(y,f).
\]
Divide \eqref{eq:v080-old-flag-normal-cap} by
\(\mathcal N_{\rm flag}(z)\ge\lambda_*\), proving
\eqref{eq:v080-conditional-normal-cap}.  For a frozen coarse layer,
\(H_i/\lambda_*\) is a fixed finite constant and hence is \(r^{-o(1)}\) as \(r\downarrow0\).
For the diagonal choice, all retained quantities lie on fixed \(W\)-power layers and
\(|\log W|/|\log r|\to0\), exactly as in
\cref{cor:v078-fine-scale-absorbs-coarse-tax}.
\end{proof}

\begin{proposition}[Four old flags at the free cycle endpoint exclude the fully coherent
subunit-defect return]
\label{prop:v080-four-flag-coherent-exclusion}
Let \(h_r\) be the normalized fine-scale Maslov graphon on the occurrence probability space
\((\mathfrak f,z,\widehat\rho_*)\).  Its kernel depends only on the physical endpoints \(z,z'\);
write its edge mass as
\[
M=r^{\chi+o(1)}.
\]
Consider the determinant-small four-cycles returned by the collision-preserving Maslov routing,
and write an ordered cycle as
\[
p,y,q,y',p.
\]
Restrict to the non-line packet output and, before sampling any old flag marks over the free
endpoint \(y'\), use the physical variables \(p,q,y\) to define
\[
E_{\rm new}
:=
\Span\{a_y-a_p,a_q-a_y\}.
\]
The determinant-small condition places \(y'\) in the associated direct packet, while
\(E_{\rm new}\) and its normal are independent of the old flags sampled over \(y'\).  Extend the
physical cycle conditionally by four independent samples from \(\kappa_{y'}\).  Let
\(\mathcal Z_{\rm all\text{-}coh}(\vartheta)\) be the portion on
which all four sampled old-flag normals lie within \(\vartheta\) of the new normal
\([n_{\rm new}]\).  Then
\begin{equation}
\label{eq:v080-four-flag-coherent-cycle-cost}
\boxed{
\mathcal Z_{\rm all\text{-}coh}(\vartheta)
\lesssim_U
\left(\frac{H_i}{\lambda_*}\vartheta\right)^4.
}
\end{equation}
Consequently, if this fully coherent class carries a fixed portion of a rank-loss output of mass
\(\gtrsim M^4\), then
\begin{equation}
\label{eq:v080-coherent-cycle-necessary-angular-scale}
\boxed{
M
\lesssim_U
\frac{H_i}{\lambda_*}\vartheta.
}
\end{equation}
In the fixed-coarse/all-fine regime, with
\(\vartheta=r^{\alpha+o(1)}\), a fully coherent survivor must satisfy
\begin{equation}
\label{eq:v080-coherent-cycle-power-region}
\boxed{\chi\ge\alpha-o(1).}
\end{equation}

In particular, every total-defect layer with \(\chi<1\) admits parameters
\[
\chi<\alpha<d-\beta<1,
\qquad
\Delta=r^d,
\qquad
\theta_{\rm hr}=r^\beta,
\]
for which the direct new-packet thickness obeys
\begin{equation}
\label{eq:v080-new-packet-thinner-than-angle}
\boxed{
t_{\rm new}
\lesssim
r+\frac{\Delta}{\theta_{\rm hr}}
=
o(r^\alpha)
=
o(\vartheta).
}
\end{equation}
The fully coherent rank-loss return is then impossible by
\eqref{eq:v080-coherent-cycle-power-region}; every remaining old/new comparison contains a
\(\vartheta\)-transverse flag at the free endpoint \(y'\) and returns, with the positive fine-scale
thickness ratio \(t_{\rm new}/\vartheta=o(1)\), to the transverse-polarization ledger.  The
close-anchor, horizontal-line, and prior physical-bush outputs retain their existing destinations.
\end{proposition}

\begin{proof}
Condition on the four physical vertices \(p,y,q,y'\).  The new plane
\(E_{\rm new}\) is already fixed by \(p,q,y\), while the four old flags over \(y'\) are independent
under \(\kappa_{y'}^{\otimes4}\).  Apply
\eqref{eq:v080-conditional-normal-cap} four times with the prescribed normal
\([n_{\rm new}]\).  Their conditional product is at most
\[
C_U
\left(\frac{H_i}{\lambda_*}\vartheta\right)^4.
\]
Integrate the physical four-cycle kernel, which is nonnegative and bounded by one.  Since each
\(\kappa_{y'}\) has total mass one, the fourfold conditional extension leaves the unmarked
four-cycle mass unchanged.  This proves \eqref{eq:v080-four-flag-coherent-cycle-cost}.  Comparison
with a coherent lower bound \(\gtrsim M^4\) and taking fourth roots gives
\eqref{eq:v080-coherent-cycle-necessary-angular-scale}; now use
\eqref{eq:v080-coarse-angular-tax-subpower} to obtain
\eqref{eq:v080-coherent-cycle-power-region}.

If \(\chi<1\), choose \(\alpha\in(\chi,1)\), then choose
\(d\in(\alpha,1)\) and
\(\beta\in(0,d-\alpha)\).  Since \(d<1\), one has \(r\ll\Delta\), as required by the
arbitrary-threshold Maslov routing.  Since both \(1\) and \(d-\beta\) exceed \(\alpha\),
\eqref{eq:v080-new-packet-thinner-than-angle} follows from
\cref{lem:v066-direct-horizontal-packet-thickness}.  The coherent necessary condition
\(\chi\ge\alpha-o(1)\) contradicts the strict choice \(\alpha>\chi\).  On the complement, at least
one of the four old-flag normals over \(y'\) is separated from \(n_{\rm new}\) by
\(\vartheta\).  Choose the first such flag and integrate the remaining conditional flag variables
to one.  The unequal-thickness intersection estimate of
\cref{lem:v078-old-flag-new-cycle-polarization,prop:v066-rank-packet-transverse-return} supplies
the factor \(t_{\rm new}/\vartheta=o(1)\) in the genuinely free \(y'\)-variable, with the frozen
old thickness absorbed into the coarse constant.
\end{proof}

\begin{lemma}[Conditional tensoring of old flags is lossless for every physical collision kernel]
\label{lem:v081-lossless-conditional-flag-tensor}
Use the disintegration from
\cref{cor:v080-conditional-flag-normal-cap},
\[
d\widehat\rho_*(\mathfrak f,z)
=
d\kappa_z(\mathfrak f)\,d\widehat\mu_*(z).
\]
For an integer \(k\ge1\), define the \(k\)-flag extension
\begin{equation}
\label{eq:v081-k-flag-extension}
d\widehat\rho_*^{[k]}(\mathfrak f_1,\ldots,\mathfrak f_k,z)
:=
\left(\prod_{j=1}^k d\kappa_z(\mathfrak f_j)\right)
d\widehat\mu_*(z).
\end{equation}
Then \(\widehat\rho_*^{[k]}\) is a probability measure with physical marginal
\(\widehat\mu_*\).  For every \(m\ge1\) and every nonnegative kernel
\(K(z_1,\ldots,z_m)\) depending only on the physical endpoints,
\begin{equation}
\label{eq:v081-physical-kernel-tensor-invariance}
\boxed{
\int K(z_1,\ldots,z_m)
\prod_{\ell=1}^m
d\widehat\rho_*^{[k]}(\boldsymbol{\mathfrak f}_\ell,z_\ell)
=
\int K(z_1,\ldots,z_m)
\prod_{\ell=1}^m d\widehat\mu_*(z_\ell).
}
\end{equation}
In particular the fine Maslov graph edge mass, its four-cycle mass, and every
collision-preserving restriction are unchanged by the lift.  No new geometric incidence is
assumed: the \(k\) marks are independent samples from the actual conditional old-flag
occurrence law at the same physical point.
\end{lemma}

\begin{proof}
For every physical point \(z\), \(\kappa_z\) is a probability measure.  Integrating the product
in \((\mathfrak f_1,\ldots,\mathfrak f_k)\) therefore gives one.  This proves the physical marginal
identity, and applying it independently at the \(m\) physical endpoints proves
\eqref{eq:v081-physical-kernel-tensor-invariance}.
\end{proof}

\begin{proposition}[Eight old flags at the free cycle endpoint close the coherent defect range
below two]
\label{prop:v081-eight-flag-full-defect-exclusion}
Lift the normalized fine Maslov graphon \(h_r\) to
\(\widehat\rho_*^{[8]}\) by
\cref{lem:v081-lossless-conditional-flag-tensor}; write its unchanged edge mass as
\[
M=r^{\chi+o(1)}.
\]
Write an ordered determinant-small cycle as
\[
p,y,q,y',p.
\]
Before the flag lift, use the physical variables \(p,q,y\) to set
\[
E_{\rm new}
=
\Span\{a_y-a_p,a_q-a_y\};
\]
the determinant-small condition places the remaining endpoint \(y'\) in the corresponding direct
packet, while \(E_{\rm new}\) is independent of the old flags sampled over \(y'\).
Let \(\mathcal Z_{\rm eight\text{-}coh}(\vartheta)\) be the cycle mass for which all eight old
flag normals sampled conditionally over this free endpoint lie within \(\vartheta\) of
\([n_{\rm new}]\).  Then
\begin{equation}
\label{eq:v081-eight-flag-coherent-cycle-cost}
\boxed{
\mathcal Z_{\rm eight\text{-}coh}(\vartheta)
\lesssim_U
\left(\frac{H_i}{\lambda_*}\vartheta\right)^8.
}
\end{equation}
If this class carries a fixed portion of a rank-loss output of mass
\(\gtrsim M^4\), it follows that
\begin{equation}
\label{eq:v081-eight-flag-coherent-necessary}
\boxed{
M
\lesssim_U
\left(\frac{H_i}{\lambda_*}\right)^2\vartheta^2.
}
\end{equation}
Thus, in the fixed-coarse/all-fine regime, a fully coherent survivor at
\(\vartheta=r^{\alpha+o(1)}\) must satisfy
\begin{equation}
\label{eq:v081-eight-flag-power-region}
\boxed{\chi\ge2\alpha-o(1).}
\end{equation}

For every bad total-defect exponent \(\chi<2\), choose
\begin{equation}
\label{eq:v081-full-defect-parameter-window}
\boxed{
\frac{\chi}{2}<\alpha<d-\beta<1,
\qquad
\vartheta=r^\alpha,
\qquad
\Delta=r^d,
\qquad
\theta_{\rm hr}=r^\beta.
}
\end{equation}
Then the direct new-packet thickness is \(o(\vartheta)\), the fully coherent output contradicts
the eight-flag power-region inequality above, and every complementary packet has at least one of the
eight old flags transverse to the new polarization at the free endpoint \(y'\).  It therefore
returns with a positive fine-scale thickness gain to
\cref{lem:v078-old-flag-new-cycle-polarization}.  Hence no determinant-small,
non-line, fully coherent rank-loss return survives at any exponent \(\chi<2\).
\end{proposition}

\begin{proof}
Condition on \(p,q,y,y'\), using the displayed ordering only to designate \(y'\) as the free
packet endpoint.  The new normal is already fixed by \(p,q,y\).  The eight old flags over \(y'\)
are independent under \(\kappa_{y'}^{\otimes8}\).  By
\eqref{eq:v080-conditional-normal-cap}, their simultaneous coherence has conditional probability
at most
\[
C_U
\left(\frac{H_i}{\lambda_*}\vartheta\right)^8.
\]
Integrate the physical four-cycle kernel.  It is nonnegative and bounded by one, so
\eqref{eq:v081-eight-flag-coherent-cycle-cost} follows.  The tensor invariance
\eqref{eq:v081-physical-kernel-tensor-invariance} ensures that the rank-loss lower bound remains
\(\gtrsim M^4\).  Compare the two bounds and take fourth roots to obtain
\eqref{eq:v081-eight-flag-coherent-necessary}; then use
\eqref{eq:v080-coarse-angular-tax-subpower} to get
\eqref{eq:v081-eight-flag-power-region}.

If \(\chi<2\), the interval \((\chi/2,1)\) is nonempty.  Choose
\(\alpha\) there, then \(d\in(\alpha,1)\) and
\(\beta\in(0,d-\alpha)\).  Exactly as in
\eqref{eq:v080-new-packet-thinner-than-angle},
\[
t_{\rm new}
\lesssim
r+r^{d-\beta}
=
o(r^\alpha).
\]
The strict inequality \(2\alpha>\chi\) contradicts
\eqref{eq:v081-eight-flag-power-region}.  On the complement, one of the eight actual old flags
over the free endpoint \(y'\) is \(\vartheta\)-transverse to the new packet.  Choose the first such
flag, integrate the other seven conditional probability variables to one, and apply the
unequal-thickness transverse intersection bound in the genuinely free \(y'\)-variable, as recorded in
\cref{lem:v078-old-flag-new-cycle-polarization,prop:v066-rank-packet-transverse-return}.  Since
\(t_{\rm new}/\vartheta=o(1)\) and the old layer is frozen, this gives the required positive
fine-scale gain.
\end{proof}

\begin{corollary}[Finite flag amplification excludes coherent rank loss at any prescribed
mesoscopic threshold]
\label{cor:v081-k-flag-arbitrary-threshold}
Let
\[
M=r^{\chi+o(1)},
\qquad
0\le\chi<2,
\]
and prescribe any \(0<d<1\), so that \(\Delta=r^d\).  Choose an integer
\[
k>\frac{4\chi}{d},
\]
then choose
\begin{equation}
\label{eq:v081-k-flag-parameter-window}
\boxed{
\frac{4\chi}{k}<\alpha<d-\beta<d,
\qquad
\vartheta=r^\alpha,
\qquad
\theta_{\rm hr}=r^\beta.
}
\end{equation}
On the \(k\)-flag extension
\(\widehat\rho_*^{[k]}\), a fully coherent rank-loss output of mass
\(\gtrsim M^4\) would force
\begin{equation}
\label{eq:v081-k-flag-coherent-survivor}
\boxed{
M
\lesssim_U
\left(\frac{H_i}{\lambda_*}\vartheta\right)^{k/4},
\qquad
\chi\ge\frac{k\alpha}{4}-o(1),
}
\end{equation}
which contradicts \eqref{eq:v081-k-flag-parameter-window}.  Since
\[
t_{\rm new}
\lesssim
r+r^{d-\beta}
=o(r^\alpha),
\]
every complementary rank-loss packet has a transverse old flag at the free endpoint and pays the
fine-scale thickness gain.  Thus, after the close-anchor, horizontal-line, and transverse ledgers
are removed, the arbitrary-threshold Maslov routing produces a physical
\(O(r^{1-d})\)-Lagrangian bush for every prescribed \(d\in(0,1)\).
\end{corollary}

\begin{proof}
By
\cref{lem:v081-lossless-conditional-flag-tensor}, the \(k\)-flag lift leaves the physical
four-cycle mass unchanged.  Applying
\eqref{eq:v080-conditional-normal-cap} independently to the \(k\) flags over the free endpoint gives
\[
\mathcal Z_{k\text{-}\rm coh}(\vartheta)
\lesssim_U
\left(\frac{H_i}{\lambda_*}\vartheta\right)^k.
\]
Comparison with \(M^4\), followed by fourth roots, proves
\eqref{eq:v081-k-flag-coherent-survivor}.  The parameter inequalities give both its contradiction
and \(t_{\rm new}=o(\vartheta)\).  The first transverse flag is then treated exactly as in the proof
of \cref{prop:v081-eight-flag-full-defect-exclusion}.
\end{proof}

\begin{corollary}[The final compatibility is either transverse-paid or one coherent Maslov diagonal]
\label{cor:v078-final-two-generation-routing}
After dyadic regularization of \(t_{\rm old},t_{\rm new}\) and a slowly vanishing choice
\(\vartheta=W^{o(1)}\), every old-flag/new-cycle contribution from
\cref{pr:v078-final-six-point-compatibility} has one of two destinations:
\begin{enumerate}[label=\textnormal{(\roman*)}]
\item it pays the transverse factor
\((t_{\rm old}+r)(t_{\rm new}+r)/\vartheta\) from
\eqref{eq:v078-two-generation-transverse-cost};
\item all retained old and new packet data lie in one cell-locally aligned
\(E\oplus E\) polarization, and the new collision kernel factors exactly as
\begin{equation}
\label{eq:v078-final-coherent-normal-tangential-factorization}
\boxed{
\mathbf 1_{\{|\Pi_s(z)-\Pi_s(z')|\le r\}}
\le
\mathbf 1_{\{|N_s(z,z')|\le r\}}
\mathbf 1_{\{|T_s(z,z')|\le r\}}.
}
\end{equation}
\end{enumerate}
Hence the sole coherent remainder is the relative tangential diagonal conditioned by the normal
Reeb window, now for the low old-flag-reuse measure.  This is precisely the symplectic
normal/tangential interface of \cref{lem:v056-polarized-collision-factorization}; no direct use of
\cref{lem:v046-shared-two-edge-flag} is made without its containment hypothesis.
Moreover, after collision-preserving degree regularization, its rank-loss return is subject to the
quantitative anchor-diagonal constraint
\eqref{eq:v079-coherent-survivor-power-debt}; any four-cycle mass exceeding that bound is already
in the transverse or physical-output alternatives.
The lossless eight-flag extension at the free endpoint is stronger still.  By
\cref{lem:v081-lossless-conditional-flag-tensor}, it leaves the physical collision graph and its
four-cycle lower bound unchanged, while
\cref{prop:v081-eight-flag-full-defect-exclusion} removes the fully coherent rank-loss return for
every bad fine total-defect exponent \(\chi<2\).  Thus failure of
\eqref{eq:v078-final-fine-conditional-tangential} has no independent coherent incidence
destination: it returns to a transverse old/new polarization, a horizontal-line or close-anchor
chart, or a physical Lagrangian bush.
\end{corollary}

\begin{proposition}[Eight-flag amplification reduces every bad coherent rank-loss return to a
physical bush]
\label{prop:v081-rank-loss-to-bush-reduction}
Fix one coarse low-reuse occurrence layer \(G_*\), normalize its physical occurrence marginal, and
let \(r\downarrow0\) range over fine collision scales.  Suppose that collision-preserving pruning
and two-sided degree regularization give a normalized Maslov graph with
\[
M=r^{\chi+o(1)},
\qquad
0\le\chi<2.
\]
Choose
\[
\frac{\chi}{2}<\alpha<d-\beta<1,
\qquad
\vartheta=r^\alpha,
\qquad
\Delta=r^d,
\qquad
\theta_{\rm hr}=r^\beta.
\]
Then the arbitrary-threshold routing has one of the following four destinations:
\begin{enumerate}[label=\textnormal{(\roman*)}]
\item a close-anchor/angular or horizontal-line packet;
\item a transverse old/new polarization packet, with thickness gain
\[
\frac{t_{\rm new}}{\vartheta}
\lesssim
r^{1-\alpha}+r^{d-\beta-\alpha}
=o(1);
\]
\item a physical \(O(r^{1-d})\)-Lagrangian point-probe bush of normalized occurrence mass
\(\gtrsim M^3\);
\item a fully coherent determinant-small rank-loss packet.
\end{enumerate}
The fourth alternative is impossible.  Hence, after the first two destinations have been paid in
their weighted ledgers, every remaining bad collision layer produces the physical bush in
\textnormal{(iii)}.

Writing \(m_*=\nu_i(G_*)>0\), that bush has unrepeated selector-direction mass
\(\gtrsim_{G_*}m_*M^3\), and the exact Reeb escape map gives
\begin{equation}
\label{eq:v081-forced-bush-neighborhood-credit}
\boxed{
\mathcal L^4\!\left(N_{C_{G_*}r^{1-d}}(K_b)\right)
\gtrsim_{G_*}
m_*M^3
=
r^{3\chi+o(1)}.
}
\end{equation}
\end{proposition}

\begin{proof}
Apply
\cref{prop:v057-total-defect-four-cycle-return,prop:v050-arbitrary-threshold-maslov-routing}.
Its first output is \textnormal{(iii)}.  Split its determinant-small output by the anchored
horizontal geometry.  The close-anchor and horizontal-line pieces give \textnormal{(i)}; the
remaining cycles lie in the direct contact-polarized packet of
\cref{prop:v067-no-time-anchored-packet}.  Order such a cycle as \(p,y,q,y',p\), with \(y'\) the
free packet endpoint, and lift that endpoint to eight conditional old flags by
\cref{lem:v081-lossless-conditional-flag-tensor}.  The fully coherent part is impossible by
\cref{prop:v081-eight-flag-full-defect-exclusion}.  On its complement, the first transverse flag at
\(y'\) gives \textnormal{(ii)} through
\cref{lem:v078-old-flag-new-cycle-polarization}; the displayed gain is the direct packet thickness
divided by \(\vartheta\).

On the frozen occurrence layer,
\(\lambda_*\le\mathcal N_{\rm flag}<2\lambda_*\), so the normalized physical occurrence marginal
is comparable, with constants depending only on \(G_*\), to
\(m_*^{-1}\nu_i|_{G_*}\).  Therefore a normalized bush of mass \(\gtrsim M^3\) contains actual
selector-direction mass \(\gtrsim_{G_*}m_*M^3\).  Its radius is
\(O(r/\Delta)=O(r^{1-d})\).  The inclusion of its ideal Reeb escape in the corresponding
neighborhood of \(K_b\), followed by the one-bush volume formula
\cref{prop:v014-near-K,thm:v014-one-volume}, proves
\eqref{eq:v081-forced-bush-neighborhood-credit}.
\end{proof}

\subsection{Maximal bush extraction and vector Frostman escape}

\begin{proposition}[Maximal extraction gives a paid induced branch or an \(M\)-mass bush family]
\label{prop:v081-maximal-bush-family}
Keep the notation of
\cref{prop:v081-rank-loss-to-bush-reduction}, fix any \(d\in(0,1)\), and choose the finite flag
number and angular parameters from
\cref{cor:v081-k-flag-arbitrary-threshold}.  Let \(h\) be the normalized Maslov graph and suppose
\[
e(h)=M=r^{\chi+o(1)},
\qquad
0\le\chi<2,
\]
Then at least one of the following holds:
\begin{enumerate}[label=\textnormal{(\roman*)}]
\item there is an induced measurable subgraph \(A\) whose absolute edge mass is at least \(M/2\)
and whose arbitrary-threshold routing places a fixed part of its normalized rank-loss output in a
close-anchor/angular, horizontal-line, or transverse-polarization ledger;
\item there is a finite family of pairwise disjoint selector sets
\[
H_1,\ldots,H_N\subset G_*,
\qquad
N\lesssim M^{-3},
\]
such that every \(H_j\) is an \(O(r^{1-d})\)-Lagrangian point-probe bush and
\begin{equation}
\label{eq:v081-maximal-bush-direction-mass}
\boxed{
\sum_{j=1}^N\widehat\mu_*(H_j)
\gtrsim M.
}
\end{equation}
Every \(H_j\) in \textnormal{(ii)} may be chosen with the conditional determinant-nondegenerate
origin witness of \cref{lem:v081-origin-marked-bush}; the origin-witness kernel has physical
marginal \(\widehat\mu_*|_{H_j}\) and survives every later restriction.
\end{enumerate}
In the second alternative, let \(E_j\) be the corresponding ideal Reeb escape, trimmed to a fixed
outer Reeb wing, and use
the original selector-direction normalization on \(G_*\).  Then
\begin{equation}
\label{eq:v081-maximal-bush-escaped-L1}
\boxed{
\sum_{j=1}^N |E_j|
\gtrsim_{G_*}
m_*M.
}
\end{equation}
Writing
\[
\mathcal O_{\rm bush}
:=
\sum_{1\le j<k\le N}|E_j\cap E_k|,
\]
the exact escaped second moment gives
\begin{equation}
\label{eq:v081-maximal-bush-physical-second-moment}
\boxed{
\mathcal L^4\!\left(N_{C_{G_*}r^{1-d}}(K_b)\right)
\gtrsim_{G_*}
\frac{(m_*M)^2}
{m_*M+\mathcal O_{\rm bush}}.
}
\end{equation}
Thus the loss from \(M\) to \(M^3\) is not intrinsic to the physical output.  After maximal
de-multiplication, the only remaining loss is the overlap of the actual Reeb-escaped bushes.
\end{proposition}

\begin{proof}
Work with the normalized physical occurrence marginal
\(\widehat\mu_*\).  Make a finite maximal selection of disjoint bush leaf sets.  If at some stage
the complement \(A\) has induced edge mass
\[
\iint_{A\times A}h\,d\widehat\mu_*d\widehat\mu_*
\ge
\frac M2,
\]
normalize \(\widehat\mu_*|_A\).  Its normalized edge mass is at least \(M/2\).  The same
free-endpoint conditional flag law is available on \(A\), and the exponent of its edge mass is at
most \(\chi+o(1)\).  Apply the arbitrary-threshold routing and
\cref{cor:v081-k-flag-arbitrary-threshold}.  If a paid geometric output occurs, alternative
\textnormal{(i)} holds and the selection stops.  Otherwise use the origin-marked form
\cref{lem:v081-origin-marked-bush} of the \(O(r^{1-d})\)-bush.  Its unnormalized
\(\widehat\mu_*\)-mass is at least
\[
\widehat\mu_*(A)
\left(
\frac{M/2}{\widehat\mu_*(A)^2}
\right)^3
\gtrsim
M^3.
\]
Assume alternative \textnormal{(i)} never occurs.  Consequently a maximal bush family is finite
and has \(N\lesssim M^{-3}\).  Its complement has induced
edge mass below \(M/2\).  If \(U=\bigcup_jH_j\), then \(0\le h\le1\) gives
\[
\frac M2
<
e(h)-e(h|_{U^c\times U^c})
\le
2\widehat\mu_*(U).
\]
This proves \eqref{eq:v081-maximal-bush-direction-mass}.

The dyadic occurrence bounds on \(G_*\) convert the normalized mass in
\eqref{eq:v081-maximal-bush-direction-mass} into original selector-direction mass
\(\gtrsim_{G_*}m_*M\).  The \(H_j\) are disjoint, and the outer-wing one-bush Jacobian is bounded
above and below.  The one-bush volume formula therefore gives
\eqref{eq:v081-maximal-bush-escaped-L1}.  Finally put
\[
\mathcal M_{\rm bush}:=\sum_j\mathbf 1_{E_j}.
\]
Then
\[
\int\mathcal M_{\rm bush}=\sum_j|E_j|,
\qquad
\int\mathcal M_{\rm bush}^2
=
\sum_j|E_j|+2\mathcal O_{\rm bush}.
\]
Since every \(E_j\subset N_{C_{G_*}r^{1-d}}(K_b)\), Cauchy--Schwarz proves
\eqref{eq:v081-maximal-bush-physical-second-moment}.
\end{proof}

\begin{principle}[The remaining escaped overlap is again a contact--Maslov collision]
\label{pr:v081-bush-overlap-contact-return}
For every pair \(j\ne k\), the exact affine homothety formula
\cref{thm:v014-pair-overlap} writes
\[
|E_j\cap E_k|
=
|U|\int |s-s_j|^3
\sigma\!\left(H_j\cap\Psi_{jk,s}^{-1}(H_k)\right)\,ds.
\]
Every direction pair counted by this integral obeys the secondary physical collision
\eqref{eq:v014-secondary-collision} at thickness \(O(r^{1-d})\).  Splitting those pairs by their
angular separation gives exactly two destinations: the fixed-angle part is a coarser
Reeb-window Maslov graph, and the near-angle part is an angular Darboux refinement.  Hence the
denominator \(\mathcal O_{\rm bush}\) in
\eqref{eq:v081-maximal-bush-physical-second-moment} contains no unstructured set-overlap term.  Its
weighted summation across the maximal family, with the old flag occurrence measures retained, is
the contact-overlap estimate treated below.
\end{principle}

\begin{theorem}[A direction-Frostman Lagrangian bush escapes to a front-Frostman measure]
\label{thm:v081-bush-frostman-escape}
Let \(r_n\downarrow0\).  For each \(n\), suppose \(H_n\subset U\) is an
\(\varepsilon_n\)-approximate Lagrangian point-probe bush centered at
\((c_n,s_n)\), where
\[
\varepsilon_n\downarrow0,
\qquad
s_n\in J\Subset I,
\]
and let \(\eta_n\) be the normalized selector-direction measure on \(H_n\).  Assume that for some
\(0<\delta<3\), some \(\rho_n\downarrow0\) with
\(\varepsilon_n\le\rho_n\), and one constant \(C_\delta\),
\begin{equation}
\label{eq:v081-bush-direction-frostman}
\boxed{
\eta_n(B(a,R))
\le
C_\delta R^{3-\delta}
\qquad
(a\in\mathbb R^3,\ \rho_n\le R\le1).
}
\end{equation}
Then the actual Sticky Reeb front carries a nonzero
\((4-\delta)\)-Frostman measure.  In particular,
\begin{equation}
\label{eq:v081-bush-frostman-dimension}
\boxed{
\dim_HK_b\ge4-\delta.
}
\end{equation}
If such bushes exist for every \(\delta>0\), then
\(\dim_HK_b=4\).
\end{theorem}

\begin{proof}
Choose once and for all an interval \(I_{\rm out}\Subset I\) with positive length and
\[
\dist(I_{\rm out},J)>0.
\]
For \(a\in H_n\), write the actual selector front as
\[
F_n(a,s):=(b(a)+sa,s).
\]
The approximate bush relation gives, uniformly for \(s\in I_{\rm out}\),
\begin{equation}
\label{eq:v081-actual-to-ideal-bush}
\left|
b(a)+sa-\bigl(c_n+(s-s_n)a\bigr)
\right|
\lesssim
\varepsilon_n.
\end{equation}
Let
\[
\Lambda_n
:=
(F_n)_\#
\left(
\eta_n\otimes\frac{\mathbf1_{I_{\rm out}}\,ds}{|I_{\rm out}|}
\right).
\]
These are probability measures supported on \(K_b\).  Fix a physical ball
\(B((y_0,s_0),R)\) with \(R\ge\rho_n\).  Only Reeb parameters in an interval of length
\(O(R)\) contribute.  At each such parameter,
\eqref{eq:v081-actual-to-ideal-bush} and
\(\dist(I_{\rm out},J)>0\) place every contributing direction in a ball of radius
\[
\lesssim R+\varepsilon_n
\lesssim R.
\]
Therefore \eqref{eq:v081-bush-direction-frostman} gives
\begin{equation}
\label{eq:v081-approx-front-frostman}
\boxed{
\Lambda_n(B((y_0,s_0),R))
\lesssim_{\delta,I,J}
R^{4-\delta}
\qquad
(R\ge\rho_n).
}
\end{equation}

The front lies in a fixed compact region after the standing chart localization, so a subsequence
of \(\Lambda_n\) converges weakly to a probability measure \(\Lambda\) supported on \(K_b\).
For every fixed \(R>0\), eventually \(R\ge\rho_n\).  Choose a continuous cutoff equal to one on a
given closed \(R\)-ball and supported in its \((R+\zeta)\)-enlargement.  Weak convergence and
\eqref{eq:v081-approx-front-frostman} give an upper bound
\(C(R+\zeta)^{4-\delta}\); let \(\zeta\downarrow0\).  Thus \(\Lambda\) is
\((4-\delta)\)-Frostman, proving \eqref{eq:v081-bush-frostman-dimension}.
\end{proof}

\begin{theorem}[Reeb-coherent vector Frostman escape for a family of Lagrangian bushes]
\label{thm:v081-bush-family-frostman-escape}
For every \(n\), let
\[
\{H_{n,j}:1\le j\le N_n\}
\]
be pairwise disjoint selector-direction sets.  Suppose \(H_{n,j}\) is an
\(\varepsilon_n\)-approximate Lagrangian point-probe bush centered at
\((c_{n,j},s_{n,j})\), where
\[
\varepsilon_n\downarrow0,
\qquad
s_{n,j}\in J\Subset I.
\]
Put
\[
q_n:=\sum_{j=1}^{N_n}\sigma(H_{n,j})>0.
\]
Choose a fixed outer interval \(I_{\rm out}\Subset I\) with
\(\dist(I_{\rm out},J)>0\), and define
\[
\mathfrak a_{n,j}(y,s)
:=
\frac{y-c_{n,j}}{s-s_{n,j}}
\qquad
(y\in\mathbb R^3,\ s\in I_{\rm out}).
\]
Assume that for some \(0<\delta<3\), some
\(\rho_n\downarrow0\) with \(\varepsilon_n\le\rho_n\), and one constant \(C_\delta\),
\begin{equation}
\label{eq:v081-vector-bush-direction-frostman}
\boxed{
\sup_{\substack{y\in\mathbb R^3\\ s\in I_{\rm out}}}
\frac1{q_n}
\sum_{j=1}^{N_n}
\sigma\!\left(
H_{n,j}\cap B(\mathfrak a_{n,j}(y,s),R)
\right)
\le
C_\delta R^{3-\delta}
\qquad
(\rho_n\le R\le1).
}
\end{equation}
Then \(K_b\) carries a nonzero
\((4-\delta)\)-Frostman measure and
\[
\dim_HK_b\ge4-\delta.
\]
Thus \eqref{eq:v081-vector-bush-direction-frostman}, rather than a separate Frostman bound on one
arbitrarily selected bush, is the exact Reeb-coherent physical aggregation condition for the maximal family in
\cref{prop:v081-maximal-bush-family}.
\end{theorem}

\begin{proof}
Define a probability measure on the actual front by
\[
\Lambda_n
:=
\frac1{q_n|I_{\rm out}|}
\sum_{j=1}^{N_n}
(F|_{H_{n,j}\times I_{\rm out}})_\#
\bigl(\sigma|_{H_{n,j}}\otimes ds\bigr),
\qquad
F(a,s):=(b(a)+sa,s).
\]
Fix a physical \(R\)-ball centered at \((y_0,s_0)\), with \(R\ge\rho_n\).  Only a Reeb interval
of length \(O(R)\) contributes.  At each contributing \(s\), the approximate bush relation and
\(\dist(I_{\rm out},J)>0\) place every contributing direction in
\[
B\!\left(\mathfrak a_{n,j}(y_0,s),C(R+\varepsilon_n)\right)
\subset
B\!\left(\mathfrak a_{n,j}(y_0,s),C'R\right).
\]
Therefore \eqref{eq:v081-vector-bush-direction-frostman}, with the same physical \(y_0,s\) for
all components, gives after summing in \(j\), dividing by \(q_n\), and integrating the
\(O(R)\)-length time interval
\[
\Lambda_n(B((y_0,s_0),R))
\lesssim_{\delta,I,J}
R^{4-\delta}.
\]
Weak compactness and the continuous-cutoff argument from
\cref{thm:v081-bush-frostman-escape} produce a
\((4-\delta)\)-Frostman weak limit supported on \(K_b\).
\end{proof}

\subsection{Coherent hairbrushes and re-collision}

\begin{corollary}[Failure of bush-family escape is a Reeb-coherent hairbrush concentration]
\label{cor:v081-bush-family-angular-concentration}
Fix \(0<\delta<3\).  If no subsequence of the maximal bush families satisfies
\eqref{eq:v081-vector-bush-direction-frostman} with this \(\delta\), then there are a subsequence,
\(A_n\to\infty\), and radii
\[
\varepsilon_n\le R_n\downarrow0,
\]
and points \(y_n\in\mathbb R^3\), \(s_n^{\rm out}\in I_{\rm out}\) such that, with
\[
a_{n,j}:=\mathfrak a_{n,j}(y_n,s_n^{\rm out}),
\]
\begin{equation}
\label{eq:v081-bush-family-angular-concentration}
\boxed{
\sum_{j=1}^{N_n}
\sigma\!\left(H_{n,j}\cap B(a_{n,j},R_n)\right)
>
A_nq_nR_n^{3-\delta}.
}
\end{equation}
The caps retain their bush-center labels, selected Reeb windows, and conditional old-flag laws.
More importantly, all cap centers are generated by the same physical outer-wing point
\((y_n,s_n^{\rm out})\).  Thus the sole non-Frostman output of the maximal physical family is a
genuine Reeb-coherent multi-bush hairbrush, not an arbitrary simultaneous cap choice.
\end{corollary}

\begin{proof}
Take the contrapositive of
\eqref{eq:v081-vector-bush-direction-frostman} with constants \(A_n\to\infty\).  The left side of
\eqref{eq:v081-bush-family-angular-concentration} is at most \(q_n\); hence
\[
R_n^{3-\delta}<A_n^{-1},
\]
which forces \(R_n\to0\).
\end{proof}

\begin{corollary}[The \(M^{-4}\) debt is exactly a lower bound for the bad hairbrush scale]
\label{cor:v081-hairbrush-bad-scale-lower-bound}
For the maximal families in alternative \textnormal{(ii)} of
\cref{prop:v081-maximal-bush-family}, every Reeb-coherent failure
\eqref{eq:v081-bush-family-angular-concentration} satisfies
\begin{equation}
\label{eq:v081-hairbrush-bad-scale-lower-bound}
\boxed{
R_n^\delta
\gtrsim_{U,G_*}
A_n\frac{q_n}{N_n}
\gtrsim_{U,G_*}
A_nM_n^4.
}
\end{equation}
Equivalently, at every radius satisfying
\begin{equation}
\label{eq:v081-hairbrush-automatic-frostman-range}
R^\delta\le c_Uq_n/N_n,
\end{equation}
the Reeb-coherent family Frostman bound holds automatically.  Thus the factor \(M_n^{-4}\) in
\eqref{eq:v081-selected-family-cap-density} is not an artifact of selecting one component: it is
the exact family cardinality-to-mass scale barrier.
\end{corollary}

\begin{proof}
Every direction cap of radius \(R\) has normalized Lebesgue mass \(O_U(R^3)\).  Hence, uniformly in
the common outer point,
\[
\sum_{j=1}^{N_n}
\sigma\!\left(H_{n,j}\cap B(\mathfrak a_{n,j}(y,s),R)\right)
\lesssim_U N_nR^3.
\]
Comparison with \eqref{eq:v081-bush-family-angular-concentration} gives
\(A_nq_nR_n^{3-\delta}\lesssim_UN_nR_n^3\), which is the first inequality in
\eqref{eq:v081-hairbrush-bad-scale-lower-bound}.  Use
\(q_n\gtrsim_{G_*}M_n\) and \(N_n\lesssim M_n^{-3}\) for the second.  The same cap-volume bound
divided by \(q_n\) proves
\eqref{eq:v081-hairbrush-automatic-frostman-range}.
\end{proof}

\begin{lemma}[A Reeb-coherent cap family is one exact re-collision cocycle orbit]
\label{lem:v081-hairbrush-recollision-cocycle}
Fix one outer point \((y,s)\in\mathbb R^3\times I_{\rm out}\).  For every pair of distinct bush labels
\(j,k\),
\begin{equation}
\label{eq:v081-hairbrush-center-transport}
\boxed{
\Psi_{jk,s}\!\left(\mathfrak a_{n,j}(y,s)\right)
=
\mathfrak a_{n,k}(y,s).
}
\end{equation}
Moreover, uniformly for \(s\in I_{\rm out}\),
\begin{equation}
\label{eq:v081-hairbrush-cap-transport}
\Psi_{jk,s}\!\left(
B(\mathfrak a_{n,j}(y,s),R)
\right)
\subset
B(\mathfrak a_{n,k}(y,s),C_{I,J}R).
\end{equation}
If \(|u-s|\le R\) and the cap on the left meets \(U\), the same uniform bound gives
\begin{equation}
\label{eq:v081-hairbrush-moving-cap-transport}
\Psi_{jk,u}\!\left(
B(\mathfrak a_{n,j}(y,s),R)
\right)
\subset
B(\mathfrak a_{n,k}(y,s),C_{I,J,U}R).
\end{equation}
Thus the caps in
\eqref{eq:v081-bush-family-angular-concentration} are not independently centered angular
patches: they are a single orbit of the exact contact re-collision maps.
\end{lemma}

\begin{proof}
By \eqref{eq:v014-Psi},
\[
\begin{aligned}
\Psi_{jk,s}\!\left(\mathfrak a_{n,j}(y,s)\right)
&=
\frac{c_{n,j}-c_{n,k}+(s-s_{n,j})(y-c_{n,j})/(s-s_{n,j})}
{s-s_{n,k}}\\
&=
\frac{y-c_{n,k}}{s-s_{n,k}}
=
\mathfrak a_{n,k}(y,s).
\end{aligned}
\]
The linear part of \(\Psi_{jk,s}\) is the scalar
\((s-s_{n,j})/(s-s_{n,k})\).  Since \(s\in I_{\rm out}\),
\(s_{n,j},s_{n,k}\in J\), and \(\dist(I_{\rm out},J)>0\), its absolute value is uniformly
bounded.  This proves \eqref{eq:v081-hairbrush-cap-transport}.
For \(|u-s|\le R\), subtract the two explicit affine formulas and use that every cap meeting
\(U\) has bounded center.  The coefficient denominators remain separated from zero, so the center
error is \(O_{I,J,U}(R)\), proving
\eqref{eq:v081-hairbrush-moving-cap-transport}.
\end{proof}

\begin{lemma}[The common-hairbrush blow-up lands exactly in the Lagrangian pencil]
\label{lem:v081-hairbrush-maslov-blowup}
Suppress \(n\), write
\[
s_0:=s_n^{\rm out},
\qquad
y_0:=y_n,
\qquad
d_j:=s_0-s_{n,j},
\qquad
a_j:=\frac{y_0-c_{n,j}}{d_j},
\]
and put
\[
a=a_j+R_n\xi,
\qquad
s=s_0+R_n\tau.
\]
The ideal outer-wing front of the \(j\)-th bush, translated by \((y_0,s_0)\) and dilated by
\(R_n^{-1}\), is exactly
\begin{equation}
\label{eq:v081-hairbrush-rescaled-front}
\boxed{
R_n^{-1}
\left(
c_{n,j}+(s-s_{n,j})a-y_0
\right)
=
d_j\xi+\tau a_j+R_n\tau\xi.
}
\end{equation}
For two labels \(j\ne k\), directions
\(a_j+R_n\xi\) and \(a_k+R_n\eta\) re-collide at the scaled time \(\tau\) if and only if
\begin{equation}
\label{eq:v081-hairbrush-rescaled-maslov}
\boxed{
\beta_{jk}+\tau\alpha_{jk}
+R_n\tau(\xi-\eta)
=0,
\qquad
\alpha_{jk}:=a_j-a_k,
\qquad
\beta_{jk}:=d_j\xi-d_k\eta.
}
\end{equation}
Consequently every ideal blow-up limit \(R_n\to0\) obeys
\begin{equation}
\label{eq:v081-hairbrush-limit-maslov}
\boxed{
\beta_{jk}+\tau\alpha_{jk}=0,
\qquad
(\alpha_{jk},\beta_{jk})\in L_\tau.
}
\end{equation}
Thus a Reeb-coherent Frostman failure returns, under its canonical physical blow-up, to the same
Maslov/Lagrangian incidence pencil used in the angular-collapse routing.
For the actual selector front, the rescaled error relative to
\eqref{eq:v081-hairbrush-rescaled-front} is at most
\[
\varepsilon_n/R_n.
\]
Hence \eqref{eq:v081-hairbrush-limit-maslov} is an exact actual-front limit when
\(\varepsilon_n/R_n\to0\); without that separation one obtains the corresponding finite-scale
Maslov incidence at thickness \(O(\varepsilon_n/R_n)\), and no zero-error limit is claimed.
\end{lemma}

\begin{proof}
The identity \(c_{n,j}=y_0-d_ja_j\) gives
\[
\begin{aligned}
c_{n,j}+(s_0+R_n\tau-s_{n,j})(a_j+R_n\xi)-y_0
&=
R_n(d_j\xi+\tau a_j)+R_n^2\tau\xi,
\end{aligned}
\]
which is \eqref{eq:v081-hairbrush-rescaled-front}.  Equality of the two rescaled horizontal
fronts is exactly \eqref{eq:v081-hairbrush-rescaled-maslov}.  Letting \(R_n\to0\) gives
\eqref{eq:v081-hairbrush-limit-maslov}, and the last membership is the definition of the
Lagrangian graph \(L_\tau=\{(\alpha,\beta):\beta+\tau\alpha=0\}\).
\end{proof}

\begin{proposition}[A Reeb-coherent bad cap family merges into one larger physical Lagrangian bush]
\label{prop:v081-hairbrush-merges-to-one-bush}
For the caps in \eqref{eq:v081-bush-family-angular-concentration}, put
\[
A_{n,j}:=H_{n,j}\cap B(a_{n,j},R_n),
\qquad
S_n:=\sum_j\sigma(A_{n,j}),
\]
and
\[
H_n^{\rm hair}
:=
\bigcup_{j=1}^{N_n}A_{n,j}.
\]
Then the disjointness of the \(H_{n,j}\) gives
\begin{equation}
\label{eq:v081-hairbrush-merged-mass}
\boxed{
\sigma(H_n^{\rm hair})=S_n
>
A_nq_nR_n^{3-\delta}.
}
\end{equation}
Moreover \(H_n^{\rm hair}\) is one
\(C_{I,J}(R_n+\varepsilon_n)\)-Lagrangian point-probe bush centered at the actual outer-wing point
\((y_n,s_n^{\rm out})\).  Since \(\varepsilon_n\le R_n\),
\begin{equation}
\label{eq:v081-hairbrush-merged-bush}
\boxed{
|b(a)+s_n^{\rm out}a-y_n|
\lesssim_{I,J}R_n
\qquad(a\in H_n^{\rm hair}).
}
\end{equation}
Consequently the exact Reeb escape map gives the physical credit
\begin{equation}
\label{eq:v081-hairbrush-merged-volume}
\boxed{
\mathcal L^4\!\left(N_{C_{I,J}R_n}(K_b)\right)
\gtrsim_{I,J,U}
S_n
>
A_nq_nR_n^{3-\delta}.
}
\end{equation}
Thus a failure of the Reeb-coherent family Frostman bound is itself a larger physical bush, not
merely an angular label attached to the old bush family.
\end{proposition}

\begin{proof}
For \(a\in A_{n,j}\), the \(\varepsilon_n\)-bush relation and the definition of
\(a_{n,j}\) give
\[
\begin{aligned}
|b(a)+s_n^{\rm out}a-y_n|
&\le
|b(a)-(c_{n,j}-s_{n,j}a)|
+
|c_{n,j}+(s_n^{\rm out}-s_{n,j})a-y_n|\\
&\le
\varepsilon_n
+
|s_n^{\rm out}-s_{n,j}|\,|a-a_{n,j}|\\
&\lesssim_{I,J}
\varepsilon_n+R_n.
\end{aligned}
\]
This proves \eqref{eq:v081-hairbrush-merged-bush};
\eqref{eq:v081-hairbrush-merged-mass} is the cap concentration inequality.  Because
\(I_{\rm out}\Subset I\), one can choose, uniformly in \(s_n^{\rm out}\), a Reeb subinterval of
fixed positive length inside \(I\) and a fixed positive distance from \(s_n^{\rm out}\).  Apply
\cref{prop:v014-near-K,thm:v014-one-volume} to the merged bush on that interval.  Its ideal escape
has volume comparable to \(S_n\) and lies in the \(C_{I,J}R_n\)-neighborhood of the actual front,
proving \eqref{eq:v081-hairbrush-merged-volume}.
\end{proof}

\begin{corollary}[At the hairbrush scale the physical product collision is already large]
\label{cor:v081-hairbrush-coarse-product-collision}
There is an interval \(W_n'\subset I\) of length \(c_{I,J}R_n\), containing or adjacent to
\(s_n^{\rm out}\), such that
\begin{equation}
\label{eq:v081-hairbrush-coarse-product-collision}
\boxed{
\int_{W_n'}
\iint_{H_n^{\rm hair}\times H_n^{\rm hair}}
\mathbf 1_{\{|\Pi_u(z(a))-\Pi_u(z(a'))|\le C_{I,J,U}R_n\}}
\,d\sigma(a)d\sigma(a')du
\gtrsim_{I,J,U}
R_nS_n^2.
}
\end{equation}
Hence the ideal-coupling thickening issue does not occur at the coarse hairbrush scale
\(R_n\); it arises only when one asks for a collision-preserving transfer to a scale
\(o(R_n)\).
\end{corollary}

\begin{proof}
By \eqref{eq:v081-hairbrush-merged-bush}, every point
\(\Pi_{s_n^{\rm out}}(z(a))\), \(a\in H_n^{\rm hair}\), lies in one
\(O(R_n)\)-ball.  Since \(U\) is bounded, the same holds with a fixed enlargement for
\(|u-s_n^{\rm out}|\le c_{I,J}R_n\).  The collision indicator is therefore one on the full product
for every \(u\in W_n'\).  Integrate its product mass \(S_n^2\) over \(W_n'\).
\end{proof}

\begin{corollary}[Family Frostman failure reduces to one physical-bush Frostman problem]
\label{cor:v081-family-to-merged-bush-frostman}
Let
\[
\eta_n^{\rm hair}
:=
S_n^{-1}\sigma|_{H_n^{\rm hair}}.
\]
For each fixed \(0<\delta<3\), exactly one of the following occurs after passing to a subsequence:
\begin{enumerate}[label=\textnormal{(\roman*)}]
\item the measures \(\eta_n^{\rm hair}\) satisfy the uniform direction-Frostman hypothesis of
\cref{thm:v081-bush-frostman-escape} above the bush scale \(C_{I,J}R_n\), and therefore
\(\dim_HK_b\ge4-\delta\);
\item there are \(\widetilde A_n\to\infty\), radii
\(\widetilde R_n\ge C_{I,J}R_n\) with \(\widetilde R_n\downarrow0\), and direction caps such that
\begin{equation}
\label{eq:v081-merged-bush-angular-concentration}
\boxed{
\eta_n^{\rm hair}\!\left(B(\widetilde a_n,\widetilde R_n)\right)
>
\widetilde A_n\widetilde R_n^{3-\delta}.
}
\end{equation}
The restricted direction set in \textnormal{(ii)} remains one
\(C_{I,J}R_n\)-bush centered at \((y_n,s_n^{\rm out})\) and retains all original bush labels and
old conditional Maslov flags.
\end{enumerate}
\end{corollary}

\begin{proof}
Apply \cref{thm:v081-bush-frostman-escape,%
cor:v081-bush-frostman-or-angular-concentration} to the merged bushes from
\cref{prop:v081-hairbrush-merges-to-one-bush}.  Restricting directions does not alter the pointwise
bush inequality or the conditional flag laws.
\end{proof}

\begin{corollary}[Merging removes the artificial component-selection factor \(M^{-4}\)]
\label{cor:v081-merged-bush-cap-debt}
In alternative \textnormal{(ii)} of
\cref{cor:v081-family-to-merged-bush-frostman}, put
\[
\widetilde p_n
:=
\sigma\!\left(
H_n^{\rm hair}\cap B(\widetilde a_n,\widetilde R_n)
\right).
\]
Then
\begin{equation}
\label{eq:v081-merged-bush-cap-mass}
\boxed{
\widetilde p_n
>
\widetilde A_nS_n\widetilde R_n^{3-\delta}.
}
\end{equation}
After translating and dilating this cap by \(\widetilde R_n^{-1}\), its normalized direction
measure has density
\begin{equation}
\label{eq:v081-merged-bush-cap-density}
\boxed{
\widetilde D_{n,\rm cap}
\lesssim_U
\frac{\widetilde R_n^3}{\widetilde p_n}
\lesssim_U
\frac{\widetilde R_n^\delta}{\widetilde A_nS_n}.
}
\end{equation}
Necessarily
\begin{equation}
\label{eq:v081-merged-bush-cap-scale}
\boxed{
\widetilde R_n^\delta
\gtrsim_U
\widetilde A_nS_n.
}
\end{equation}
Thus the \(M_n^{-4}\) factor in
\eqref{eq:v081-selected-family-cap-density} is only the cost of selecting one original component.
After the Reeb-coherent caps are merged, the correct remaining normalization debt is
\(S_n^{-1}\), with no bush-cardinality factor.
\end{corollary}

\begin{proof}
Equation \eqref{eq:v081-merged-bush-cap-mass} is
\eqref{eq:v081-merged-bush-angular-concentration} multiplied by \(S_n\).  The affine rescaling
calculation of \cref{lem:v081-bush-cap-normalization-ledger} gives
\eqref{eq:v081-merged-bush-cap-density}.  Since a probability measure supported in a fixed bounded
three-dimensional ball has density norm bounded below by a positive dimensional constant,
\eqref{eq:v081-merged-bush-cap-scale} follows from
\eqref{eq:v081-merged-bush-cap-density}.
\end{proof}

\begin{lemma}[Inside one Reeb hairbrush, clustered homothety centers force a rank-one direction locus]
\label{lem:v081-hairbrush-cluster-rank-one}
Keep the notation of \cref{lem:v081-hairbrush-maslov-blowup}.  If
\(s_{n,j}\ne s_{n,k}\), denote the distinguished homothety center of the pair by
\[
a_{jk}^{\rm hom}
:=
\frac{c_{n,j}-c_{n,k}}{s_{n,j}-s_{n,k}}.
\]
Then
\begin{equation}
\label{eq:v081-hairbrush-hom-center-formula}
\boxed{
a_{jk}^{\rm hom}
=
\frac{d_ka_k-d_ja_j}{d_k-d_j},
\qquad
d_j(a_j-a^\sharp)-d_k(a_k-a^\sharp)
=
-(d_k-d_j)(a_{jk}^{\rm hom}-a^\sharp).
}
\end{equation}
Let \(0<R,\eta\le1\), and let \(G=(V,E)\) be a connected graph of bush labels with at least one
edge such that
\begin{equation}
\label{eq:v081-hairbrush-cluster-hypotheses}
|s_{n,j}-s_{n,k}|\ge\delta_s,
\qquad
|a_{jk}^{\rm hom}-a^\sharp|\le\eta
\qquad(jk\in E).
\end{equation}
If the graph diameter is at most \(L\), there is a vector \(v_G\in\mathbb R^3\) such that
\begin{equation}
\label{eq:v081-hairbrush-rank-one-curve}
\boxed{
\dist\!\left(
a_j,
\left\{a^\sharp+\frac{v_G}{d}:d\in s_0-J\right\}
\right)
\lesssim_{I,J}L\eta
\qquad(j\in V).
}
\end{equation}
In particular, if all caps \(B(a_j,R)\), \(j\in V\), meet the bounded direction chart \(U\), then
\begin{equation}
\label{eq:v081-hairbrush-cluster-cap-mass}
\boxed{
\sum_{j\in V}
\sigma\!\left(H_{n,j}\cap B(a_j,R)\right)
\lesssim_{I,J,U,\delta_s}
(R+L\eta)^2.
}
\end{equation}
Thus homothety-center clustering inside a common Reeb hairbrush is already a quantitative
direction-rank-loss branch; it is not an additional unstructured fixed-angle remainder.
\end{lemma}

\begin{proof}
Since \(c_{n,j}=y_0-d_ja_j\) and \(s_{n,j}=s_0-d_j\), direct substitution gives the first identity
in \eqref{eq:v081-hairbrush-hom-center-formula}; rearranging gives the second.  Hence the quantities
\[
v_j:=d_j(a_j-a^\sharp)
\]
differ by \(O_{I,J}(\eta)\) across every edge.  Fix a root \(j_0\) and put \(v_G=v_{j_0}\).
Summation along a path of length at most \(L\), followed by division by
\(|d_j|\gtrsim_{I,J}1\), proves
\eqref{eq:v081-hairbrush-rank-one-curve}.

The parameter set \(s_0-J\) is a compact interval separated from zero.  Under the cluster
hypotheses above, the fact that the caps meet \(U\) and that one edge
has \(|d_j-d_k|\ge\delta_s\) bounds \(a^\sharp\) and \(v_G\) in terms of
\(I,J,U,\delta_s\).  Therefore the curve
\(d\mapsto a^\sharp+v_G/d\) has uniformly bounded length.  Its
\((R+L\eta)\)-neighborhood in \(\mathbb R^3\) has volume
\(O_{I,J,U,\delta_s}((R+L\eta)^2)\).  The sets \(H_{n,j}\) are pairwise disjoint, so their cap
masses sum to at most the normalized Lebesgue measure of that tube.  This proves the claimed
cluster cap-mass bound.
\end{proof}

\begin{corollary}[A heavy Reeb-hairbrush star cannot have one clustered homothety center]
\label{cor:v081-hairbrush-heavy-star-center-separation}
Fix a root label \(j_0\) and a set of leaves \(\mathcal N\) satisfying
\[
|s_{n,j_0}-s_{n,k}|\ge\delta_s
\qquad(k\in\mathcal N).
\]
If all \(a_{j_0k}^{\rm hom}\), \(k\in\mathcal N\), lie in one ball of radius \(\eta\), then
\begin{equation}
\label{eq:v081-hairbrush-clustered-star-mass}
\boxed{
\sum_{k\in\mathcal N\cup\{j_0\}}
\sigma\!\left(H_{n,k}\cap B(a_k,R)
\right)
\lesssim_{I,J,U,\delta_s}
(R+\eta)^2.
}
\end{equation}
Consequently every star whose retained cap mass is larger than the right side contains two leaves
whose distinguished homothety centers are \(\eta\)-separated.  This is the geometric
separated-center input for the next old flag-polarization comparison; no transverse flag conclusion
is asserted before that comparison is made.
\end{corollary}

\begin{proof}
Apply \cref{lem:v081-hairbrush-cluster-rank-one} to the star graph, whose diameter is at most two.
The final assertion is the contrapositive.
\end{proof}

\begin{proposition}[A coherent bad cap family forces exact local re-collision energy]
\label{prop:v081-hairbrush-local-recollision-energy}
Keep the cap notation \(A_{n,j},S_n\) from
\cref{prop:v081-hairbrush-merges-to-one-bush}.
For all large \(n\), choose an interval \(W_n\subset I_{\rm out}\) of length comparable to
\(R_n\), having \(s_n^{\rm out}\) as an endpoint or midpoint, and define the localized ideal
outer-wing pieces
\[
\mathcal E_{n,j}
:=
\left\{
\bigl(c_{n,j}+(u-s_{n,j})a,u\bigr):
a\in A_{n,j},\ u\in W_n
\right\}.
\]
Then all \(\mathcal E_{n,j}\) lie in one \(O(R_n)\)-ball around
\((y_n,s_n^{\rm out})\),
\[
|\mathcal E_{n,j}|\asymp_{I,J}|U|R_n\sigma(A_{n,j}),
\]
and
\begin{equation}
\label{eq:v081-hairbrush-local-energy}
\boxed{
\sum_{j\ne k}
\int_{W_n}
|u-s_{n,j}|^3
\sigma\!\left(
A_{n,j}\cap\Psi_{jk,u}^{-1}(A_{n,k})
\right)du
\gtrsim_{I,J}
\frac{S_n^2}{R_n^2}-R_nS_n.
}
\end{equation}
In particular, the bad-cap inequality gives the explicit lower input
\begin{equation}
\label{eq:v081-hairbrush-local-energy-input}
S_n>A_nq_nR_n^{3-\delta}
\end{equation}
to the exact the preceding stage re-collision ledger.  The first term on the right of
\eqref{eq:v081-hairbrush-local-energy} dominates the diagonal precisely when
\(S_n\gg R_n^3\), equivalently when the coherent hairbrush has genuine angular density above
ambient cap volume.
\end{proposition}

\begin{proof}
For \(a\in A_{n,j}\), the defining identity
\(c_{n,j}+(s_n^{\rm out}-s_{n,j})a_{n,j}=y_n\) gives
\[
\left|
c_{n,j}+(u-s_{n,j})a-y_n
\right|
\lesssim_{I,J,U}R_n
\qquad
(u\in W_n).
\]
This proves the common-ball containment.  The map
\[
(a,u)\longmapsto\bigl(c_{n,j}+(u-s_{n,j})a,u\bigr)
\]
is injective and has Jacobian \(|u-s_{n,j}|^3\asymp_{I,J}1\) on
\(A_{n,j}\times W_n\), proving the asserted volume formula.  Cauchy--Schwarz inside the common
four-dimensional \(O(R_n)\)-ball yields
\[
\sum_{j,k}|\mathcal E_{n,j}\cap\mathcal E_{n,k}|
\gtrsim_{I,J}
\frac{(R_nS_n)^2}{R_n^4}
=
\frac{S_n^2}{R_n^2}.
\]
The diagonal is
\(\sum_j|\mathcal E_{n,j}|\lesssim_{I,J}|U|R_nS_n\).  For \(j\ne k\), the localized version
of the exact pair-overlap calculation in
\cref{thm:v014-pair-overlap} gives
\[
|\mathcal E_{n,j}\cap\mathcal E_{n,k}|
=
|U|\int_{W_n}|u-s_{n,j}|^3
\sigma\!\left(A_{n,j}\cap\Psi_{jk,u}^{-1}(A_{n,k})\right)du.
\]
Subtract the diagonal and divide by \(|U|\).
\end{proof}

\begin{proposition}[Dense coherent hairbrush energy forces high Reeb contact-label reuse]
\label{prop:v081-hairbrush-high-label-reuse}
For \(a\in A_{n,j}\) and \(u\in W_n\), define the exact target-label multiplicity
\begin{equation}
\label{eq:v081-hairbrush-label-multiplicity}
\mathfrak m_{n,j}(a,u)
:=
\sum_{k\ne j}
\mathbf 1_{A_{n,k}}\!\left(\Psi_{jk,u}(a)\right).
\end{equation}
Every term counted by \(\mathfrak m_{n,j}(a,u)\) is a distinct bush label whose ideal Reeb front
passes through the same point
\[
\bigl(c_{n,j}+(u-s_{n,j})a,u\bigr).
\]
Moreover the off-diagonal energy in
\eqref{eq:v081-hairbrush-local-energy} has the exact representation
\begin{equation}
\label{eq:v081-hairbrush-energy-as-reuse}
\boxed{
\mathcal R_n
:=
\sum_{j\ne k}\int_{W_n}|u-s_{n,j}|^3
\sigma\!\left(A_{n,j}\cap\Psi_{jk,u}^{-1}(A_{n,k})\right)du
\asymp_{I,J}
\sum_j\int_{W_n}\int_{A_{n,j}}
\mathfrak m_{n,j}(a,u)\,d\sigma(a)du.
}
\end{equation}
For every \(B\ge1\), the contribution of
\(\{\mathfrak m_{n,j}<B\}\) to the right side is at most
\begin{equation}
\label{eq:v081-hairbrush-low-reuse-bound}
\boxed{BR_nS_n.}
\end{equation}
Consequently, whenever \(S_n/R_n^3\to\infty\), a fixed portion of \(\mathcal R_n\) is carried by
\begin{equation}
\label{eq:v081-hairbrush-high-reuse-threshold}
\boxed{
\mathfrak m_{n,j}(a,u)
\gtrsim_{I,J}
\frac{S_n}{R_n^3}
\gtrsim
A_nq_nR_n^{-\delta}.
}
\end{equation}
Thus the high-density half of
\cref{cor:v081-hairbrush-density-recollision-dichotomy} is an actual common Reeb contact-star
reuse event, with all bush labels still present.
\end{proposition}

\begin{proof}
For fixed \(j,u\), Tonelli gives
\[
\int_{A_{n,j}}\mathfrak m_{n,j}(a,u)\,d\sigma(a)
=
\sum_{k\ne j}
\sigma\!\left(A_{n,j}\cap\Psi_{jk,u}^{-1}(A_{n,k})\right).
\]
On the fixed outer wing, \(|u-s_{n,j}|^3\asymp_{I,J}1\), proving
\eqref{eq:v081-hairbrush-energy-as-reuse}.  On the set where
\(\mathfrak m_{n,j}<B\), delete all other restrictions and use
\(|W_n|\asymp R_n\) and \(\sum_j\sigma(A_{n,j})=S_n\); this gives
\eqref{eq:v081-hairbrush-low-reuse-bound}.

If \(S_n/R_n^3\to\infty\),
\eqref{eq:v081-hairbrush-local-energy} gives
\(\mathcal R_n\gtrsim S_n^2/R_n^2\).  Choose a sufficiently small fixed
\(c_{I,J}>0\) and put \(B=c_{I,J}S_n/R_n^3\).  Then
\eqref{eq:v081-hairbrush-low-reuse-bound} is at most half of \(\mathcal R_n\), so the complementary
high-reuse set carries a fixed portion.  Finally
\eqref{eq:v081-hairbrush-local-energy-input} gives the final comparison in
\eqref{eq:v081-hairbrush-high-reuse-threshold} after decreasing the implicit constant.
\end{proof}

\begin{lemma}[Ideal re-collision energy is the fine product-collision limit]
\label{lem:v081-hairbrush-ideal-product-limit}
For \(u\in W_n\), write
\[
T_{n,j,u}(a):=c_{n,j}+(u-s_{n,j})a
\]
and define
\[
\mathcal C^{\rm id}_{n,\rho}
:=
\sum_{j\ne k}\int_{W_n}
\iint_{A_{n,j}\times A_{n,k}}
\mathbf 1_{\{|T_{n,j,u}(a)-T_{n,k,u}(a')|\le\rho\}}
\,d\sigma(a)d\sigma(a')du.
\]
For every fixed \(n\),
\begin{equation}
\label{eq:v081-hairbrush-ideal-product-limit}
\boxed{
\lim_{\rho\downarrow0}
\rho^{-3}\mathcal C^{\rm id}_{n,\rho}
\asymp_{I,J,U}
\mathcal R_n.
}
\end{equation}
Moreover
\begin{equation}
\label{eq:v081-hairbrush-actual-ideal-inclusion}
\boxed{
|T_{n,j,u}(a)-T_{n,k,u}(a')|\le\rho
\Longrightarrow
|\Pi_u(z(a))-\Pi_u(z(a'))|
\le\rho+2\varepsilon_n.
}
\end{equation}
Consequently, if a scale \(\rho_n=o(R_n)\) can be chosen in the differentiation regime of
\eqref{eq:v081-hairbrush-ideal-product-limit} with \(\varepsilon_n=o(\rho_n)\), then
\eqref{eq:v081-hairbrush-coupling-to-product-target} holds with
\(\widetilde r_n=2\rho_n\).  The only missing issue in that transfer is uniform separation between
the bush-error scale and the direction-set differentiation scale.
\end{lemma}

\begin{proof}
Let \(f_{n,j,u}\) be the density of
\((T_{n,j,u})_\#(\sigma|_{A_{n,j}})\) with respect to three-dimensional Lebesgue measure.  On the
outer wing,
\[
f_{n,j,u}(x)
=
\frac{1}{|U|\,|u-s_{n,j}|^3}
\mathbf 1_{T_{n,j,u}(A_{n,j})}(x).
\]
If
\[
\phi_\rho(x):=\frac{\mathbf 1_{B(0,\rho)}(x)}{|B(0,\rho)|},
\]
then the normalized \((j,k,u)\)-summand of
\(\rho^{-3}\mathcal C^{\rm id}_{n,\rho}\) is, up to the dimensional ball constant,
\[
\int f_{n,j,u}(x)(f_{n,k,u}*\phi_\rho)(x)\,dx.
\]
The approximate-identity theorem gives convergence in \(L^2\) to
\(\int f_{n,j,u}f_{n,k,u}\).  The \(L^2\) norms are bounded on the outer wing by constants times
\(\sigma(A_{n,j})^{1/2}\) and \(\sigma(A_{n,k})^{1/2}\), so dominated convergence applies in
\(u\); the label family is finite for fixed \(n\).  Finally the exact change of variables from
\cref{thm:v014-pair-overlap} identifies
\[
\int f_{n,j,u}f_{n,k,u}
\asymp_{I,J,U}
\sigma\!\left(A_{n,j}\cap\Psi_{jk,u}^{-1}(A_{n,k})\right).
\]
Summation proves \eqref{eq:v081-hairbrush-ideal-product-limit}.  Each actual bush point differs
from its ideal point by at most \(\varepsilon_n\), so the triangle inequality proves
\eqref{eq:v081-hairbrush-actual-ideal-inclusion} and the last assertion.
\end{proof}

\begin{proposition}[Finite-scale cell transfer: product collision, contact-star reuse, or a density budget]
\label{prop:v081-hairbrush-finite-cell-transfer}
Fix \(n\), a scale
\[
\varepsilon_n\le r\le R_n,
\]
and a half-open direction grid \(\mathcal Q_r\) of side \(r\).  Put
\[
w_{n,k,Q}:=\sigma(A_{n,k}\cap Q),
\]
and call \((k,Q)\) \(\omega\)-heavy when
\[
w_{n,k,Q}\ge\omega r^3,
\qquad 0<\omega\le C_U.
\]
Define the reverse exact label multiplicity
\begin{equation}
\label{eq:v081-hairbrush-reverse-label-multiplicity}
\mathfrak m^{\rm rev}_{n,k}(a',u)
:=
\sum_{j\ne k}
\mathbf 1_{A_{n,j}}\!\left(\Psi_{kj,u}(a')\right).
\end{equation}
For every \(B\ge1\), at least one of the following holds:
\begin{enumerate}[label=\textnormal{(\roman*)}]
\item a fixed portion of \(\mathcal R_n\) lands in \(\omega\)-heavy target cells, and the actual
finite-thickness product collision obeys
\begin{equation}
\label{eq:v081-hairbrush-heavy-cell-product-transfer}
\boxed{
\begin{aligned}
&\sum_{j\ne k}\int_{W_n}
\iint_{A_{n,j}\times A_{n,k}}
\mathbf 1_{\{|\Pi_u(z(a))-\Pi_u(z(a'))|\le C_{I,J,U}r\}}
\,d\sigma(a)d\sigma(a')du\\
&\hspace{30mm}\gtrsim_{I,J,U}
\omega r^3\mathcal R_n;
\end{aligned}
}
\end{equation}
\item a fixed portion of \(\mathcal R_n\) lands where
\begin{equation}
\label{eq:v081-hairbrush-reverse-high-reuse}
\boxed{\mathfrak m^{\rm rev}_{n,k}(a',u)\ge B,}
\end{equation}
which is a common ideal Reeb contact star rooted at the target endpoint;
\item the complete ideal energy satisfies
\begin{equation}
\label{eq:v081-hairbrush-light-cell-energy-budget}
\boxed{
\mathcal R_n
\lesssim_{I,J,U}
BN_n\omega R_n^4.
}
\end{equation}
\end{enumerate}
No limiting differentiation scale is used in this trichotomy.
\end{proposition}

\begin{proof}
For a target cell \(Q\), set
\[
E_{n,jk,Q}(u)
:=
\left\{
a\in A_{n,j}:
\Psi_{jk,u}(a)\in A_{n,k}\cap Q
\right\}.
\]
If \(a\in E_{n,jk,Q}(u)\) and \(a'\in A_{n,k}\cap Q\), then the ideal source point equals the
ideal target point generated by \(\Psi_{jk,u}(a)\), while
\[
|T_{n,k,u}(a')-T_{n,k,u}(\Psi_{jk,u}(a))|
\lesssim_{I,J,U}r.
\]
The two actual points introduce at most \(2\varepsilon_n\le2r\) additional error.  Hence the
left side of \eqref{eq:v081-hairbrush-heavy-cell-product-transfer} is bounded below by
\begin{equation}
\label{eq:v081-hairbrush-cell-weighted-coupling}
\sum_{j\ne k}\sum_{Q\in\mathcal Q_r}
\int_{W_n}
w_{n,k,Q}\,
\sigma(E_{n,jk,Q}(u))\,du.
\end{equation}
On the heavy cells, \(w_{n,k,Q}\ge\omega r^3\).  Since the outer-wing Jacobians are bounded above
and below, a fixed portion of \(\mathcal R_n\) in those cells proves
\eqref{eq:v081-hairbrush-heavy-cell-product-transfer}.

It remains to estimate the part landing in light target cells and satisfying
\(\mathfrak m^{\rm rev}_{n,k}<B\).  Change variables
\(a'=\Psi_{jk,u}(a)\).  All Jacobians are uniformly comparable on the outer wing, so this energy is
bounded by a constant times
\[
\sum_k\int_{W_n}
\int_{\substack{a'\in A_{n,k}\cap Q\\
Q\ {\rm light}}}
\mathfrak m^{\rm rev}_{n,k}(a',u)\,d\sigma(a')du.
\]
Each \(A_{n,k}\) lies in a direction ball of radius \(R_n\), which meets
\(O_U((R_n/r)^3)\) grid cells.  Therefore its total mass in light cells is at most
\[
C_U(R_n/r)^3\omega r^3
\lesssim_U\omega R_n^3.
\]
Using \(\mathfrak m^{\rm rev}_{n,k}<B\), summing the \(N_n\) target labels, and integrating over
\(|W_n|\asymp R_n\) gives \(C_{I,J,U}BN_n\omega R_n^4\).  If neither of the first two alternatives
carries a fixed portion of \(\mathcal R_n\), this proves
\eqref{eq:v081-hairbrush-light-cell-energy-budget}.
\end{proof}

\subsection{Fine Maslov returns and stopping leaves}

\begin{corollary}[The balanced finite-cell choice removes the unknown differentiation threshold]
\label{cor:v081-hairbrush-balanced-cell-transfer}
Assume \(S_n/R_n^3\) is larger than a sufficiently large constant and set
\[
B_n:=c_{I,J}\frac{S_n}{R_n^3},
\qquad
\omega_n:=c_{I,J,U}\frac{S_n}{N_nR_n^3},
\]
with the constants chosen sufficiently small.  Then alternative \textnormal{(iii)} of
\cref{prop:v081-hairbrush-finite-cell-transfer} is impossible.  Consequently, at every scale
\(\varepsilon_n\le r\le R_n\), either
\begin{equation}
\label{eq:v081-hairbrush-balanced-product-transfer}
\boxed{
\mathcal C^{\rm phys}_{n,r}
\gtrsim_{I,J,U}
\frac{S_n}{N_nR_n^3}\,r^3\mathcal R_n
\gtrsim_{I,J,U}
\frac{S_n^3}{N_nR_n^5}\,r^3,
}
\end{equation}
where \(\mathcal C^{\rm phys}_{n,r}\) denotes the physical product-collision expression in
\eqref{eq:v081-hairbrush-heavy-cell-product-transfer}, or a fixed portion of
\(\mathcal R_n\) is carried by reverse contact stars of label degree
\begin{equation}
\label{eq:v081-hairbrush-balanced-reuse}
\boxed{
\mathfrak m^{\rm rev}_{n,k}
\gtrsim_{I,J}
\frac{S_n}{R_n^3}.
}
\end{equation}
Moreover the Frostman-failure lower bound and maximal-family estimates give
\begin{equation}
\label{eq:v081-hairbrush-average-cell-density}
\boxed{
\omega_n
\gtrsim_{I,J,U,G_*}
A_nM_n^4R_n^{-\delta}.
}
\end{equation}
Thus the finite-scale transfer loses exactly the average relative density of one original labeled
cap; if that density is too small, the complementary output is high Reeb contact-star reuse.
\end{corollary}

\begin{proof}
By \eqref{eq:v081-hairbrush-local-energy} and \(S_n/R_n^3\gg1\),
\[
\mathcal R_n\gtrsim_{I,J}\frac{S_n^2}{R_n^2}.
\]
With the displayed choices,
\[
B_nN_n\omega_nR_n^4
=
c_{I,J,U}^2\frac{S_n^2}{R_n^2}.
\]
Choose the product of the constants below the implicit constant in the preceding energy lower
bound.  Then \eqref{eq:v081-hairbrush-light-cell-energy-budget} cannot hold.  The two remaining
alternatives give \eqref{eq:v081-hairbrush-balanced-product-transfer} and
\eqref{eq:v081-hairbrush-balanced-reuse}.  Finally
\[
\frac{S_n}{N_nR_n^3}
>
A_n\frac{q_n}{N_n}R_n^{-\delta}
\gtrsim_{G_*}
A_nM_n^4R_n^{-\delta}
\]
by \eqref{eq:v081-hairbrush-local-energy-input},
\(q_n\gtrsim_{G_*}M_n\), and \(N_n\lesssim M_n^{-3}\).
\end{proof}

\begin{proposition}[The fixed-angle physical branch is a normalized fine Maslov graph]
\label{prop:v081-hairbrush-fine-maslov-graph}
Let
\[
d\widehat\mu_n^{\rm hair}(z(a))
:=
S_n^{-1}\mathbf 1_{H_n^{\rm hair}}(a)\,d\sigma(a),
\]
and fix \(\tau_0>0\).  Suppose that, for some
\(\varepsilon_n\le r\le R_n\), a fixed portion of the physical product collision in
\eqref{eq:v081-hairbrush-heavy-cell-product-transfer} is carried by pairs satisfying
\(|a-a'|\ge\tau_0\).  Define
\begin{equation}
\label{eq:v081-hairbrush-fine-maslov-kernel}
h_{n,r}(z(a),z(a'))
:=
\mathbf 1_{\{|a-a'|\ge\tau_0\}}
\frac{1}{C_{\tau_0}r}
\int_{W_n}
\mathbf 1_{\{|\Pi_u(z(a))-\Pi_u(z(a'))|\le C r\}}\,du,
\end{equation}
where \(C_{\tau_0}\) is chosen sufficiently large.  Then
\(0\le h_{n,r}\le1\), every edge of this graph is an \(O(r)\)-thickened Maslov
incidence, and its normalized edge mass satisfies
\begin{equation}
\label{eq:v081-hairbrush-fine-maslov-edge-mass}
\boxed{
\widetilde M_{n,r}
:=
\iint h_{n,r}(z,z')\,
d\widehat\mu_n^{\rm hair}(z)d\widehat\mu_n^{\rm hair}(z')
\gtrsim_{\tau_0,I,J,U}
\omega r^2\frac{\mathcal R_n}{S_n^2}.
}
\end{equation}
In the dense regime \(S_n/R_n^3\gg1\), this gives
\begin{equation}
\label{eq:v081-hairbrush-fine-maslov-edge-scale}
\widetilde M_{n,r}
\gtrsim_{\tau_0,I,J,U}
\omega\Bigl(\frac r{R_n}\Bigr)^2.
\end{equation}
For the balanced choice of
\cref{cor:v081-hairbrush-balanced-cell-transfer},
\begin{equation}
\label{eq:v081-hairbrush-balanced-fine-maslov-edge}
\boxed{
\widetilde M_{n,r}
\gtrsim_{\tau_0,I,J,U,G_*}
A_nM_n^4R_n^{-\delta}
\Bigl(\frac r{R_n}\Bigr)^2.
}
\end{equation}
The conditional old Maslov flags on the original occurrences lift to this graph without
additional mass loss.  If the required fixed-angle portion is absent, the discarded collision
mass is instead concentrated in the \(\tau_0\)-Darboux branch.
\end{proposition}

\begin{proof}
For fixed \(a,a'\), the difference
\[
\Pi_u(z(a))-\Pi_u(z(a'))
\]
is affine in \(u\), with \(u\)-slope \(a-a'\).  Hence, on
\(|a-a'|\ge\tau_0\), the set of \(u\in W_n\) for which its norm is at most
\(Cr\) has length \(O_{\tau_0}(r)\).  Enlarging \(C_{\tau_0}\) therefore gives
\(h_{n,r}\le1\).

The collision relation itself is
\[
\bigl(b(a)-b(a')\bigr)+u(a-a')=O(r),
\]
so each edge lies in the \(O(r)\)-neighborhood of the Maslov plane \(L_u\).
After restricting to the assumed fixed-angle portion, dividing the physical product-collision
lower bound in \cref{prop:v081-hairbrush-finite-cell-transfer} by the endpoint normalization
\(S_n^2\) and the time-window normalization \(C_{\tau_0}r\) gives
\[
\widetilde M_{n,r}
\gtrsim_{\tau_0,I,J,U}
\frac{\omega r^3\mathcal R_n}{rS_n^2},
\]
which is the asserted fine Maslov edge-mass bound.  The local-energy bound
\(\mathcal R_n\gtrsim S_n^2/R_n^2\) proves its scale form.  The average-cell-density estimate then
proves the balanced fine-edge form.
Finally, the finite-cell product consists of actual selector endpoints
\(z(a),z(a')\); the second endpoint need not equal the exact transported endpoint
\(z(\Psi_{jk,u}(a))\).  Nevertheless each endpoint carries its own conditional probability law
\(\kappa_{z(a)}\) or \(\kappa_{z(a')}\).  Integrating any fixed finite tensor product of those laws
gives \(1\), by \cref{lem:v081-lossless-conditional-flag-tensor}.  Preservation on the hairbrush
energy is \cref{cor:v081-hairbrush-energy-flag-lift}.
Therefore decorating the physical edge measure by the old flags costs no additional mass factor.
\end{proof}

\begin{corollary}[At the bush scale the fixed-angle Maslov return has constant edge mass]
\label{cor:v081-hairbrush-coarse-maslov-graph}
Fix \(\tau_0>0\).  If a fixed portion of the coarse physical collision from
\cref{cor:v081-hairbrush-coarse-product-collision} is carried by
\(|a-a'|\ge\tau_0\), then
\begin{equation}
\label{eq:v081-hairbrush-coarse-maslov-edge}
\boxed{
\iint h^{\rm coarse}_{n,R_n}(z,z')\,
d\widehat\mu_n^{\rm hair}(z)d\widehat\mu_n^{\rm hair}(z')
\gtrsim_{\tau_0,I,J,U}1,
}
\end{equation}
where
\[
h^{\rm coarse}_{n,R_n}(z(a),z(a'))
:=
\mathbf 1_{\{|a-a'|\ge\tau_0\}}
\frac1{C_{\tau_0}R_n}
\int_{W_n'}
\mathbf 1_{\{|\Pi_u(z(a))-\Pi_u(z(a'))|\le C_{I,J,U}R_n\}}\,du
\]
and \(0\le h^{\rm coarse}_{n,R_n}\le1\).  Hence the density factor
\(\omega_n\) is a strict fine-scale transfer cost; it is absent at the intrinsic bush scale.
If the fixed-angle portion is absent, a fixed portion of the coarse collision lies in the
\(\tau_0\)-Darboux branch.
\end{corollary}

\begin{proof}
On \(|a-a'|\ge\tau_0\), the collision-time window has length
\(O_{\tau_0}(R_n)\), so the kernel is bounded by \(1\).  Divide the fixed-angle part of
\eqref{eq:v081-hairbrush-coarse-product-collision} by
\(R_nS_n^2\) and use
\(\sigma(H_n^{\rm hair})=S_n\).
\end{proof}

\begin{corollary}[The exact two-scale scalar frontier after the Maslov return]
\label{cor:v081-hairbrush-two-scale-maslov-frontier}
Assume the fixed-angle branch of
\cref{prop:v081-hairbrush-fine-maslov-graph}, use the balanced cell density
\(\omega_n\), and suppose that the resulting merged bush has a further bad direction cap of radius
\(\widetilde R_n\), excess \(\widetilde A_n\), and normalized density constant
\(\widetilde D_{n,\rm cap}\) from
\cref{cor:v081-merged-bush-cap-debt}.  Then the quotient of the next cap-density tax by the
available normalized Maslov edge mass obeys
\begin{equation}
\label{eq:v081-hairbrush-two-scale-maslov-frontier}
\boxed{
\frac{\widetilde D_{n,\rm cap}}{\widetilde M_{n,r}}
\lesssim_{\tau_0,I,J,U,G_*}
\frac{\widetilde R_n^\delta R_n^\delta}
{\widetilde A_nS_nA_nM_n^4}
\Bigl(\frac{R_n}{r}\Bigr)^2.
}
\end{equation}
At the intrinsic bush scale, if the fixed-angle alternative of
\cref{cor:v081-hairbrush-coarse-maslov-graph} holds, the sharper quotient is
\begin{equation}
\label{eq:v081-hairbrush-coarse-two-scale-frontier}
\boxed{
\frac{\widetilde D_{n,\rm cap}}
{\widetilde M^{\rm coarse}_{n,R_n}}
\lesssim_{\tau_0,I,J,U}
\frac{\widetilde R_n^\delta}{\widetilde A_nS_n}.
}
\end{equation}
Thus the finite-cell and cap-normalization steps introduce no hidden factor: their complete
two-scale loss is the right side of
\eqref{eq:v081-hairbrush-two-scale-maslov-frontier}.  The already proved estimates make this
loss subpower whenever that displayed right side is \(r^{-o(1)}\).  Without a scale-coherent
relation among \(r,R_n,\widetilde R_n,A_n,\widetilde A_n,M_n,S_n\), however,
\eqref{eq:v081-hairbrush-two-scale-maslov-frontier} alone does not imply such absorption.
\end{corollary}

\begin{proof}
Divide \eqref{eq:v081-merged-bush-cap-density} by
\eqref{eq:v081-hairbrush-fine-maslov-edge-scale} to obtain
\[
\frac{\widetilde D_{n,\rm cap}}{\widetilde M_{n,r}}
\lesssim
\frac{\widetilde R_n^\delta}{\widetilde A_nS_n\omega_n}
\Bigl(\frac{R_n}{r}\Bigr)^2.
\]
Now use the average-cell-density estimate, equivalently
\(\omega_n^{-1}\lesssim A_n^{-1}M_n^{-4}R_n^\delta\).
For the coarse two-scale frontier, divide the merged-bush cap-density bound by the coarse Maslov
edge bound.
\end{proof}

\begin{proposition}[Intrinsic-scale Maslov renormalization preserves a fixed mass fraction]
\label{prop:v081-hairbrush-intrinsic-renormalization}
Assume the fixed-angle alternative of
\cref{cor:v081-hairbrush-coarse-maslov-graph}, and suppose the close-anchor, horizontal-line, and
transverse old/new-polarization outputs are routed through their already assigned ledgers.  Fix
\(d\in(0,1)\).  Then, after passing to a subsequence, either one of those paid outputs occurs, or
there are pairwise disjoint sets
\[
H^{(1)}_{n,1},\ldots,H^{(1)}_{n,N_n^{(1)}}
\subset H_n^{\rm hair}
\]
such that
\begin{equation}
\label{eq:v081-hairbrush-renormalized-family}
\boxed{
N_n^{(1)}\lesssim_{\tau_0,I,J,U,d}1,
\qquad
\sum_{\ell=1}^{N_n^{(1)}}\sigma(H^{(1)}_{n,\ell})
\gtrsim_{\tau_0,I,J,U,d}S_n,
}
\end{equation}
and every \(H^{(1)}_{n,\ell}\) is an
\(O(R_n^{1-d})\)-Lagrangian point-probe bush.  All conditional old-flag laws are retained.

Consequently, for every fixed \(0<\delta<3\), either this new family satisfies the vector
Frostman hypothesis of
\cref{thm:v081-bush-family-frostman-escape}, with the roles of the two separated Reeb windows
\(J\) and \(I_{\rm out}\) interchanged, along a subsequence and
\(\dim_HK_b\ge4-\delta\), or it has a further Reeb-coherent hairbrush concentration at radii
\[
R_n^{1-d}\lesssim \widehat R_n\downarrow0.
\]
The latter concentration again merges into one physical bush by
\cref{prop:v081-hairbrush-merges-to-one-bush}.  This entire intrinsic-scale return uses neither
\(\omega_n^{-1}\) nor the cap-density factor
\(\widetilde D_{n,\rm cap}\).
\end{proposition}

\begin{proof}
Let \(h_n=h^{\rm coarse}_{n,R_n}\) and use the probability measure
\(\widehat\mu_n^{\rm hair}\).  By
\eqref{eq:v081-hairbrush-coarse-maslov-edge}, its edge mass \(M_n^{\rm hair}\) is bounded below by
a fixed positive constant.  In particular
\[
M_n^{\rm hair}=R_n^{o(1)},
\]
so it lies in the exponent range \(\chi=0<2\) of
\cref{cor:v081-k-flag-arbitrary-threshold}.  Apply the graph-theoretic maximal extraction proof of
\cref{prop:v081-maximal-bush-family} with collision scale \(R_n\) and threshold
\(\Delta_n=R_n^d\).  Unless one of the stated paid ledgers occurs, it gives
\[
N_n^{(1)}\lesssim (M_n^{\rm hair})^{-3}\lesssim1,
\qquad
\sum_\ell
\widehat\mu_n^{\rm hair}(H^{(1)}_{n,\ell})
\gtrsim M_n^{\rm hair}\gtrsim1,
\]
and bush error \(O(R_n/\Delta_n)=O(R_n^{1-d})\).
Multiplying the last mass inequality by \(S_n\) proves
\eqref{eq:v081-hairbrush-renormalized-family}.  The lossless conditional tensor identity
\cref{lem:v081-lossless-conditional-flag-tensor} retains the old flags.

The new bush centers lie in the compact window containing \(W_n'\subset I_{\rm out}\).
Apply \cref{thm:v081-bush-family-frostman-escape} after relabeling that window as the center
window and choosing a compact subinterval of the original \(J\) as the next outer window.  Their
separation is bounded below by a fixed positive constant.  If the vector Frostman hypothesis
fails, the proof of \cref{cor:v081-bush-family-angular-concentration} supplies
\(\widehat R_n\ge cR_n^{1-d}\) tending to zero and a common outer-point cap family whose outer
level is back in \(J\).  \Cref{prop:v081-hairbrush-merges-to-one-bush} merges it.  No direction cap is normalized before
the coarse Maslov graph is formed, which is why neither of the two displayed density factors
appears.
\end{proof}

\begin{corollary}[One alternating return forces a two-probe angular lock]
\label{cor:v081-hairbrush-alternating-two-probe-lock}
Keep the non-Frostman alternative in
\cref{prop:v081-hairbrush-intrinsic-renormalization}.  Let
\[
(y_n^{(0)},s_n^{(0)})\in\mathbb R^3\times I_{\rm out}
\]
be the center of the incoming merged bush \(H_n^{\rm hair}\), and let
\[
(y_n^{(1)},s_n^{(1)})\in\mathbb R^3\times J
\]
be the common outer point which centers the returned merged hairbrush
\(\widehat H_n^{\rm hair}\subset H_n^{\rm hair}\).  Then
\begin{equation}
\label{eq:v081-hairbrush-two-probe-direction-lock}
\boxed{
\widehat H_n^{\rm hair}
\subset
B\!\left(
\frac{y_n^{(1)}-y_n^{(0)}}{s_n^{(1)}-s_n^{(0)}},
C_{I,J}\bigl(R_n+\widehat R_n\bigr)
\right)
\subset
B(a_n^\sharp,C_{I,J,d}\widehat R_n).
}
\end{equation}
Consequently, for every fixed \(\tau_0>0\), all pairs in
\(\widehat H_n^{\rm hair}\times\widehat H_n^{\rm hair}\) lie in the
\(\tau_0\)-angular Darboux branch for large \(n\).  In particular an indefinitely repeated
fixed-angle hairbrush-renormalization chain is impossible: after one alternating-window return,
the next geometric destination is necessarily an angular Darboux chart.

More quantitatively, if a blow-up at direction scale \(T_n\downarrow0\) has a noncollapsed
horizontal limit, then
\begin{equation}
\label{eq:v081-hairbrush-two-probe-scale-lock}
\boxed{
T_n\lesssim_{I,J}R_n+\widehat R_n
\asymp_{I,J,d}\widehat R_n.
}
\end{equation}
Thus the alternating return cannot hide a third, substantially larger noncollapsed angular
scale.  In particular there is no noncollapsed angular blow-up for which the inherited physical
error is \(o(T_n)\); the remaining cap is at the physical resolution scale and cannot be credited
with the error-free tangent Maslov relation without a further weighted coarse/fine argument.
\end{corollary}

\begin{proof}
Every \(a\in\widehat H_n^{\rm hair}\subset H_n^{\rm hair}\) obeys the inherited point-probe
relation
\[
|b(a)+s_n^{(0)}a-y_n^{(0)}|
\lesssim_{I,J}R_n.
\]
The second hairbrush merge gives
\[
|b(a)+s_n^{(1)}a-y_n^{(1)}|
\lesssim_{I,J}\widehat R_n.
\]
Subtracting and using
\(\dist(J,I_{\rm out})>0\) proves the first containment in
\eqref{eq:v081-hairbrush-two-probe-direction-lock}.  Since
\(\widehat R_n\gtrsim R_n^{1-d}\ge R_n\), the second containment follows.  Its diameter tends to
zero, which proves the fixed-angle assertion.

For the last claim, write \(a=a_n+T_n\xi\) on a putative noncollapsed blow-up.  After subtracting
the two point-probe equations and dividing by \(T_n\), the variation in \(\xi\) is
\(O((R_n+\widehat R_n)/T_n)\).  If
\((R_n+\widehat R_n)/T_n\to0\), every limiting horizontal secant vanishes, contradicting
noncollapse.  This proves \eqref{eq:v081-hairbrush-two-probe-scale-lock}.
\end{proof}

\begin{proposition}[A resolution-scale angular cap is either a good leaf or a tax-free fine return]
\label{prop:v081-resolution-cap-freeze-or-return}
Let \(H\subset G_*\) be any fixed measurable selector set of mass
\[
p:=\sigma(H)>0
\]
contained in one direction ball of radius \(0<T<1\), and retain at every point of \(H\) the
conditional old-flag laws inherited from \(G_*\).  Then exactly one of the following applies.
\begin{enumerate}[label=\textnormal{(\roman*)}]
\item the relative residual estimate
\[
\mathfrak Z_{Cr,\ge r}^{J^+}(\sigma|_H)
\lesssim
p^2r^{2-o(1)}
\]
holds uniformly for all sufficiently small \(r\), so \(H\) is a good terminal leaf in the
fixed-coarse/all-fine decomposition;
\item there is a sequence of failure scales \(r_k\downarrow0\) satisfying
\begin{equation}
\label{eq:v081-resolution-cap-diagonal-scale}
\boxed{
r_k\le\min\{T^k,p^k\},
\qquad
p=r_k^{o(1)},
\qquad
T=r_k^{o(1)},
\qquad
\frac{T^3}{p}=r_k^{-o(1)}.
}
\end{equation}
After translating the cap and dilating directions by \(T^{-1}\), every density or mass
normalization introduced by this fixed cap is therefore subpower at the new collision scale.
The normalized contact--Maslov graph, the conditional flag amplification, and the subsequent
bush/rank-loss routing may be applied on \(H\) without a fixed power loss.
\end{enumerate}
In particular the resolution-scale cap from
\cref{cor:v081-hairbrush-alternating-two-probe-lock} is not itself a terminal normalization
obstruction.  What remains after this proposition is the aggregation of the resulting nested
good-leaf/fine-return stopping tree, not an unpaid scalar density factor at one fixed node.
\end{proposition}

\begin{proof}
If \textnormal{(i)} fails, there are arbitrarily small failure scales.  Choose the \(k\)-th one
below \(\min\{T^k,p^k\}\).  Then
\[
\frac{|\log p|}{|\log r_k|}
\le\frac1k,
\qquad
\frac{|\log T|}{|\log r_k|}
\le\frac1k,
\]
which proves \eqref{eq:v081-resolution-cap-diagonal-scale}.  The rescaled normalized direction
measure has density at most \(C_UT^3/p\).  Hence this density, the restriction normalization
\(p^{-1}\), and every fixed old-flag layer constant are \(r_k^{-o(1)}\), exactly as in
\cref{lem:v078-coarse-fine-mass-freezing,cor:v078-fine-scale-absorbs-coarse-tax}.
Restriction changes which physical endpoints are retained but not their conditional flag laws, by
\cref{lem:v081-lossless-conditional-flag-tensor}.  Thus all hypotheses used by the fine
contact--Maslov return survive with only subpower constants.
\end{proof}

\begin{lemma}[Quadratic aggregation of terminal stopping leaves]
\label{lem:v081-stopping-leaf-quadratic-aggregation}
For one of the finitely many shifted physical \(r\)-grids, define
\[
\mathcal E_r^J(\lambda)
:=
\int_J\sum_{Q\in\mathcal Q_r}
\bigl((\Pi_s)_\#\lambda(Q)\bigr)^2\,ds.
\]
Let \(\lambda\) have mass \(m>0\).  Suppose that for every sufficiently small \(r\) there is a
decomposition
\[
\lambda=\sum_{\gamma\in\mathcal L(r)}\lambda_\gamma+\lambda_{\rm tail}^{(r)}
\]
into positive measures, where \(\mathcal L(r)\) is finite, such that
\begin{equation}
\label{eq:v081-stopping-leaf-input}
\boxed{
\begin{aligned}
\mathcal E_r^J(\lambda_\gamma)
&\le
C(r)\|\lambda_\gamma\|^2r^{3-o(1)}
\qquad(\gamma\in\mathcal L(r)),\\
C(r)&=r^{-o(1)},\\
\|\lambda_{\rm tail}^{(r)}\|
&\le
mr^{3/2-o(1)}.
\end{aligned}
}
\end{equation}
Then
\begin{equation}
\label{eq:v081-stopping-leaf-output}
\boxed{
\mathcal E_r^J(\lambda)
\lesssim
m^2r^{3-o(1)}.
}
\end{equation}
If \eqref{eq:v081-stopping-leaf-input} holds for the finite family of shifted grids needed to
dominate physical \(Cr\)-balls, then the full collision estimate
\eqref{eq:v078-final-coherent-full-collision}, and hence the relative residual target
\eqref{eq:v078-final-low-reuse-relative-residual}, follows.
\end{lemma}

\begin{proof}
For each \(s,Q\), write
\[
F_\lambda(s,Q):=(\Pi_s)_\#\lambda(Q).
\]
Then \(\mathcal E_r^J(\lambda)=\|F_\lambda\|_{L^2(ds\times\#Q)}^2\).
Minkowski and positivity give
\[
\begin{aligned}
\mathcal E_r^J(\lambda)^{1/2}
&\le
\sum_{\gamma\in\mathcal L(r)}
\mathcal E_r^J(\lambda_\gamma)^{1/2}
+
\mathcal E_r^J(\lambda_{\rm tail}^{(r)})^{1/2}\\
&\lesssim
C(r)^{1/2}r^{3/2-o(1)}
\sum_\gamma\|\lambda_\gamma\|
+
|J|^{1/2}\|\lambda_{\rm tail}^{(r)}\|\\
&\lesssim
mr^{3/2-o(1)}.
\end{aligned}
\]
Squaring proves \eqref{eq:v081-stopping-leaf-output}.  A fixed finite shifted-grid covering
dominates every pair with
\(|\Pi_s(z)-\Pi_s(z')|\le Cr\) by a same-cube pair at a comparable scale.  Summing those grid
energies and applying
\cref{prop:v048-energy-residual-equivalence} gives the final assertion.
\end{proof}

\begin{corollary}[The remaining stopping-tree condition is a quantitative tail estimate]
\label{cor:v081-stopping-tree-tail-criterion}
Build a nested stopping tree by applying
\cref{prop:v081-resolution-cap-freeze-or-return} at every resolution-scale cap.  Suppose that at
depth \(L\) the nonterminal nodes are pairwise disjoint and have total mass \(t_L\downarrow0\),
while the finitely many terminal leaves above depth \(L\) obey their relative residual estimates.
Let \(C_L<\infty\) dominate both their residual constants and the reciprocals of their positive
masses.  By the global bounded-direction-density estimate for \(|a-a'|<r\), each such leaf then
obeys the full shifted-grid estimate in
\eqref{eq:v081-stopping-leaf-input} with constant \(O(C_L)\).  If there is a depth choice
\(L=L(r)\to\infty\) such that
\begin{equation}
\label{eq:v081-stopping-tree-tail-target}
\boxed{
t_{L(r)}
\le
mr^{3/2-o(1)},
\qquad
C_{L(r)}
=
r^{-o(1)},
}
\end{equation}
then the Maslov-bush-to-Frostman bridge
\cref{prob:v081-maslov-bush-frostman-bridge} closes.
\end{corollary}

\begin{proof}
Take the terminal leaves above depth \(L(r)\) as
\(\mathcal L(r)\) and the sum of the nonterminal nodes at that depth as
\(\lambda_{\rm tail}^{(r)}\).  Their supports form a measurable partition of the root measure up
to null sets.  For a terminal leaf of mass \(m_\gamma\), the collision/residual comparison gives
the required fixed-angle grid bound, while its near-direction contribution is
\(O(m_\gamma r^3)\le O(C_Lm_\gamma^2r^3)\).  Apply
\cref{lem:v081-stopping-leaf-quadratic-aggregation}.
\end{proof}

\begin{corollary}[Absence of a good leaf forces arbitrarily deep alternating endpoint-star chains]
\label{cor:v081-no-good-leaf-endpoint-chain}
Assume \(\dim_HK_b<4\).  Then no positive-direction-mass stopping node can terminate in
alternative \textnormal{(i)} of
\cref{prop:v081-resolution-cap-freeze-or-return}: after normalization, its relative residual
estimate would close by
\cref{thm:v048-residual-frostman-criterion}.  Consequently, unless one of the already paid
geometric ledgers occurs, for every integer \(L\) there are nested positive-mass selector sets
\[
H^{(0)}\supset H^{(1)}\supset\cdots\supset H^{(L)}
\]
and direction radii
\[
T_0>T_1>\cdots>T_L>0,
\qquad
T_{\ell+1}\le\frac12T_\ell,
\]
such that every transition retains the conditional old flags and supplies two point-probe centers
in the alternating separated windows \(J\) and \(I_{\rm out}\).  Each returned
\(H^{(\ell+1)}\) is contained in one \(O(T_{\ell+1})\)-direction cap and carries the physical
endpoint-star witnesses of the fine Maslov graph which produced it.

Thus the geometric data required by the multiscale endpoint-star side of
the optional Maslov secant--to--curve promotion are present to arbitrary finite depth.  This statement does
\emph{not} by itself synchronize the continuous center, angular-chart, and polarization labels
from one depth to the next.  \Cref{lem:v081-power-excess-logarithmic-block} below removes
mere scalar scale sparsity by thickening every power-excess failure to a positive logarithmic
block.  The missing alternative to
\eqref{eq:v081-stopping-tree-tail-target} is therefore coherent selection of the physical
endpoint-star data across those blocks, not the existence of enough bad thickness parameters.
\end{corollary}

\begin{proof}
A positive-mass good node, normalized to a probability measure, satisfies the codimension-two
residual hypothesis and therefore gives a four-dimensional front subset.  Under the contrary
dimension assumption, every positive-mass node must have arbitrarily fine failure scales.  Apply
\cref{prop:v081-resolution-cap-freeze-or-return} and then the bush-family/hairbrush route recorded
in
\cref{prop:v081-hairbrush-intrinsic-renormalization,%
cor:v081-hairbrush-alternating-two-probe-lock}.  The returned cap radius tends to zero along the
chosen failure sequence, so, after going sufficiently far out, it is at most half the parent
radius.  Restriction preserves positive mass and the conditional flag laws.  Induction gives every
finite \(L\).  The last paragraph is the scale-density hypothesis explicitly required in
the optional Maslov secant--to--curve promotion; finite-depth selection alone does not synchronize the
continuous geometric labels.
\end{proof}

\begin{lemma}[A power-excess residual scale fills a positive logarithmic block]
\label{lem:v081-power-excess-logarithmic-block}
Let \(H\) have mass \(p>0\), fix \(0<\eta<2\), and suppose
\begin{equation}
\label{eq:v081-power-excess-one-scale}
\mathfrak Z_{r_n,\ge\tau_0}^{J^+}(\sigma|_H)
\ge
A_np^2r_n^{2-\eta},
\qquad
A_n\longrightarrow\infty.
\end{equation}
Put
\[
\alpha_\eta
:=
\frac{2-\eta}{2-\eta/2}
\in(0,1).
\]
Then, for every
\[
r_n\le R\le r_n^{\alpha_\eta},
\]
one has
\begin{equation}
\label{eq:v081-power-excess-block}
\boxed{
\mathfrak Z_{R,\ge\tau_0}^{J^+}(\sigma|_H)
\ge
A_np^2R^{2-\eta/2}.
}
\end{equation}
The interval of logarithmic scales has relative length
\begin{equation}
\label{eq:v081-power-excess-log-density}
\boxed{
\frac{\log(r_n^{\alpha_\eta}/r_n)}
{\log(1/r_n)}
=
1-\alpha_\eta
=
\frac{\eta}{4-\eta}>0.
}
\end{equation}
Hence a Hausdorff counterexample cannot force only isolated residual-failure scales.  At every
nonterminal positive-mass node it supplies positive-relative-length logarithmic blocks of
normalized Maslov graphs, with exponent still strictly below two.
\end{lemma}

\begin{proof}
The residual is monotone in its thickness:
\[
\mathfrak Z_{R,\ge\tau_0}^{J^+}(\sigma|_H)
\ge
\mathfrak Z_{r_n,\ge\tau_0}^{J^+}(\sigma|_H)
\qquad(R\ge r_n).
\]
Write \(R=r_n^\beta\), where
\(\alpha_\eta\le\beta\le1\).  The definition of \(\alpha_\eta\) gives
\[
\beta(2-\eta/2)\ge2-\eta.
\]
Since \(0<r_n<1\),
\[
r_n^{2-\eta}
\ge
r_n^{\beta(2-\eta/2)}
=
R^{2-\eta/2}.
\]
Insert this in \eqref{eq:v081-power-excess-one-scale}.  The logarithmic-length calculation is
immediate.
\end{proof}

\subsection{Coherence across logarithmic blocks}

\begin{lemma}[Finite-dimensional geometric labels synchronize with only subpower loss]
\label{lem:v081-subpower-label-synchronization}
Let \(\mathcal Y\) be the compact label space consisting of the two bounded physical point-probe
centers in \(J\) and \(I_{\rm out}\), one of the finitely many angular charts, the old/new
polarization planes in the relevant compact Grassmannians, and the angular-scale exponent
\[
\varkappa_R
:=
\frac{\log(1/\tau_R)}{\log(1/R)}
\in[0,1],
\]
where \(R\le\tau_R\le1\) is the selected direction-shell scale.  There are constants
\(D_{\mathcal Y},C_{\mathcal Y}<\infty\) such that \(\mathcal Y\) is covered by at most
\[
C_{\mathcal Y}\zeta^{-D_{\mathcal Y}}
\]
balls of radius \(\zeta\).

On the logarithmic block from
\cref{lem:v081-power-excess-logarithmic-block}, make any measurable choice
\(R\mapsto\mathfrak y_R\in\mathcal Y\) of the geometric output of the flagged Maslov routing.
Put \(L_n=\log(1/r_n)\) and choose
\[
\zeta_n:=\exp(-\sqrt{L_n})=r_n^{o(1)}.
\]
Then there are a label ball \(B_{\mathcal Y}(\mathfrak y_n,\zeta_n)\) and a measurable subset
\(\mathcal S_n\subset[r_n,r_n^{\alpha_\eta}]\) such that
\begin{equation}
\label{eq:v081-subpower-label-scale-content}
\boxed{
\int_{\mathcal S_n}\frac{dR}{R}
\gtrsim_{\eta,\mathcal Y}
\zeta_n^{D_{\mathcal Y}}L_n
=
r_n^{o(1)}L_n,
\qquad
\mathfrak y_R\in B_{\mathcal Y}(\mathfrak y_n,\zeta_n)
\quad(R\in\mathcal S_n).
}
\end{equation}
After a subsequence, \(\mathfrak y_n\to\mathfrak y_\infty\), so all retained centers, angular
charts, and polarization planes converge on \(\mathcal S_n\).  The loss in
\eqref{eq:v081-subpower-label-scale-content} is subpower and therefore does not consume the fixed
gap between the residual exponent \(2-\eta/2\) and the endpoint exponent \(2\).
\end{lemma}

\begin{proof}
Cover \(\mathcal Y\) by the stated number of \(\zeta_n\)-balls and partition the logarithmic block
according to the first ball containing \(\mathfrak y_R\).  Its \(dR/R\)-length is
\((1-\alpha_\eta)L_n\).  One cell therefore has measure at least
\[
\frac{(1-\alpha_\eta)L_n}
{C_{\mathcal Y}\zeta_n^{-D_{\mathcal Y}}},
\]
which is \eqref{eq:v081-subpower-label-scale-content}.  Compactness gives the convergent
subsequence.  Finally
\[
\zeta_n^{D_{\mathcal Y}}
=
\exp(-D_{\mathcal Y}\sqrt{L_n})
=
r_n^{D_{\mathcal Y}/\sqrt{L_n}}
=
r_n^{o(1)},
\]
whereas the exponent gain \(\eta/2\) is fixed.
\end{proof}

\begin{proposition}[Positive-density coherent label or separated label variation]
\label{prop:v081-log-block-label-measure-dichotomy}
Start from the logarithmic block
\[
\mathcal B_n^0=[r_n,r_n^{\alpha_\eta}],
\qquad
|\mathcal B_n^0|_{\log}
:=
\int_{\mathcal B_n^0}\frac{dR}{R}.
\]
First restrict to one of the finitely many angular charts on a measurable scale subset
\(\mathcal B_n\subset\mathcal B_n^0\) satisfying
\[
|\mathcal B_n|_{\log}
\gtrsim
|\mathcal B_n^0|_{\log}.
\]
Let
\(R\mapsto\mathfrak y_R\in\mathcal Y\) be the remaining continuous geometric-label choice from
\cref{lem:v081-subpower-label-synchronization}, and define the label probability
\[
\varpi_n
:=
\frac1{|\mathcal B_n|_{\log}}
(\mathfrak y_\bullet)_\#
\left(\mathbf1_{\mathcal B_n}\frac{dR}{R}\right).
\]
After passing to a subsequence, \(\varpi_n\rightharpoonup\varpi\).  According as \(\varpi\) has
an atom or is non-atomic, the corresponding one of the following two alternatives is available.
\begin{enumerate}[label=\textnormal{(\roman*)}]
\item There are \(\mathfrak y_\infty\in\mathcal Y\), \(\theta>0\), radii
\(\zeta_n\downarrow0\), and scale sets \(\mathcal S_n\subset\mathcal B_n\) such that
\begin{equation}
\label{eq:v081-positive-density-label-atom}
\boxed{
\int_{\mathcal S_n}\frac{dR}{R}
\ge
\theta|\mathcal B_n|_{\log},
\qquad
d_{\mathcal Y}(\mathfrak y_R,\mathfrak y_\infty)
\le\zeta_n
\quad(R\in\mathcal S_n).
}
\end{equation}
Thus the center/polarization labels converge, in one fixed angular chart, on a fixed positive proportion of every
retained logarithmic block.
\item There are compact sets \(Y_1,Y_2\subset\mathcal Y\) and constants
\(\theta,c>0\) such that
\begin{equation}
\label{eq:v081-positive-density-label-variation}
\boxed{
\dist(Y_1,Y_2)\ge c,
\qquad
\int_{\{R\in\mathcal B_n:\mathfrak y_R\in Y_i\}}\frac{dR}{R}
\ge
\theta|\mathcal B_n|_{\log}
\quad(i=1,2)
}
\end{equation}
for all large \(n\).  Hence a fixed positive proportion of scales carries genuinely separated
physical-center or polarization data.
\end{enumerate}
The first alternative supplies the positive-density finite-dimensional label coherence required
on the endpoint-star side of
the optional Maslov secant--to--curve promotion.  In the second, the corresponding implication is
the variation-to-four-cluster projection alternative stated there; it is no longer hidden inside a
scale or entropy loss.
\end{proposition}

\begin{proof}
Compactness of \(\mathcal Y\) gives a weakly convergent subsequence.  If \(\varpi\) has an atom
\(\mathfrak y_\infty\) of mass \(2\theta>0\), choose a decreasing sequence of continuity radii
\(\zeta^{(k)}\downarrow0\) for the balls about \(\mathfrak y_\infty\).  Weak convergence gives
\[
\varpi_n(B(\mathfrak y_\infty,\zeta^{(k)}))\ge\theta
\]
for all \(n\) sufficiently large depending on \(k\).  A diagonal choice
\(k=k(n)\to\infty\) proves \eqref{eq:v081-positive-density-label-atom}.

If \(\varpi\) has no atoms, its support contains two distinct points.  Choose disjoint compact
neighborhoods \(Y_1,Y_2\), separated by some \(c>0\), whose \(\varpi\)-masses are both larger than
\(2\theta>0\) and whose boundaries have zero \(\varpi\)-mass.  Weak convergence then gives
\(\varpi_n(Y_i)\ge\theta\) for \(i=1,2\), proving
\eqref{eq:v081-positive-density-label-variation}.
\end{proof}

\begin{corollary}[The positive-density dichotomy retains the physical endpoint-star measures]
\label{cor:v081-endpoint-measure-label-dichotomy}
Let \(\mathcal X\) be the fixed compact localization of the finite physical endpoint/flag tuples
used by one routed endpoint star, and let
\(\Gamma_R\in\mathcal P(\mathcal X)\) be its normalized retained occurrence measure at scale
\(R\).  Replace the finite-dimensional label space in
\cref{prop:v081-log-block-label-measure-dichotomy} by
\[
\mathcal Y^\sharp
:=
\mathcal Y\times\mathcal P(\mathcal X),
\]
where \(\mathcal P(\mathcal X)\) has any metric inducing weak convergence.  Then the same
atomic/non-atomic conclusion holds.

In the atomic alternative, on a fixed positive proportion of each logarithmic block,
\begin{equation}
\label{eq:v081-positive-density-endpoint-measure-coherence}
\boxed{
\mathfrak y_R\longrightarrow\mathfrak y_\infty,
\qquad
\Gamma_R\rightharpoonup\Gamma_\infty,
}
\end{equation}
uniformly in the selected scale sets in the diagonal sense of
\eqref{eq:v081-positive-density-label-atom}.  Thus the actual physical edge-chain measures, not
only their center and polarization labels, satisfy the endpoint-measure coherence part of
the optional Maslov secant--to--curve promotion on positive-density \emph{thickness} blocks.  Passing from that
thickness density to the exact angular-scale-density hypothesis of that problem remains a separate
scale-map check.  In the non-atomic alternative, two positive-proportion scale families remain a
fixed positive distance apart in either geometric labels or weak endpoint-star measure.
\end{corollary}

\begin{proof}
The space \(\mathcal P(\mathcal X)\) is compact metrizable because \(\mathcal X\) is compact
metrizable.  Hence \(\mathcal Y^\sharp\) is compact.  Apply the proof of
\cref{prop:v081-log-block-label-measure-dichotomy} verbatim to the pushforward logarithmic scale
probabilities on \(\mathcal Y^\sharp\).  No covering-number estimate is used in that proof.
\end{proof}

\begin{corollary}[The atomic label branch has an exact angular-scale exponent trichotomy]
\label{cor:v081-atomic-angular-scale-trichotomy}
In the atomic alternative
\eqref{eq:v081-positive-density-label-atom}, let
\(\varkappa_\infty\in[0,1]\) be the limiting angular-scale exponent.  Uniformly on the selected
positive-density scale sets,
\begin{equation}
\label{eq:v081-atomic-angular-scale-law}
\boxed{
\tau_R
=
R^{\varkappa_\infty+o(1)},
\qquad
\frac R{\tau_R}
=
R^{1-\varkappa_\infty+o(1)}.
}
\end{equation}
Consequently:
\begin{enumerate}[label=\textnormal{(\roman*)}]
\item if \(0<\varkappa_\infty<1\), the physical error is \(o(\tau_R)\), so this is a genuine
noncollapsed angular Darboux scale;
\item if \(\varkappa_\infty=1\), the angular and physical scales are power-comparable up to
subpower factors; this is the only exponent compatible with the resolution-scale cap branch, and
the two-probe scale lock of
\cref{cor:v081-hairbrush-alternating-two-probe-lock,%
prop:v081-resolution-cap-freeze-or-return} forces precisely this endpoint;
\item if \(\varkappa_\infty=0\), then
\(\tau_R=R^{o(1)}\); this is the sole remaining degenerate thickness-to-angular scale map.
\end{enumerate}
Thus the atomic branch has no unidentified continuum of scale ratios.
\end{corollary}

\begin{proof}
The exponent coordinate is part of the compact label and converges uniformly on the selected
sets in \eqref{eq:v081-positive-density-label-atom}.  Hence
\[
\log(1/\tau_R)
=
(\varkappa_\infty+o(1))\log(1/R),
\]
which is \eqref{eq:v081-atomic-angular-scale-law}.  The three alternatives are the three possible
positions of \(\varkappa_\infty\) in \([0,1]\).
\end{proof}

\begin{lemma}[A compatible alternating chain has one limiting Reeb orbit]
\label{lem:v081-nested-probe-reeb-orbit}
Let
\[
H^{(0)}\supset H^{(1)}\supset\cdots
\]
be one compatible nested branch of the endpoint-star tree in
\cref{cor:v081-no-good-leaf-endpoint-chain}.  Suppose that \(H^{(\ell)}\) is contained in a
direction ball of radius \(T_\ell\downarrow0\).  For \(j=0,1\), let
\[
(y_\ell^{(j)},s_\ell^{(j)})
\in
\mathbb R^3\times(I_{\rm out}\cup J)
\]
be the two alternating point probes retained on \(H^{(\ell+1)}\), and write
\(\epsilon_\ell^{(j)}\to0\) for their physical errors:
\begin{equation}
\label{eq:v081-nested-probe-error}
\boxed{
|b(a)+s_\ell^{(j)}a-y_\ell^{(j)}|
\le
\epsilon_\ell^{(j)}
\qquad
(a\in H^{(\ell+1)},\ j=0,1).
}
\end{equation}
After passage to a subsequence there are
\[
a_*,b_*\in\mathbb R^3,
\qquad
s_*^{(j)}\in I_{\rm out}\cup J,
\]
such that
\begin{equation}
\label{eq:v081-nested-probe-reeb-orbit}
\boxed{
s_\ell^{(j)}\longrightarrow s_*^{(j)},
\qquad
y_\ell^{(j)}
\longrightarrow
b_*+s_*^{(j)}a_*
\qquad(j=0,1).
}
\end{equation}
Thus every limiting physical probe center on this branch belongs to the Euclidean image
\begin{equation}
\label{eq:v081-single-reeb-orbit}
\boxed{
\mathcal O_{a_*,b_*}
:=
\{(b_*+sa_*,s):s\in\mathbb R\}
}
\end{equation}
of one Reeb orbit above the limiting line \((a_*,b_*)\).  In particular, continuous transverse
drift of the physical centers is not an independent non-atomic label on a compatible infinite
branch; only the longitudinal Reeb level can still drift.  Moreover, because
\(\dist(J,I_{\rm out})>0\),
\begin{equation}
\label{eq:v081-nested-probe-slope}
\boxed{
\frac{y_\ell^{(1)}-y_\ell^{(0)}}
{s_\ell^{(1)}-s_\ell^{(0)}}
\longrightarrow a_*.
}
\end{equation}
\end{lemma}

\begin{proof}
Choose \(a_\ell\in H^{(\ell+1)}\).  If \(m>\ell\), then both \(a_m\) and \(a_\ell\) belong to
the \(O(T_\ell)\)-cap containing \(H^{(\ell)}\).  Hence \((a_\ell)\) is Cauchy and converges to
some \(a_*\).  The fixed localization of \(G_*\) makes \(b(a_\ell)\) bounded.  Pass to a
subsequence on which
\[
b(a_\ell)\to b_*,
\qquad
s_\ell^{(j)}\to s_*^{(j)}
\quad(j=0,1).
\]
Substitution of \(a=a_\ell\) in
\eqref{eq:v081-nested-probe-error} gives
\[
y_\ell^{(j)}
=
b(a_\ell)+s_\ell^{(j)}a_\ell+o(1),
\]
which proves \eqref{eq:v081-nested-probe-reeb-orbit}.  Subtract the two identities and divide by
the uniformly separated denominator \(s_\ell^{(1)}-s_\ell^{(0)}\) to obtain
\eqref{eq:v081-nested-probe-slope}.
\end{proof}

\begin{proposition}[Four Reeb times are a Maslov firewall for a coherent endpoint star]
\label{prop:v081-four-time-polarization-firewall}
Let \(R_n\downarrow0\), let \(R_n\le\tau_n\le1\), and suppose
\begin{equation}
\label{eq:v081-four-time-error-negligible}
\boxed{
\frac{R_n}{\tau_n}\longrightarrow0.
}
\end{equation}
For \(j=0,1,2,3\), suppose that the retained physical endpoint-star measures furnish
\[
s_{n,j}\to s_j,
\qquad
V_{n,j}\to V_\infty\in\Gr(3,6),
\qquad
q_{n,j}=(\alpha_{n,j},\beta_{n,j})\in S^5,
\]
where
\begin{equation}
\label{eq:v081-four-time-flag-incidence}
\boxed{
\dist(q_{n,j},V_{n,j})=o(1),
\qquad
|\beta_{n,j}+s_{n,j}\alpha_{n,j}|
\lesssim
\frac{R_n}{\tau_n}+o(1).
}
\end{equation}
If \(s_0,s_1,s_2,s_3\) are distinct, then
\begin{equation}
\label{eq:v081-four-time-front-null}
\boxed{
P_{V_\infty}\equiv0,
\qquad
V_\infty\in
\mathsf K_{101}\cup\mathsf K_{110}\cup\mathsf K_{200},
\qquad
\rank(\pi_a|_{V_\infty})\le2.
}
\end{equation}
Consequently this configuration cannot support a noncollapsed three-dimensional horizontal
blow-up.

More generally, if the horizontal frame matrices of the coherent planes approach the rank-\(r\)
locus, then \(r+1\) distinct limiting Reeb times suffice.  This is exactly the rank-adapted
interpolation of \cref{lem:v046-rank-adapted-interpolation}.
\end{proposition}

\begin{proof}
After a further subsequence, \(q_{n,j}\to q_j\in S^5\).  The two displayed incidence and
error-negligibility relations in the proposition give
\[
q_j\in V_\infty\cap L_{s_j},
\qquad
q_j\ne0.
\]
The matrix-pencil/Maslov equivalence
\eqref{eq:v003-maslov-equivalence} yields
\[
P_{V_\infty}(s_j)=0
\qquad(0\le j\le3).
\]
This polynomial has degree at most three, so four distinct roots force
\(P_{V_\infty}\equiv0\).  The front-null classification and direction-rank law used in
\cref{thm:v046-planar-angular-exclusion} give
\eqref{eq:v081-four-time-front-null}.  The final assertion follows from
\cref{lem:v046-rank-adapted-interpolation} after letting the incidence and rank errors tend to zero.
\end{proof}

\begin{corollary}[Finite-scale compactness form of the four-time firewall]
\label{cor:v081-finite-scale-four-time-firewall}
Let
\[
\mathfrak N
:=
\{V\in\Gr(3,6):P_V\equiv0\}.
\]
For every \(\gamma>0\) and \(\eta>0\) there is
\(\varepsilon_0=\varepsilon_0(\gamma,\eta,I^\sharp)>0\) with the following property.  If
\[
\max_{i,j}d_{\Gr}(V_i,V_j)\le\varepsilon_0,
\qquad
|s_i-s_j|\ge\gamma\quad(i\ne j),
\]
and unit vectors \(q_i=(\alpha_i,\beta_i)\) satisfy
\[
\dist(q_i,V_i)\le\varepsilon_0,
\qquad
|\beta_i+s_i\alpha_i|\le\varepsilon_0
\qquad(0\le i\le3),
\]
then
\[
\dist_{\Gr}(V_i,\mathfrak N)<\eta
\qquad(0\le i\le3).
\]
Thus, after fixing the Reeb packet separation and the rank-loss neighborhood, choosing the
physical incidence and polarization-cell errors below \(\varepsilon_0\) sends every coherent
four-time packet to the quantitative front-null/rank-loss ledger.
\end{corollary}

\begin{proof}
Otherwise take a violating sequence with both errors tending to zero while \(\gamma,\eta\) remain
fixed.  Compactness of \(I^\sharp\), \(S^5\), and \(\Gr(3,6)\) gives limits with one common plane
\(V\), four pairwise \(\gamma\)-separated times, and nonzero vectors in
\(V\cap L_{s_i}\).  \Cref{prop:v081-four-time-polarization-firewall} gives \(V\in\mathfrak N\), contradicting the
\(\eta\)-separation.  The classification in
\eqref{eq:v081-four-time-front-null} identifies this neighborhood with the stated horizontal
rank-loss output.
\end{proof}

\begin{lemma}[One coherent polarization is front-null or has three shrinking Reeb packets]
\label{lem:v081-coherent-plane-three-time-packets}
Let \(I^\sharp\) be a fixed compact Reeb interval, let \(V_n\to V_\infty\) in
\(\Gr(3,6)\), and let
\[
q_n=(\alpha_n,\beta_n)\in S^5,
\qquad
s_n\in I^\sharp,
\]
satisfy
\[
\dist(q_n,V_n)\le\varepsilon_n,
\qquad
|\beta_n+s_n\alpha_n|\le\varepsilon_n,
\qquad
\varepsilon_n\downarrow0.
\]
Exactly one of the following holds.
\begin{enumerate}[label=\textnormal{(\roman*)}]
\item \(P_{V_\infty}\equiv0\), and hence
\[
V_\infty\in
\mathsf K_{101}\cup\mathsf K_{110}\cup\mathsf K_{200},
\qquad
\rank(\pi_a|_{V_\infty})\le2.
\]
\item The real zero set
\[
Z_\infty
:=
\{s\in I^\sharp:P_{V_\infty}(s)=0\}
\]
has at most three points and
\begin{equation}
\label{eq:v081-coherent-three-time-packet-width}
\boxed{
\dist(s_n,Z_\infty)
\lesssim_{I^\sharp,V_\infty}
\left(
\varepsilon_n+d_{\Gr}(V_n,V_\infty)
\right)^{1/3}.
}
\end{equation}
If \(Z_\infty=\varnothing\), no such approximate incidences exist for all large \(n\).
\end{enumerate}
\end{lemma}

\begin{proof}
Choose orthonormal frame matrices \(A_n,B_n\) for \(V_n\), converging after an orthogonal change of
frame to matrices \(A_\infty,B_\infty\) for \(V_\infty\).  Project \(q_n\) to \(V_n\).  In frame
coordinates this gives a vector \(c_n\) with \(|c_n|=1+O(\varepsilon_n)\) and
\[
|(B_n+s_nA_n)c_n|
\lesssim_{I^\sharp}
\varepsilon_n.
\]
The other two singular values of \(B_n+s_nA_n\) are uniformly bounded, so
\[
|P_{V_n}(s_n)|
=
|\det(B_n+s_nA_n)|
\lesssim_{I^\sharp}
\varepsilon_n.
\]
Continuity of the determinant on the compact frame and Reeb-parameter sets yields
\begin{equation}
\label{eq:v081-coherent-plane-polynomial-small}
|P_{V_\infty}(s_n)|
\lesssim_{I^\sharp}
\varepsilon_n+d_{\Gr}(V_n,V_\infty).
\end{equation}
If \(P_{V_\infty}\equiv0\), the classification and rank law give \textnormal{(i)}.  Otherwise factor
the fixed nonzero polynomial over \(\mathbb C\).  On \(I^\sharp\), its nonreal factors and its
distance from zero away from \(Z_\infty\) are uniformly bounded below.  Since its degree is at most
three,
\[
|P_{V_\infty}(s)|
\gtrsim_{I^\sharp,V_\infty}
\min\{1,\dist(s,Z_\infty)^3\}.
\]
Combine this with \eqref{eq:v081-coherent-plane-polynomial-small}.  If \(Z_\infty\) is empty, the
same compactness bound contradicts the right side tending to zero.
\end{proof}

\begin{corollary}[The three Reeb packets return to the existing flagged Maslov routing]
\label{cor:v081-three-time-packets-reroute}
Work on a positive-density logarithmic block with \(\varkappa_\infty<1\).  Suppose the atomic
endpoint-measure alternative retains a fixed portion of the current determinant-small four-cycle
occurrence measure, including its actual cycle-plane label, and that these plane labels converge to
\(V_\infty\).  Then this coherent-polarization part has no new multiscale endpoint.  At the same
occurrence-mass scale it returns to one of:
\begin{enumerate}[label=\textnormal{(\roman*)}]
\item a physical Lagrangian point-probe bush;
\item a smaller angular packet or a close-anchor/horizontal-line ledger;
\item a transverse old/new polarization ledger.
\end{enumerate}
In particular a coherent plane together with an at-most-three-level Reeb packet is not a fourth
alternative to the local contact--Maslov trichotomy.
\end{corollary}

\begin{proof}
If \(P_{V_\infty}\equiv0\), then
\(\rank(\pi_a|_{V_\infty})\le2\).  The current cycle planes therefore lie in the determinant-small
horizontal-rank branch.  Apply the arbitrary-threshold route and the lossless finite-flag
amplification
\cref{cor:v081-k-flag-arbitrary-threshold,prop:v081-rank-loss-to-bush-reduction}: its fully coherent
return is excluded, while its complement gives one of the three stated outputs.

Suppose \(P_{V_\infty}\not\equiv0\).  By
\cref{lem:v081-coherent-plane-three-time-packets}, every retained collision label lies in one of at
most three intervals whose common width tends to zero and whose centers are the real roots of
\(P_{V_\infty}\).  Discard empty root packets.  If two root centers coalesce, combine their
intervals; otherwise their nonzero limiting gaps are bounded below.

Apply the weighted cycle-relative time trichotomy
\cref{lem:v050-weighted-relative-time-trichotomy} with a time resolution larger than the shrinking
packet widths and smaller than the fixed gaps between distinct roots.  Its close-anchor output is
angular.  Its two-interval output returns, without loss of the \(Z/T\) occurrence-mass scale, to a
shrinking bush or angular packet by
\cref{cor:v050-two-time-scale-persistence}.  In the remaining output, three pairwise separated
labels must lie in the three distinct root packets.  Their Vandermonde product is therefore bounded
below by a positive constant depending only on \(V_\infty\).  The weighted continuous anchored rank
packet
\cref{prop:v050-continuous-anchored-rank-packet} applies.  Its nondegenerate determinant complement
is already the physical bush output of
\cref{prop:v050-arbitrary-threshold-maslov-routing}; its determinant-small part is again eliminated
or sent to a transverse ledger by
\cref{cor:v081-k-flag-arbitrary-threshold,prop:v081-rank-loss-to-bush-reduction}.  These cases exhaust
the retained four-cycle occurrence measure.
\end{proof}

\begin{proposition}[Same-vertex non-atomic polarization variation is already transverse-paid]
\label{prop:v081-same-vertex-polarization-variation-paid}
Let a normalized fine Maslov graph have edge mass
\[
M=r^{\chi+o(1)},
\qquad
0\le\chi<2,
\]
and retain the conditional old flags by
\cref{lem:v081-lossless-conditional-flag-tensor}.  Suppose a fixed portion of its determinant-small
four-cycle occurrence mass has, at the decorated free endpoint, an old polarization normal in a
compact set \(N_1\subset\mathbb P^2\) and the current new-cycle polarization normal in a compact set
\(N_2\subset\mathbb P^2\), where
\[
\dist(N_1,N_2)\ge c_0>0.
\]
Then this portion belongs to the transverse old/new polarization ledger; it cannot be an unpaid
coherent endpoint-star return.  More precisely, after choosing the packet threshold
\(\vartheta=r^\alpha<c_0/4\), every such decorated endpoint obeys the selector-mass bound
\[
\lesssim_{U,c_0}
(t_{\rm old}+r)(t_{\rm new}+r),
\]
and the finite-flag parameter choice in
\cref{cor:v081-k-flag-arbitrary-threshold} makes the complementary class in which every sampled old
flag is coherent with the new plane smaller than the required \(M^4\) rank-loss mass.
\end{proposition}

\begin{proof}
Orient representatives of the projective normals locally.  Their separation gives
\[
|n_{\rm old}\times n_{\rm new}|
\gtrsim c_0.
\]
The old and new packets meet the same graph product cell at the decorated endpoint.  Alternative
\textnormal{(i)} of
\cref{lem:v078-old-flag-new-cycle-polarization}, used with any
\(\vartheta<c_0/4\), therefore gives
\[
\sigma(\text{possible decorated directions})
\lesssim_{U,c_0}
(t_{\rm old}+r)(t_{\rm new}+r).
\]
For the quantitative four-cycle occurrence, lift the free endpoint to the number \(k\) of
independent conditional old flags prescribed in
\cref{cor:v081-k-flag-arbitrary-threshold}.  If at least one flag lies in \(N_1\) while the new
normal lies in \(N_2\), the preceding transverse estimate applies.  If all flags lie within the
\(\vartheta\)-cell of the new normal, their occurrence mass is bounded by
\[
\left(\frac{H_i}{\lambda_*}\vartheta\right)^k
\]
as in the proof of
\cref{prop:v081-eight-flag-full-defect-exclusion}.  The inequalities
\[
\frac{4\chi}{k}<\alpha<d-\beta
\]
make this \(o(M^4)\), while
\(t_{\rm new}=o(\vartheta)\).  Hence a fixed portion of the rank-loss mass cannot remain in the
fully coherent class, and every surviving portion is transverse-paid.
\end{proof}

\begin{corollary}[Residual polarization variation is purely cross-branch]
\label{cor:v081-polarization-variation-cross-branch}
In the non-atomic polarization alternative of
\cref{prop:v081-log-block-label-measure-dichotomy}, any positive occurrence mass on which the two
separated polarization families are joined by a retained old-flag/current-cycle incidence at one
decorated endpoint is disposed of by
\cref{prop:v081-same-vertex-polarization-variation-paid}.  Consequently the only polarization
variation not already present in the transverse ledger is segregated variation between distinct
stopping-tree branches: no positive-mass decorated endpoint carries both sides of the separation.

Thus the four-cluster alternative in
the optional Maslov secant--to--curve promotion is needed only for this cross-branch remainder.  It is not needed
to handle local drift of a coherent Maslov polarization, which is already excluded by the
conditional multi-flag symplectic comparison.
\end{corollary}

\begin{proposition}[Separated polarization branches either decouple or create a transverse
Maslov four-cycle]
\label{prop:v081-cross-branch-polarization-decoupling}
Let \(\nu_1,\nu_2\) be two restrictions of one normalized selector occurrence measure, with
\[
\nu_i(G)\ge\theta>0
\qquad(i=1,2),
\]
and suppose every retained old polarization normal over \(\nu_i\) lies in a compact set
\(N_i\subset\mathbb P^2\), where
\[
\dist(N_1,N_2)\ge c_0>0.
\]
Let \(h_r\), \(0\le h_r\le1\), be the physical normalized Maslov graph at scale \(r\), with all old
conditional flags retained, and put
\[
X_r
:=
\iint h_r(z_1,z_2)\,d\nu_1(z_1)d\nu_2(z_2).
\]
Exactly one of the following occurs.
\begin{enumerate}[label=\textnormal{(\roman*)}]
\item \(X_r\le r^{2-o(1)}\); the cross-branch contribution is already at the quadratic residual
endpoint and may be deleted before treating the two diagonal graphs separately.
\item Along a subsequence
\[
X_r=r^{\chi+o(1)}
\qquad\text{for some }0\le\chi<2.
\]
Then the cross graph returns, at the same power-excess scale, to a physical Lagrangian bush, a
close-anchor/angular or horizontal-line output, or the transverse old/new polarization ledger.
\end{enumerate}
Thus separated cross-branch polarization mixing is not an additional coherent endpoint.
\end{proposition}

\begin{proof}
Only \textnormal{(ii)} requires proof.  Normalize \(\nu_i\) to probabilities and symmetrize the
bipartite graph on their disjoint union.  Because both masses are bounded below by \(\theta\), this
changes \(X_r\) only by a constant.  Two applications of Cauchy--Schwarz give alternating
bipartite four-cycle mass \(\gtrsim_\theta X_r^4\).  Apply the arbitrary-threshold Maslov routing
\cref{prop:v050-arbitrary-threshold-maslov-routing}.  Its determinant-nondegenerate output is a
physical bush.  Its close-anchor and horizontal-line outputs are already among the stated ledgers.

It remains to consider a determinant-small alternating cycle.  Its vertices alternate between
\(\nu_1\) and \(\nu_2\), and all retain their old conditional polarization flags.  Let
\(n_{\rm new}\) be the normal of the current cycle packet.  If
\(\vartheta<c_0/4\), then \(n_{\rm new}\) cannot lie within \(\vartheta\) of both \(N_1\) and
\(N_2\).  Hence at least one of the four vertices has an old flag
\(\vartheta\)-transverse to the new cycle polarization.  Cyclically re-root the four-cycle so that
this vertex is the decorated free endpoint.  The rotation costs only a factor four.  \Cref{lem:v078-old-flag-new-cycle-polarization} places that portion in the transverse ledger.
Therefore a fixed part of the alternating four-cycle mass reaches one of the listed outputs.
\end{proof}

\begin{corollary}[Non-atomic polarization is reduced to coherent components without a power loss]
\label{cor:v081-nonatomic-polarization-component-reduction}
At a fine scale \(r\), choose \(\zeta_r\downarrow0\) with
\(\zeta_r^{-1}=r^{-o(1)}\), and split the normalized Maslov graph according as the two retained old
polarization normals have distance at most \(C\zeta_r\) or greater than \(C\zeta_r\).  Cover the
compact polarization space by \(r^{-o(1)}\) \(\zeta_r\)-cells and sum the far graph over cell pairs.
If its total edge mass exceeds \(r^{2-o(1)}\), one separated cell pair still has edge exponent
strictly below two, and
\cref{prop:v081-cross-branch-polarization-decoupling} returns it to a physical or transverse
output.  Otherwise delete the far graph at total quadratic cost.  A bounded-overlap net
decomposition of the near graph leaves \(r^{-o(1)}\) graph pieces on each of which both endpoint
polarizations lie within \(O(\zeta_r)\) of one common plane.

Consequently the non-atomic polarization alternative does not require a new four-cluster theorem
at the collision-energy level.  After a subpower partition it is either already transverse-paid or
graph-decoupled into the coherent components handled by
\cref{cor:v081-three-time-packets-reroute}.  The only branch not covered by this conclusion is the
resolution endpoint \(\varkappa_\infty=1\), where the physical incidence error is not
Maslov-negligible.
\end{corollary}

\begin{proof}
The number of polarization cells and ordered cell pairs is \(r^{-o(1)}\).  Hence pigeonholing a
far contribution loses only a subpower factor and preserves every fixed exponent gap below two.
Apply
\cref{prop:v081-cross-branch-polarization-decoupling} to a separated color class.  If no such
contribution is excessive, positivity gives a total deleted edge mass \(r^{2-o(1)}\).  Within each
remaining near piece the two old polarization labels lie in one \(O(\zeta_r)\)-ball.  Compare the
current cycle plane with that ball.  A separated current plane is same-vertex transverse-paid by
\cref{prop:v081-same-vertex-polarization-variation-paid}; a coherent current plane enters
\cref{cor:v081-three-time-packets-reroute}.  The piece edge masses add before normalization, and
the subpower net count is absorbed by the fixed exponent gap.
\end{proof}

\begin{corollary}[Symplectic reduction of the remaining label drift]
\label{cor:v081-nested-label-symplectic-reduction}
Apply the endpoint-measure compactness of
\cref{cor:v081-endpoint-measure-label-dichotomy} along one compatible nested branch, and include in
the occurrence label an actual normalized Maslov secant \(q\), its Reeb parameter \(s\), and its
polarization plane \(V\).  On every exponent branch
\[
\varkappa_\infty<1
\]
one has
\[
R/\tau_R=R^{1-\varkappa_\infty+o(1)}\to0,
\]
so the limiting occurrence measure is supported on the exact incidence relation
\[
q\in V\cap L_s,
\qquad q\ne0.
\]
Then:
\begin{enumerate}[label=\textnormal{(\roman*)}]
\item the limiting physical point-probe centers lie on the single Reeb orbit
\(\mathcal O_{a_*,b_*}\) from
\cref{lem:v081-nested-probe-reeb-orbit};
\item if the polarization marginal is non-atomic, two separated compact sets of polarization
planes carry positive occurrence mass; their same-vertex part is transverse-paid by
\cref{prop:v081-same-vertex-polarization-variation-paid}, leaving only cross-branch variation as
the polarized four-cluster input;
\item if the polarization marginal has an atom \(V_\infty\), then, on any noncollapsed
three-dimensional horizontal branch, the conditional Reeb-time support above \(V_\infty\) has at
most three points; at finite scale all flags whose planes converge to \(V_\infty\) lie in the
three shrinking packets of
\eqref{eq:v081-coherent-three-time-packet-width}.
\end{enumerate}
Indeed four points in the conditional time support allow four disjoint time neighborhoods and
four convergent endpoint-star flags, contradicting
\cref{prop:v081-four-time-polarization-firewall}.  Thus, below the resolution endpoint
\(\varkappa_\infty=1\), the non-atomic alternative of
\cref{prop:v081-log-block-label-measure-dichotomy} is no longer arbitrary physical-center drift:
it is either genuine polarization variation or a finite, at-most-three-level Reeb packet over one
coherent polarization.  Whenever that polarization is the retained current cycle-plane label, the
latter packet is returned to the existing bush/angular/transverse routes by
\cref{cor:v081-three-time-packets-reroute}.  The endpoint \(\varkappa_\infty=1\) remains separate because
\eqref{eq:v081-four-time-error-negligible} fails there; the two-probe scale lock forbids replacing
its resolution-size error by an exact tangent Maslov incidence.
\end{corollary}

\begin{proof}
The first assertion is \cref{lem:v081-nested-probe-reeb-orbit}.  If the polarization marginal has
no atoms, the compact-set separation argument in
\cref{prop:v081-log-block-label-measure-dichotomy} applies to that marginal itself.  If it has an
atom \(V_\infty\) and the conditional time support contains four distinct points, choose disjoint
small neighborhoods of those points.  Each has positive conditional mass, so weak convergence of
the retained occurrence measures supplies one actual normalized flag from each neighborhood.
Their plane labels converge to \(V_\infty\), and for
\(\varkappa_\infty<1\) their normalized physical incidence errors tend to zero.  \Cref{prop:v081-four-time-polarization-firewall} gives horizontal rank at most two, contrary to the
noncollapse hypothesis.  \Cref{lem:v081-coherent-plane-three-time-packets} supplies the
finite-scale packet width.  This proves the trichotomy and also explains why the same inference is
unavailable at the resolution endpoint.
\end{proof}

\begin{corollary}[Every below-resolution label branch returns to the existing geometric outputs]
\label{cor:v081-below-resolution-label-closure}
Assume the actual endpoint-star tuple measure, including its current cycle plane and normalized
Maslov secants, is retained as in
\cref{cor:v081-endpoint-measure-label-dichotomy}.  On every positive-density exponent branch with
\[
\varkappa_\infty<1,
\]
the label compactness alternative has no independent output:
\begin{enumerate}[label=\textnormal{(\roman*)}]
\item physical-center drift is confined to one limiting Reeb orbit by
\cref{lem:v081-nested-probe-reeb-orbit};
\item separated polarization variation is transverse-paid or graph-decoupled by
\cref{prop:v081-cross-branch-polarization-decoupling,%
cor:v081-nonatomic-polarization-component-reduction};
\item over one coherent polarization, four limiting Reeb levels force front-null horizontal rank
loss by \cref{prop:v081-four-time-polarization-firewall}, while at most three levels return to the
flagged bush/angular/transverse routing by
\cref{cor:v081-three-time-packets-reroute}.
\end{enumerate}
Consequently any endpoint-star chain which avoids all previously paid geometric outputs must, after
passing to a further positive-density subsequence, satisfy
\[
\boxed{\varkappa_\infty=1.}
\]
Thus the sole remaining multiscale endpoint is the resolution-scale angular cap; no
below-resolution Maslov or polarization branch remains.
\end{corollary}

\begin{proof}
For \(\varkappa_\infty<1\),
\[
R/\tau_R=R^{1-\varkappa_\infty+o(1)}\to0.
\]
Hence the physical endpoint-star measures converge to exact Maslov incidences by
\cref{lem:v046-rescaled-contact-secant} and
\eqref{eq:v046-limit-maslov-support}.  Apply the three alternatives just listed.  They exhaust,
respectively, the physical-center coordinate, the polarization coordinate, and the
Reeb-time/current-cycle part of the compact endpoint-star label.  Therefore an unpaid branch cannot
have \(\varkappa_\infty<1\).  The exponent trichotomy
\cref{cor:v081-atomic-angular-scale-trichotomy} leaves only
\(\varkappa_\infty=1\).
\end{proof}

\subsection{Resolution shells and the density cocycle}

\begin{lemma}[The resolution shell is automatically relative-quadratic]
\label{lem:v081-resolution-shell-relative-quadratic}
Let \(H\subset U\) have direction mass \(p>0\), and let
\(\mu_H=(\id,b)_\#(\sigma|_H)\).  For \(0<\rho\le\tau<1\), the angular-shell residual satisfies
\begin{equation}
\label{eq:v081-resolution-shell-basic-bound}
\boxed{
\mathfrak Z_{\rho,\tau}^{J^+}(\mu_H)
\lesssim_U
p\tau^2.
}
\end{equation}
Consequently, along the tax-free failure sequence of
\cref{prop:v081-resolution-cap-freeze-or-return}, where
\[
p=\rho^{o(1)},
\]
every shell with angular exponent
\[
\frac{\log(1/\tau)}{\log(1/\rho)}\longrightarrow1
\]
obeys the relative endpoint estimate
\begin{equation}
\label{eq:v081-resolution-shell-relative-bound}
\boxed{
\mathfrak Z_{\rho,\tau}^{J^+}(\mu_H)
\lesssim
p^2\rho^{2-o(1)}.
}
\end{equation}
In particular a fixed power-excess residual failure cannot be carried by
\(\varkappa_\infty=1\).
\end{lemma}

\begin{proof}
Discard the two residual indicators in the definition
\eqref{eq:v048-maslov-residual-content}.  On the shell
\(\tau\le|a-a'|<2\tau\), the weight is at most \(\tau^{-1}\).  The bounded Lebesgue direction
density gives, for every \(a\in H\),
\[
\sigma\{a'\in H:|a-a'|<2\tau\}
\lesssim_U\tau^3.
\]
Integration in \(a\) proves
\eqref{eq:v081-resolution-shell-basic-bound}.  At exponent one,
\[
\tau=\rho^{1+o(1)},
\qquad
\frac1p=\rho^{-o(1)}.
\]
Therefore
\[
p\tau^2
=
p^2\rho^2
\left(\frac{\tau}{\rho}\right)^2\frac1p
=
p^2\rho^{2-o(1)},
\]
which is \eqref{eq:v081-resolution-shell-relative-bound}.  This is incompatible with a lower bound
\(\gtrsim p^2\rho^{2-\eta}\) for any fixed \(\eta>0\).
\end{proof}

\begin{corollary}[Every genuine fine return has Maslov-negligible physical error]
\label{cor:v081-fine-return-error-negligible}
At a resolution cap, decompose any relative power-excess residual failure into the dyadic angular
shells of \cref{lem:v048-residual-shell-ledger}.  The \(r^{-o(1)}\) number of shells costs no fixed
exponent.  By
\cref{lem:v081-resolution-shell-relative-quadratic}, a selected bad shell cannot have
\(\varkappa_\infty=1\).  Hence, after a subsequence,
\[
\varkappa_\infty<1,
\qquad
\frac{\rho}{\tau_\rho}\longrightarrow0.
\]
Thus every actual fine return from
\cref{prop:v081-resolution-cap-freeze-or-return} lies in the error-negligible contact--Maslov
regime and is subject to
\cref{cor:v081-below-resolution-label-closure}.  The resolution cap introduces no additional Maslov endpoint and reduces the proof to the mass aggregation of successive physical bush returns.
\end{corollary}

\begin{proof}
The ultra-near term in
\cref{lem:v048-residual-shell-ledger} is already below the target after the standard choice of its
cutoff.  If the sum of the remaining \(O(\log(2/\rho))=\rho^{-o(1)}\) shells has a fixed power
excess, one shell retains that excess up to a subpower factor.  The preceding lemma excludes
exponent one for this shell.  The displayed scale separation and the last assertion follow from
\cref{lem:v046-rescaled-contact-secant}.
\end{proof}

\begin{lemma}[Mass loss and non-Frostman concentration form one density cocycle]
\label{lem:v081-bush-return-density-cocycle}
Let a stopping node \(H\) have direction mass \(p\) and be contained in a direction cap of radius
\(T\).  Put
\[
D(H):=\frac{p}{T^3}.
\]
Suppose one complete fine-Maslov/bush/non-Frostman return first retains a proportion at least
\(cM\) of the node mass in its maximal physical bush family.  At the normalized Frostman-failure
scale \(R\), suppose the Reeb-coherent mergers then retain a proportion at least
\[
cA R^{3-\delta}
\]
and, after the intrinsic Maslov renormalization and alternating Reeb return, place that mass inside
a child direction cap of radius at most \(CT\widehat R\).  Here \(R\) and \(\widehat R\) are not
identified; in the returned hairbrush branch one only has the previously proved lower scale lock
\(\widehat R\gtrsim R^{1-d}\).  If the child is \(H'\), then
\begin{equation}
\label{eq:v081-bush-return-density-cocycle}
\boxed{
D(H')
\gtrsim_{I,J,U}
A M R^{3-\delta}\widehat R^{-3}D(H).
}
\end{equation}
For a chain of such returns,
\begin{equation}
\label{eq:v081-bush-return-density-product}
\boxed{
\prod_{\ell<L}
\left(c_{I,J,U}A_\ell M_\ell
R_\ell^{3-\delta}\widehat R_\ell^{-3}\right)
\lesssim_U
\frac1{D(H^{(0)})}.
}
\end{equation}
Consequently, an infinite branch cannot have
\[
A_\ell M_\ell
R_\ell^{3-\delta}\widehat R_\ell^{-3}
\ge C_{I,J,U}>1
\]
at every sufficiently large generation.  Apart from density-increasing generations whose
multiplicative gains have bounded total product, every infinite survivor must lie in the
compensation regime
\begin{equation}
\label{eq:v081-bush-return-compensation-regime}
\boxed{
M_\ell
\lesssim_{I,J,U}
A_\ell^{-1}\widehat R_\ell^3R_\ell^{\delta-3}.
}
\end{equation}
\end{lemma}

\begin{proof}
The hypotheses give
\[
p'
\gtrsim
A M R^{3-\delta}p,
\qquad
T'\lesssim T\widehat R.
\]
Divide the first inequality by \((T')^3\) to obtain
\eqref{eq:v081-bush-return-density-cocycle}.  Iteration gives the lower bound by the product on the
left of \eqref{eq:v081-bush-return-density-product}.  On the other hand the direction marginal is
bounded Lebesgue measure, so every selector subset contained in a \(T'\)-cap has mass
\(\lesssim_U(T')^3\); hence \(D(H')\lesssim_U1\).  Comparing the iterated lower and uniform upper
bounds proves \eqref{eq:v081-bush-return-density-product}.  The last two assertions are immediate.
\end{proof}

\begin{corollary}[The stopping-tree tail is confined to the compensation regime]
\label{cor:v081-stopping-tail-compensation}
After
\cref{cor:v081-fine-return-error-negligible}, every nonterminal node first passes through the exact
contact--Maslov routing and then through the physical maximal-bush/Frostman alternative.  For the
non-Frostman children, denote the normalized Frostman-failure radius by \(R_\ell\), the final
relative radius of the returned child cap by \(\widehat R_\ell\), and combine all fixed retained
fractions and excess constants into \(A_\ell\).  The mass and radius
hypotheses of
\cref{lem:v081-bush-return-density-cocycle} follow from
\cref{prop:v081-maximal-bush-family,%
cor:v081-bush-family-angular-concentration,%
prop:v081-hairbrush-merges-to-one-bush,%
cor:v081-family-to-merged-bush-frostman,%
prop:v081-hairbrush-intrinsic-renormalization}.

Therefore the part of the normalized stopping tree on which
\[
A_\ell M_\ell
R_\ell^{3-\delta}\widehat R_\ell^{-3}
\]
is uniformly larger than one has finite branch depth.  The density cocycle therefore confines the
normalized tail to nodes in the displayed compensation regime.  Equivalently, its remaining
normalized scalar is the explicit domination of the bad Maslov edge mass
\(M_\ell\) by the exact two-scale cap-density debt
\(A_\ell^{-1}\widehat R_\ell^3R_\ell^{\delta-3}\).  This is the same
\(R_\ell/\widehat R_\ell\) cubic enlargement already isolated in the fine and coarse two-scale
frontier estimates.  The cocycle is paired with the absolute-flow exhaustion below, which removes this term.
\end{corollary}

\begin{lemma}[Lossless exhaustion of a hereditary Markov route]
\label{lem:v081-hereditary-markov-exhaustion}
Let \(\Gamma\) be a finite occurrence measure with distinguished source map
\(\pi_{\rm src}\), and allow any finite or countable collection of conditional probability marks
over that source.  Suppose that, for every nonzero restriction \(\Gamma'\le\Gamma\), one application
of a routing alternative produces an output submeasure \(\Theta\le\Gamma'\) with
\[
 \|\Theta\|\ge c\|\Gamma'\|,
 \qquad c>0,
\]
and assigns \(\Theta\) to finitely or countably many measurable destinations by source restrictions
and extensions by probability kernels.  Then \(\Gamma\) has a countable destination decomposition
\begin{equation}
\label{eq:v081-hereditary-markov-exhaustion}
 \boxed{
 \Gamma=\sum_{j\ge1}\Theta_j,
 \qquad
 \sum_{j\ge1}\|\Theta_j\|=\|\Gamma\|,
 }
\end{equation}
up to a null remainder.  If the source marginal of \(\Gamma\) is
\(f\,d\sigma\), then the destination source marginals are
\(f_j,d\sigma\) with
\[
 0\le f_j,
 \qquad
 \sum_j f_j\le f
 \quad\text{a.e.}
\]
Every inherited mark is unchanged at the retained source point.
\end{lemma}

\begin{proof}
Apply the route to \(\Gamma_0:=\Gamma\), remove the produced output, and apply it again to
\(\Gamma_1:=\Gamma_0-\Theta_0\).  More generally set
\(\Gamma_{n+1}:=\Gamma_n-\Theta_n\).  Positivity and the fixed retention give
\[
 \|\Gamma_n\|\le(1-c)^n\|\Gamma\|,
\]
so the remainder tends to zero in total variation and the destination measures sum to
\(\Gamma\).  A source restriction replaces a density by a smaller density, and a probability
kernel has source marginal equal to the measure from which it is extended.  Summing the resulting
Radon--Nikodym densities proves \(\sum_jf_j\le f\).  The same observation shows that an inherited
conditional mark is never resampled or moved to another endpoint.
\end{proof}

\begin{lemma}[Cap-square summation pays fractional sparse edge tails]
\label{lem:v081-cap-square-edge-tail}
Let
\[
d\nu_\gamma=f_\gamma\,d\sigma,
\qquad
0\le f_\gamma\le1,
\qquad
\sum_\gamma f_\gamma\le1,
\]
be finite selector-direction submeasures, each supported in a direction cap of radius at most
\(T\), and put
\[
p_\gamma:=\|\nu_\gamma\|,
\qquad
\sum_\gamma p_\gamma\le m,
\]
where \(m\le\|\sigma\|\).  For measurable kernels \(0\le h_\gamma\le1\),
\begin{equation}
\label{eq:v081-cap-square-edge-tail}
\boxed{
\sum_\gamma
\iint
h_\gamma(a,a')\,d\nu_\gamma(a)d\nu_\gamma(a')
\lesssim_U
mT^3.
}
\end{equation}
Consequently, if
\[
T^3\le mr^{2-o(1)},
\]
the complete within-node edge tail is at most \(m^2r^{2-o(1)}\).  No bound on the total tail mass
\(\sum_\gamma p_\gamma\) stronger than \(m\) is required.
\end{lemma}

\begin{proof}
Bounded Lebesgue direction density and \(f_\gamma\le1\) give
\(p_\gamma\lesssim_UT^3\), while \(h_\gamma\le1\) gives
\[
\iint h_\gamma(a,a')\,d\nu_\gamma(a)d\nu_\gamma(a')
\le p_\gamma^2.
\]
Therefore
\[
\sum_\gamma p_\gamma^2
\le
\left(\sup_\gamma p_\gamma\right)\sum_\gamma p_\gamma
\lesssim_U
mT^3.
\]
\end{proof}

\clearpage
\section{Cross-generation edge-flow closure}

\subsection{Absolute occurrences and the edge-Carleson reduction}

\begin{proposition}[Exact edge-Carleson criterion for closing the bush-return tree]
\label{prop:v081-bush-tree-edge-carleson-criterion}
Keep the hypotheses of
\cref{prob:v081-maslov-bush-frostman-bridge}.  Suppose the collision-preserving bush-return
construction can be saturated to depth \(L=L(r)\) so that, for the incoming absolute normalized
Maslov edge mass \(\mathcal M_{\rm in}(r)\),
\begin{equation}
\label{eq:v081-bush-tree-edge-carleson}
\boxed{
\mathcal M_{\rm in}(r)
\lesssim
r^{-o(1)}
\left(
\mathcal M_{\rm paid}(r)
 +
\mathcal M_{\rm cross}(r)
 +
 \sum_{\gamma\in\mathcal S_L}p_\gamma T_\gamma^3
\right).
}
\end{equation}
Here \(\mathcal S_L\) is the family of sparse terminal hairbrush nodes, \(T_\gamma\) is the
radius of the terminal direction cap, and
\(\mathcal M_{\rm paid}\lesssim m^2r^{2-o(1)}\), and
\(\mathcal M_{\rm cross}\lesssim m^2r^{2-o(1)}\).  If every node in
\(\mathcal S_L\) lies in a direction cap of radius \(T_L\) and
\[
T_L^3\le mr^{2-o(1)},
\]
then the relative residual estimate
\eqref{eq:v078-final-low-reuse-relative-residual} follows and the Maslov--Frostman bridge closes.
The terminal nodes may be fractional selector submeasures
\[
d\nu_\gamma=f_\gamma\,d\sigma,
\qquad
\sum_{\gamma\in\mathcal S_L}f_\gamma\le1,
\qquad
p_\gamma:=\|\nu_\gamma\|;
\]
thus the statement applies to the target-cell/time-packet occurrence partitions below, not only to
pairwise disjoint selector sets.
\end{proposition}

\begin{proof}
Bounded direction density and the fractional source domination give
\[
\sum_{\gamma\in\mathcal S_L}p_\gamma T_\gamma^3
\lesssim_U
mT_L^3
\lesssim
m^2r^{2-o(1)}.
\]
Insert this and the two assumed edge bounds into
\eqref{eq:v081-bush-tree-edge-carleson}.  The collision/residual comparison
\cref{prop:v048-energy-residual-equivalence} gives
\eqref{eq:v078-final-low-reuse-relative-residual}; then
\cref{prop:v078-final-coherent-factor-budget,thm:v048-residual-frostman-criterion} close the stated
bridge.
\end{proof}

\begin{lemma}[The incoming Maslov edge flow survives maximal bush extraction linearly]
\label{lem:v081-incoming-edge-flow-bush-extraction}
Let \(h\) be a symmetric measurable kernel with \(0\le h\le1\) on a probability space
\((G,\mu)\), put
\[
d_h(z):=\int_Gh(z,z')\,d\mu(z'),
\qquad
M:=\iint_{G\times G}h\,d\mu d\mu,
\]
and define the incoming edge-flow measure
\[
d\lambda_h(z):=d_h(z)\,d\mu(z).
\]
Suppose the maximal bush extraction gives pairwise disjoint leaf sets \(H_j\), with
\(H:=\bigcup_jH_j\), such that
\[
\iint_{H^c\times H^c}h\,d\mu d\mu<\frac M2.
\]
Then
\begin{equation}
\label{eq:v081-incoming-edge-flow-bush-mass}
\boxed{
\lambda_h(H)>\frac M4,
\qquad
\sum_j\lambda_h(H_j)=\lambda_h(H),
\qquad
0\le\lambda_h\le\mu.
}
\end{equation}
Moreover, on \(\{d_h>0\}\),
\[
d\kappa_z^h(z')
:=
\frac{h(z,z')}{d_h(z)}\,d\mu(z')
\]
is a probability kernel and the old physical Maslov edge measure disintegrates exactly as
\begin{equation}
\label{eq:v081-incoming-edge-flow-disintegration}
\boxed{
h(z,z')\,d\mu(z)d\mu(z')
=
d\lambda_h(z)d\kappa_z^h(z').
}
\end{equation}
Thus the maximal-bush passage retains a fixed fraction of the incoming \emph{edge} mass, not merely
a comparable amount of unweighted vertex mass.  The neighbor \(z'\), together with all older
conditional Maslov flags, may be carried as a probability mark over the bush endpoint \(z\) with
no further normalization factor.
\end{lemma}

\begin{proof}
Write
\[
E_{11}:=\iint_{H\times H}h,
\qquad
E_{10}:=\iint_{H\times H^c}h,
\qquad
E_{00}:=\iint_{H^c\times H^c}h,
\]
where all integrals use \(d\mu d\mu\).  Symmetry gives
\[
M=E_{11}+2E_{10}+E_{00},
\qquad
\lambda_h(H)=E_{11}+E_{10}
\ge\frac{M-E_{00}}2
>\frac M4.
\]
Disjointness gives the middle identity in
\eqref{eq:v081-incoming-edge-flow-bush-mass}, and \(d_h\le1\) gives the domination
\(\lambda_h\le\mu\).  The definition of \(\kappa_z^h\) proves both its normalization and
\eqref{eq:v081-incoming-edge-flow-disintegration}.  Since all other old flags are already
conditional probability laws over \(z\), their tensor product with \(\kappa_z^h\) again has total
mass one.
\end{proof}

\begin{lemma}[Dyadic degree layers retain the linear edge flow]
\label{lem:v081-edge-flow-degree-layers}
Keep \(h,\mu,d_h,\lambda_h,M\) from
\cref{lem:v081-incoming-edge-flow-bush-extraction}.  Fix \(A>2\), a scale \(0<r<1/2\), put
\(K:=\lceil A\log_2(1/r)\rceil\), and define
\[
D_k:=\{2^{-k-1}<d_h\le2^{-k}\}
\qquad
\left(0\le k<K\right).
\]
Then
\begin{equation}
\label{eq:v081-edge-flow-degree-layer-tail}
\boxed{
\lambda_h\!\left(\left\{d_h\le2^{-K}\right\}\right)
\le2^{-K}\le r^A,
\qquad
2^{-k-1}\mu|_{D_k}
<
\lambda_h|_{D_k}
\le
2^{-k}\mu|_{D_k}.
}
\end{equation}
Thus all but an \(O(r^A)\) part of the incoming edge flow is a disjoint union of only
\(O_A(\log(1/r))\) pieces on which edge flow and selector mass are pointwise comparable.  No
single degree layer has to be selected.
\end{lemma}

\begin{proof}
The second assertion is the definition of \(D_k\).  Since \(d_h\le2^{-K}\) on the discarded set
and \(\mu(G)=1\), its \(\lambda_h\)-mass is at most \(2^{-K}\le r^A\).  The displayed layers and
this remainder form a disjoint partition up to null boundaries.
\end{proof}

\begin{proposition}[A physical Reeb bush transports every comparable incoming edge flow linearly]
\label{prop:v081-one-bush-linear-edge-flow-transfer}
Let \(H\) be an \(R\)-approximate Lagrangian point-probe bush, let \(\nu\) be a probability
selector measure supported on \(H\), and let \(\lambda\) be a finite measure satisfying
\begin{equation}
\label{eq:v081-comparable-bush-edge-flow}
\theta\,d\nu
\le d\lambda
\le2\theta\,d\nu
\end{equation}
for some \(\theta>0\).  Choose the outer Reeb interval \(W\) of length
\(|W|\asymp R\) from
\cref{cor:v081-hairbrush-coarse-product-collision}, and define the normalized physical collision
kernel
\[
g_R(z,z')
:=
\frac1{|W|}\int_W
\mathbf1_{\{|\Pi_u(z)-\Pi_u(z')|\le C_{I,J,U}R\}}\,du.
\]
After increasing the fixed geometric constant, \(g_R=1\) on \(H\times H\).  For any angular
threshold \(0<\tau<1\), split
\[
g_R=g_R^{\rm ang}+g_R^{\rm fix}
\]
by \(|a-a'|<\tau\) and \(|a-a'|\ge\tau\), and put
\[
d_{\rm ang}(z):=\int g_R^{\rm ang}(z,z')\,d\nu(z'),
\qquad
d_{\rm fix}(z):=\int g_R^{\rm fix}(z,z')\,d\nu(z').
\]
Then \(d_{\rm ang}+d_{\rm fix}=1\).  Use the measurable partition
\[
X_{\rm ang}:=\{d_{\rm ang}\ge1/2\},
\qquad
X_{\rm fix}:=H\setminus X_{\rm ang}.
\]
Thus
\begin{equation}
\label{eq:v081-one-bush-linear-edge-retention}
\boxed{
\lambda(X_{\rm ang})+\lambda(X_{\rm fix})=\lambda(H),
\qquad
d_\star(z)\ge\frac12
\quad(z\in X_\star,\ \star\in\{{\rm ang},{\rm fix}\}).
}
\end{equation}
For \(z\in X_\star\), define
\begin{equation}
\label{eq:v081-one-bush-reeb-markov-kernel}
dQ_z^\star(z')
:=
\frac{g_R^\star(z,z')}{d_\star(z)}\,d\nu(z')
\end{equation}
and use \(Q_z^{\rm ang}\) on \(X_{\rm ang}\) and \(Q_z^{\rm fix}\) on
\(X_{\rm fix}\).  These are probability kernels, with density at most \(2\) relative to \(\nu\),
supported on actual physical \(O(R)\)-collisions.  In the fixed-angle case it is supported on the normalized coarse
Maslov graph of
\cref{cor:v081-hairbrush-coarse-maslov-graph}; in the angular case it is supported on the angular
Darboux branch.  If \(d\kappa_z(\zeta)\) is any conditional old-neighbor or old-polarization law,
then
\[
\sum_{\star\in\{{\rm ang},{\rm fix}\}}
\mathbf1_{X_\star}(z)
d\lambda(z)d\kappa_z(\zeta)\,dQ_z^\star(z')
\]
has total mass \(\lambda(H)\).  Hence the Reeb collision step transports all of
the incoming edge flow without replacing it by the quadratic product mass \(\lambda(H)^2\).

More precisely, for either \(\star\) with \(\lambda(X_\star)>0\), normalize
\[
\alpha_\star:=\lambda(X_\star)^{-1}\lambda|_{X_\star}
\]
and take disjoint source and target copies with probability measure
\[
\widetilde\nu_\star
:=
\frac12\alpha_\star^{(0)}+\frac12\nu^{(1)}.
\]
Writing \(dQ_z^\star(z')=q_\star(z,z')d\nu(z')\), define the symmetric two-copy kernel
\[
\widetilde h_\star((z,0),(z',1))
=
\widetilde h_\star((z',1),(z,0))
:=
\frac12q_\star(z,z'),
\]
and set it to zero on equal copies.  Then
\begin{equation}
\label{eq:v081-one-bush-two-copy-graph}
\boxed{
0\le\widetilde h_\star\le1,
\qquad
e_{\widetilde\nu_\star}(\widetilde h_\star)=\frac14.
}
\end{equation}
Under \eqref{eq:v081-comparable-bush-edge-flow}, the source marginal
\(\alpha_\star\) is comparable, within a factor two, to the normalized restriction of \(\nu\) to
\(X_\star\).  Thus \eqref{eq:v081-one-bush-two-copy-graph} is a normalized physical Maslov graph
with fixed edge mass and no hidden cap-density factor.  In the four-cycle routing, use the target
copy for the anchors and the source copy for the free/common-neighbor vertices; the conditional old
neighbor and polarization marks therefore remain attached to the transported edge flow.
\end{proposition}

\begin{proof}
The point-probe relation places all projected bush endpoints in one \(O(R)\)-ball throughout the
interval \(W\), exactly as in the proof of
\cref{cor:v081-hairbrush-coarse-product-collision}; hence \(g_R=1\) after fixing the displayed
constant.  The angular and fixed-angle kernels partition it, so their degrees sum to one at every
source point.  On \(X_{\rm ang}\), the angular degree is at least \(1/2\) by definition.  On
\(X_{\rm fix}\), one has
\[
d_{\rm fix}=1-d_{\rm ang}>\frac12.
\]
This proves \eqref{eq:v081-one-bush-linear-edge-retention}.  Equation
\eqref{eq:v081-one-bush-reeb-markov-kernel} is therefore a probability kernel and its density is at
most two.  The support assertions follow from the two
angular pieces of the defining physical collision kernel.  Finally integrate first in
\(\kappa_z(\zeta)\) and then in the appropriate \(Q_z^\star\); both are probability laws.
The density \(q_\star\) is at most two.  Hence the two-copy kernel is at most one, and direct
integration over its two oriented cross-copy pieces gives
\[
e_{\widetilde\nu_\star}(\widetilde h_\star)
=
2\cdot\frac14\cdot\frac12
\int_{X_\star}\int_Hq_\star(z,z')\,d\nu(z')d\alpha_\star(z)
=
\frac14.
\]
Finally, \eqref{eq:v081-comparable-bush-edge-flow} implies
\[
\frac12
\frac{\nu|_{X_\star}}{\nu(X_\star)}
\le
\alpha_\star
\le
2
\frac{\nu|_{X_\star}}{\nu(X_\star)}
\]
after harmlessly weakening the exact normalization constants.  This proves the last assertion.
\end{proof}

\begin{corollary}[The two-copy Maslov routing may be oriented toward the edge-flow source]
\label{cor:v081-oriented-two-copy-maslov-routing}
For either two-copy graph
\(\widetilde h_\star\) in
\eqref{eq:v081-one-bush-two-copy-graph}, the arbitrary-threshold routing of
\cref{prop:v050-arbitrary-threshold-maslov-routing} may be performed with both anchors in the
target copy and both common-neighbor/free vertices in the source copy.  Its bush output therefore
has source marginal mass bounded below by a fixed positive constant, while its determinant-small
four-cycle output has fixed positive four-cycle mass.  In both cases the conditional old-neighbor
and polarization marks remain attached to the source vertices.
\end{corollary}

\begin{proof}
The proof of \cref{prop:v050-arbitrary-threshold-maslov-routing} is bipartite.  With
\(\alpha_\star\) on the source side, \(\nu\) on the target side, and
\[
k(z,z'):=\frac12q_\star(z,z'),
\]
one has \(0\le k\le1\) and
\[
\iint k(z,z')\,d\alpha_\star(z)d\nu(z')=\frac12.
\]
Apply the first Cauchy--Schwarz in the source variable and the second in the target variable.
The resulting four-cycle has the ordered form
\[
p^{(1)},y^{(0)},q^{(1)},(y')^{(0)},p^{(1)}.
\]
Thus the anchors \(p,q\) lie in the target copy and the common neighbors \(y,y'\) lie in the
source copy, with total four-cycle content at least \(2^{-4}\).  Splitting by the determinant
threshold gives the rank-loss output with a fixed constant.  On its complement, the Cramer-rule
and Reeb-time pigeonhole argument in the cited proposition selects its Lagrangian bush from the
common-neighbor variable, hence from the source copy, again with fixed positive mass.  All added
marks are conditional probability variables over that source endpoint and integrate to one.
\end{proof}

\begin{corollary}[One bush-return generation preserves a fixed portion of the incoming Maslov edge]
\label{cor:v081-one-generation-edge-flow-conservation}
Let the incoming normalized Maslov kernel satisfy
\[
e(h)=M=r^{\chi+o(1)},
\qquad 0\le\chi<2,
\]
and let \(H=\bigcup_jH_j\) be the maximal physical bush family of
\cref{prop:v081-maximal-bush-family} in its un-paid alternative.  Decompose the edge-flow measure
\(\lambda_h|_H\) simultaneously by the disjoint bushes \(H_j\) and the degree layers of
\cref{lem:v081-edge-flow-degree-layers}, with any fixed \(A>2\).  On every nonempty piece normalize
the selector measure and apply
\cref{prop:v081-one-bush-linear-edge-flow-transfer}.  Then, for all sufficiently small \(r\), the
sum of the transported angular or fixed-angle physical occurrence measures is at least
\begin{equation}
\label{eq:v081-one-generation-edge-flow-conservation}
\boxed{
\sum_{j,k}\|\Gamma_{j,k}^{(1)}\|
=
\lambda_h(H)-O(r^A)
\gtrsim M.
}
\end{equation}
Every \(\Gamma_{j,k}^{(1)}\) consists of actual physical Reeb collisions, retains the old Maslov
neighbor and all conditional polarization flags, and has transition density bounded by a fixed
constant.  Thus a complete bush-return generation has no power loss in incoming edge mass.
\end{corollary}

\begin{proof}
\Cref{lem:v081-incoming-edge-flow-bush-extraction} gives
\(\lambda_h(H)>M/4\).  The discarded low-degree part has mass at most \(r^A\) by
\cref{lem:v081-edge-flow-degree-layers}.  Since \(A>2>\chi\), one has \(r^A=o(M)\).  On each
remaining disjoint piece, the two measures satisfy
\eqref{eq:v081-comparable-bush-edge-flow}; sum the two branches in
\eqref{eq:v081-one-bush-linear-edge-retention}.  The support and flag assertions are those of
\cref{prop:v081-one-bush-linear-edge-flow-transfer} and
\eqref{eq:v081-incoming-edge-flow-disintegration}.
\end{proof}

\subsection{Angular and separated-time strict progress}

\begin{corollary}[Superpolynomially small bush--degree pieces carry negligible edge flow]
\label{cor:v081-polynomial-edge-flow-pieces}
In \cref{cor:v081-one-generation-edge-flow-conservation}, fix \(A>2\).  Since
\[
N\lesssim M^{-3}=r^{-3\chi+o(1)}
\]
and there are \(O_A(\log(1/r))\) degree layers, the total edge-flow mass of all pieces
\(H_j\cap D_k\) having
\[
\lambda_h(H_j\cap D_k)<r^{A+7}
\]
is at most \(r^{A+1-o(1)}\).  Every retained piece has selector mass
\[
\mu(H_j\cap D_k)\ge\lambda_h(H_j\cap D_k)\ge r^{A+7}.
\]
Thus, after a negligible relative-quadratic discard, every bush entering the same-window
amplification argument has the polynomial mass lower bound required by
\cref{cor:v081-same-window-amplification-depth}.
\end{corollary}

\begin{proof}
There are at most
\[
O_A\!\left(M^{-3}\log(1/r)\right)
\le
r^{-6-o(1)}
\]
pieces because \(\chi<2\).  Multiply this bound by \(r^{A+7}\).  The selector-mass lower bound
uses \(d_h\le1\), hence \(\lambda_h\le\mu\).
\end{proof}

\begin{lemma}[Adaptive source-cap disintegration preserves the old-neighbor occurrence]
\label{lem:v081-adaptive-source-cap-occurrence}
Keep the notation and hypotheses of
\cref{prop:v081-one-bush-linear-edge-flow-transfer}.  For any
\(0<\tau<1\), partition the direction chart by half-open cubes
\(Q\in\mathcal Q_\tau\) of side \(\tau\), and put
\[
X_{\star,Q}:=X_\star\cap\{a_z\in Q\},
\qquad
\Lambda_{\star,Q}:=\lambda(X_{\star,Q})
\quad
(\star\in\{{\rm ang},{\rm fix}\}).
\]
For every nonnull piece normalize
\[
\alpha_{\star,Q}:=\Lambda_{\star,Q}^{-1}\lambda|_{X_{\star,Q}}
\]
and retain the same Markov kernel \(Q_z^\star\) from
\eqref{eq:v081-one-bush-reeb-markov-kernel}.  Then
\begin{equation}
\label{eq:v081-adaptive-source-cap-mass}
\boxed{
\sum_{\star,Q}\Lambda_{\star,Q}=\lambda(H),
\qquad
\sum_{\star,Q}
\Lambda_{\star,Q}\,
\alpha_{\star,Q}(dz)\,\kappa_z(d\zeta)\,Q_z^\star(dz')
=
\sum_\star
\mathbf1_{X_\star}(z)\lambda(dz)\kappa_z(d\zeta)Q_z^\star(dz').
}
\end{equation}
Thus no cap is selected and no factor depending on
\(\#\mathcal Q_\tau\) is lost.  On each normalized piece the two-copy construction of
\eqref{eq:v081-one-bush-two-copy-graph} still has edge mass \(1/4\).

Moreover:
\begin{enumerate}[label=\textnormal{(\roman*)}]
\item on an angular piece, both the source direction and the Markov target direction lie in one
fixed enlargement \(Q^*\) of \(Q\), of diameter \(O(\tau)\);
\item on a fixed-angle piece, \(|a_z-a_{z'}|\ge\tau\), so the normalized physical Maslov error is
\(O(R/\tau)\), and every physical-bush output of the oriented two-copy routing has its direction
support in \(Q^*\).
\end{enumerate}
All old neighbors and all old polarization flags remain conditional probability marks over the
source copy on every cap piece.
\end{lemma}

\begin{proof}
The half-open cubes partition the source direction chart.  Hence the sets
\(X_{\star,Q}\) are pairwise disjoint and their union over \(\star,Q\) is \(H\), up to null
boundaries.  Integrating the two probability kernels \(\kappa_z\) and \(Q_z^\star\) proves
\eqref{eq:v081-adaptive-source-cap-mass}.  The proof of
\eqref{eq:v081-one-bush-two-copy-graph} uses only that \(Q_z^\star\) is a probability kernel with
density at most two; restricting and renormalizing its source marginal does not change that
calculation.

If \(\star={\rm ang}\), the definition of \(g_R^{\rm ang}\) gives
\(|a_z-a_{z'}|<\tau\).  Since \(a_z\in Q\), both directions lie in a fixed enlargement \(Q^*\).
If \(\star={\rm fix}\), the opposite angular inequality holds.  The physical collision error is
\(O(R)\), hence division by the angular scale gives \(O(R/\tau)\).  Finally
\cref{cor:v081-oriented-two-copy-maslov-routing} chooses the free/common-neighbor variables, and
therefore its bush output, from the source copy.  Those directions lie in \(Q^*\).  Conditional
marks integrate to one before and after the restriction.
\end{proof}

\begin{corollary}[Angular fine returns have unconditional occurrence-level strict progress]
\label{cor:v081-angular-old-occurrence-strict-progress}
Let a parent node be contained in a direction cap of radius \(T\).  Prescribe any
\(\tau=o(T)\).  Along the fine-failure sequence of
\cref{prop:v081-resolution-cap-freeze-or-return}, choose the next physical collision scale so far
out that the returned bush error \(R\) satisfies
\[
R=o(\tau).
\]
Then \cref{lem:v081-adaptive-source-cap-occurrence} sends the entire incoming old-neighbor
occurrence measure, with all polarization flags attached, into a disjoint sum of the following two
forms:
\begin{enumerate}[label=\textnormal{(\roman*)}]
\item angular children supported in direction caps of radius \(O(\tau)=o(T)\);
\item fixed-angle Maslov graphs with vanishing normalized error, whose physical-bush outputs are
again supported in direction caps of radius \(O(\tau)\).
\end{enumerate}
In particular an angular fine return never has to replace the incoming occurrence by the degree
measure of a newly sampled residual graph.  The second requirement formerly listed in
\cref{prob:v081-strict-progress-edge-flow-gate} is therefore satisfied.
\end{corollary}

\begin{proof}
The diagonal choice in
\eqref{eq:v081-resolution-cap-diagonal-scale} permits the new collision scale, and hence the bush
error furnished by the arbitrary-threshold Maslov routing, to be chosen below any prescribed
positive function of the fixed parent data.  Choose it so that \(R/\tau\to0\), and apply
\cref{lem:v081-adaptive-source-cap-occurrence}.  Its mass identity is linear in the absolute weights
\(\Lambda_{\star,Q}\), so keeping all cap pieces introduces neither a pigeonhole loss nor a cap-count
factor.  The two support conclusions give the asserted strict cap decrease.
\end{proof}

\begin{lemma}[The coarse self-collision of one bush is time-local]
\label{lem:v081-one-bush-self-return-time-local}
Let \(H\) be an \(R\)-approximate point-probe bush centered at Reeb time \(s_0\), and let
\(z=(a,b),z'=(a',b')\in H\) satisfy \(|a-a'|\ge\tau>0\).  If
\[
|\Pi_u(z)-\Pi_u(z')|\lesssim R
\qquad
(|u-s_0|\lesssim R),
\]
then their closest collision time and transverse residual obey
\begin{equation}
\label{eq:v081-one-bush-self-return-time-local}
\boxed{
|s_*(z,z')-s_0|
\lesssim
\frac R\tau,
\qquad
|r(z,z')|\lesssim R.
}
\end{equation}
Consequently the fixed-angle two-copy graph
\eqref{eq:v081-one-bush-two-copy-graph} need not create a new separated Reeb probe: its bush output
may be the original bush \(H\) itself.
\end{lemma}

\begin{proof}
Write \(\alpha=a-a'\) and \(\beta=b-b'\).  At \(u=s_0\), the bush relation gives
\(|\beta+s_0\alpha|\lesssim R\).  Since
\[
s_*=-\frac{\alpha\cdot\beta}{|\alpha|^2},
\qquad
\beta+s_0\alpha
=
r+(s_0-s_*)\alpha,
\qquad
r\perp\alpha,
\]
orthogonal projection onto \(\mathbb R\alpha\) gives
\(|s_*-s_0||\alpha|\lesssim R\), and projection onto \(\alpha^\perp\) gives
\(|r|\lesssim R\).  Use \(|\alpha|\ge\tau\).
\end{proof}

\begin{lemma}[A separated old Reeb time gives a lossless fiberwise two-probe cap]
\label{lem:v081-separated-old-time-fiber-cap}
Let \(H\) be an \(R\)-approximate point-probe bush centered at \((c,s_0)\).  Let
\(\mathcal G_{\rm sep}\subset H\times G\) be an old physical Maslov edge measure such that
\[
s_*(z,\zeta)\in[t_0-\delta_s,t_0+\delta_s],
\qquad
|t_0-s_0|\ge g_0>0,
\qquad
|r(z,\zeta)|\le r
\]
on its support.  For a target endpoint \(\zeta\), let \(N_\zeta\subset H\) be its source fiber.
Then there is a center \(a_\zeta^\sharp\) such that
\begin{equation}
\label{eq:v081-separated-old-time-fiber-cap}
\boxed{
N_\zeta
\subset
B\!\left(
a_\zeta^\sharp,
C_{I,U}g_0^{-1}(R+r+\delta_s)
\right)
}
\end{equation}
for almost every nonempty fiber.  Moreover the disintegration over \(\zeta\) is lossless:
\begin{equation}
\label{eq:v081-separated-old-time-fiber-mass}
\boxed{
\int
\left(
\int_{N_\zeta}h(z,\zeta)\,d\mu(z)
\right)d\mu(\zeta)
=
\iint_{\mathcal G_{\rm sep}}h(z,\zeta)\,d\mu(z)d\mu(\zeta).
}
\end{equation}
Thus separated old time supplies a genuinely new, fiberwise direction cap while retaining the
entire old-neighbor occurrence measure; no single Reeb label or target endpoint is selected.
\end{lemma}

\begin{proof}
For \(z\in N_\zeta\), the collision identity and the time-packet condition give
\[
|\Pi_{t_0}(z)-\Pi_{t_0}(\zeta)|
\lesssim_{I,U}
r+\delta_s.
\]
For \(z,z'\in N_\zeta\), subtract this relation and the two current bush relations
\[
|\Pi_{s_0}(z)-c|\lesssim R,
\qquad
|\Pi_{s_0}(z')-c|\lesssim R.
\]
Since
\[
\bigl(\Pi_{t_0}(z)-\Pi_{t_0}(z')\bigr)
-
\bigl(\Pi_{s_0}(z)-\Pi_{s_0}(z')\bigr)
=
(t_0-s_0)(a_z-a_{z'}),
\]
one obtains
\[
|a_z-a_{z'}|
\lesssim_{I,U}
g_0^{-1}(R+r+\delta_s).
\]
Choose any measurable representative of each nonempty fiber as its cap center; the usual rational
grid selection makes this choice measurable.  Equation
\eqref{eq:v081-separated-old-time-fiber-mass} is Tonelli's theorem applied to the original
edge occurrence measure.
\end{proof}

\begin{corollary}[The moving fiber caps aggregate without a continuous-label loss]
\label{cor:v081-separated-fiber-cap-aggregation}
Keep \cref{lem:v081-separated-old-time-fiber-cap} and put
\[
T:=C_{I,U}g_0^{-1}(R+r+\delta_s).
\]
Write \(a_\zeta\) for the direction coordinate of the old target endpoint.
Partition direction space by a half-open grid of side \(T\), and assign an old edge occurrence
\((z,\zeta)\) to the unique grid cell \(Q\) containing the fiber-cap center
\(a_\zeta^\sharp\).  Let \(\Gamma_Q\) be the resulting occurrence measures.  Then
\begin{equation}
\label{eq:v081-separated-fiber-cap-aggregation}
\boxed{
\sum_Q\|\Gamma_Q\|
=
\|\mathcal G_{\rm sep}\|,
\qquad
\operatorname{supp}(\pi_z)_\#\Gamma_Q
\subset B_Q^*,
}
\end{equation}
where \(B_Q^*\) is a \(C_{I,U}T\)-direction cap and the family
\(\{B_Q^*\}_Q\) has bounded overlap.  If the selector measure has total mass \(m\) and bounded
three-dimensional direction density, the part for which the old target \(\zeta\) also belongs to
the same enlarged cap obeys
\begin{equation}
\label{eq:v081-separated-fiber-internal-edge}
\boxed{
\sum_Q
\Gamma_Q\{(z,\zeta):a_\zeta\in B_Q^*\}
\lesssim_U
mT^3.
}
\end{equation}
The complementary occurrence measure is one aggregate cross-cap physical Maslov graph.  Therefore
the separated-time branch either reaches the cap-square threshold \(T^3\le mr^{2-o(1)}\), or
continues in caps of radius \(T\), with its entire incoming edge mass accounted for.
\end{corollary}

\begin{proof}
Equation \eqref{eq:v081-separated-old-time-fiber-cap} shows that if
\(a_\zeta^\sharp\in Q\), every source direction belongs to a fixed enlargement \(B_Q^*\) of \(Q\).
The half-open assignment partitions the original occurrence measure, proving
\eqref{eq:v081-separated-fiber-cap-aggregation}.  Fixed enlargements of equal grid cells have
bounded overlap.  Since \(0\le h\le1\), the internal part is bounded by
\[
\sum_Q\mu(B_Q^*)^2
\le
\left(\sup_Q\mu(B_Q^*)\right)\sum_Q\mu(B_Q^*)
\lesssim_U
T^3m.
\]
All remaining occurrences have source and old target in different caps, so their sum, not each
individual pair of caps, is the aggregate cross-cap graph.
\end{proof}

\subsection{Same-window amplification}

\begin{corollary}[Separated-window old edges satisfy strict progress]
\label{cor:v081-separated-window-strict-progress}
If the old collision-time packet and the current bush center lie in the two fixed separated Reeb
windows, then \(g_0\gtrsim_{I,J}1\).  Along every genuinely bad fine return,
\[
R+r+\delta_s=o(T_{\rm parent})
\]
after the time-packet width is chosen between the physical error and the parent cap scale.  Hence
the child radius in
\cref{cor:v081-separated-fiber-cap-aggregation} satisfies
\[
T=o(T_{\rm parent}).
\]
The separated-time part of
\cref{prob:v081-strict-progress-edge-flow-gate} is therefore closed: its full old-neighbor
occurrence mass is either internal-cap paid, aggregate-cross routed, or placed in strictly smaller
child caps.
\end{corollary}

\begin{proof}
The window separation gives the lower bound for \(g_0\).  \Cref{cor:v081-fine-return-error-negligible} supplies the physical-error separation, and the
weighted time-packet routing permits
\[
R+r\ll\delta_s\ll T_{\rm parent}
\]
on a genuine fine shell.  Substitute these estimates in the definition of \(T\) in
\cref{cor:v081-separated-fiber-cap-aggregation}.  Its mass and internal/cross conclusions are
unchanged because no time label was selected.
\end{proof}

\begin{lemma}[A same-window old-neighbor return is quadratic or amplifies the physical bush]
\label{lem:v081-same-window-edge-flow-amplification}
Let \(H\) be an \(R\)-approximate point-probe bush centered at \((c,s_0)\), let
\(q:=\mu(H)>0\), and suppose that on \(H\)
\[
\theta\le d_h(z)\le2\theta.
\]
Let \(\mathcal G_{\rm same}\subset H\times G\) be any part of the old physical Maslov graph on
which
\[
|s_*(z,\zeta)-s_0|\le\delta_s,
\qquad
|r(z,\zeta)|\le r,
\]
and write
\[
E_{\rm same}
:=
\iint_{\mathcal G_{\rm same}}h(z,\zeta)\,d\mu(z)d\mu(\zeta).
\]
Let \(H^+\) be the set of target endpoints \(\zeta\) occurring in
\(\mathcal G_{\rm same}\).  Then \(H\cup H^+\) is a
\(C_{I,U}(R+\delta_s+r)\)-approximate Lagrangian bush centered at \((c,s_0)\), and
\begin{equation}
\label{eq:v081-same-window-target-mass}
\boxed{
\mu(H^+)\ge\frac{E_{\rm same}}q.
}
\end{equation}
In particular, if \(E_{\rm same}\ge c_0\lambda_h(H)\), then for every \(B\ge1\) one has the
dichotomy
\begin{equation}
\label{eq:v081-same-window-quadratic-or-amplified}
\boxed{
\theta\le Bq
\ \Longrightarrow\
\lambda_h(H)\le2Bq^2,
\qquad
\theta>Bq
\ \Longrightarrow\
\mu(H^+)\ge c_0Bq.
}
\end{equation}
Thus a large same-window portion cannot simply repeat an unchanged small bush: it is either already
quadratic in the bush mass or enlarges that same physical bush by the factor
\(\gtrsim\theta/q\).
\end{lemma}

\begin{proof}
For \((z,\zeta)\in\mathcal G_{\rm same}\), the Pythagorean collision identity gives
\[
|\Pi_{s_*(z,\zeta)}(z)-\Pi_{s_*(z,\zeta)}(\zeta)|
=
|r(z,\zeta)|
\le r.
\]
Since the direction chart is bounded,
\[
\begin{aligned}
|\Pi_{s_0}(\zeta)-c|
&\le
|\Pi_{s_0}(\zeta)-\Pi_{s_*}(\zeta)|
+|\Pi_{s_*}(\zeta)-\Pi_{s_*}(z)|\\
&\qquad
+|\Pi_{s_*}(z)-\Pi_{s_0}(z)|
+|\Pi_{s_0}(z)-c|\\
&\lesssim_{I,U}
\delta_s+r+R.
\end{aligned}
\]
This is the asserted common point-probe bush relation.  Since \(0\le h\le1\),
\[
E_{\rm same}
\le
\mu(H)\mu(H^+)
=
q\mu(H^+),
\]
which proves \eqref{eq:v081-same-window-target-mass}.  On the degree layer,
\[
\theta q\le\lambda_h(H)\le2\theta q.
\]
If \(\theta\le Bq\), the upper bound gives the first alternative in
\eqref{eq:v081-same-window-quadratic-or-amplified}.  If \(\theta>Bq\), then
\[
\mu(H^+)
\ge
\frac{E_{\rm same}}q
\ge
c_0\frac{\lambda_h(H)}q
\ge
c_0\theta
>
c_0Bq,
\]
which proves the second.
\end{proof}

\begin{corollary}[Same-window edge flow of a disjoint bush family is amplified or square-summable]
\label{cor:v081-same-window-family-square-sum}
Let \(\{H_j\}\) be pairwise disjoint physical bushes, put \(q_j:=\mu(H_j)\), and decompose each
\(H_j\) into the degree layers \(D_{j,k}\) of
\cref{lem:v081-edge-flow-degree-layers}.  On a layer for which the same-window part carries at least
one half of \(\lambda_h(D_{j,k})\), apply
\cref{lem:v081-same-window-edge-flow-amplification} with \(c_0=1/2\).  For every \(B\ge1\), either
an amplified same-center bush is produced, or the total old edge flow of all non-amplified
same-window layers satisfies
\begin{equation}
\label{eq:v081-same-window-family-square-sum}
\boxed{
\mathcal M_{\rm same,quad}
\lesssim
B\sum_{j,k}\mu(D_{j,k})^2
\le
B\sum_jq_j^2.
}
\end{equation}
Layers on which the same-window part carries less than one half send a fixed fraction of their
edge flow to the separated-time branch and hence to
\cref{cor:v081-separated-window-strict-progress}.  Thus the entire old-neighbor occurrence is
partitioned into separated-time strict progress, physical bush amplification, and the single
square-sum remainder in \eqref{eq:v081-same-window-family-square-sum}.
\end{corollary}

\begin{proof}
On a non-amplified same-window layer,
\eqref{eq:v081-same-window-quadratic-or-amplified} gives
\[
\lambda_h(D_{j,k})
\lesssim
B\mu(D_{j,k})^2.
\]
The degree layers are disjoint inside each \(H_j\), so
\[
\sum_k\mu(D_{j,k})^2
\le
\left(\sum_k\mu(D_{j,k})\right)^2
\le q_j^2.
\]
Sum in \(j\).  The last assertion is the complementary half of the same/separated time split and
uses the lossless occurrence disintegration
\eqref{eq:v081-separated-old-time-fiber-mass}.
\end{proof}

\begin{corollary}[The same-window square sum is paid at the cap threshold]
\label{cor:v081-same-window-cap-square-payment}
If, in \cref{cor:v081-same-window-family-square-sum}, every \(H_j\) is contained in a direction cap
of radius \(T\), then
\[
\mathcal M_{\rm same,quad}
\lesssim_U
BmT^3,
\qquad
m:=\sum_jq_j.
\]
Consequently this entire same-window remainder is
\(\lesssim m^2r^{2-o(1)}\) whenever
\[
T^3\le B^{-1}mr^{2-o(1)}.
\]
\end{corollary}

\begin{proof}
\Cref{lem:v081-cap-square-edge-tail} gives the cap-square estimate
\(\sum_jq_j^2\lesssim_UmT^3\).  Insert it into the preceding same-window square sum.
\end{proof}

\begin{corollary}[A same-window amplification chain has sublogarithmic depth]
\label{cor:v081-same-window-amplification-depth}
Let \(r\downarrow0\), put
\[
B(r):=\exp\!\sqrt{\log(1/r)}=r^{-o(1)},
\]
and suppose the initial bush mass satisfies \(q_0\ge r^{C_0}\) for one fixed \(C_0<\infty\).
Iterate only the amplified alternative of
\eqref{eq:v081-same-window-quadratic-or-amplified}, with a fixed retained fraction
\(c_0>0\).  Then either a quadratic alternative
\[
\lambda_h(H_\ell)\lesssim B(r)q_\ell^2
\]
occurs, or after at most
\begin{equation}
\label{eq:v081-same-window-amplification-depth}
\boxed{
L_{\rm amp}
\le
\frac{\log(1/q_0)+O(1)}
{\log(c_0B(r))}
=
O_{C_0,c_0}\!\left(\sqrt{\log(1/r)}\right)
=
o(\log(1/r))
}
\end{equation}
steps the physical bush has direction mass bounded below by a fixed positive constant.
If the bush errors tend to zero, the latter alternative gives
\(\dim_HK_b=4\).  Hence, under the contrary dimension hypothesis, every polynomial-mass
same-window chain reaches its quadratic branch in sublogarithmic depth and incurs only subpower
fixed-constant loss.
\end{corollary}

\begin{proof}
At every amplified step,
\[
q_{\ell+1}\ge c_0B(r)q_\ell.
\]
Since \(q_\ell\le1\), iteration gives
\[
(c_0B(r))^\ell q_0\le1
\]
until a fixed positive mass is reached.  Taking logarithms and using
\(\log B(r)=\sqrt{\log(1/r)}\) proves
\eqref{eq:v081-same-window-amplification-depth}.

For a bush of direction mass \(q_\ell\ge q_*>0\), its normalized direction measure has density at
most \(C_U/q_*\).  It therefore satisfies the uniform direction-Frostman estimate of
\cref{thm:v081-bush-frostman-escape} for every \(0<\delta<3\).  Along vanishing bush errors that
theorem gives \((4-\delta)\)-Frostman measures on \(K_b\) for every \(\delta>0\), hence
\(\dim_HK_b=4\).
\end{proof}

\subsection{Common-target Markovization}

\begin{lemma}[Amplified same-window targets either grow or form a weighted common-target star]
\label{lem:v081-amplified-target-overlap-star}
Let \(\{S_i\}\) be pairwise disjoint source pieces of selector masses
\(q_i:=\mu(S_i)>0\), and let
\(\Gamma_i\le\mu|_{S_i}\otimes\mu\) be same-window old-neighbor occurrence measures.  Write
\[
e_i(\zeta)
:=
\frac{d(\pi_\zeta)_\#\Gamma_i}{d\mu}(\zeta),
\qquad
w_i(\zeta):=\frac{e_i(\zeta)}{q_i},
\qquad
n(\zeta):=\sum_iw_i(\zeta).
\]
Then \(0\le w_i\le1\).  Suppose the amplified alternative of
\cref{lem:v081-same-window-edge-flow-amplification} holds on every \(S_i\), so that
\begin{equation}
\label{eq:v081-amplified-target-weight}
\int w_i\,d\mu
=
\frac{\|\Gamma_i\|}{q_i}
\ge c_0Bq_i.
\end{equation}
Put \(Q:=\sum_iq_i\) and \(H^+:=\bigcup_i\{w_i>0\}\).  For all sufficiently large \(B\), at least
one of the following two quantitative conclusions is available:
\begin{enumerate}[label=\textnormal{(\roman*)}]
\item the target union genuinely amplifies the available selector mass,
\begin{equation}
\label{eq:v081-amplified-target-union-growth}
\mu(H^+)\ge\frac{c_0}{2}B^{1/2}Q;
\end{equation}
\item the off-diagonal weighted common-target content satisfies
\begin{equation}
\label{eq:v081-amplified-target-common-star}
\boxed{
\sum_{i\ne j}\int w_i(\zeta)w_j(\zeta)\,d\mu(\zeta)
\gtrsim_{c_0}
B^{3/2}Q.
}
\end{equation}
\end{enumerate}
The second quantity is an actual common-old-target occurrence, not an overlap of unweighted
supports.  More precisely, where \(e_i(\zeta)>0\), define
\[
d\Gamma_i(z,\zeta)
=
e_i(\zeta)\,dK_{i,\zeta}(z)\,d\mu(\zeta).
\]
Then \(K_{i,\zeta}\) is a probability kernel and every summand in
\eqref{eq:v081-amplified-target-common-star} lifts losslessly to
\begin{equation}
\label{eq:v081-amplified-target-two-star-lift}
w_i(\zeta)w_j(\zeta)\,d\mu(\zeta)\,
dK_{i,\zeta}(z_i)\,dK_{j,\zeta}(z_j).
\end{equation}
All old polarization flags over \(z_i,z_j\) may be tensored into this lift.

If \(S_i\) is an \(R_i\)-bush centered at \((c_i,s_i)\), and the \(i\)-th occurrence has collision
time within \(\delta_i\) of \(s_i\) and residual at most \(r\), then every common target
\(\zeta=(a_\zeta,b_\zeta)\) in \eqref{eq:v081-amplified-target-two-star-lift} satisfies
\begin{equation}
\label{eq:v081-amplified-target-carrier-line}
\left|
c_i-\bigl(b_\zeta+s_i a_\zeta\bigr)
\right|
\lesssim_{I,U}
R_i+\delta_i+r.
\end{equation}
Thus the large-overlap alternative is precisely a weighted Reeb carrier-line star through actual
selector line \(\zeta\).
\end{lemma}

\begin{proof}
Because \(\Gamma_i\le\mu|_{S_i}\otimes\mu\), one has
\[
0\le e_i(\zeta)\le\mu(S_i)=q_i,
\]
and hence \(0\le w_i\le1\).  Summing
\eqref{eq:v081-amplified-target-weight} gives
\[
W:=\int n\,d\mu\ge c_0BQ.
\]
If \eqref{eq:v081-amplified-target-union-growth} fails, Cauchy--Schwarz on \(H^+\) gives
\[
\int n^2\,d\mu
\ge
\frac{W^2}{\mu(H^+)}
>
2B^{1/2}W.
\]
On the other hand \(w_i^2\le w_i\), so
\[
\sum_i\int w_i^2\,d\mu
\le W.
\]
Subtracting this diagonal term from \(\int n^2\,d\mu\), and then using
\(W\ge c_0BQ\), proves the asserted common-star bound.  The definition of \(K_{i,\zeta}\) is the
disintegration of the actual \(i\)-th occurrence; integrating the two probability kernels in
the displayed two-star lift returns its weight.  Conditional flag
laws again integrate to one.

Finally, the same-window estimate in the proof of
\cref{lem:v081-same-window-edge-flow-amplification} applies to every
\((z_i,\zeta)\in\operatorname{supp}\Gamma_i\) and gives
\[
|\Pi_{s_i}(\zeta)-c_i|
\lesssim_{I,U}
R_i+\delta_i+r.
\]
Since \(\Pi_{s_i}(\zeta)=b_\zeta+s_i a_\zeta\), this is
\eqref{eq:v081-amplified-target-carrier-line}.
\end{proof}

\begin{corollary}[High common-target multiplicity has a linear different-source Markovization]
\label{cor:v081-common-target-linear-markovization}
In \cref{lem:v081-amplified-target-overlap-star}, restrict to one dyadic source-mass layer
\[
q\le q_i<2q.
\]
Let
\[
e(\zeta):=\sum_i e_i(\zeta),
\qquad
Y:=\{\zeta:n(\zeta)\ge B^{1/2}\}.
\]
If the union-growth alternative
\eqref{eq:v081-amplified-target-union-growth} fails, then
\begin{equation}
\label{eq:v081-common-target-high-edge-mass}
\boxed{
\sum_i\Gamma_i(S_i\times Y)
\ge
\frac14\sum_i\|\Gamma_i\|.
}
\end{equation}
For \(\zeta\in Y\), first disintegrate the incoming \(i\)-th edge by \(K_{i,\zeta}\), and then choose
a second source index and endpoint by
\begin{equation}
\label{eq:v081-common-target-second-source-kernel}
\mathsf P_\zeta(j,dz_j)
:=
\frac{e_j(\zeta)}{e(\zeta)}\,dK_{j,\zeta}(z_j).
\end{equation}
This is a probability kernel and, conditionally on an incoming edge from \(S_i\),
\begin{equation}
\label{eq:v081-common-target-different-source-probability}
\mathsf P_\zeta\{j=i\}
\le
\frac{2}{n(\zeta)}
\le
2B^{-1/2}.
\end{equation}
Consequently, after discarding the event \(j=i\), the measure
\begin{equation}
\label{eq:v081-common-target-linear-two-star}
\boxed{
\sum_i
\mathbf1_Y(\zeta)\,
d\Gamma_i(z_i,\zeta)
\sum_{j\ne i}
\frac{e_j(\zeta)}{e(\zeta)}
dK_{j,\zeta}(z_j)
}
\end{equation}
has mass at least a fixed positive fraction of
\(\sum_i\|\Gamma_i\|\), for all sufficiently large \(B\).  Thus the common-target star is linear in
the incoming old-neighbor occurrence mass.  Both source endpoints retain independent conditional
old-polarization flags, and both of their bush centers satisfy
\eqref{eq:v081-amplified-target-carrier-line} for the same actual selector line \(\zeta\).
\end{corollary}

\begin{proof}
On the dyadic layer,
\[
qn(\zeta)\le e(\zeta)<2qn(\zeta).
\]
Failure of \eqref{eq:v081-amplified-target-union-growth}, together with
\(W=\int n\,d\mu\ge c_0BQ\), gives
\[
\int_{\{n<B^{1/2}\}}n\,d\mu
\le
B^{1/2}\mu(H^+)
<
\frac{c_0}{2}BQ
\le
\frac W2.
\]
Hence \(\int_Yn\,d\mu\ge W/2\).  Since
\[
\sum_i\|\Gamma_i\|
=
\int e\,d\mu
<
2qW,
\qquad
\int_Ye\,d\mu
\ge
q\int_Yn\,d\mu
\ge
\frac{qW}{2},
\]
one obtains \eqref{eq:v081-common-target-high-edge-mass}.  Equation
\eqref{eq:v081-common-target-second-source-kernel} has total mass one.  Moreover
\[
\frac{e_i(\zeta)}{e(\zeta)}
\le
\frac{q_i}{qn(\zeta)}
<
\frac2{n(\zeta)},
\]
which proves \eqref{eq:v081-common-target-different-source-probability}.  Integrate the two
probability kernels and combine the \(1/4\) retained high-target fraction with
\(1-2B^{-1/2}\).  Tensoring conditional flag laws changes none of these masses.
\end{proof}

\begin{corollary}[A fixed number of distinct source flags can be sampled linearly]
\label{cor:v081-common-target-four-source-star}
Keep \cref{cor:v081-common-target-linear-markovization}.  For every fixed integer \(k\ge2\), extend
each incoming edge over \(Y\) by \(k-1\) conditionally independent samples from
\(\mathsf P_\zeta\).  The probability that the original source index and the \(k-1\) new source
indices are not all distinct is at most
\begin{equation}
\label{eq:v081-common-target-distinct-loss}
\binom{k}{2}\,2B^{-1/2}.
\end{equation}
After restricting to distinct indices, the resulting \(k\)-source star therefore retains a fixed
fraction of the original incoming old-neighbor occurrence mass.  In particular \(k=4\) gives, over
one actual selector target \(\zeta\), four distinct physical source bushes together with four actual
old collision times and four independently retained old polarization flags.  This four-flag star
is obtained by Markov extension of the original occurrence; it incurs no fourth power of its edge
mass.
\end{corollary}

\begin{proof}
At every \(\zeta\in Y\), each source-index atom of
\(\mathsf P_\zeta\) is at most \(2B^{-1/2}\) by the different-source probability bound.  For each of the
\(\binom{k}{2}\) pairs of samples, condition on the first member of the pair and use this maximal
atom bound for the second.  The union bound proves the asserted distinct-source loss.  Integrating
the \(k-1\) probability kernels leaves the
incoming occurrence mass unchanged before the collision event is discarded.  Old flag laws are
additional conditional probability kernels and hence are retained as well.
\end{proof}

\subsection{Flagged common-target routing}

\begin{lemma}[Every nondegenerate Maslov bush may retain a lossless origin flag]
\label{lem:v081-origin-marked-bush}
In the determinant-nondegenerate branch of
\cref{prop:v050-arbitrary-threshold-maslov-routing}, write the collision thickness as
\(\epsilon_{\rm col}\), the horizontal determinant threshold as \(\Delta\), and
\[
R_{\rm org}:=C_{U,J}\epsilon_{\rm col}/\Delta.
\]
The physical bush may be chosen so that every one of its source endpoints \(z\) carries a
conditional origin witness
\[
\mathfrak w=(p,q,z^\dagger,s,V,\mathbf q),
\qquad
V\in\Gr(3,6),
\qquad
\mathbf q=(\alpha,\beta)\in S^5,
\]
with the following properties:
\begin{equation}
\label{eq:v081-origin-marked-bush}
\boxed{
\dist(\mathbf q,V)=0,
\qquad
|\beta+s\alpha|
\lesssim
\epsilon_{\rm col}/\tau,
\qquad
|s-\bar s|\lesssim R_{\rm org}.
}
\end{equation}
Here \(\tau\) is the retained fixed-angle lower bound, \(\bar s\) is the bush-center Reeb time,
\((p,z,q,z^\dagger,p)\) is an actual determinant-nondegenerate physical Maslov four-cycle, and
\[
V
:=
\Span\{
\mathbf z_z-\mathbf z_p,\,
\mathbf z_q-\mathbf z_p,\,
\mathbf z_{z^\dagger}-\mathbf z_p
\},
\qquad
\mathbf q
:=
\frac{\mathbf z_z-\mathbf z_p}{|\mathbf z_z-\mathbf z_p|},
\]
with \(\mathbf z_x:=(a_x,b_x)\).  The origin-witness law is a probability kernel over \(z\);
therefore attaching it changes neither the physical bush measure nor any later restriction of that
measure.
\end{lemma}

\begin{proof}
Use the notation in the proof of
\cref{thm:v046-common-neighbor-dichotomy,prop:v050-arbitrary-threshold-maslov-routing}.  For the
selected anchors \(p,q\), let
\[
d\lambda_{p,q}(z)=h(p,z)h(q,z)\,d\nu(z),
\qquad
K(p,q)=\|\lambda_{p,q}\|,
\]
and let \(\mathcal G_{p,q}\) be the determinant-nondegenerate set of ordered common-neighbor pairs
\((z,z^\dagger)\).  The chosen anchor pair satisfies
\[
\frac{(\lambda_{p,q}\otimes\lambda_{p,q})(\mathcal G_{p,q})}{K(p,q)}
\gtrsim M^3.
\]
Partition the selected edge label \(s_{pz}\) into half-open intervals \(I\) of length comparable to
\(R_{\rm org}\), and let \(A_I\) be the set of \(z\) in the \(I\)-packet which have positive
\(\mathcal G_{p,q}\)-degree.  Cramer's rule gives
\[
|s_{pz}-s_{pz^\dagger}|\lesssim R_{\rm org}
\qquad
((z,z^\dagger)\in\mathcal G_{p,q}),
\]
so a good pair meets only \(O(1)\) neighboring time packets.  If \(I^*\) denotes the union of those
neighbors, then
\[
\begin{aligned}
(\lambda_{p,q}\otimes\lambda_{p,q})(\mathcal G_{p,q})
&=
\sum_I
\int_{A_I}\int
\mathbf1_{\mathcal G_{p,q}}(z,z^\dagger)\,
d\lambda_{p,q}(z^\dagger)d\lambda_{p,q}(z)\\
&\le
\sum_I\lambda_{p,q}(A_I)\lambda_{p,q}(I^*)\\
&\lesssim
K(p,q)\sup_I\lambda_{p,q}(A_I).
\end{aligned}
\]
Choose \(I\) attaining the displayed lower bound and call its center \(\bar s\).  Then
\(\lambda_{p,q}(A_I)\gtrsim M^3\); since \(h\le1\), also
\(\nu(A_I)\gtrsim M^3\).  Every \(z\in A_I\) has a good partner \(z^\dagger\).  Disintegrating the
restricted good-pair measure in \(z^\dagger\), and making an arbitrary measurable choice on a null
set, gives the asserted probability kernel of witnesses.  Thus \(H:=A_I\) has exactly the same
mass lower bound as the usual bush output and every point of \(H\) is origin-marked.

The horizontal determinant condition implies that the three displayed phase-space differences
are linearly independent, so \(V\in\Gr(3,6)\), and the definition gives
\(\mathbf q\in V\).  It also gives \(|a_z-a_p|\gtrsim_U\Delta\), so one may take
\(\tau\gtrsim_U\Delta\); on the already separated fixed-angle branch use its sharper retained
\(\tau\).  Hence
\(|\mathbf z_z-\mathbf z_p|\gtrsim\tau\).  Dividing the actual contact-edge estimate
\[
|(b_z-b_p)+s_{pz}(a_z-a_p)|\lesssim\epsilon_{\rm col}
\]
by this norm proves the second assertion in
\eqref{eq:v081-origin-marked-bush}.  Finally \(s=s_{pz}\in I\), so
\(|s-\bar s|\lesssim R_{\rm org}\), and the contact edge \(pz\) gives
\[
|\Pi_{\bar s}(z)-\Pi_{\bar s}(p)|
\le
|\Pi_s(z)-\Pi_s(p)|
+|s-\bar s|\,|a_z-a_p|
\lesssim_{U,J}
R_{\rm org}.
\]
Thus \(A_I\) is the asserted point-probe bush centered at
\((\Pi_{\bar s}(p),\bar s)\).  Integrating the conditional partner law first proves the
physical-marginal assertion.
\end{proof}

\begin{lemma}[An origin-marked bush plane stays quantitatively outside the front-null locus]
\label{lem:v081-origin-plane-away-front-null}
For the origin plane \(V\) in
\cref{lem:v081-origin-marked-bush}, suppose the phase-space coordinates remain in the fixed bounded
chart and the generating horizontal determinant is at least \(\Delta\).  Then
\begin{equation}
\label{eq:v081-origin-plane-away-front-null}
\boxed{
\dist_{\Gr}(V,\mathfrak N)
\gtrsim_U
\Delta,
\qquad
\mathfrak N=\{W\in\Gr(3,6):P_W\equiv0\}.
}
\end{equation}
\end{lemma}

\begin{proof}
Let \(D\) be the \(6\times3\) matrix whose columns are the three phase-space differences defining
\(V\), and let \(D_a\) be its horizontal \(3\times3\) block.  The horizontal nondegeneracy gives
\[
|\det D_a|\ge\Delta.
\]
If \(Q=D(D^{\mathsf T}D)^{-1/2}\) is an orthonormal frame for \(V\), then its horizontal frame
\(Q_a\) satisfies
\[
|\det Q_a|
=
\frac{|\det D_a|}{\sqrt{\det(D^{\mathsf T}D)}}
\gtrsim_U
\Delta,
\]
because the three phase-space differences are uniformly bounded.  The distance of \(Q_a\) from
the rank-at-most-two matrices is its least singular value, which is at least
\(|\det Q_a|\) since the other two singular values are at most one.  Every
\(W\in\mathfrak N\) has \(\rank(\pi_a|_W)\le2\) by
\eqref{eq:v081-four-time-front-null}.  The Grassmannian distance controls the distance between
orthonormal horizontal frame matrices up to a fixed factor, proving
\eqref{eq:v081-origin-plane-away-front-null}.
\end{proof}

\begin{corollary}[The linear four-source star has the carrier-rooted polarized lift]
\label{prob:v081-common-target-carrier-rooted-lift}
Let \(\Gamma^{(4)}\) be the linear four-source common-target occurrence of
\cref{cor:v081-common-target-four-source-star}.  Retain on each source bush the origin flag of
\cref{lem:v081-origin-marked-bush}.  Choose two Reeb resolutions
\[
\max_i\left\{
R_{{\rm org},i},
\epsilon_{{\rm col},i}/\tau_i
\right\}
\ll
\delta_{\rm edge}
\ll
\delta_s
\ll
\tau.
\]
On the same-window part
\[
|s_*(z_i,\zeta)-\bar s_i|\le\delta_{\rm edge},
\]
the fourfold Markov extension carries actual triples
\[
(\mathbf q_i,V_i,s_i),
\qquad
\mathbf q_i=(\alpha_i,\beta_i),
\qquad
0\le i\le3,
\]
such that
\begin{equation}
\label{eq:v081-carrier-rooted-polarized-lift}
\boxed{
\dist(\mathbf q_i,V_i)=0,
\qquad
|\beta_i+s_i\alpha_i|=o(\delta_s),
\qquad
|s_i-s_*(z_i,\zeta)|=o(\delta_s).
}
\end{equation}
The lift preserves the physical four-source occurrence exactly.  Hence the four Reeb roots used by
\cref{prop:v081-four-time-polarization-firewall} are the origin-cycle roots of the actual source
bushes and are synchronized with the common selector carrier; no independently sampled target
flag is substituted for them.
\end{corollary}

\begin{proof}
In the maximal-bush construction
\cref{prop:v081-maximal-bush-family}, use the origin-marked refinement
\cref{lem:v081-origin-marked-bush} whenever the arbitrary-threshold routing produces a physical
bush.  The proof of maximality is unchanged because the refined bush has the same
\(\gtrsim M^3\) selector-mass lower bound.  Restrictions to maximal bushes, degree layers, adaptive
source caps, and same-window old edges preserve the mark: each is a physical restriction, while
the later Reeb and common-target extensions are probability kernels.  The same applies to every
descendant after its next fine residual failure produces a fresh physical Maslov bush.

For the \(i\)-th marked source,
\eqref{eq:v081-origin-marked-bush} gives
\[
\dist(\mathbf q_i,V_i)=0,
\qquad
|\beta_i+s_i\alpha_i|
\lesssim
\epsilon_{{\rm col},i}/\tau_i,
\qquad
|s_i-\bar s_i|
\lesssim
R_{{\rm org},i}.
\]
The defining same-window restriction gives
\[
|s_*(z_i,\zeta)-\bar s_i|\le\delta_{\rm edge}.
\]
The triangle inequality and the displayed hierarchy prove
\eqref{eq:v081-carrier-rooted-polarized-lift}.  Finally, tensoring the four conditional
origin-witness kernels with the Markov extension integrates to one in the new variables, so its
physical marginal is exactly \(\Gamma^{(4)}\).
\end{proof}

\begin{lemma}[Separated origin polarizations are routed at the source endpoint]
\label{lem:v081-common-target-source-polarization-cycle}
Keep one dyadic source-mass layer of
\cref{cor:v081-common-target-linear-markovization}, retain the origin marks from
\cref{prob:v081-common-target-carrier-rooted-lift}, and also lift every source endpoint \(z\) by its
actual conditional old flags.  Suppose a positive portion of the retained common-target occurrence
has two source origin planes lying in Grassmannian cells
\(\mathcal V_1,\mathcal V_2\subset\Gr(3,6)\) with
\[
\dist(\mathcal V_1,\mathcal V_2)\ge\vartheta,
\qquad
\vartheta^{-1}=r^{-o(1)}.
\]
Then that portion admits an occurrence-preserving sourcewise routing into:
\begin{enumerate}[label=\textnormal{(\roman*)}]
\item transverse old/new polarization, close-anchor, horizontal-line, or coherent rank-loss
outputs already assigned to the paid ledgers;
\item angular fine-return or physical Lagrangian-bush children retaining their inherited source
caps and all common-target marks.
\end{enumerate}
No flag is moved from a source endpoint to the common target, and no power of the incoming
occurrence mass is lost.
\end{lemma}

\begin{proof}
Disintegrate the four-source extension over its distinguished source occurrence.  By symmetry
among the four source samples, on a fixed fraction of the separated-plane event the distinguished
source carries one member of a separated pair.  Keep its origin plane, the other three sources, and
the common target as conditional probability marks.  The resulting physical source marginal is therefore a
submeasure of the original incoming occurrence; all such marginals add to at most that occurrence.

On every source bush and degree layer apply
\cref{prop:v081-one-bush-linear-edge-flow-transfer}.  Its pointwise angular/fixed split preserves
the whole restricted source occurrence, including the three extra source marks and the common
target.  On the fixed part the symmetric two-copy graph has edge mass \(1/4\) by
\eqref{eq:v081-one-bush-two-copy-graph}.  Apply
\cref{cor:v081-oriented-two-copy-maslov-routing,cor:v081-k-flag-arbitrary-threshold} with the source
copy as the decorated free side.  The old flag has therefore remained at the same source endpoint
as the current cycle plane.  A transverse plane is paid by
\cref{lem:v078-old-flag-new-cycle-polarization,%
prop:v081-same-vertex-polarization-variation-paid}; the close-anchor/angular, horizontal-line, and
coherent rank-loss outputs have their established destinations.  The remaining output is the
physical source-side Lagrangian bush of
\cref{prop:v081-rank-loss-to-bush-reduction}.  It retains the inherited cap and all conditional
common-target marks.  The angular part is sent to the smaller-cap child of
\cref{cor:v081-angular-old-occurrence-strict-progress}.

Every operation above is either a restriction of the distinguished source occurrence or an
extension by a probability kernel.  Summing the source bushes, degree layers, and polarization
cells is therefore linear.  The fixed edge mass \(1/4\) is taken only after normalizing each
restricted source piece, so no power of its absolute mass is introduced.
\end{proof}

\begin{corollary}[The linear four-source star reaches the polarization/time trichotomy]
\label{cor:v081-common-target-polarized-four-star}
On the four-source occurrence from
\cref{cor:v081-common-target-four-source-star}, use the carrier-rooted lift of
\cref{prob:v081-common-target-carrier-rooted-lift}, and fix a polarization resolution
\(\vartheta\downarrow0\) and a Reeb-time resolution \(\delta_s\downarrow0\).  The origin cycles use
the same determinant threshold \(\Delta\) at one routed generation; equivalently first retain its
dyadic layer.  Fix
\[
\eta_{\rm rk}<c_U\Delta/2
\]
with the constant from
\cref{lem:v081-origin-plane-away-front-null}, and choose the next fine failure scale so that the
normalized incidence error and the \(O(\vartheta)\)-coherence diameter are below
\[
\varepsilon_0(\delta_s,\eta_{\rm rk},I^\sharp)
\]
from \cref{cor:v081-finite-scale-four-time-firewall}.  Up to a subpower polarization-cell decomposition, every
four-source occurrence belongs to one of:
\begin{enumerate}[label=\textnormal{(\roman*)}]
\item two of its origin Maslov planes \(V_i\) are \(\vartheta\)-separated;
\item all four origin planes are \(O(\vartheta)\)-coherent and their four origin Reeb times contain four
pairwise \(\delta_s\)-separated values;
\item all four origin planes are \(O(\vartheta)\)-coherent and their origin Reeb times lie in at most
three intervals of length \(O(\delta_s)\).
\end{enumerate}
By \cref{lem:v081-common-target-source-polarization-cycle}, the first alternative is routed at its
actual source endpoints into paid outputs or source-side physical bush children, with no power loss
in the incoming occurrence.  In the second alternative the four origin-cycle Maslov incidences,
synchronized with the same selector target, give the Vandermonde input of
\cref{prop:v081-four-time-polarization-firewall,cor:v081-finite-scale-four-time-firewall}.  That
corollary places every \(V_i\) within \(\eta_{\rm rk}\) of \(\mathfrak N\), contradicting
\cref{lem:v081-origin-plane-away-front-null}; hence this branch is empty.  In the third
alternative, at least two source bushes
belong to the same shrinking Reeb-time packet and, by
\eqref{eq:v081-amplified-target-carrier-line}, merge at one physical point of the common selector
line up to error \(O(R+\delta_{\rm edge}+\delta_s+r)\).  Consequently the final gate is reduced to the lossless
occurrence-level aggregation of the merged packets in \textnormal{(iii)}; both the transverse
branch and the four-separated-time coherent branch are already routed.
\end{corollary}

\begin{proof}
Cover the compact polarization space by a bounded-overlap
\(\vartheta\)-net.  The number of cells is subpower for the usual choice
\(\vartheta^{-1}=r^{-o(1)}\), and all cells are retained in the absolute occurrence sum.  Either
two retained cell labels are separated or all four lie in one fixed enlargement of a single cell.
In the latter case, four real numbers are either pairwise \(\delta_s\)-separated or are covered by
at most three intervals of length \(2\delta_s\).  This proves the trichotomy without selecting one
label cell or one time packet.

For \textnormal{(i)}, apply
\cref{lem:v081-common-target-source-polarization-cycle}.  For \textnormal{(ii)}, apply the
finite-scale four-time firewall and then
\cref{lem:v081-origin-plane-away-front-null}, using the parameter hierarchy in the statement.  In
the last alternative choose a time interval containing two
source labels.  Their common target \(\zeta\) satisfies
\eqref{eq:v081-amplified-target-carrier-line} for both centers; replacing their two times by the
interval representative changes the front point by \(O_U(\delta_s)\).  The union of their
disjoint source subsets is therefore an
\(O(R+\delta_s+r)\)-point-probe bush at that representative point.
\end{proof}

\begin{lemma}[A four-sample Reeb test gives a lossless three-packet stopping]
\label{lem:v081-four-sample-three-packet-stopping}
Let
\[
d\Gamma(z,\zeta)
=
e(\zeta)\,dP_\zeta(z)\,d\mu(\zeta)
\]
be a finite common-target occurrence measure, where \(P_\zeta\) is a probability kernel and every
source occurrence has a measurable Reeb label \(s(z,\zeta)\in I^\sharp\).  Fix
\(\delta_s>0\).  At each target put
\[
\mathfrak q(\zeta)
:=
P_\zeta^{\otimes4}
\left\{
(z_0,z_1,z_2,z_3):
|s(z_i,\zeta)-s(z_j,\zeta)|>\delta_s
\text{ for all }i\ne j
\right\}.
\]
There is a measurable stopping decomposition
\begin{equation}
\label{eq:v081-three-packet-stopping-decomposition}
\boxed{
\Gamma
=
\Gamma_{\rm four}
+
\sum_{\ell\ge0}\sum_{m=1}^3\Gamma_{\ell,m}
}
\end{equation}
up to a null limiting tail, with the following properties.
\begin{enumerate}[label=\textnormal{(\roman*)}]
\item Over every part assigned to \(\Gamma_{\rm four}\), the four-fold Markov extension contains a
pairwise \(\delta_s\)-separated four-time tuple with conditional probability at least \(1/8\).
\item For every \((\ell,m)\), there is a measurable interval center
\(t_{\ell,m}(\zeta)\) such that
\[
|s(z,\zeta)-t_{\ell,m}(\zeta)|\le2\delta_s
\qquad
\text{on }\operatorname{supp}\Gamma_{\ell,m}.
\]
\item After \(L\) stopping rounds, the as-yet unassigned occurrence mass is at most
\[
2^{-L}\|\Gamma\|.
\]
\end{enumerate}
Thus a common-target occurrence is exhausted linearly by the four-separated-time output and at
most three coherent Reeb packets per stopping round.  No target and no time packet is selected.
\end{lemma}

\begin{proof}
We first use the following elementary probability fact.  If \(\nu\) is a probability measure on a
real interval and
\[
\nu^{\otimes4}\{ |s_i-s_j|>\delta_s\ (i\ne j)\}<\frac18,
\]
then three intervals of radius \(\delta_s\) cover at least one half of \(\nu\).  Indeed, if no three
such intervals do so, then for every choice of at most three centers the complement of their union
has mass greater than \(1/2\).  Draw four independent samples successively.  After the first
sample, each of the next three has conditional probability greater than \(1/2\) of lying outside
the \(\delta_s\)-neighborhoods of all previous samples.  The probability of four pairwise separated
samples is therefore greater than \(2^{-3}=1/8\), a contradiction.

Apply this fact to the pushforward of \(P_\zeta\) by \(s(\,\cdot\,,\zeta)\).  On
\(\{\mathfrak q\ge1/8\}\), put the current occurrence into \(\Gamma_{\rm four}\).  On the
complement, first take the three radius-\(\delta_s\) intervals supplied by the probability fact.
Approximating their centers by a fixed countable rational grid enlarges them to radius
\(2\delta_s\).  Choose the lexicographically first such rational triple whose union has conditional
mass at least \(1/2\), assign each covered source to the first interval containing its time, and
remove the three resulting packet measures.  This choice is measurable.  Repeat the
construction with the conditional remainder.  Each low-four-time round removes at least half of
the remaining occurrence mass, while every high-four-time part is assigned to
\(\Gamma_{\rm four}\).  The remainder after \(L\) rounds is therefore at most \(2^{-L}\|\Gamma\|\);
continuity from above gives a null limiting tail and
\eqref{eq:v081-three-packet-stopping-decomposition}.
\end{proof}

\begin{lemma}[Target cells turn synchronized Reeb packets into a fractional physical bush family]
\label{lem:v081-target-cell-three-packet-bushes}
Let one of the packet measures \(\Gamma_{\ell,m}\) from
\cref{lem:v081-four-sample-three-packet-stopping} be supported on physical old edges satisfying
\[
|\Pi_u(z)-\Pi_u(\zeta)|\le \rho,
\qquad
|s-u|\le\varepsilon_{\rm sync},
\qquad
|s-t_{\ell,m}(\zeta)|\le2\delta_s.
\]
Partition the bounded selector line chart of the target \(\zeta=(a_\zeta,b_\zeta)\) into half-open
cells \(C\) of diameter \(\eta\), and quantize \(t_{\ell,m}(\zeta)\) into half-open intervals of
length \(\delta_s\).  Assign every occurrence to its unique pair \((C,J)\), and let
\(\Gamma_{\ell,m,C,J}\) be the resulting restrictions.  Then
\begin{equation}
\label{eq:v081-target-cell-packet-mass}
\boxed{
\sum_{C,J}\|\Gamma_{\ell,m,C,J}\|
=
\|\Gamma_{\ell,m}\|.
}
\end{equation}
For representatives \(\zeta_C\in C\) and \(t_J\in J\), the source marginal of every nonnull piece
is supported on the physical point-probe bush
\begin{equation}
\label{eq:v081-target-cell-packet-bush}
\boxed{
\left|
\Pi_{t_J}(z)-\Pi_{t_J}(\zeta_C)
\right|
\lesssim_{I,U}
\rho+\varepsilon_{\rm sync}+\delta_s+\eta.
}
\end{equation}
If the incoming source was already restricted by
\cref{lem:v081-adaptive-source-cap-occurrence} to a direction cap of radius \(\tau\), every such
bush remains in that same \(\tau\)-cap.  Moreover the sum of all source marginals in
\eqref{eq:v081-target-cell-packet-mass} is exactly the source marginal of
\(\Gamma_{\ell,m}\), even when their physical supports overlap.  Hence this is a fractional
occurrence partition, not an unweighted family requiring a choice of one target cell.
\end{lemma}

\begin{proof}
The half-open target cells and time intervals give a measurable partition, proving
\eqref{eq:v081-target-cell-packet-mass}.  For an occurrence assigned to \((C,J)\), write
\(s=s(z,\zeta)\) and let \(u\) be its actual old-edge collision time.  Then
\[
\begin{aligned}
|\Pi_{t_J}(z)-\Pi_{t_J}(\zeta_C)|
&\le
|\Pi_u(z)-\Pi_u(\zeta)|
+|s-u|\,|a_z-a_\zeta|\\
+|t_J-s|\,|a_z-a_\zeta|\\
&\qquad
+|\Pi_{t_J}(\zeta)-\Pi_{t_J}(\zeta_C)|\\
&\lesssim_{I,U}
\rho+\varepsilon_{\rm sync}+\delta_s+\eta.
\end{aligned}
\]
Here the direction chart and the relevant Reeb interval are bounded, and the target cell controls
both \(a_\zeta\) and \(b_\zeta\).  This proves
\eqref{eq:v081-target-cell-packet-bush}.  Direction-cap support is inherited from the source
restriction.  Finally, pushforward commutes with the disjoint sum of the occurrence restrictions,
which proves the last assertion.
\end{proof}

\begin{corollary}[The same-window selector-carrier gate closes at occurrence level]
\label{cor:v081-same-window-carrier-reduction}
Let \(r_0\) be the fixed root collision scale in the desired edge-Carleson estimate, and let
\(\rho\) be the later fine failure scale which produces the current physical bush.  Put
\[
B_0:=B(r_0)=\exp\sqrt{\log(1/r_0)}.
\]
\Cref{lem:v081-same-window-edge-flow-amplification} is used here with its old-edge time-packet
width renamed \(\delta_{\rm edge}\); the larger symbol \(\delta_s\) is reserved for the
four-sample Maslov/Reeb stopping.  \Cref{prob:v081-common-target-carrier-rooted-lift} supplies a synchronization error
\[
\varepsilon_{\rm sync}
:=
\max_i|s_i-s_*(z_i,\zeta)|
=o(\delta_s).
\]
Apply the adaptive source-cap disintegration at a radius \(\tau=o(T_{\rm parent})\), with the physical collision error chosen
\(o(\tau)\).  For every polynomial-mass same-window source family,
its entire incoming
old-neighbor occurrence, up to an \(O(r_0^A)\) tail for arbitrary fixed \(A\), is a disjoint absolute
sum of:
\begin{enumerate}[label=\textnormal{(\roman*)}]
\item quadratic cap pieces of total mass
\[
\lesssim_U B_0m\tau^3;
\]
\item transverse-polarization, excluded four-separated-time, angular, horizontal-line, or
front-nondegenerate pieces already assigned to the paid ledgers;
\item physical point-probe bush children of error
\[
O(\rho+\varepsilon_{\rm sync}+\delta_s+\eta)=o(\tau)
\]
whose source directions remain in \(O(\tau)\)-caps and whose occurrence masses add linearly to the
unpaid input.
\end{enumerate}
If
\[
\tau^3\le B_0^{-1}mr_0^{2-o(1)},
\]
the first part is at the relative quadratic target.  Hence every unpaid same-window occurrence
makes strict physical/cap progress.
\end{corollary}

\begin{proof}
In every local lemma used below, relabel its collision thickness \(r\) as the present fine scale
\(\rho\).  First choose \(\tau\) from the root-scale cap-square target, and then use
\cref{prop:v081-resolution-cap-freeze-or-return} to choose \(\rho\) so far out on the failure
sequence that all physical and synchronization errors are \(o(\tau)\).

The non-amplified part is
\cref{cor:v081-same-window-family-square-sum,%
cor:v081-same-window-cap-square-payment}, because
\cref{lem:v081-adaptive-source-cap-occurrence} places every oriented source bush in an
\(O(\tau)\)-cap.  This gives \textnormal{(i)}.

For the amplified part apply
\cref{lem:v081-amplified-target-overlap-star}.  In its union-growth alternative, reverse the old
edges and use their target marginal.  Since the source pieces are disjoint,
\[
\sum_i e_i(\zeta)\le\sum_iq_i\le1,
\]
so the reversed edge-flow marginal is still dominated by \(\mu\), and its total occurrence mass is
unchanged.  Degree-layering therefore restarts the same physical-bush/Frostman alternative without
a power loss.  If union growth repeats \(L\) times from initial selector mass \(q_0\), then
\[
\left(cB_0^{1/2}\right)^Lq_0\le1.
\]
For \(q_0\ge r_0^{C_0}\), taking logarithms gives
\[
L
\lesssim_{C_0}
\frac{\log(1/r_0)}{\log B_0}
=
O_{C_0}\!\left(\sqrt{\log(1/r_0)}\right).
\]
A fixed positive-mass vanishing-error bush family gives the vector-Frostman output by
\cref{lem:v081-mass-conserving-vector-frostman-stopping,%
thm:v081-bush-family-frostman-escape} and is recorded in the paid ledger.  Otherwise this branch
terminates in sublogarithmic depth.

In the common-target alternative, first dyadically layer the source masses; the number of layers is
subpower after
\cref{cor:v081-polynomial-edge-flow-pieces}.  \Cref{cor:v081-common-target-linear-markovization} retains a fixed fraction of the incoming edge
mass in a different-source Markov star, and
\cref{cor:v081-common-target-four-source-star} supplies four distinct source flags without taking
an edge-mass power.  Separated source polarizations are routed at their actual source endpoints by
\cref{lem:v081-common-target-source-polarization-cycle}.  On the coherent part, the
carrier-rooted coupling
\cref{prob:v081-common-target-carrier-rooted-lift} synchronizes the old-edge times with the Maslov
planes, and
\cref{cor:v081-common-target-polarized-four-star} excludes the four-separated-time output: the
finite-scale Reeb-front firewall would place every origin plane near the front-null locus, contrary
to its determinant-nondegenerate origin cycle.

Apply
\cref{lem:v081-four-sample-three-packet-stopping} to the remaining conditional Reeb-time law.
Choose \(L_{\rm pkt}=\lceil A\log_2(1/r_0)\rceil\).  Its unassigned tail is at most
\(r_0^A\) times the incoming occurrence, while all removed mass is contained in three time packets
per stopping round.  This internal stopping has exact mass addition and introduces no factor
\(C^{L_{\rm pkt}}\).  Finally apply
\cref{lem:v081-target-cell-three-packet-bushes} with
\[
\rho+\varepsilon_{\rm sync}+\delta_s+\eta=o(\tau).
\]
Every retained packet becomes the physical child in \textnormal{(iii)}, remains in its inherited
\(O(\tau)\)-source cap, and all target cells and all three packets are summed rather than selected.
This exhausts the common-target alternative and proves the corollary.
\end{proof}

\subsection{Cross-generation conservation and deterministic saturation}

\begin{proposition}[Cross-generation edge-flow conservation]
\label{thm:v081-cross-generation-edge-flow-conservation}
In the bush-return tree of
\cref{prop:v081-bush-tree-edge-carleson-criterion}, retain the incoming edge-flow disintegration
\eqref{eq:v081-incoming-edge-flow-disintegration} at every nonterminal node.  Route every extracted
physical bush by the two pointwise branches and the two-copy graph of
\cref{prop:v081-one-bush-linear-edge-flow-transfer}.  Assign transverse polarization, horizontal,
front-nondegenerate, angular-good, and vector-Frostman outputs to
\(\mathcal M_{\rm paid}\); assign every occurrence whose old neighbor leaves the current child cap
to \(\mathcal M_{\rm cross}\).  The strict progress needed for this routing is supplied by
\cref{cor:v081-separated-window-strict-progress} on separated old-time windows,
\cref{cor:v081-angular-old-occurrence-strict-progress} on angular fine returns, and
\cref{cor:v081-same-window-carrier-reduction} on the same-window selector-carrier branch.  Then,
for every fixed \(A>2\) and every finite depth \(L\), the routing may be exhausted so that
\begin{equation}
\label{eq:v081-finite-depth-edge-flow-conservation}
\boxed{
\mathcal M_{\rm in}(r)
\le
\mathcal M_{\rm paid}(r)
+
\mathcal M_{\rm cross}(r)
+
\sum_{\gamma\in\mathcal S_L}\mathcal M_{\gamma}^{\rm same}(r)
+
O_A(Lr^A)\mathcal M_{\rm in}(r),
}
\end{equation}
where \(\mathcal M_{\gamma}^{\rm same}\) is the terminal old-edge flow whose two endpoints lie in
the cap \(\gamma\).  If the terminal source marginal is
\(d\nu_\gamma=f_\gamma d\sigma\), has mass \(p_\gamma\), and is supported in a cap of radius
\(T_\gamma\), then
\begin{equation}
\label{eq:v081-terminal-markov-cap-bound}
\boxed{
\mathcal M_{\gamma}^{\rm same}(r)
\lesssim_U p_\gamma T_\gamma^3,
\qquad
\sum_{\gamma\in\mathcal S_L}f_\gamma\le f_{\rm root}.
}
\end{equation}
Consequently \eqref{eq:v081-bush-tree-edge-carleson} holds for every
\(L=r^{-o(1)}\); no factor \(C^L\) occurs.
\end{proposition}

\begin{proof}
Represent the node by its ordered old-edge occurrence measure
\[
d\Gamma_v(z,\zeta)
:=h_v(z,\zeta)\,d\mu_v(z)d\mu_v(\zeta)
=d\lambda_{h_v}(z)d\kappa_z^{h_v}(\zeta).
\]
The first coordinate is now the source once and for all; no later orientation factor is needed.
On every nonzero unused restriction of \(\Gamma_v\), maximal bush extraction and
\cref{lem:v081-incoming-edge-flow-bush-extraction} assign a uniform positive fraction either to a
paid output or to disjoint physical bushes.  \Cref{lem:v081-hereditary-markov-exhaustion} therefore exhausts \(\Gamma_v\) by countably many
such pieces.  Their source marginals add to \(\lambda_{h_v}\), and the old endpoint \(\zeta\)
remains the conditional mark \(\kappa_z^{h_v}\) throughout.

On every bush piece use the degree partition of
\cref{lem:v081-edge-flow-degree-layers}.  The normalized low-degree remainder costs at most \(r^A\)
times the incoming absolute node weight.  Every other layer satisfies
\eqref{eq:v081-comparable-bush-edge-flow}.  \Cref{cor:v081-one-generation-edge-flow-conservation} and, more precisely, the pointwise
angular/fixed partition in
\cref{prop:v081-one-bush-linear-edge-flow-transfer}, transport its entire remaining occurrence
mass through actual \(O(R)\)-Reeb collisions.  No unweighted square of its selector mass is taken.

Before routing either source branch, apply the adaptive source-cap disintegration
\cref{lem:v081-adaptive-source-cap-occurrence}, with a cap radius prescribed below the parent
radius and the bush error chosen still smaller.  The cap pieces retain the whole absolute
occurrence mass and all old marks.  For a fixed-angle source piece, the symmetric two-copy graph
\eqref{eq:v081-one-bush-two-copy-graph} has edge mass \(1/4\).  Apply
\cref{cor:v081-oriented-two-copy-maslov-routing,cor:v081-k-flag-arbitrary-threshold} with the target
copy as the anchor side and the source copy as the free/common-neighbor side.  Its close-anchor,
line, and transverse-polarization outputs are paid.  Split the retained old-neighbor time relative
to the current bush center.  The separated-time part enters a genuinely smaller alternating-window
child by \cref{cor:v081-separated-window-strict-progress}.  The same-window part is partitioned,
with its absolute occurrence mass unchanged, into paid outputs or genuinely smaller-cap physical
children by \cref{cor:v081-same-window-carrier-reduction}; this includes the amplified
common-target branch.  The old neighbor, every old polarization, and the origin-cycle marks remain
attached by
\cref{lem:v081-lossless-conditional-flag-tensor,lem:v081-origin-marked-bush,%
prob:v081-common-target-carrier-rooted-lift}.  Every fixed-proportion output of this routing is
exhausted on the unused source occurrence by
\cref{lem:v081-hereditary-markov-exhaustion}; the two-copy edge mass remains \(1/4\) on each
remainder.  The coherent rank-loss alternative is excluded by the same finite-flag argument
because the graph edge mass is fixed.  For an angular source
piece, \cref{cor:v081-angular-old-occurrence-strict-progress} places both copies directly in a
strictly smaller cap; at a further failure scale it repeats the same lossless source-cap
disintegration instead of sampling an unrelated fine residual graph.  Consequently every
nonterminal node has the absolute identity
\[
\mathfrak m_v
=
\mathfrak m_{v,{\rm paid}}
+
\mathfrak m_{v,{\rm cross}}
+
\sum_{w\in{\rm ch}(v)}\mathfrak m_w
+\mathfrak e_v,
\qquad
0\le\mathfrak e_v\lesssim_A r^A\mathfrak m_v,
\]
where \(\mathfrak m_v=\|\Gamma_v\|\).  The child source marginals are fractional restrictions of
the parent marginal, and their Radon--Nikodym densities add pointwise to at most the parent density.

Sum this identity over all nodes above depth \(L\).  At each depth the node occurrences are
submeasures of a partition of the root occurrence, so
\[
\sum_{|v|=k}\mathfrak m_v\le\mathcal M_{\rm in}(r),
\qquad
\sum_{|v|<L}\mathfrak e_v
\lesssim_A Lr^A\mathcal M_{\rm in}(r).
\]
All intermediate child terms telescope with coefficient one.  This proves the finite-depth
identity and removes the former \(C^L\).

At a terminal cap \(\gamma\), split the old-neighbor mark according as \(\zeta\) lies inside or
outside that cap.  The outside part is \(\mathcal M_{\rm cross}\).  For the inside part, use the
root edge density \(h\le1\), the source density \(f_\gamma\), and bounded Lebesgue direction
density:
\[
\mathcal M_\gamma^{\rm same}
\le
\int f_\gamma(z)\mathbf1_\gamma(z)
\left(\int\mathbf1_\gamma(\zeta)\,d\sigma(\zeta)\right)d\sigma(z)
\lesssim_U p_\gamma T_\gamma^3.
\]
This proves \eqref{eq:v081-terminal-markov-cap-bound}.  If \(L=r^{-o(1)}\), choose \(A>2\)
larger than the fixed terminal exponent reserve.  Then \(Lr^A=r^{A-o(1)}\) is absorbed, and the
result is \eqref{eq:v081-bush-tree-edge-carleson}.
\end{proof}

\begin{proposition}[Lossless deterministic saturation of the return tree]
\label{prop:v081-lossless-cap-saturation}
Fix a coarse occurrence layer of source mass \(m>0\), a root collision thickness \(r\), and put
\[
B(r):=\exp\!\sqrt{\log(1/r)}.
\]
At a node in a cap of radius \(T_v\), prescribe
\[
\tau_v:=\frac{T_v}{C_*B(r)},
\]
where \(C_*\) dominates the fixed cap enlargement in
\cref{lem:v081-adaptive-source-cap-occurrence}.  Every non-paid child can be chosen with
\begin{equation}
\label{eq:v081-deterministic-cap-contraction}
\boxed{T_w\le B(r)^{-1}T_v\qquad(w\in{\rm ch}(v)).}
\end{equation}
The choices preserve the entire absolute occurrence measure and all inherited marks.  The least
depth at which \(T_L^3\le mr^{2-o(1)}\) satisfies
\begin{equation}
\label{eq:v081-deterministic-saturation-depth}
\boxed{L(r)=O_m\!\left(\sqrt{\log(1/r)}\right)=o(\log(1/r)).}
\end{equation}
The assertion is uniform in the number and shape of the children and in every carrier label.
\end{proposition}

\begin{proof}
Since \(B(r)\to\infty\), one has \(\tau_v=o(T_v)\).  On a positive-mass node which is not a good
residual leaf, \cref{prop:v081-resolution-cap-freeze-or-return} supplies arbitrarily fine failure
scales.  Choose the node's fine scale only after \(\tau_v\) is fixed, so that every physical,
synchronization, polarization, and target-cell error is at most
\(T_v/(C_*^2B(r))\).  The angular, separated-window, and same-window support statements in
\cref{cor:v081-angular-old-occurrence-strict-progress,cor:v081-separated-window-strict-progress,%
cor:v081-same-window-carrier-reduction} then give
\eqref{eq:v081-deterministic-cap-contraction}.

The fine scales are auxiliary certificates attached to individual occurrence nodes; they are not
identified with one common generation scale.  The tree is countable because it is built from
half-open cap, target-cell, degree, and rational time-packet partitions, so all nodewise choices can
be made.  More importantly, all pieces are retained.  The identities in
\cref{lem:v081-hereditary-markov-exhaustion,lem:v081-adaptive-source-cap-occurrence,%
lem:v081-four-sample-three-packet-stopping,lem:v081-target-cell-three-packet-bushes} show that no
minimum node mass enters the construction.  Each cap normalization is subpower at its own chosen
fine scale by \eqref{eq:v081-resolution-cap-diagonal-scale}; such constants are not multiplied
across nodes or generations because
\cref{thm:v081-cross-generation-edge-flow-conservation} uses absolute occurrence measures.

Finally \(T_L^3\le T_0^3B(r)^{-3L}\).  Since \(m\) and \(T_0\) belong to the fixed coarse layer,
the least admissible integer satisfies
\[
L\le
1+\frac{\log(T_0^3/m)+(2+o(1))\log(1/r)}{3\log B(r)}
=O_m\!\left(\sqrt{\log(1/r)}\right).
\]
Carrier labels are passive conditional marks in every cited mass identity, proving uniformity.
\end{proof}

\begin{corollary}[Maslov--bush--Frostman bridge]
\label{cor:v081-maslov-bush-frostman-bridge-closed}
For every fixed coarse low-reuse occurrence layer \(G_*\), either the relative residual estimate
\eqref{eq:v078-final-low-reuse-relative-residual} holds uniformly at all sufficiently fine scales,
or \(K_b\) carries a nonzero
\((4-\delta)\)-Frostman measure for every \(0<\delta<3\).  In particular, under the contrary
hypothesis \(\dim_HK_b<4\), the relative residual estimate holds.  Thus the physical bush output
cannot remain as an unassigned subquadratic Maslov residual.
\end{corollary}

\begin{proof}
If a positive-mass good leaf occurs, use
\cref{thm:v081-bush-family-frostman-escape} or the good alternative of
\cref{lem:v081-mass-conserving-vector-frostman-stopping}.  Otherwise every positive-mass node has
arbitrarily fine residual-failure scales.  \Cref{prop:v081-lossless-cap-saturation} supplies a simultaneous mass-preserving saturation with
\(L=O(\sqrt{\log(1/r)})\) and \(T_L^3\le mr^{2-o(1)}\).  Unlike the former diagonal choice, this
uses the predetermined contraction \(B(r)^{-1}\) and the exact absolute-flow identities, so neither
the retained mass nor a minimum node mass enters the depth estimate.
\Cref{thm:v081-cross-generation-edge-flow-conservation} supplies
\eqref{eq:v081-bush-tree-edge-carleson}.  Applying the edge--Carleson criterion once more now
gives the relative residual estimate and closes the stated bridge.
\end{proof}

\begin{theorem}[Contact--symplectic closure of the four-dimensional Sticky Kakeya selector]
\label{thm:v081-sticky-kakeya-closure}
For every compact full-direction Sticky selector in the contactification used in this manuscript,
the Euclidean Reeb front \(K_b\subset\mathbb R^4\) satisfies
\[
\boxed{\dim_HK_b=4.}
\]
\end{theorem}

\begin{proof}
Assume \(\dim_HK_b<4\).  The finite Darboux, angular-sector, Reeb-window, and occurrence-layer
decompositions lose only fixed or subpower factors.  Fix one coarse occurrence layer and suppose
that its relative residual target fails along arbitrarily fine scales.
\Cref{prop:v078-relative-residual-maslov-return} converts that failure, with the quadratic layer
mass exactly normalized out, into a normalized Maslov graph.  The old-flag/new-cycle comparison
and finite conditional flag amplification route this graph into the transverse-polarization,
horizontal-line, front-nondegenerate, angular, or physical Lagrangian-bush outputs; see
\cref{cor:v078-final-two-generation-routing,%
cor:v081-k-flag-arbitrary-threshold,%
prop:v081-rank-loss-to-bush-reduction}.  The first four outputs obey the required quadratic
relative estimates in their assigned ledgers and therefore cannot carry the assumed power-excess
failure.  For the last output,
\cref{cor:v081-maslov-bush-frostman-bridge-closed} gives either that same relative estimate or a
nonzero \((4-\delta)\)-Frostman measure for every \(0<\delta<3\); the latter contradicts
\(\dim_HK_b<4\).  Hence the relative residual target holds on every fixed coarse layer.

Summing the finite/coarse layers and absorbing all fixed coarse constants at the fine scale as in
\cref{lem:v078-coarse-fine-mass-freezing,cor:v078-fine-scale-absorbs-coarse-tax} gives, for every
\(\eta>0\),
\[
\mathfrak Z_\rho^{J^+}(\mu)
\lesssim_\eta
\rho^{2-\eta}
\qquad(0<\rho<1).
\]
The codimension-two residual criterion
\cref{thm:v048-residual-frostman-criterion} then yields
\(\dim_HK_b=4\), a contradiction.
\end{proof}

\subsection{The marked source-hereditary interface}

\begin{corollary}[The polarized selector-carrier strict-progress gate is closed]
\label{prob:v081-strict-progress-edge-flow-gate}
Let a returned bush \(H\), centered at \((c,s_0)\), carry the conserved occurrence measure
\[
d\lambda(z)d\kappa_z^h(\zeta)
\]
from \eqref{eq:v081-incoming-edge-flow-disintegration}, and let its physical coarse Reeb collision
be represented by the Markov kernels
\eqref{eq:v081-one-bush-reeb-markov-kernel}.  For a fixed root scale \(r_0\), choose the current
fine failure scale \(\rho\) after the child-cap target.  Its entire incoming occurrence is
partitioned, up to an arbitrarily prescribed \(O(r_0^A)\) tail, into paid outputs and physical children lying in
direction caps \(o(T_{\rm parent})\).  In particular, the strict-progress input used in
\cref{thm:v081-cross-generation-edge-flow-conservation} holds with no additional hypothesis.
\end{corollary}

\begin{proof}
The separated old-time part has strict progress by
\cref{lem:v081-separated-old-time-fiber-cap,%
cor:v081-separated-fiber-cap-aggregation,%
cor:v081-separated-window-strict-progress}.  The angular all-fine part also has occurrence-level
strict progress by
\cref{lem:v081-adaptive-source-cap-occurrence,%
cor:v081-angular-old-occurrence-strict-progress}: it uses the coarse Reeb Markov transition on
every adaptive source cap and never substitutes the degree measure of a newly sampled fine graph.

For the same-window part apply \cref{cor:v081-same-window-carrier-reduction}.  Its non-amplified
pieces are paid by the cap-square ledger.  In its amplified common-target branch,
\cref{cor:v081-common-target-linear-markovization,cor:v081-common-target-four-source-star} preserve
a fixed fraction of the original occurrence linearly.  Separated planes are routed sourcewise by
\cref{lem:v081-common-target-source-polarization-cycle}.  On the coherent part,
\cref{prob:v081-common-target-carrier-rooted-lift} makes the old-edge Reeb times and the Maslov
plane incidences simultaneous;
four-separated-time flags are then impossible by the finite-scale Reeb-front firewall and the
quantitative separation of every determinant-nondegenerate origin plane from the front-null locus,
as recorded in \cref{cor:v081-common-target-polarized-four-star}.  Finally,
\cref{lem:v081-four-sample-three-packet-stopping} partitions the remaining conditional time law
exactly into three-packet pieces with \(O(r_0^A)\) tail, and
\cref{lem:v081-target-cell-three-packet-bushes} turns every such piece into a fractional physical
bush family in \(O(\tau)\)-source caps.  Choosing
\(\rho+\varepsilon_{\rm sync}+\delta_s+\eta=o(\tau)=o(T_{\rm parent})\) gives strict cap progress.  All cap, target-cell, and
time-packet pieces are summed rather than selected, so the absolute incoming occurrence is
conserved.
\end{proof}

\begin{lemma}[Mass-conserving vector-Frostman stopping for a maximal bush family]
\label{lem:v081-mass-conserving-vector-frostman-stopping}
Let
\[
\{H_j:1\le j\le N\}
\]
be a finite pairwise disjoint family of \(\varepsilon\)-bushes.  Let \(\lambda\) be a finite measure
supported on their disjoint union and assume, for some fixed \(C_\lambda<\infty\),
\[
\lambda|_{H_j}\le C_\lambda\sigma|_{H_j},
\qquad
q:=\|\lambda\|>0.
\]
Use the common outer window \(I_{\rm out}\) and the vector centers
\(\mathfrak a_j(y,s)\) from
\cref{thm:v081-bush-family-frostman-escape}.  Fix \(0<\delta<3\), a lower radius
\(\rho\ge\varepsilon\), and \(C_0\ge1\).  There is a pairwise disjoint measurable decomposition of
a subset carrying at least \(q/2\) of the family mass into exactly one of the following forms.
\begin{enumerate}[label=\textnormal{(\roman*)}]
\item A restriction \(\lambda_{\rm good}\le\lambda\), of mass at least \(q/2\), whose normalization
satisfies
\[
\sup_{\substack{y\in\mathbb R^3\\s\in I_{\rm out}}}
\frac1{\|\lambda_{\rm good}\|}
\sum_j
\lambda_{\rm good}\!\left(
H_j\cap B(\mathfrak a_j(y,s),R)
\right)
\le
2C_0R^{3-\delta}
\qquad(\rho\le R\le1).
\]
\item Countably many disjoint Reeb-coherent hairbrush packets
\(\mathcal H_\ell\), with masses \(S_\ell\), witness points
\((y_\ell,s_\ell)\), and radii \(\rho\le R_\ell\le1\), such that
\begin{equation}
\label{eq:v081-mass-conserving-hairbrush-total}
\boxed{
\sum_\ell S_\ell\ge\frac q2,
\qquad
S_\ell
>
C_0q_\ell R_\ell^{3-\delta}.
}
\end{equation}
Here \(q_\ell\) is the \(\lambda\)-mass remaining immediately before \(\mathcal H_\ell\) is removed.
Every \(\mathcal H_\ell\) merges into one physical \(O(R_\ell)\)-bush centered at
\((y_\ell,s_\ell)\), its underlying selector-direction mass is at least
\(C_\lambda^{-1}S_\ell\), and all old conditional Maslov flags remain attached.
\end{enumerate}
Thus failure of the vector-Frostman output need not be represented by one arbitrarily small
hairbrush: a fixed fraction of the maximal-family mass is retained across the full stopping
collection.
\end{lemma}

\begin{proof}
Start with \(\lambda_0=\lambda\), of mass \(q_0=q\).  Given
\(\lambda_\ell\), stop with alternative \textnormal{(i)} if its normalization satisfies the
displayed vector-Frostman inequality with constant \(2C_0\).  Otherwise choose a triple
\((y_\ell,s_\ell,R_\ell)\) whose violating cap mass is at least one half of the supremal violation,
and remove the corresponding physical selector subsets from every \(H_j\).  Denote their union by
\(\mathcal H_\ell\), its mass by \(S_\ell\), and the remainder by
\(\lambda_{\ell+1}\).  Then
\[
S_\ell>C_0q_\ell R_\ell^{3-\delta}.
\]
The removed selector sets are disjoint by construction.

Continue while the removed mass is below \(q/2\).  If the procedure stops, the remaining mass is
at least \(q/2\) and gives \textnormal{(i)}.  If it does not stop and the accumulated removed mass
reaches \(q/2\), retain the first countable initial segment which does so and obtain
\textnormal{(ii)}.  The only other formal possibility is that the remainders decrease to a measure
\(\lambda_\infty\) of mass greater than \(q/2\) while the removed masses have sum below \(q/2\).
Continuity from above gives convergence of every fixed cap mass.  A strict vector-Frostman
violation for \(\lambda_\infty\) would then, after a small enlargement of its ball and rational
approximation of \((y,s,R)\), give a uniformly positive near-supremal violation for all sufficiently
large remainders, contradicting convergence of the removed masses to zero.  Hence
\(\lambda_\infty\) satisfies the bound and gives \textnormal{(i)}.

For every removed packet, all cap centers are generated by its single witness
\((y_\ell,s_\ell)\).  The calculation in
\cref{prop:v081-hairbrush-merges-to-one-bush} therefore merges it into one physical bush.
The domination by \(C_\lambda\sigma\) gives the asserted lower bound for its underlying direction
mass.  In alternative \textnormal{(i)}, the proof of
\cref{thm:v081-bush-family-frostman-escape} applies with the component restrictions of
\(\lambda_{\rm good}\) in place of \(\sigma|_{H_j}\): push forward
\(\lambda_{\rm good}\otimes ds\), and use the displayed vector Frostman bound in the same
\(O(R)\)-long Reeb window.  Thus this weighted good alternative has the same
\((4-\delta)\)-Frostman conclusion.
Restrictions preserve the conditional flag laws by
\cref{lem:v081-lossless-conditional-flag-tensor}.
\end{proof}

\begin{principle}[Finite-scale routing of ideal reuse into the flagged contact-star trichotomy]
\label{pr:v081-hairbrush-reuse-routing}
At the coarse scale \(R_n\),
\cref{prop:v081-hairbrush-merges-to-one-bush,%
cor:v081-hairbrush-coarse-product-collision} already give a genuine physical bush and full product
collision mass.  \Cref{cor:v081-hairbrush-coarse-maslov-graph} therefore gives, on its
fixed-angle part, a normalized thickened Maslov graph of constant edge mass; no
ideal-to-physical conversion or cell-density factor is needed there.  The issue below concerns only
a strictly finer Maslov scale.
\Cref{prop:v081-hairbrush-intrinsic-renormalization} sends that constant-edge graph
through the arbitrary-threshold route and returns, unless a paid ledger occurs, an \(O(1)\)-member
bush family carrying a fixed fraction of the normalized hairbrush mass.  Thus a single
intrinsic-scale Maslov return is lossless.  \Cref{cor:v081-hairbrush-alternating-two-probe-lock} then uses the required alternation of the two
separated Reeb windows to show that a second merged hairbrush lies in one shrinking direction cap.
Hence fixed-angle hairbrush renormalization does not iterate indefinitely; its only return is the
angular Darboux chart at a scale comparable to the inherited physical error.
\Cref{prop:v081-resolution-cap-freeze-or-return} freezes that actual positive-mass
cap.  It is either already a good all-fine-scales residual leaf, or every cap-normalization
constant becomes subpower along a new, much finer failure sequence and the same flagged
contact--Maslov route restarts without a power tax.

The measure in \eqref{eq:v081-hairbrush-energy-as-reuse} is an ideal Monge-type coupling: for each
source direction \(a\) and target label \(k\), it uses the single transported direction
\(\Psi_{jk,u}(a)\).  \Cref{cor:v081-hairbrush-energy-flag-lift} shows that its actual
selector endpoints retain their old conditional Maslov flags.  \Cref{lem:v081-hairbrush-ideal-product-limit} shows that, under uniform differentiation-scale
separation, the lossless benchmark would be to find
\[
C\varepsilon_n\le\widetilde r_n=o(R_n)
\]
such that
\begin{equation}
\label{eq:v081-hairbrush-coupling-to-product-target}
\boxed{
\begin{aligned}
&\sum_{j\ne k}\int_{W_n}
\iint_{A_{n,j}\times A_{n,k}}
\mathbf 1_{\{|\Pi_u(z(a))-\Pi_u(z(a'))|\le\widetilde r_n\}}
\,d\sigma(a)d\sigma(a')du\\
&\hspace{30mm}\gtrsim_{I,J,U}
\widetilde r_n^3\mathcal R_n.
\end{aligned}
}
\end{equation}
Here \(\widetilde r_n^3\) is the transverse three-dimensional thickening of the Monge graph.

Uniform differentiation is not assumed.  \Cref{prop:v081-hairbrush-finite-cell-transfer} proves at every
\(\varepsilon_n\le r\le R_n\) the unconditional replacement: either physical product collision
retains \(\omega r^3\mathcal R_n\), reverse Reeb contact-star reuse is high, or the complete energy
is bounded by \(BN_n\omega R_n^4\).  \Cref{cor:v081-hairbrush-balanced-cell-transfer} chooses \(B,\omega\) so that the third case is
absent whenever \(S_n\gg R_n^3\).  Thus the unknown differentiation threshold has been eliminated.
On the fixed-angle portion, \Cref{prop:v081-hairbrush-fine-maslov-graph} divides out only the genuine
\(O_{\tau_0}(r)\) collision-time window and returns the physical product branch to a normalized
\(O(r)\)-Maslov graph.  Its edge mass is
\[
\widetilde M_{n,r}
\gtrsim
\omega_n(r/R_n)^2,
\]
and the old conditional polarizations remain attached without a second normalization loss.
The complementary small-angle portion is the Darboux branch.
\Cref{cor:v081-hairbrush-two-scale-maslov-frontier} compares this edge mass with the
next merged-cap density tax and records the exact remaining two-scale quotient; no relation between
the two scales is inserted by hand.
The remaining scalar in the product branch is exactly
\[
\omega_n^{-1}
\lesssim_{I,J,U,G_*}
A_n^{-1}M_n^{-4}R_n^\delta,
\]
while the complementary branch is an actual high-reuse contact star.  After the original
occurrence weights are restored, these are precisely the two inputs to
\cref{prop:v074-continuous-high-reuse-trichotomy}; its outputs are angular collapse, a tetrahedral
Reeb-front seed, or one contact-polarized \(E\oplus E\) packet.  The last output is handled by the
old/new flag comparison.  No label count is identified with direction mass before this weighted
regularization.
\end{principle}

\begin{corollary}[Dyadic label reuse has an exact clustered-star mass threshold]
\label{cor:v081-hairbrush-dyadic-star-threshold}
Restrict the target labels of one reuse star from
\eqref{eq:v081-hairbrush-high-reuse-threshold} to a dyadic cap-mass layer
\[
p\le\sigma(A_{n,k})<2p,
\qquad
|s_{n,j_0}-s_{n,k}|\ge\delta_s
\qquad(k\in\mathcal N),
\]
where the near-equal center-level pairs have already been routed separately, and suppose
\(\#\mathcal N\ge B\).  Then either two distinguished centers
\(a_{j_0k}^{\rm hom}\) are \(\eta\)-separated, or
\begin{equation}
\label{eq:v081-hairbrush-clustered-reuse-threshold}
\boxed{
Bp
\lesssim_{I,J,U,\delta_s}
(R_n+\eta)^2.
}
\end{equation}
Hence, at the level of the ideal coupling, failure to obtain separated homothety centers is an
explicit scalar cap-mass loss, not an unlabeled moderate-recurrence remainder.  Passage to the
physical weighted star still requires
\eqref{eq:v081-hairbrush-coupling-to-product-target}.
\end{corollary}

\begin{proof}
If no two distinguished centers are \(\eta\)-separated, they lie in one \(\eta\)-ball.  Apply
\cref{cor:v081-hairbrush-heavy-star-center-separation}; its left side is at least
\(\sum_{k\in\mathcal N}\sigma(A_{n,k})\ge Bp\).
\end{proof}

\begin{corollary}[The local re-collision energy retains the old Maslov flags losslessly]
\label{cor:v081-hairbrush-energy-flag-lift}
Write \(z(a):=(a,b(a))\) for the selector point over direction \(a\).  For every \(k\ge1\), every
pair \(j\ne k'\), and every \(u\in W_n\), the corresponding summand in
\eqref{eq:v081-hairbrush-local-energy} is unchanged after decorating both physical selector
endpoints independently by \(k\) old flags:
\begin{equation}
\label{eq:v081-hairbrush-flagged-energy}
\boxed{
\begin{aligned}
&\int_{A_{n,j}\cap\Psi_{jk',u}^{-1}(A_{n,k'})}
\left(\prod_{\ell=1}^k d\kappa_{z(a)}(\mathfrak f_\ell)\right)
\left(\prod_{\ell=1}^k
d\kappa_{z(\Psi_{jk',u}(a))}(\mathfrak g_\ell)\right)d\sigma(a)\\
&\hspace{28mm}=
\sigma\!\left(A_{n,j}\cap\Psi_{jk',u}^{-1}(A_{n,k'})\right).
\end{aligned}
}
\end{equation}
Consequently neither the coherent cap restriction nor passage to the exact re-collision energy
discards the conditional polarization normals used by the Maslov routing.
\end{corollary}

\begin{proof}
Each \(\kappa_z\) is a probability measure.  Integrate the flag variables first, exactly as in
\cref{lem:v081-lossless-conditional-flag-tensor}.  The cap restriction changes the retained
physical endpoints but not their conditional laws, by
\cref{lem:v081-bush-cap-normalization-ledger}.
\end{proof}

\begin{corollary}[The coherent hairbrush has an exact density--re-collision dichotomy]
\label{cor:v081-hairbrush-density-recollision-dichotomy}
For the maximal families in alternative \textnormal{(ii)} of
\cref{prop:v081-maximal-bush-family}, let \(M_n\) denote the associated normalized collision edge
mass and set
\[
\Theta_n:=A_nq_nR_n^{-\delta}.
\]
After passing to a subsequence, one of the following holds:
\begin{enumerate}[label=\textnormal{(\roman*)}]
\item \(\Theta_n\to\infty\), and the off-diagonal local re-collision energy satisfies
\begin{equation}
\label{eq:v081-hairbrush-positive-recollision-energy}
\boxed{
\sum_{j\ne k}
\int_{W_n}|u-s_{n,j}|^3
\sigma\!\left(
A_{n,j}\cap\Psi_{jk,u}^{-1}(A_{n,k})
\right)du
\gtrsim_{I,J}
A_n^2q_n^2R_n^{4-2\delta}.
}
\end{equation}
\item \(\Theta_n\lesssim1\), and therefore
\begin{equation}
\label{eq:v081-hairbrush-small-mass-debt}
\boxed{
q_n\lesssim A_n^{-1}R_n^\delta,
\qquad
M_n\lesssim_{G_*}A_n^{-1}R_n^\delta.
}
\end{equation}
\end{enumerate}
The first alternative is precisely an input to the fixed-angle contact-homothety recurrence ledger;
only its balanced sign-reversing subcase may be sent to the two-time rigidity theorem.  The second
records, without hiding it in an angular label, the remaining small bush-mass debt.
\end{corollary}

\begin{proof}
By \eqref{eq:v081-hairbrush-local-energy-input},
\(S_n/R_n^3>\Theta_n\).  In the first alternative the first term in
\eqref{eq:v081-hairbrush-local-energy} eventually dominates the diagonal term, and substitution of
\(S_n>A_nq_nR_n^{3-\delta}\) proves
\eqref{eq:v081-hairbrush-positive-recollision-energy}.  In the second alternative the definition
of \(\Theta_n\) gives the first inequality in
\eqref{eq:v081-hairbrush-small-mass-debt}; the second follows from
\(q_n\gtrsim_{G_*}M_n\).
\end{proof}

\begin{corollary}[A Reeb-coherent bad cap family selects one labeled bush with an explicit \(M^{-4}\)
normalization debt]
\label{cor:v081-family-cap-single-bush-debt}
Assume the second alternative of
\cref{prop:v081-maximal-bush-family}, so that
\[
q_n:=\sum_j\sigma(H_{n,j})\gtrsim_{G_*}M_n,
\qquad
N_n\lesssim M_n^{-3}.
\]
If the Reeb-coherent cap concentration
\eqref{eq:v081-bush-family-angular-concentration} holds, then some index \(j_n\) satisfies
\begin{equation}
\label{eq:v081-selected-family-cap-mass}
\boxed{
p_n
:=
\sigma\!\left(H_{n,j_n}\cap B(a_{n,j_n},R_n)\right)
\gtrsim_{G_*}
A_nM_n^4R_n^{3-\delta}.
}
\end{equation}
After translating that cap and dilating directions by \(R_n^{-1}\), its normalized direction
measure has density
\begin{equation}
\label{eq:v081-selected-family-cap-density}
\boxed{
D_{n,\rm cap}
\lesssim_{G_*}
\frac{R_n^3}{p_n}
\lesssim_{G_*}
\frac{R_n^\delta}{A_nM_n^4}.
}
\end{equation}
The selected component retains one actual Lagrangian-bush center, its Reeb labels, the old
conditional Maslov flags, and the common outer-wing hairbrush point
\((y_n,s_n^{\rm out})\).
\end{corollary}

\begin{proof}
The left side of
\eqref{eq:v081-bush-family-angular-concentration} is a sum of \(N_n\) nonnegative cap masses.
One is therefore at least the sum divided by \(N_n\).  Use
\(q_n\gtrsim M_n\) and \(N_n\lesssim M_n^{-3}\) to obtain
\eqref{eq:v081-selected-family-cap-mass}.  The rescaling Jacobian gives
\(D_{n,\rm cap}\lesssim R_n^3/p_n\), proving
\eqref{eq:v081-selected-family-cap-density}.  All labels belong to the selected bush and are
unchanged by this restriction.
\end{proof}

\begin{corollary}[The physical bush output has one exact remaining failure mode]
\label{cor:v081-bush-frostman-or-angular-concentration}
For any sequence of approximate Lagrangian point-probe bushes satisfying the geometric hypotheses
of \cref{thm:v081-bush-frostman-escape}, fix \(0<\delta<3\).  If no subsequence satisfies its
uniform direction-Frostman hypothesis with this \(\delta\), then there are a subsequence, numbers
\(A_n\to\infty\), and radii
\[
\varepsilon_n\le R_n\downarrow0,
\]
and direction balls \(B(a_n,R_n)\) for which
\begin{equation}
\label{eq:v081-bush-angular-concentration}
\boxed{
\eta_n(B(a_n,R_n))
>
A_n
R_n^{3-\delta}.
}
\end{equation}
Thus the non-Frostman bush output is an actual angular Darboux concentration of the normalized
bush-leaf measure.  It must be passed to the angular-collapse/Maslov-pencil analysis with its bush
center and old flags retained.
\end{corollary}

\begin{proof}
This is the contrapositive of
\eqref{eq:v081-bush-direction-frostman}.  If one constant \(C_\delta\) worked after a subsequence,
\cref{thm:v081-bush-frostman-escape} would apply.  Otherwise take \(A_n\to\infty\) and make the
corresponding diagonal choice.  Since the left side of
\eqref{eq:v081-bush-angular-concentration} is at most one,
\[
R_n^{3-\delta}<A_n^{-1},
\]
so \(R_n\to0\).
\end{proof}

\begin{lemma}[Exact normalization ledger for a non-Frostman bush cap]
\label{lem:v081-bush-cap-normalization-ledger}
Let \(H\subset U\) be a bush leaf set of selector-direction mass
\[
q:=\sigma(H)>0,
\qquad
\eta:=q^{-1}\sigma|_H.
\]
Suppose that for some direction cap \(B(a_0,R)\),
\[
m_R
:=
\eta(B(a_0,R))
>
A R^{3-\delta}.
\]
Let
\[
S_R(a):=\frac{a-a_0}{R},
\qquad
\eta_R
:=
(S_R)_\#
\left(
\frac{\sigma|_{H\cap B(a_0,R)}}{qm_R}
\right).
\]
Then \(\eta_R\) is a probability measure on a fixed unit ball, absolutely continuous with respect
to three-dimensional Lebesgue measure, and
\begin{equation}
\label{eq:v081-bush-cap-rescaled-density}
\boxed{
\left\|\frac{d\eta_R}{d\mathcal L^3}\right\|_\infty
\lesssim_U
\frac{R^3}{qm_R}
\lesssim_U
\frac{R^\delta}{qA}.
}
\end{equation}
If the original bush error is \(\varepsilon\le R\), its error in the rescaled direction chart is
\(O(\varepsilon/R)\).  Moreover conditioning on the cap does not alter the pointwise old-flag law
\(\kappa_z\); consequently the prescribed-normal estimate
\eqref{eq:v080-conditional-normal-cap} remains available at every retained physical endpoint.
\end{lemma}

\begin{proof}
The unnormalized cap has direction mass \(qm_R\).  Under the dilation \(S_R\), Lebesgue density
acquires the Jacobian factor \(R^3\).  Division by \(qm_R\) gives the first bound in
\eqref{eq:v081-bush-cap-rescaled-density}; the assumed lower bound for \(m_R\) gives the second.
The bush-error statement is the same affine rescaling.  Finally \(\kappa_z\) is the conditional
distribution of flag marks over a fixed physical endpoint \(z\).  Restricting the physical
marginal to a direction cap changes which \(z\)'s are retained but does not change that conditional
distribution at any retained \(z\).
\end{proof}

\begin{principle}[Maslov--bush--to--Frostman bridge: reduction and discharge]
\label{prob:v081-maslov-bush-frostman-bridge}
Assume a fixed coarse low-reuse occurrence layer \(G_*\) violates the relative residual target
\eqref{eq:v078-final-low-reuse-relative-residual} along arbitrarily fine scales.  The maximal
\(M\)-mass family of point-probe bushes from \cref{prop:v081-maximal-bush-family}, with its old
Maslov flags and Reeb windows still attached, reduces the two contact-overlap sources in
\cref{pr:v081-bush-overlap-contact-return} to either of the following equivalent closing outputs:
\begin{enumerate}[label=\textnormal{(\roman*)}]
\item the weighted conditional tangential estimate
\eqref{eq:v078-final-fine-conditional-tangential} after all bush outputs are aggregated;
\item on one fixed positive-measure set \(G\subset J\), probability measures
\(\nu_s\) supported on \(K_{b,s}\) satisfying, for every \(\varepsilon>0\),
\[
\nu_s(B(y,R))
\lesssim_\varepsilon
R^{3-\varepsilon}
\qquad
(s\in G,\ 0<R<1).
\]
\end{enumerate}
The second output closes by
\cref{lem:v047-slice-frostman-lift}; the first closes by
\cref{prop:v078-final-coherent-factor-budget,thm:v048-residual-frostman-criterion}.
The required step is therefore a scale-coherent aggregation theorem for the actual
Reeb-escaped Lagrangian bushes, not another exclusion of coherent polarization and not a
categorical symplectic invariant.  For a strict fine-scale treatment, after
\cref{thm:v081-bush-family-frostman-escape,%
cor:v081-bush-frostman-or-angular-concentration,%
prop:v081-hairbrush-merges-to-one-bush,cor:v081-merged-bush-cap-debt}, it is enough to show that the
fine-scale fixed-angle ledger pays the merged cap-density debt
\eqref{eq:v081-merged-bush-cap-density}; the componentwise \(M^{-4}\) debt
\eqref{eq:v081-selected-family-cap-density} is avoidable.  The non-Frostman input is now exactly the
Reeb-coherent hairbrush of
\cref{cor:v081-bush-family-angular-concentration}; its scale is constrained by
\cref{cor:v081-hairbrush-bad-scale-lower-bound}.  \Cref{prop:v081-hairbrush-merges-to-one-bush} merges all coherent caps into one actual physical
bush, and \cref{cor:v081-family-to-merged-bush-frostman} closes whenever the normalized directions
of these merged bushes have a uniform Frostman subsequence.  In the remaining single-bush
non-Frostman branch, by
\cref{lem:v081-hairbrush-recollision-cocycle,lem:v081-hairbrush-maslov-blowup,%
prop:v081-hairbrush-local-recollision-energy,cor:v081-hairbrush-energy-flag-lift}, the retained
labeled caps form an exact \(\Psi\)-cocycle orbit, their canonical physical blow-up returns to the
Lagrangian pencil, their whole mass enters the preceding re-collision integral, and all old
conditional Maslov flags survive.  What remains is to split that localized integral
into the near-angle Darboux branch, the balanced sign-reversing two-time subcase controlled by
\cref{thm:v016-two-time-self,prop:v017-sign-reversal-rigid}, and the genuinely moderate
fixed-angle recurrence that must pay the explicit density debt
\eqref{eq:v081-merged-bush-cap-density}.  No generic moderate overlap is being promoted to a
small symmetric-difference defect.  \Cref{cor:v081-hairbrush-density-recollision-dichotomy} isolates the alternative small-mass debt
from the genuinely off-diagonal recurrence input.  On the latter input,
\cref{lem:v081-hairbrush-cluster-rank-one,%
cor:v081-hairbrush-heavy-star-center-separation} remove the clustered-center star unless its cap
mass already lies in a rank-one direction tube.  The remaining fixed-angle star therefore has
separated homothety centers before the old/new flag-polarization split is invoked.  \Cref{prop:v081-hairbrush-high-label-reuse} identifies the high-density energy with a common Reeb
contact-label reuse integral, while
\cref{prop:v081-hairbrush-finite-cell-transfer,%
cor:v081-hairbrush-balanced-cell-transfer,prop:v081-hairbrush-fine-maslov-graph,%
pr:v081-hairbrush-reuse-routing} remove the unknown
differentiation threshold: at every finite scale they give either a physical product collision with
the explicit cell-density factor \(\omega_n\), or high reverse contact-star reuse.  In the
fixed-angle physical branch the collision measure is now an honest normalized thickened Maslov
graph with edge mass
\(\widetilde M_{n,r}\gtrsim\omega_n(r/R_n)^2\); the small-angle complement is routed to the
Darboux chart.  At \(r=R_n\), the full coarse collision instead gives constant normalized edge
mass by \cref{cor:v081-hairbrush-coarse-maslov-graph}.  \Cref{prop:v081-hairbrush-intrinsic-renormalization} then bypasses both density factors for one
complete generation: it returns a bounded-cardinality bush family carrying a fixed normalized
mass fraction, which either satisfies vector Frostman or forms another common outer-point
hairbrush.  The latter alternates back to the original separated Reeb window, so
\cref{cor:v081-hairbrush-alternating-two-probe-lock} places its entire direction set in one
\(O(\widehat R_n)\)-cap.  Thus the strict fine-scale scalar obstruction is
\eqref{eq:v081-hairbrush-two-scale-maslov-frontier}, while the intrinsic-scale formulation leaves
no repeated fixed-angle chain.  Its remaining endpoint is the nonplanar angular Darboux blow-up at
the scale locked by \eqref{eq:v081-hairbrush-two-probe-scale-lock}.  Since an error-free
noncollapsed tangent limit would require the opposite scale separation, the planar exclusion
\cref{thm:v046-planar-angular-exclusion} cannot simply be cited here.  The remaining endpoint is
handled nodewise by \cref{prop:v081-resolution-cap-freeze-or-return}: the cap is either a good
residual leaf or a tax-free fine return.  The remaining global endpoint is therefore the weighted
aggregation of the nested good-leaf/fine-return stopping tree.  \Cref{lem:v081-stopping-leaf-quadratic-aggregation} shows that terminal-leaf cross terms create no
additional incidence loss; by \cref{cor:v081-stopping-tree-tail-criterion}, it is enough to prove
the quantitative tail/leaf-constant condition
\eqref{eq:v081-stopping-tree-tail-target}.  If that condition fails under a null-front
hypothesis, \cref{cor:v081-no-good-leaf-endpoint-chain} supplies alternating physical
endpoint-star chains to arbitrary finite depth.  The exact remaining substitute is to upgrade
the positive-logarithmic failure blocks from
\cref{lem:v081-power-excess-logarithmic-block}.  \Cref{lem:v081-subpower-label-synchronization} already gives a common shrinking label cell for
the physical centers, angular charts, and polarization planes with only subpower loss.  Thus the
stronger label-measure dichotomy
\cref{prop:v081-log-block-label-measure-dichotomy} leaves exactly two cases: a fixed positive
proportion of scales with convergent center/polarization labels, or two positive-proportion scale
families with separated labels.  The edge-flow argument below discharges both cases by producing the corresponding geometric gain or, equivalently, direct Carleson tail decay.  \Cref{cor:v081-endpoint-measure-label-dichotomy} retains the actual physical endpoint-star
measures in this alternative.  In the atomic case,
\cref{cor:v081-atomic-angular-scale-trichotomy} further reduces the thickness-to-angle map to
\(0<\varkappa_\infty<1\), the resolution endpoint \(\varkappa_\infty=1\), or the sole scale-degenerate
endpoint \(\varkappa_\infty=0\).  The contact--symplectic information now removes two more apparent
degrees of freedom.  \Cref{lem:v081-nested-probe-reeb-orbit} shows that, on a compatible
nested branch, all alternating physical probe centers converge on the Euclidean image of one Reeb
orbit.  On every branch with \(\varkappa_\infty<1\), the retained physical flags have vanishing
normalized incidence error; hence
\cref{prop:v081-four-time-polarization-firewall,cor:v081-nested-label-symplectic-reduction} show that
four distinct Reeb levels over one coherent polarization force the front-null K\"ahler locus and
horizontal rank at most two.  Thus the remaining noncollapsed part below the resolution endpoint
is reduced to genuine polarization variation or an at-most-three-level Reeb packet over one
polarization; \cref{lem:v081-coherent-plane-three-time-packets} gives those three packets the
finite-scale width \(O((R/\tau_R+d_{\Gr}(V_R,V_\infty))^{1/3})\).
\Cref{cor:v081-three-time-packets-reroute} sends the latter, when it is the retained
current cycle-plane label, back to the existing bush/angular/transverse alternatives.  A further
non-Frostman bush return is then precisely the resolution cap already isolated above.  Thus the
genuinely new below-resolution endpoint is polarization variation.  \Cref{prop:v081-same-vertex-polarization-variation-paid} and
\cref{cor:v081-polarization-variation-cross-branch} remove its same-vertex part by the old/new
symplectic transverse ledger, while
\cref{prop:v081-cross-branch-polarization-decoupling,%
cor:v081-nonatomic-polarization-component-reduction} either charge its cross-branch graph or
decouple it into coherent components.  At \(\varkappa_\infty=1\) the inherited coarse error is
resolution-sized, so the exact-root argument is deliberately not invoked.  However
\cref{lem:v081-resolution-shell-relative-quadratic,cor:v081-fine-return-error-negligible} show that
this lowest angular shell is already relative-quadratic: every genuine fine residual failure has
\(\rho/\tau_\rho\to0\) and therefore re-enters the exact contact--Maslov routing.  The final
endpoint is the mass aggregation of successive physical bush returns.  \Cref{lem:v081-bush-return-density-cocycle} combines the loss \(M_\ell\), the Frostman-failure
radius \(R_\ell\), and the returned child-cap radius \(\widehat R_\ell\) into the exact density
multiplier
\(A_\ell M_\ell R_\ell^{3-\delta}\widehat R_\ell^{-3}\).  Since Lebesgue direction density is
bounded, generations with a multiplier uniformly above one have finite depth.  \Cref{cor:v081-stopping-tail-compensation} confines the final tail to
\(M_\ell\lesssim A_\ell^{-1}\widehat R_\ell^3R_\ell^{\delta-3}\).  \Cref{lem:v081-mass-conserving-vector-frostman-stopping} removes the former one-cap selection loss:
either a positive fraction of the maximal family is vector-Frostman, or a disjoint collection of
Reeb-coherent hairbrushes retains total mass \(\gtrsim m_*M\).  \Cref{lem:v081-incoming-edge-flow-bush-extraction} retains the stronger linear edge mass as the
conditional flow measure \(\lambda_h\), and
\cref{prop:v081-one-bush-linear-edge-flow-transfer} transports it through the physical Reeb
re-collision graph without the false \(p^2\) normalization.  Its coarse return can nevertheless be
time-local by \cref{lem:v081-one-bush-self-return-time-local}.  The shared-target obstruction is
closed at the conserved-occurrence level by
\cref{lem:v081-origin-marked-bush,prob:v081-common-target-carrier-rooted-lift,%
prob:v081-strict-progress-edge-flow-gate}; consequently
\cref{thm:v081-cross-generation-edge-flow-conservation,%
cor:v081-maslov-bush-frostman-bridge-closed} close the edge-Carleson estimate and the present
bridge.
\Cref{cor:v081-hairbrush-dyadic-star-threshold} exposes the exact scalar loss if the
distinguished centers do not separate.  Thus this displayed reduction is discharged, rather than
retained as an additional hypothesis.
\end{principle}

\begin{proposition}[Exact factor budget in the final coherent polarization]
\label{prop:v078-final-coherent-factor-budget}
Let \(\lambda\) be the restriction of the low-reuse selector measure to one fixed-angle,
cell-locally coherent \(E\oplus E\) class, put \(m=\|\lambda\|\), and use the coordinates
\((x,q,y,p)\) of \eqref{eq:v056-polarized-coordinates}.  Suppose its \(q\)-marginal satisfies
\begin{equation}
\label{eq:v078-relative-normal-offset-frostman}
\lambda_q(I)
\le
D(|I|+r)
\qquad
\text{for every interval }I.
\end{equation}
Let
\[
d\Lambda_r(z,z',s)
:=
\mathbf 1_{\{|N_s(z,z')|\le r\}}
\,d\lambda(z)d\lambda(z')ds.
\]
Then
\begin{equation}
\label{eq:v078-final-normal-window-budget}
\boxed{
\|\Lambda_r\|
\lesssim_J
Dmr\log(2/r).
}
\end{equation}
If, in addition,
\begin{equation}
\label{eq:v078-final-relative-normal-density}
\boxed{D\lesssim mr^{-o(1)}}
\end{equation}
and the collision-preserved tangential shading satisfies
\begin{equation}
\label{eq:v078-final-conditional-tangential-factor}
\boxed{
\int
\mathbf 1_{\{|T_s(z,z')|\le r\}}\,d\Lambda_r(z,z',s)
\lesssim
r^{2-o(1)}\|\Lambda_r\|,
}
\end{equation}
then
\begin{equation}
\label{eq:v078-final-coherent-full-collision}
\boxed{
\int_J\iint
\mathbf 1_{\{|\Pi_s(z)-\Pi_s(z')|\le r\}}
\,d\lambda(z)d\lambda(z')ds
\lesssim
m^2r^{3-o(1)}.
}
\end{equation}
Equivalently, the relative Maslov residual bound
\eqref{eq:v078-final-low-reuse-relative-residual} holds on this coherent class.
\end{proposition}

\begin{proof}
For intervals of length at least \(r\),
\eqref{eq:v078-relative-normal-offset-frostman} gives
\(\lambda_q(I)\le 2D|I|\).  Thus
\eqref{eq:v078-final-normal-window-budget} is
\cref{lem:v055-parallel-offset-window-sum} with \(W=r\); intervals shorter than \(r\) occur only
in the first shell of that proof and are absorbed by the \(+r\) convention.  The kernel factorization
\eqref{eq:v078-final-coherent-normal-tangential-factorization}, followed by
\eqref{eq:v078-final-conditional-tangential-factor}, bounds the full collision integral by
\[
r^{2-o(1)}\|\Lambda_r\|
\lesssim
Dmr^{3-o(1)}.
\]
Use \eqref{eq:v078-final-relative-normal-density} to obtain
\eqref{eq:v078-final-coherent-full-collision}.  Finally apply the collision/residual equivalence
\eqref{eq:v048-energy-residual-equivalence} after dividing the collision integral by its universal
window factor \(r\).
\end{proof}

\begin{lemma}[Occurrence-density transfer exposes the sharp ledger offset tax]
\label{lem:v078-occurrence-density-relative-offset-tax}
Retain the low alternative in \cref{lem:v078-coherent-flag-rerooting}, and assume, as in its proof,
that this alternative carries a fixed portion of the first moment.  Write
\[
H_{\rm lo}
:=
\{0<\mathcal N_{\rm flag},\ P_i^{(H)}\le L\},
\qquad
A_{\rm lo}
:=
\int_{H_{\rm lo}}\mathcal N_{\rm flag}\,d\nu_i.
\]
Then
\begin{equation}
\label{eq:v078-low-occurrence-retained-first-moment}
\boxed{
A_{\rm lo}\gtrsim K_iH_iP_i.
}
\end{equation}
After discarding occurrence values carrying \(o(A_{\rm lo})\) and taking a joint logarithmic
dyadic layer in \(\mathcal N_{\rm flag}\) and \(P_i^{(H)}\), there are
\(\lambda_*>0\), \(0<\mathsf P_*\le L\), and
\[
G_*
\subset
\left\{
z\in H_{\rm lo}:
\lambda_*\le\mathcal N_{\rm flag}(z)<2\lambda_*,
\quad
\mathsf P_*\le P_i^{(H)}(z)<2\mathsf P_*
\right\}
\]
such that, with \(m_*:=\nu_i(G_*)\),
\begin{equation}
\label{eq:v078-occurrence-layer-mass-identity}
\boxed{
\lambda_*m_*
\asymp
\int_{G_*}\mathcal N_{\rm flag}\,d\nu_i
=
K_iH_iP_i\,W^{o(1)},
\qquad
\lambda_*\le 4H_i\mathsf P_*.
}
\end{equation}
For every interval \(I\) in the coherent normal coordinate \(q\),
\begin{equation}
\label{eq:v078-occurrence-to-physical-q-slab}
\boxed{
(\nu_i|_{G_*})_q(I)
\le
\lambda_*^{-1}
\int_{G_*\cap\{q\in I\}}
\mathcal N_{\rm flag}\,d\nu_i.
}
\end{equation}
Consequently the relative normal estimate needed in
\eqref{eq:v078-relative-normal-offset-frostman}--\eqref{eq:v078-final-relative-normal-density}
follows on this layer from the single weighted slab estimate
\begin{equation}
\label{eq:v078-weighted-occurrence-q-slab-target}
\boxed{
\int_{G_*\cap\{q\in I\}}
\mathcal N_{\rm flag}\,d\nu_i
\lesssim
\left(\int_{G_*}\mathcal N_{\rm flag}\,d\nu_i\right)
(|I|+r)r^{-o(1)}
\quad\text{for every }I.
}
\end{equation}
The pointwise low-reuse cap alone proves only an absolute estimate.  Indeed bounded direction
density and \eqref{eq:v078-flag-density-by-rerooted-pair-mass} give
\begin{equation}
\label{eq:v078-trivial-occurrence-q-slab}
\int_{G_*\cap\{q\in I\}}
\mathcal N_{\rm flag}\,d\nu_i
\lesssim_U
H_i\mathsf P_*(|I|+r),
\end{equation}
so after \eqref{eq:v078-occurrence-layer-mass-identity} its relative-density loss is bounded by
the layer-sharp tax
\begin{equation}
\label{eq:v078-exact-relative-offset-tax}
\boxed{
\mathfrak g_{\rm rel}
:=
\frac{\mathsf P_*}{K_iP_i}\,W^{-o(1)}.
}
\end{equation}
The weighted dense/sparse argument below completes the coarse low old-flag reuse bound by proving
\eqref{eq:v078-weighted-occurrence-q-slab-target}, or gain the reciprocal factor
\(\mathfrak g_{\rm rel}^{-1}\) in the conditional tangential estimate before the normal and
tangential factors are multiplied.
\end{lemma}

\begin{proof}
The first assertion is exactly the low-half case in the proof of
\cref{lem:v078-coherent-flag-rerooting}.  Dyadically decompose the nonzero values of
\(\mathcal N_{\rm flag}\) and \(P_i^{(H)}\); values below sufficiently small fixed powers of \(W\)
carry \(o(A_{\rm lo})\), and the remaining number of joint layers is \(W^{-o(1)}\).  One joint
layer therefore gives the lower bound in
\eqref{eq:v078-occurrence-layer-mass-identity}; its upper bound uses
\eqref{eq:v078-flag-lift-first-moment}, while its pointwise upper bound uses
\eqref{eq:v078-flag-density-by-rerooted-pair-mass}.  On \(G_*\),
\(1\le\lambda_*^{-1}\mathcal N_{\rm flag}\), which proves
\eqref{eq:v078-occurrence-to-physical-q-slab}.  Substitution of
\eqref{eq:v078-weighted-occurrence-q-slab-target} and
\eqref{eq:v078-occurrence-layer-mass-identity} yields the required relative marginal estimate.

For the final assertion, the map \(a\mapsto q=n\cdot a\) sends a \(q\)-interval of length
\(|I|\) to a direction slab of measure \(O_U(|I|+r)\) after scale-\(r\) enlargement.  Multiply this
by the pointwise bound
\(\mathcal N_{\rm flag}\le2H_iP_i^{(H)}\le4H_i\mathsf P_*\) to obtain
\eqref{eq:v078-trivial-occurrence-q-slab}.  Dividing its coefficient by the total weighted mass in
\eqref{eq:v078-occurrence-layer-mass-identity} gives
\(\mathsf P_* /(K_iP_i)\) up to subpower factors, which is
\eqref{eq:v078-exact-relative-offset-tax}.
\end{proof}

\begin{proposition}[Exact cancellation of the relative offset tax]
\label{prop:v078-relative-offset-tax-cancellation}
Use the occurrence layer \(G_*\) from
\cref{lem:v078-occurrence-density-relative-offset-tax}, put
\(\lambda:=\nu_i|_{G_*},\ m_*:=\|\lambda\|\), and let \(\Lambda_r\) be its normal-window pair
measure.  The absolute slab estimate \eqref{eq:v078-trivial-occurrence-q-slab} gives an admissible
normal-offset constant
\begin{equation}
\label{eq:v078-absolute-to-relative-normal-constant}
\boxed{
D_*
\lesssim_U
\frac{H_i\mathsf P_*}{\lambda_*}
\lesssim
m_*\frac{\mathsf P_*}{K_iP_i}\,W^{-o(1)}.
}
\end{equation}
Consequently
\begin{equation}
\label{eq:v078-taxed-normal-window-budget}
\boxed{
\|\Lambda_r\|
\lesssim
m_*^2\frac{\mathsf P_*}{K_iP_i}\,rW^{-o(1)}.
}
\end{equation}
If the collision-preserved tangential shading satisfies the tax-cancelling estimate
\begin{equation}
\label{eq:v078-tax-cancelling-tangential-target}
\boxed{
\int
\mathbf 1_{\{|T_s(z,z')|\le r\}}\,d\Lambda_r(z,z',s)
\lesssim
\frac{K_iP_i}{\mathsf P_*}\,r^{2-o(1)}\|\Lambda_r\|,
}
\end{equation}
then the full coherent collision bound
\eqref{eq:v078-final-coherent-full-collision}, and hence the relative residual target
\eqref{eq:v078-final-low-reuse-relative-residual}, holds on \(G_*\).
Thus relative \(q\)-equidistribution is sufficient but not necessary: the existing
contact-polarized dense/sparse route may instead close the branch by producing exactly the factor
\(K_iP_i/\mathsf P_*\) in \eqref{eq:v078-tax-cancelling-tangential-target}; its sparse alternative remains a
physical Lagrangian-bush exit.
\end{proposition}

\begin{proof}
Divide \eqref{eq:v078-trivial-occurrence-q-slab} by \(\lambda_*\) and use
\eqref{eq:v078-occurrence-to-physical-q-slab}; this gives the first bound in
\eqref{eq:v078-absolute-to-relative-normal-constant}.  The second follows from
\(\lambda_*m_*\gtrsim K_iH_iP_iW^{o(1)}\).  Apply
\cref{lem:v055-parallel-offset-window-sum} at scale \(r\) to obtain
\eqref{eq:v078-taxed-normal-window-budget}.  Multiplying it by the conditional fraction in
\eqref{eq:v078-tax-cancelling-tangential-target} cancels
\(\mathsf P_* /(K_iP_i)\) and gives
\[
\int_J\iint
\mathbf 1_{\{|\Pi_s(z)-\Pi_s(z')|\le r\}}
\,d\lambda(z)d\lambda(z')ds
\lesssim
m_*^2r^{3-o(1)}.
\]
Now use \cref{prop:v048-energy-residual-equivalence} exactly as in
\cref{prop:v078-final-coherent-factor-budget}.
\end{proof}

\begin{proposition}[The final collision lifts exactly to two old-flag occurrences]
\label{prop:v078-double-occurrence-collision-lift}
Keep the notation of \cref{lem:v078-occurrence-density-relative-offset-tax}.  Write
\(\mathfrak f=(y,f)\) for a rooted flag and define the restricted occurrence measure
\begin{equation}
\label{eq:v078-restricted-flag-occurrence-measure}
d\rho_*(\mathfrak f,z)
:=
\mathbf 1_{G_*}(z)w_f(z)\,d\Lambda(y,f)d\nu_i(z).
\end{equation}
Its physical marginal and total mass are
\begin{equation}
\label{eq:v078-occurrence-physical-marginal}
\boxed{
(\pi_z)_\#\rho_*
=
\mathbf 1_{G_*}\mathcal N_{\rm flag}\,d\nu_i,
\qquad
\|\rho_*\|
=
\int_{G_*}\mathcal N_{\rm flag}\,d\nu_i
\asymp
\lambda_*m_*.
}
\end{equation}
Its rooted-flag marginal retains the actual packet weights and obeys only the honest bound
\begin{equation}
\label{eq:v078-occurrence-flag-marginal}
(\pi_{\mathfrak f})_\#\rho_*
=
\left(\int_{G_*}w_f(z)\,d\nu_i(z)\right)d\Lambda(y,f)
\le
2H_i\,d\Lambda(y,f).
\end{equation}
For every nonnegative kernel \(K(z,z',s)\),
\begin{equation}
\label{eq:v078-double-occurrence-kernel-comparison}
\boxed{
\int_J\iint_{G_*\times G_*}K(z,z',s)\,d\nu_i(z)d\nu_i(z')ds
\asymp
\lambda_*^{-2}
\int_J\iint K(z,z',s)\,d\rho_*(\mathfrak f,z)d\rho_*(\mathfrak f',z')ds.
}
\end{equation}
In particular, with the normal kernel retained, the tax-cancelling target
\eqref{eq:v078-tax-cancelling-tangential-target} is equivalent, up to constants and subpower
losses, to
\begin{equation}
\label{eq:v078-six-point-weighted-tangential-target}
\boxed{
\int
\mathbf 1_{\{|N_s(z,z')|\le r\}}
\mathbf 1_{\{|T_s(z,z')|\le r\}}
\,d\rho_*d\rho_*'ds
\lesssim
\frac{K_iP_i}{\mathsf P_*}\,r^{2-o(1)}
\int
\mathbf 1_{\{|N_s(z,z')|\le r\}}
\,d\rho_*d\rho_*'ds.
}
\end{equation}
This is the collision-preserved old-flag/new-cycle six-point estimate with all occurrence weights
still present.  No unweighted selector measure and no categorical symplectic invariant has been
substituted for it.
\end{proposition}

\begin{proof}
The first identity follows from the definition of
\(\mathcal N_{\rm flag}\); the second is Tonelli.  Equation
\eqref{eq:v078-occurrence-flag-marginal} uses
\(\int w_f\,d\nu_i<2H_i\).  Since
\(\lambda_*\le\mathcal N_{\rm flag}<2\lambda_*\) on \(G_*\), multiplying two copies of this
comparison and integrating against \(K\) proves
\eqref{eq:v078-double-occurrence-kernel-comparison}.  Apply it first to the product of the normal
and tangential kernels and then to the normal kernel alone; the common factor
\(\lambda_*^{-2}\) cancels, yielding
\eqref{eq:v078-six-point-weighted-tangential-target}.
\end{proof}

\begin{corollary}[Fine collision scales absorb every coarse flag-power tax]
\label{cor:v078-fine-scale-absorbs-coarse-tax}
Let \(W_n\downarrow0\) be coarse old-flag scales, let \(G_{*,n}\) be occurrence layers as above, and
choose the new collision scales \(r_n\) by
\cref{lem:v078-coarse-fine-mass-freezing}.  All retained quantities have been placed on fixed
\(W_n\)-power layers.  In particular, for some fixed \(C<\infty\),
\[
W_n^{o(1)}
\lesssim
\frac{\mathsf P_{*,n}}{K_{i,n}P_{i,n}}
\le
W_n^{-C+o(1)}.
\]
Therefore
\begin{equation}
\label{eq:v078-coarse-tax-is-fine-subpower}
\boxed{
\frac{\mathsf P_{*,n}}{K_{i,n}P_{i,n}}
=
r_n^{-o(1)},
\qquad
\frac{K_{i,n}P_{i,n}}{\mathsf P_{*,n}}
=
r_n^{o(1)}.
}
\end{equation}
The normal constant in
\eqref{eq:v078-absolute-to-relative-normal-constant} consequently satisfies
\begin{equation}
\label{eq:v078-fine-relative-normal-density}
\boxed{
D_{*,n}
\lesssim
m_{*,n}r_n^{-o(1)}.
}
\end{equation}
Hence the relative normal-offset gate
\eqref{eq:v078-final-relative-normal-density} is automatic on the coarse/fine scale pairs.  At
those pairs the tax-cancelling target
\eqref{eq:v078-six-point-weighted-tangential-target} is equivalent, at subpower precision, to the
single collision-preserved conditional tangential estimate
\begin{equation}
\label{eq:v078-final-fine-conditional-tangential}
\boxed{
\int
\mathbf 1_{\{|N_s|\le r_n\}}
\mathbf 1_{\{|T_s|\le r_n\}}
\,d\rho_{*,n}d\rho_{*,n}'ds
\lesssim
r_n^{2-o(1)}
\int
\mathbf 1_{\{|N_s|\le r_n\}}
\,d\rho_{*,n}d\rho_{*,n}'ds.
}
\end{equation}
Thus no fixed power inherited from the old flag generation remains in the final fine-scale
numerical ledger.  If one coarse \(W,H,G_*\) is frozen instead, the same conclusions hold, with
constants depending on that coarse layer, uniformly at the exponent level for all sufficiently
small \(r\).  This fixed-coarse/all-fine form, rather than the diagonal sequence alone, is the one
that can feed the Hausdorff Frostman passage.
\end{corollary}

\begin{proof}
The lower bound on the ratio follows, up to a subpower factor, from
\eqref{eq:v078-occurrence-layer-mass-identity},
\(\lambda_*\le4H_i\mathsf P_*\), and \(m_*\le1\).  Its upper bound follows from the fixed power
dyadic ranges of \(K_i,P_i,\mathsf P_*\).  By
\eqref{eq:v078-frozen-mass-is-subpower},
\[
\frac{|\log W_n|}{|\log r_n|}
\le
\frac1n.
\]
Taking logarithms proves \eqref{eq:v078-coarse-tax-is-fine-subpower}.  Substitute the first part
into \eqref{eq:v078-absolute-to-relative-normal-constant} to obtain
\eqref{eq:v078-fine-relative-normal-density}.  The second part is absorbed into the \(r_n^{-o(1)}\)
allowance in \eqref{eq:v078-six-point-weighted-tangential-target}, which gives
\eqref{eq:v078-final-fine-conditional-tangential}.  For a fixed coarse layer, replace the ratio
\(|\log W_n|/|\log r_n|\) by a fixed numerator divided by \(|\log r|\to\infty\).
\end{proof}

\begin{lemma}[A \(q\)-separated anchor plane is transverse to the coherent polarization]
\label{lem:v066-anchor-plane-q-transverse}
Let \(E_*\subset\mathbb R^3\) have unit normal \(n_*\), and write
\(\mathfrak q(a)=n_*\cdot a\).  If a two-plane \(E\) contains a nonzero vector \(u\) with
\[
|\mathfrak q(u)|\ge R_q
\]
and \(u\) lies in a fixed bounded set, then, for either unit normal \(n_E\) of \(E\),
\begin{equation}
\label{eq:v066-anchor-plane-angle}
\boxed{
|n_E\times n_*|
\gtrsim_U R_q.
}
\end{equation}
For the plane \(E\) in
\cref{lem:v066-direct-horizontal-packet-thickness}, every retained rank-loss cycle from
\cref{prop:v064-phase-labeled-rank-loss-cycles} satisfies this conclusion.
\end{lemma}

\begin{proof}
The norm of the restriction of \(a\mapsto n_*\cdot a\) to \(E\) is
\[
\sup_{\substack{e\in E\\|e|=1}}|n_*\cdot e|
=|n_E\times n_*|.
\]
Testing it on \(u/|u|\), and using the boundedness of \(|u|\), proves
\eqref{eq:v066-anchor-plane-angle}.  In the anchored packet take
\(u=a_{y_0}-a_p\).  This is one of the four retained edges, so
\eqref{eq:v064-four-shear-annuli} gives
\(|\mathfrak q(u)|\gtrsim R_q\).
\end{proof}

\begin{proposition}[The polarized rank packet returns to the transverse Maslov ledger]
\label{prop:v066-rank-packet-transverse-return}
Assume \(1/2\le\chi<3/5\), choose \(\alpha\) as in
\eqref{eq:v064-alpha-gap}, and put \(R_q=W^\alpha\).  Choose
\begin{equation}
\label{eq:v066-transverse-parameter-choice}
\boxed{
\alpha<d<1,\qquad
0<\beta<d-\alpha,\qquad
\Delta=W^d,\qquad
\theta_{\rm hr}=W^\beta.
}
\end{equation}
Then \(W/\Delta\to0\), as required by the arbitrary-threshold Maslov routing, while every
polarized-plane rank packet obtained from the surviving four-cycles has
\begin{equation}
\label{eq:v066-thin-transverse-packet}
\boxed{
t_{\rm dir}
\lesssim
W+W^{d-\beta}
=o(W^\alpha)=o(R_q),
\qquad
|n_E\times n_*|\gtrsim R_q.
}
\end{equation}
Thus this output cannot remain in the coherent polarization class at angular resolution
\(o(R_q)\); it is a genuinely transverse Maslov packet.

More precisely, if a reused direction cell is also carried by an affine \(E_*\)-packet of
thickness \(w_*\), then its actual direction mass is bounded by
\begin{equation}
\label{eq:v066-unequal-transverse-intersection}
\boxed{
\sigma(H_{\rm cell})
\lesssim_U
\frac{w_*t_{\rm dir}}{R_q}.
}
\end{equation}
Consequently the weighted high-reuse summation in
\cref{cor:v053-transverse-reuse-weighted-sum} holds here with
\(w^2/\vartheta\) replaced by \(w_*t_{\rm dir}/R_q\).  The formerly coherent
polarized-plane output is therefore returned to the already isolated transverse-polarization
ledger with the additional vanishing factor
\begin{equation}
\label{eq:v066-transverse-return-gain}
\boxed{
\frac{t_{\rm dir}}{R_q}
\lesssim
W^{1-\alpha}+W^{d-\beta-\alpha}
=o(1).
}
\end{equation}
\end{proposition}

\begin{proof}
The first inequality in \eqref{eq:v066-thin-transverse-packet} is
\eqref{eq:v066-actual-horizontal-packet-thickness}; the strict inequalities in
\eqref{eq:v066-transverse-parameter-choice} make both powers \(1\) and \(d-\beta\) larger than
\(\alpha\).  The angle bound is
\cref{lem:v066-anchor-plane-q-transverse}.

For completeness, the unequal-width version of
\cref{lem:v053-transverse-polarization-intersection} follows from the same coordinates:
the two normal coordinates range over intervals of lengths \(O(w_*)\) and \(O(t_{\rm dir})\), and
their Jacobian is \(|n_E\times n_*|\gtrsim R_q\).  Bounded horizontal density gives
\eqref{eq:v066-unequal-transverse-intersection}.  The proof of
\cref{cor:v053-transverse-reuse-weighted-sum} then applies verbatim, and division of
\eqref{eq:v066-thin-transverse-packet} by \(R_q\) gives
\eqref{eq:v066-transverse-return-gain}.
\end{proof}

\begin{corollary}[The final Hausdorff loss is one residual estimate]
\label{cor:v048-final-residual-endpoint}
The fixed-angle and angular analyses have the same remaining quantitative endpoint: for every
\(\eta>0\), prove
\begin{equation}
\label{eq:v048-final-residual-endpoint}
\mathfrak Z_\rho^{J^+}(\mu)
\lesssim_\eta
\rho^{2-\eta}
\end{equation}
after summing the fixed-angle shells and their Darboux-rescaled angular counterparts.  By
\cref{thm:v048-residual-frostman-criterion}, this estimate closes the full Hausdorff conclusion
directly.  By \cref{prop:v048-energy-residual-equivalence}, every violation of
\eqref{eq:v048-final-residual-endpoint} is a quantitatively excessive Maslov collision shell.  The
the preceding stage machinery routes the part above its geometric thresholds to a Lagrangian bush, angular
collapse, or a persistent nonplanar Reeb-front seed.  The precise remainder is the union
of that persistent seed branch and the diffuse exponent-gap window
\eqref{eq:v048-exact-exponent-gap}; both must be upgraded to the quadratic normal-slice estimate.
On the latter window, \cref{cor:v049-residual-excess-normal-flag} already forces a divergent staged
density in the two symplectic normal coordinates, while
\cref{prop:v050-arbitrary-threshold-maslov-routing} sends its collision graph to a shrinking bush or
determinant-small cycles at every useful \(\Delta\).  In the separated-anchor, separated-time part,
\cref{cor:v050-rank-loss-cycle-to-packet} gives the explicit saving
\eqref{eq:v050-rank-loss-saving}.  The remaining content is the mass-preserving weighted
close-anchor refinement and the final summation of the rank-packet cap to the quadratic normal-slice
estimate; the time-packet part is closed by
\cref{cor:v050-two-time-scale-persistence}, and another scalar fourth-root estimate cannot supply
the remaining summation.  the preceding flag construction retain this summation as an
old-flag/new-cycle occurrence measure.  the preceding stage uses the lossless conditional eight-flag lift
\cref{lem:v081-lossless-conditional-flag-tensor} to close its fully coherent rank-loss part, but
\cref{prop:v081-rank-loss-to-bush-reduction} shows that the other alternative is a genuine physical
point-probe bush.  Converting those bushes into scale-coherent transverse measures is precisely
\cref{prob:v081-maslov-bush-frostman-bridge}.
\end{corollary}

\begin{theorem}[Planar angular-collapse exclusion]
\label{thm:v046-planar-angular-exclusion}
Suppose a longitudinally saturated angular-collapse limit has all its secant vectors in one
three-plane \(V\).  Then
\begin{equation}
\label{eq:v046-planar-front-null}
\boxed{
P_V\equiv0,
\qquad
\|\omega|_V\|_{\rm op}=1,
\qquad
V\in
\mathsf K_{101}\cup\mathsf K_{110}\cup\mathsf K_{200},
\qquad
\rank(\pi_a|_V)\le2.
}
\end{equation}
In particular such a planar limit cannot carry a noncollapsed three-dimensional direction
blow-up.  Hence the rectifiable and quantitatively planar angular-collapse branches of a
full-direction selector are excluded.
\end{theorem}

\begin{proof}
By \eqref{eq:v046-limit-maslov-support}, for every \(s\in J\) the plane \(V\) contains a nonzero
vector of \(L_s\).  The matrix-pencil/Maslov equivalence
\eqref{eq:v003-maslov-equivalence} gives \(P_V(s)=0\) on \(J\), so the cubic polynomial vanishes
identically.  The complex-evaluation identity, the preceding classification, and the
front-null direction-rank law \eqref{eq:v005-direction-rank} give the four conclusions in
\eqref{eq:v046-planar-front-null}.  A noncollapsed three-dimensional horizontal blow-up would make
\(\pi_a|_V\) have rank three, which is impossible.
\end{proof}

\begin{lemma}[Polarized form of the characteristic collision directions]
\label{lem:v046-polarized-characteristics}
Let \(V\) be front-null.
\begin{enumerate}[label=\textnormal{(\roman*)}]
\item In the aligned strata \(\mathsf K_{101}\) and \(\mathsf K_{110}\), choose the gauge
\[
u=(e_1,0),
\qquad
Ju=(0,e_1),
\qquad
w=(c,d),
\qquad
c,d\perp e_1.
\]
At every parameter at which the pencil has rank two,
\[
\boxed{
\kappa_V(s)
=
\mathbb P\operatorname{span}\{u-sJu\}.
}
\]
Thus the entire generic Maslov collision curve lies in the unique complex line
\(\mathcal C_V=\mathbb Cu\); the symplectic-null direction \(w\) is invisible to the generic
collision pairs.

\item In \(\mathsf K_{110}\), the rank is two for every real \(s\), while the cokernel line also
moves with projective degree one.  Hence this is the only stratum with simultaneous moving
right- and left-incidence data.

\item In \(\mathsf K_{101}\), writing
\(w=(\mu e_2,\nu e_2)\), the third pencil column is
\((\nu+s\mu)e_2\).  There is at most one finite exceptional parameter
\(s_*=-\nu/\mu\) at which the rank drops to one; at that parameter the null direction \(w\) enters
\(V\cap L_{s_*}\).  This is the only aligned mechanism that can add a third collision direction at
one Reeb level.

\item In \(\mathsf K_{200}\), there is a fixed physical output plane \(W_V\subset\R^3\) such that
\[
\boxed{
\operatorname{im}(B+sA)\subset W_V
\qquad(s\in\R),
}
\]
and hence one fixed projective cokernel normal
\(n_V\in\mathbb P((\R^3)^*)\).  Its right characteristic curve has degree two.
\end{enumerate}
\end{lemma}

\begin{proof}
In the aligned gauge the pencil is
\[
B+sA=[s e_1,\ e_1,\ d+s c].
\]
Whenever it has rank two, its right kernel is generated by \((1,-s,0)\), which represents
\(u-sJu\) in \(V\).  For \(\mathsf K_{110}\), the vectors \(c,d\) are independent, so
\(d+sc\ne0\) and the rank never drops.  In \(\mathsf K_{101}\), substitute
\(c=\mu e_2\), \(d=\nu e_2\).  The exceptional statement follows immediately.  The final item is
\cref{prop:v005-forced-output-plane,thm:v005-transverse-K200} together with the characteristic-degree
theorem of the preceding stage.
\end{proof}

\begin{proposition}[Density-bearing Maslov--Alberti branch closes]
\label{prop:v046-alberti-gate}
Suppose there is a positive-\(\mu\)-mass set \(E\subset G_b\) on which
\[
0<\liminf_{r\downarrow0}\frac{\mu(B(x,r))}{r^3}
\le
\limsup_{r\downarrow0}\frac{\mu(B(x,r))}{r^3}
<\infty,
\]
and \(\mu|_E\) admits three independent Alberti representations.  Then \(E\) has a positive-mass
countably three-rectifiable part.  Its horizontal projection has positive Lebesgue measure, and the
affine Reeb sweep of that part has positive four-dimensional measure.  Hence it cannot occur in a
null-front counterexample.
\end{proposition}

\begin{proof}
The density and independent-representation hypotheses give countable three-rectifiability by the
rectifiability and Alberti-representation results of Bate and Bate--Li
\cite{Bate2015,BateLi2017}.  Because the graph projection sends
\(\mu|_E\) to Lebesgue measure on a positive-measure direction set, the area formula gives a
positive-measure rectifiable subset on which the horizontal tangent map has rank three.  On that
subset the leading coefficient \(\det A\) of the Reeb-front polynomial is nonzero.  The rectifiable
Reeb-sweep argument of
\cref{thm:v003-rectifiable-polynomial,thm:v005-no-rectifiable-selector} then gives positive
four-volume.
\end{proof}

\begin{proof}[Proof of \cref{thm:main}]
By \cref{prop:selector}, pass to one marked line in every direction.  The selected
line graph still has packing dimension three and its Euclidean front is contained in
the original front.  The contact--symplectic closure theorem
\cref{thm:v081-sticky-kakeya-closure} applies to this selector and gives dimension four
for its Reeb front.  The containing front therefore has Hausdorff dimension four as well.
\end{proof}

\section{Discussion}

The reduction uses three manifestations of the same contact--symplectic structure. The evaluation
map turns Reeb motion into Euclidean line motion. The matrix pencil $B+sA$ measures infinitesimal
front rank and identifies total collapse with Maslov trapping. Finally, the rescaled collision
equation is the same pencil equation, so the polarization attached to a collision survives the
passage from local geometry to the multiscale edge flow.

Two safeguards are important. First, being three-dimensional in line space does not make a
rough selector Lagrangian or rectifiable. Second, a high-degree graph vertex is not credited as a
Euclidean bush until the contact evaluation produces a common physical point-probe. These are
the places where a purely combinatorial reformulation would lose the geometry needed for the
Hausdorff endpoint.  The compensation region is closed at the level at which it actually occurs:
the ordered, absolute edge-flow measure.  The same identity supplies the weighted marked
finite-scale estimate.  What it does not supply is any typed Hall/Radon--Nikodym comparison; no
such algebraic conclusion is claimed here.

\appendix

\section{Measurable disintegration and passive marks}
\label{app:measure}

This appendix collects the measure-theoretic facts used in the main iteration; the
disintegration conventions are standard, as in \cite{Bogachev2007}.  They are grouped
by function rather than by the stage at which they first arose.

\begin{lemma}[Restriction--kernel calculus]
\label{lem:restriction-kernel}
Let \(d\Gamma(x,m)=d\lambda(x)dK_x(m)\), where \(K_x\) is a probability kernel.  If
\(0\le f\le1\) and \(Q_{x,m}\) is another probability kernel, then
\[
 d\widetilde\Gamma(x,m,n)=f(x,m)d\lambda(x)dK_x(m)dQ_{x,m}(n)
\]
has source marginal dominated by \(\lambda\).  For a countable disjoint family of such
restrictions, the source Radon--Nikodym densities add pointwise to at most one.
\end{lemma}

\begin{proof}
Integrating first in \(n\) removes \(Q_{x,m}\).  Integrating in \(m\) leaves the density
\(x\mapsto\int f(x,m)dK_x(m)\), which lies in \([0,1]\).  Additivity for disjoint restrictions
follows from monotone convergence.
\end{proof}

\begin{lemma}[Countable adaptive cap partition]
\label{lem:adaptive-cap-partition}
Let \(X\) be a separable metric space and let \(t(x)>0\) be measurable.  There is a countable
half-open partition \(X=\bigsqcup_jE_j\) and centers \(x_j\) such that
\(E_j\subset B(x_j,Ct(x))\) on a measurable layer where \(t\) is comparable to \(t(x_j)\).
Applying the partition to the source coordinate of an occurrence measure preserves its total
mass and every conditional mark.
\end{lemma}

\begin{proof}
Split into dyadic layers of \(t\).  On each layer choose a countable maximal net and order its
balls.  Subtract all earlier balls from the current one.  The resulting half-open cells are
disjoint, measurable, and cover the layer.  The last assertion is
\cref{lem:restriction-kernel} with indicator densities.
\end{proof}

\begin{remark}
The affine fibre label and the complete carrier history are represented by a Dirac kernel
\(x\mapsto\delta_{(\mathfrak f(x),\mathscr Q(x))}\).  Since its mass is one, it survives every
restriction and kernel extension.  This is the precise meaning of a passive mark in the proof.
\end{remark}


\section{Uniform finite-scale parameter hierarchy}
\label{app:parameters}

Fix a physical thickness \(\delta\), a bush exponent \(0<d_0<1/3\), and
\[
 B(\delta)=\exp\sqrt{\log(1/\delta)}.
\]
At a node of radius \(T_v\), first prescribe
\(\tau_v=T_v/(C_*B(\delta))\).  Only then choose the available failure scale so that the
physical, synchronization, polarization, target-cell, and time-packet errors are bounded by
\(T_v/(C_*^2B(\delta))\).  This order of choices is important: the fine scale is a local
certificate, whereas \(T_v\) is the scale carried by the absolute flow.

Before termination, \(T_v^3>\delta^{2-o(1)}\), while the physical bush error is
\(R_{\rm bush}=O(\delta^{1-d_0})\).  Hence
\[
 \frac{R_{\rm bush}}{\tau_v}
 \lesssim B(\delta)\delta^{1/3-d_0-o(1)}=o(1).
\]
Thus every local geometric error fits inside the prescribed child cap without passing below the
physical thickness.  Since \(T_w\le B(\delta)^{-1}T_v\), after
\[
 L=O(\sqrt{\log(1/\delta)})
\]
generations one has \(T_L^3\le\delta^{2-o(1)}\).  At that frontier,
\[
 \sum_\gamma p_\gamma T_\gamma^3
 \le \delta^{2-o(1)}\sum_\gamma p_\gamma,
\]
and the accumulated degree tail is
\(L\delta^A=\delta^{A-o(1)}\).  Both estimates are uniform in the number of children, their
shapes, and every passive carrier mark.

\begin{thebibliography}{99}

\bibitem{Arnold1990}
V.~I. Arnol'd,
\emph{Singularities of Caustics and Wave Fronts},
Mathematics and Its Applications (Soviet Series), vol.~62,
Kluwer Academic Publishers, Dordrecht, 1990.

\bibitem{Bate2015}
D.~Bate,
\emph{Structure of measures in Lipschitz differentiability spaces},
J. Amer. Math. Soc. \textbf{28} (2015), no.~2, 421--482.

\bibitem{BateLi2017}
D.~Bate and S.~Li,
\emph{Characterizations of rectifiable metric measure spaces},
Ann. Sci. \`Ec. Norm. Sup\'er. (4) \textbf{50} (2017), no.~1, 1--37.

\bibitem{BergerTrenn2012}
T.~Berger and S.~Trenn,
\emph{The quasi-Kronecker form for matrix pencils},
SIAM J. Matrix Anal. Appl. \textbf{33} (2012), no.~2, 336--368.

\bibitem{Bogachev2007}
V.~I. Bogachev,
\emph{Measure Theory}, vols.~I--II,
Springer, Berlin, 2007.

\bibitem{Bourgain1991}
J.~Bourgain,
\emph{Besicovitch type maximal operators and applications to Fourier analysis},
Geom. Funct. Anal. \textbf{1} (1991), no.~2, 147--187.

\bibitem{CannasDaSilva2001}
A.~Cannas da Silva,
\emph{Lectures on Symplectic Geometry},
Lecture Notes in Mathematics, vol.~1764, Springer, Berlin, 2001.

\bibitem{Federer1969}
H.~Federer,
\emph{Geometric Measure Theory},
Grundlehren der mathematischen Wissenschaften, vol.~153,
Springer-Verlag, New York, 1969.

\bibitem{Geiges2008}
H.~Geiges,
\emph{An Introduction to Contact Topology},
Cambridge Studies in Advanced Mathematics, vol.~109,
Cambridge University Press, Cambridge, 2008.

\bibitem{GohbergLancasterRodman2009}
I.~Gohberg, P.~Lancaster, and L.~Rodman,
\emph{Matrix Polynomials},
Classics in Applied Mathematics, vol.~58,
Society for Industrial and Applied Mathematics, Philadelphia, 2009.

\bibitem{GuthWangZahl2026}
L.~Guth, H.~Wang, and J.~Zahl,
\emph{A streamlined proof of the Kakeya set conjecture in \(\mathbb R^3\)},
arXiv:2601.14411 (2026).

\bibitem{KatzLabaTao2000}
N.~H. Katz, I.~\L aba, and T.~Tao,
\emph{An improved bound on the Minkowski dimension of Besicovitch sets in
\(\mathbb R^3\)},
Ann. of Math. (2) \textbf{152} (2000), no.~2, 383--446.

\bibitem{KatzTao2002}
N.~H. Katz and T.~Tao,
\emph{New bounds for Kakeya problems},
J. Anal. Math. \textbf{87} (2002), 231--263.

\bibitem{Kechris1995}
A.~S. Kechris,
\emph{Classical Descriptive Set Theory},
Graduate Texts in Mathematics, vol.~156,
Springer-Verlag, New York, 1995.

\bibitem{Mattila1995}
P.~Mattila,
\emph{Geometry of Sets and Measures in Euclidean Spaces: Fractals and Rectifiability},
Cambridge Studies in Advanced Mathematics, vol.~44,
Cambridge University Press, Cambridge, 1995.

\bibitem{VanDooren1979}
P.~Van Dooren,
\emph{The computation of Kronecker's canonical form of a singular pencil},
Linear Algebra Appl. \textbf{27} (1979), 103--140.

\bibitem{WangZahlSticky}
H.~Wang and J.~Zahl,
\emph{Sticky Kakeya sets and the sticky Kakeya conjecture},
J. Amer. Math. Soc. \textbf{39} (2026), no.~2, 515--585.

\bibitem{Wolff1995}
T.~Wolff,
\emph{An improved bound for Kakeya type maximal functions},
Rev. Mat. Iberoamericana \textbf{11} (1995), no.~3, 651--674.

\end{thebibliography}

\end{document}
